% Ordered Associative Chiral Koszul Duality
% Integration-ready chapter file (stripped from standalone amsart draft).
% Uses only \providecommand for macros that may not be in main.tex preamble.

\providecommand{\Assch}{\mathrm{Ass}^{\mathrm{ch}}}
\providecommand{\Barch}{\overline{B}^{\mathrm{ch}}}
\providecommand{\Cobar}{\Omega^{\mathrm{ch}}}
\providecommand{\coHoch}{\operatorname{coHH}}
\providecommand{\Cotor}{\operatorname{Cotor}}
\providecommand{\Coext}{\operatorname{Coext}}
\providecommand{\RHom}{R\!\operatorname{Hom}}
\providecommand{\Tot}{\operatorname{Tot}}
\providecommand{\KK}{\mathbb{K}}
\providecommand{\Dpbw}{D^{\mathrm{pbw}}}
\providecommand{\Dco}{D^{\mathrm{co}}}
\providecommand{\chotimes}{\mathbin{\otimes^{\mathrm{ch}}}}
\providecommand{\wt}{\widetilde}
\providecommand{\eps}{\varepsilon}
\providecommand{\susp}{s}
\providecommand{\coeq}{\operatorname{coeq}}
\providecommand{\id}{\mathrm{id}}
\providecommand{\op}{\mathrm{op}}
\providecommand{\cop}{\mathrm{cop}}
\providecommand{\HH}{\operatorname{HH}}
\providecommand{\Tor}{\operatorname{Tor}}
\providecommand{\Ext}{\operatorname{Ext}}
\providecommand{\gr}{\operatorname{gr}}
\providecommand{\MC}{\operatorname{MC}}
\providecommand{\Ob}{\operatorname{Ob}}
\providecommand{\eq}{\operatorname{eq}}
\providecommand{\Eone}{\mathsf{E}_1}
\providecommand{\Einf}{\mathsf{E}_{\infty}}
\providecommand{\Vlat}{V}
\providecommand{\barBgeom}{\bar{\mathbf{B}}}

\chapter{Ordered Associative Chiral Koszul Duality}
\label{ch:ordered-associative-chiral-kd}

\section{Introduction}

The ordered bar complex $\Barch^{\mathrm{ord}}(A)$ of a chiral algebra~$A$
is the cofree coalgebra $T^c(s^{-1}A)$ with differential extracting
collision residues at consecutive points and coproduct from interval-cutting.
Its linear dual, on the chirally Koszul locus, is the dg-shifted Yangian.

Three bar complexes coexist and must not be conflated.
The Francis--Gaitsgory bar $\Barch^{\mathrm{FG}}(A)$ uses only the zeroth
product $a_{(0)}b$ and sees the chiral Lie structure.
The symmetric bar $\Barch^{\Sigma}(A)$ uses all OPE products and takes
$\Sigma_n$-coinvariants --- this is the bar complex of Theorem~A.
The ordered bar $\Barch^{\mathrm{ord}}(A)$ uses all OPE products but retains
the linear ordering on configurations, with no $\Sigma_n$ quotient.
The natural map $\Barch^{\mathrm{ord}}\to \Barch^{\Sigma}$ takes
coinvariants; the associated graded of $\Barch^{\Sigma}$ recovers
$\Barch^{\mathrm{FG}}$.

Let $C = \Barch(A)$ denote the ordered bar coalgebra.
The coalgebra-side avatar of the diagonal bimodule~${}_A A_A$
is the \emph{diagonal bicomodule} $C_\Delta$, the coalgebra~$C$
equipped with its coproduct on both sides.
Interval composition on the algebra side is derived tensor over~$A$;
on the coalgebra side it is derived cotensor over~$C$.
The annulus is the self-cotensor:
\[
\HH^{\mathrm{ch}}_\bullet(A)\simeq C_\Delta \square_{C^e}^{\mathbf R} C_\Delta.
\]
The associative dual does not merely mirror the algebra; it already
remembers the open boundary formalism.
The ordered bar complex presents the $\Eone$-chiral Koszul dual
as the Steinberg variety presents the Hecke algebra.

\section{The strongly admissible regime and the foundational inputs}
\label{sec:setup}

We work over a field $k$ of characteristic $0$.  Let $X$ be a smooth curve over $k$, and let
$\mathcal D(X)$ denote the relevant derived category of $D$-modules on $X$ (or on the Ran space of
$X$) endowed with the chiral tensor product
\[
\chotimes := \boxtimes_{D_X}.
\]
An $E_1$-chiral algebra means an algebra over the non-$\Sigma$ operad $\Assch$.

\begin{definition}[Strongly admissible associative chiral algebra]
An augmented $E_1$-chiral algebra $A$ on $X$ is called \emph{strongly admissible} if the following
hold.
\begin{enumerate}[label=(\alph*)]
\item $A$ is pro-nilpotent and equipped with an exhaustive separated descending PBW-type filtration
$F^\bullet A$ which is complete for the induced topology.

\item Every associated graded piece and every finite filtration quotient is bounded below and perfect
enough that completed tensor products, totalizations, and normalized chains behave exact\-ly on these
pieces.

\item The reduced ordered bar coalgebra
\[
C:=\Barch(A)
\]
is conilpotent, and the corresponding coderived category of bicomodules admits enough coflat
resolutions so that derived cotensor products are defined and associative.

\item Whenever linear duality is invoked, each graded/conformal piece of the object under consideration
is finite-dimensional.
\end{enumerate}
\end{definition}

\begin{remark}
This is the fortified regime in which hidden homological-algebra issues are kept on the surface.
Nothing deep in the new arguments below requires more than these hypotheses; they are chosen so that
every totalization, completion, normalization, and passage to coderived categories is legitimate.
\end{remark}

\begin{hypothesis}[Repository inputs]
\label{hyp:inputs}
We assume the repository has already established the following.
\begin{enumerate}[label=(I)]
\item \textbf{Associative chiral bar--cobar.}
For every augmented pro-nilpotent strongly admissible $E_1$-chiral algebra $A$ there is a
quasi-isomorphism
\[
\Cobar \Barch(A)\xrightarrow{\sim} A,
\]
and, on the appropriate coalgebra side, the corresponding bar--cobar adjunction is an equivalence.

\item \textbf{One-sided module/comodule duality.}
For every such $A$ there is an equivalence between complete left $A$-modules and conilpotent left
$\Barch(A)$-comodules, with the usual coderived/contraderived refinement in the curved or completed
regime.

\item \textbf{Enveloping algebra formalism.}
The chiral enveloping algebra
\[
A^e := A\chotimes A^{\op}
\]
is defined, and $A$-bimodules are equivalently left $A^e$-modules.
\end{enumerate}
\end{hypothesis}

\begin{definition}[Universal twisting cochain]
Let $A$ be strongly admissible and set $C:=\Barch(A)$.  The bar--cobar equivalence produces a
canonical twisting cochain
\[
\tau:C\to A
\]
characterized by the Maurer--Cartan equation
\[
d\tau+\tau\star\tau=0.
\]
We call $\tau$ the \emph{universal twisting cochain}.
\end{definition}

\begin{definition}[Opposite and diagonal objects]
Let $A$ be an associative chiral algebra and let $C$ be a coaugmented dg coalgebra.
\begin{enumerate}[label=(\alph*)]
\item $A^{\op}$ is the same underlying object with multiplication reversed.
\item $C^{\cop}$ is the same underlying object with coproduct flipped.
\item The \emph{coalgebraic enveloping object} is
\[
C^e:=C\widehat\otimes C^{\cop}.
\]
\item The \emph{diagonal bicomodule} $C_\Delta$ is $C$ equipped with left and right coactions both
equal to $\Delta_C$.
\end{enumerate}
\end{definition}

\begin{lemma}[Bicomodules as comodules over the enveloping coalgebra]
\ClaimStatusProvedHere
\label{lem:bicom-e}
Let $C$ be a coaugmented conilpotent dg coalgebra.  Then left $C^e$-comodules are canonically the
same thing as $C$-bicomodules.
\end{lemma}

\begin{proof}
Given a left $C^e$-coaction
\[
\rho^e:M\to C\otimes C^{\cop}\otimes M,
\]
compose with the counit on the second factor and on the first factor to obtain
\[
\lambda=(\id_C\otimes \eps\otimes \id_M)\rho^e:M\to C\otimes M,
\]
\[
\rho=(\eps\otimes \id_{C^{\cop}}\otimes \id_M)\rho^e:M\to M\otimes C.
\]
These are a left and a right $C$-coaction, and coassociativity of $\rho^e$ is equivalent to the
usual bicomodule compatibility condition.

Conversely, given commuting left and right coactions $\lambda$ and $\rho$, define
\[
\rho^e
=
(\id_C\otimes \tau_{C,M}\otimes \id_M)\circ(\lambda\otimes \id_M)\circ \rho
:
M\to C\otimes C^{\cop}\otimes M,
\]
where $\tau_{C,M}$ swaps the middle two factors.  The bicomodule axioms are precisely the
coassociativity and counitality conditions for $\rho^e$.
\end{proof}

\begin{definition}[Chiral Hochschild and coHochschild theories]
Let $A$ be a strongly admissible associative chiral algebra and $C=\Barch(A)$.
For an $A$-bimodule $M$ and a $C$-bicomodule $N$, define
\[
\HH^{\mathrm{ch}}_\bullet(A,M):=\Tor_\bullet^{A^e}(A,M),
\qquad
\HH_{\mathrm{ch}}^\bullet(A,M):=\RHom_{A^e}(A,M),
\]
\[
\coHoch^{\mathrm{ch}}_\bullet(C,N):=\Cotor_\bullet^{C^e}(C_\Delta,N),
\qquad
\coHoch_{\mathrm{ch}}^\bullet(C,N):=\Coext_{C^e}^\bullet(C_\Delta,N).
\]
When $M=A$ and $N=C_\Delta$ we omit coefficients.
\end{definition}

\section{Ordered bar coalgebras and the chiral shuffle theorem}
\label{sec:shuffle}

In the associative chiral regime the basic bar object is ordered.  This is the first point at which
the theory diverges sharply from the commutative shadow: \emph{every shuffle survives as a distinct
summand.}

\begin{definition}[The simplicial ordered bar object]
Let $A$ be an augmented associative chiral algebra.  Define the simplicial object
$B^\Delta_\bullet(A)$ by
\[
B^\Delta_n(A):=A^{\chotimes n}\qquad (n\geq 0),
\]
with faces given by augmentation at the two ends and multiplication in adjacent slots, and
degeneracies given by unit insertion.  Its normalized chains
\[
N(B^\Delta_\bullet(A))
\]
are canonically identified with the reduced ordered chiral bar complex $\Barch(A)$.
\end{definition}

In suspended notation we write a bar word of length $n$ as
\[
[a_1|\cdots|a_n].
\]

\begin{definition}[EZ-admissible pair]
A pair $(A,B)$ of strongly admissible associative chiral algebras is called \emph{EZ-admissible} if
the normalization functor on bisimplicial objects built from $A$ and $B$ commutes, filtration-wise,
with the completed chiral tensor products and totalizations appearing in the proof of the
Eilenberg--Zilber theorem.
\end{definition}

\begin{theorem}[Ordered chiral shuffle theorem]
\ClaimStatusProvedHere
\label{thm:shuffle}
Let $(A,B)$ be an EZ-admissible pair of strongly admissible associative chiral algebras.
There are natural filtered chain maps
\[
\mathsf{Sh}^{\mathrm{ch}}_{A,B}:
\Barch(A)\widehat\otimes \Barch(B)\longrightarrow \Barch(A\chotimes B),
\]
\[
\mathsf{AW}^{\mathrm{ch}}_{A,B}:
\Barch(A\chotimes B)\longrightarrow \Barch(A)\widehat\otimes \Barch(B),
\]
such that
\[
\mathsf{AW}^{\mathrm{ch}}_{A,B}\circ \mathsf{Sh}^{\mathrm{ch}}_{A,B}=\id,
\qquad
\mathsf{Sh}^{\mathrm{ch}}_{A,B}\circ \mathsf{AW}^{\mathrm{ch}}_{A,B}\simeq \id.
\]
Hence there is a filtered quasi-isomorphism of conilpotent dg coalgebras
\[
\Barch(A)\widehat\otimes \Barch(B)\xrightarrow{\sim}\Barch(A\chotimes B).
\]
\end{theorem}

\begin{proof}
Set
\[
K_{\bullet,\bullet}:=B^\Delta_\bullet(A)\chotimes B^\Delta_\bullet(B).
\]
This is a bisimplicial object in the ambient abelian or stable category of $D$-modules.

Its diagonal identifies canonically with the simplicial bar object of the tensor-product algebra:
\[
\operatorname{diag}(K)_n
=
A^{\chotimes n}\chotimes B^{\chotimes n}
\cong
(A\chotimes B)^{\chotimes n}
=
B^\Delta_n(A\chotimes B).
\]
This identification is compatible with faces and degeneracies because multiplication and augmentation
on $A\chotimes B$ are componentwise.

Now apply the classical Eilenberg--Zilber theorem filtration-wise in each exact graded piece.
EZ-admissibility ensures that normalization commutes with the completed tensor constructions used
here, so one obtains a filtered quasi-isomorphism
\[
N(\operatorname{diag}K)\xrightarrow{\sim} \Tot\bigl(NK_{\bullet,\bullet}\bigr),
\]
with inverse given by Alexander--Whitney.  Since
\[
N(B^\Delta_\bullet(A))\cong \Barch(A),
\qquad
N(B^\Delta_\bullet(B))\cong \Barch(B),
\qquad
N(B^\Delta_\bullet(A\chotimes B))\cong \Barch(A\chotimes B),
\]
this yields the claimed quasi-isomorphism.

The formulas are the usual normalized Eilenberg--MacLane formulas, now interpreted in the ordered
chiral setting.  Explicitly, for bar words
\[
[a_1|\cdots|a_p]\in \Barch_p(A),\qquad [b_1|\cdots|b_q]\in \Barch_q(B),
\]
set $c_i=a_i\otimes 1$ for $1\leq i\leq p$ and $c_{p+j}=1\otimes b_j$ for $1\leq j\leq q$, and define
\[
\mathsf{Sh}^{\mathrm{ch}}_{p,q}
([a_1|\cdots|a_p]\otimes[b_1|\cdots|b_q])
=
\sum_{\sigma\in \operatorname{Sh}(p,q)}
\varepsilon(\sigma)\,
[c_{\sigma^{-1}(1)}|\cdots|c_{\sigma^{-1}(p+q)}].
\]
Similarly,
\[
\mathsf{AW}^{\mathrm{ch}}([x_1|\cdots|x_n])
=
\sum_{p=0}^n
[\pi_A(x_1)|\cdots|\pi_A(x_p)]\otimes[\pi_B(x_{p+1})|\cdots|\pi_B(x_n)],
\]
where
\[
\pi_A=\id_A\otimes \eps_B,\qquad \pi_B=\eps_A\otimes \id_B.
\]
The chain-map property uses only the simplicial identities.

Finally, the shuffle map is compatible with deconcatenation: cutting a shuffled word is equivalent to
cutting the source words and then shuffling the left halves and the right halves.  Therefore
$\mathsf{Sh}^{\mathrm{ch}}_{A,B}$ and $\mathsf{AW}^{\mathrm{ch}}_{A,B}$ are coalgebra maps.
\end{proof}

\begin{remark}
The theorem is genuinely associative.  In a commutative or $E_\infty$ setting many of the shuffle
summands are identified after the symmetric-group quotient.  Here they remain distinct.  The ordered
configuration space remembers their full combinatorics.
\end{remark}

\section{The $R$-matrix as ordered-to-unordered descent}
\label{sec:r-matrix-descent}

The ordered bar complex $\Barch^{\mathrm{ord}}(A)$ carries strictly more information
than the unordered bar complex $\Barch(A)$ whenever the OPE has poles.  The surplus
is the \emph{descent datum} for the $\Sigma_n$-covering of configuration spaces:
the $R$-matrix.  The $R$-matrix arises from the flat connection
on ordered configuration space, and the unordered bar complex is recovered
as the $R$-twisted $\Sigma_n$-descent.

\subsection*{The covering-space frame}

\begin{construction}[The covering-space frame]
\ClaimStatusProvedHere
\label{constr:r-matrix-covering-vol1}
\index{R-matrix!descent datum|textbf}
\index{descent!ordered to unordered bar complex}
The projection
\[
\pi_n\colon
\mathrm{Conf}_n^{\mathrm{ord}}(X)
\;\longrightarrow\;
\mathrm{Conf}_n(X)
\;=\;
\mathrm{Conf}_n^{\mathrm{ord}}(X)/\Sigma_n
\]
is a principal $\Sigma_n$-bundle, \'etale with fibre~$\Sigma_n$.
A local system~$\mathcal{L}$ on the base is equivalent to a
$\Sigma_n$-equivariant local system~$\pi_n^*\mathcal{L}$
on the total space: the descent correspondence
\[
\mathrm{LocSys}\bigl(\mathrm{Conf}_n(X)\bigr)
\;\simeq\;
\mathrm{LocSys}^{\Sigma_n}
\bigl(\mathrm{Conf}_n^{\mathrm{ord}}(X)\bigr).
\]
The bar complex at tensor degree~$n$ defines a local system
$\mathcal{F}_n^{\mathrm{ord}}$ on $\mathrm{Conf}_n^{\mathrm{ord}}(X)$
whose stalk at $(z_1,\dots,z_n)$ is
$(s^{-1}\bar A)^{\otimes n}$, with flat connection induced by the
residue differential of the ordered bar construction.  To descend
$\mathcal{F}_n^{\mathrm{ord}}$ to a local system on $\mathrm{Conf}_n(X)$,
one must exhibit a $\Sigma_n$-equivariant structure: isomorphisms
\[
R_\sigma\colon
\mathcal{F}_n^{\mathrm{ord}}\big|_{(z_1,\dots,z_n)}
\;\xrightarrow{\;\sim\;}
\mathcal{F}_n^{\mathrm{ord}}\big|_{(z_{\sigma(1)},\dots,z_{\sigma(n)})}
\]
for each $\sigma\in\Sigma_n$, satisfying the cocycle condition
$R_{\sigma\tau}=R_\tau\circ R_\sigma$ and compatible with the
flat connection.

For a pole-free $E_\infty$-chiral algebra, the equivariant structure
is canonical: $R_\sigma=\varepsilon(\sigma)\cdot\tau_\sigma$, where
$\tau_\sigma$ permutes tensor factors and $\varepsilon(\sigma)$ is
the Koszul sign.  The descent recovers the ordinary
$\Sigma_n$-coinvariant bar complex.

For an $E_\infty$-chiral algebra with OPE poles (every interesting vertex
algebra: $\hat\mathfrak{g}_k$, $\mathrm{Vir}_c$, $\mathcal{H}_k$, $W$-algebras),
the equivariant structure must be \emph{constructed} from the OPE data.
The construction is the $R$-matrix.
\end{construction}

\subsection*{The spectral Kohno connection}

\begin{construction}[$R$-matrix as monodromy of the Kohno connection]
\ClaimStatusProvedHere
\label{constr:r-matrix-monodromy-vol1}
\index{R-matrix!from monodromy|textbf}
\index{monodromy!R-matrix}
\index{Kohno connection}
The bar differential on $\Barch^{\mathrm{ord}}(A)$ extracts, at each
collision $z_i=z_j$, the collision residue
$r_{ij}(z_i-z_j)\in\operatorname{End}(\bar A\otimes\bar A)\otimes
\mathcal{O}(*\Delta)$.  By the bar kernel convention (the $d\log$
kernel absorbs one power), the collision residue $r(z)$ has pole
orders one less than the OPE.

These residues define a flat connection on the trivial bundle
over $\mathrm{Conf}_n^{\mathrm{ord}}(\mathbb{C})$ with fibre
$\bar A^{\otimes n}$:
\begin{equation}\label{eq:kohno-connection}
\nabla
\;=\;
d
\;-\;
\sum_{1\leq i<j\leq n}
r_{ij}(z_i-z_j)\,d\log(z_i-z_j).
\end{equation}
Here $r_{ij}$ acts in the $(i,j)$ tensor slots of~$\bar A^{\otimes n}$
and $d\log(z_i-z_j)=dz_{ij}/(z_i-z_j)$ is the logarithmic form
on $\mathrm{Conf}_n^{\mathrm{ord}}(\mathbb{C})$.

\emph{Flatness.}
The curvature
\[
\nabla^2
\;=\;
\sum_{i<j<k}
\bigl(
[r_{ij},r_{ik}]+[r_{ij},r_{jk}]+[r_{ik},r_{jk}]
\bigr)
\,d\log(z_{ij})\wedge d\log(z_{jk})
\]
vanishes if and only if each triple satisfies the classical
Yang--Baxter equation
\begin{equation}\label{eq:cybe-vol1}
[r_{12}(z_{12}),r_{13}(z_{13})]
+[r_{12}(z_{12}),r_{23}(z_{23})]
+[r_{13}(z_{13}),r_{23}(z_{23})]
\;=\;0,
\end{equation}
where $z_{ij}=z_i-z_j$.
The identity $d^2=0$ on the ordered bar complex
$\Barch^{\mathrm{ord}}(A)$ entails precisely this CYBE,
as the codimension-two strata of the ordered FM
compactification $\overline{\mathrm{FM}}_n^{\mathrm{ord}}(\mathbb{C})$
impose the three-term relation on triple collisions.
Hence $\nabla$ is flat.

\emph{The $R$-matrix.}
At $n=2$, the ordered configuration space
$\mathrm{Conf}_2^{\mathrm{ord}}(\mathbb{C})
=\{(z_1,z_2)\in\mathbb{C}^2\mid z_1\neq z_2\}$
has fundamental group $\pi_1=\mathbb{Z}$, generated by the
loop $\gamma$ in which $z_1$ circles $z_2$.
The \emph{$R$-matrix} is the monodromy of $\nabla$ around~$\gamma$:
\begin{equation}\label{eq:r-matrix-monodromy-vol1}
R(z)
\;:=\;
\operatorname{Mon}_\gamma(\nabla)
\;=\;
\mathcal{P}\exp\!\Bigl(
\oint_\gamma r(z)\,d\log z
\Bigr)
\;\in\;
\operatorname{End}(\bar A\otimes\bar A)[\![z^{-1}]\!],
\end{equation}
where $z=z_1-z_2$ is the relative coordinate and $\mathcal{P}$ denotes
path-ordering.

At general~$n$, the pure braid group
$P_n=\pi_1(\mathrm{Conf}_n^{\mathrm{ord}}(\mathbb{C}))$
has standard generators $\gamma_{ij}$ ($i<j$, the loop in which $z_i$
winds around $z_j$ with all other points fixed).  The flat connection
$\nabla$ induces a monodromy representation
\[
\rho\colon P_n\to\operatorname{GL}(\bar A^{\otimes n}),
\qquad
\gamma_{ij}\mapsto R_{ij}(z_{ij}).
\]
The braid relation $\gamma_{ij}\gamma_{ik}\gamma_{jk}
=\gamma_{jk}\gamma_{ik}\gamma_{ij}$ in $P_n$ maps to the
Yang--Baxter equation~\eqref{eq:cybe-vol1} (its exponentiated form),
so $\rho$ is a well-defined representation.
\end{construction}

\subsection*{The descent theorem}

\begin{proposition}[$R$-matrix twisted descent]
\ClaimStatusProvedHere
\label{prop:r-matrix-descent-vol1}
\index{descent!R-matrix!ordered to unordered|textbf}
\index{R-matrix!twisted descent}
Let $A$ be a strongly admissible associative chiral algebra with
ordered bar coalgebra $C=\Barch^{\mathrm{ord}}(A)$ and
$R$-matrix $R(z)$ as in
Construction~\textup{\ref{constr:r-matrix-monodromy-vol1}}.
Then the unordered bar complex is the $R$-twisted $\Sigma_n$-descent:
\begin{equation}\label{eq:descent-identification-vol1}
\Barch(A)_n
\;\simeq\;
\bigl(\Barch^{\mathrm{ord}}(A)_n\bigr)^{R\text{-}\Sigma_n},
\end{equation}
where $\Sigma_n$ acts on adjacent transpositions by
$\sigma_i\cdot(a_1\otimes\cdots\otimes a_n)
=(\tau_{i,i+1}\circ R_{i,i+1}(z_i{-}z_{i+1}))\,
(a_1\otimes\cdots\otimes a_n)$.
The $\Sigma_n$-action is well-defined if and only if the
$R$-matrix satisfies the Yang--Baxter equation.
\end{proposition}

\begin{proof}
The proof proceeds in three steps: we establish the
$\Sigma_n$-action, verify its well-definedness, and identify
the fixed locus.

\emph{Step~1: the twisted $\Sigma_n$-action.}
The covering $\pi_n\colon\mathrm{Conf}_n^{\mathrm{ord}}(\mathbb{C})
\to\mathrm{Conf}_n(\mathbb{C})$ presents the base as the orbit space
of the free $\Sigma_n$-action on ordered configurations.  An element
$\sigma_i=(i,\,i{+}1)\in\Sigma_n$ lifts to a path in
$\mathrm{Conf}_n^{\mathrm{ord}}(\mathbb{C})$ that continuously
exchanges $z_i$ and $z_{i+1}$.  The monodromy of the flat connection
$\nabla$ (eq.~\eqref{eq:kohno-connection}) along this path decomposes
into two factors:
\begin{enumerate}[label=(\roman*)]
\item the permutation $\tau_{i,i+1}$ of tensor slots in
  $\bar A^{\otimes n}$, from the coordinate transposition, and
\item the parallel transport $R_{i,i+1}(z_i-z_{i+1})$ of the fibre,
  from the flat connection on ordered configuration space.
\end{enumerate}
Define the \emph{$R$-twisted action} on generators by
\begin{equation}\label{eq:r-twisted-action}
\rho_R(\sigma_i)
\;:=\;
\tau_{i,i+1}\circ R_{i,i+1}(z_i-z_{i+1}).
\end{equation}

\emph{Step~2: well-definedness via YBE.}
To extend $\rho_R$ from the generators
$\sigma_1,\dots,\sigma_{n-1}$ to all of $\Sigma_n$, we must verify
the Coxeter relations.  The involutivity
$\rho_R(\sigma_i)^2=\id$ holds because the double monodromy
(a full loop around the branch point) is trivial in the quotient
$\Sigma_n$ (the covering sheet returns to itself).
The distant commutativity $\rho_R(\sigma_i)\rho_R(\sigma_j)
=\rho_R(\sigma_j)\rho_R(\sigma_i)$ for $|i-j|\geq 2$ is immediate
because $R_{i,i+1}$ and $R_{j,j+1}$ act in disjoint tensor slots.

The braid relation
$\rho_R(\sigma_i)\rho_R(\sigma_{i+1})\rho_R(\sigma_i)
=\rho_R(\sigma_{i+1})\rho_R(\sigma_i)\rho_R(\sigma_{i+1})$
is the critical constraint.  Expanding:
\begin{align*}
&(\tau_{i,i+1}\circ R_{i,i+1})
(\tau_{i+1,i+2}\circ R_{i+1,i+2})
(\tau_{i,i+1}\circ R_{i,i+1}) \\
={}&
(\tau_{i+1,i+2}\circ R_{i+1,i+2})
(\tau_{i,i+1}\circ R_{i,i+1})
(\tau_{i+1,i+2}\circ R_{i+1,i+2}).
\end{align*}
Conjugating the tensor permutations to the left and using the
standard identities
$\tau_{i,i+1}R_{j,j+1}\tau_{i,i+1}^{-1}=R_{\sigma_i(j),\sigma_i(j+1)}$,
this reduces to
\[
R_{12}(z_{12})\,R_{13}(z_{13})\,R_{23}(z_{23})
\;=\;
R_{23}(z_{23})\,R_{13}(z_{13})\,R_{12}(z_{12}),
\]
the Yang--Baxter equation.  This holds because the flatness
$\nabla^2=0$ implies the CYBE
(eq.~\eqref{eq:cybe-vol1}), and the path-ordered exponential
converts the infinitesimal CYBE into the exponentiated YBE
(this is the content of the Kohno--Drinfeld theorem: the
monodromy representation of a flat Knizhnik--Zamolodchikov connection
satisfies the Yang--Baxter equation).

\emph{Step~3: identification of the descent.}
The result is a genuine $\Sigma_n$-action $\rho_R$ on
$\Barch^{\mathrm{ord}}(A)_n$.  By the descent equivalence for
principal bundles, a $\Sigma_n$-equivariant local system on
the total space $\mathrm{Conf}_n^{\mathrm{ord}}(\mathbb{C})$
descends to a local system on
$\mathrm{Conf}_n(\mathbb{C})$.  The descended local system
is the equaliser of the $\rho_R$-action, which at the level of
global sections is the $R$-twisted $\Sigma_n$-invariant subspace:
\[
\bigl(\Barch^{\mathrm{ord}}(A)_n\bigr)^{R\text{-}\Sigma_n}
\;:=\;
\bigl\{
x\in\Barch^{\mathrm{ord}}(A)_n
\;\big|\;
\rho_R(\sigma)\cdot x=x
\;\text{for all }\sigma\in\Sigma_n
\bigr\}.
\]
We claim this equals $\Barch(A)_n$.  By construction, the unordered bar
complex $\Barch(A)_n$ is the complex of sections of the descended
local system $\mathcal{F}_n$ on $\mathrm{Conf}_n(\mathbb{C})$
whose pullback to $\mathrm{Conf}_n^{\mathrm{ord}}(\mathbb{C})$ is
$\mathcal{F}_n^{\mathrm{ord}}$ with the $\rho_R$-equivariant structure.
The descent-ascent adjunction for \'etale $\Sigma_n$-coverings gives
\[
\Gamma\bigl(\mathrm{Conf}_n(\mathbb{C}),\,\mathcal{F}_n\bigr)
\;=\;
\Gamma\bigl(
\mathrm{Conf}_n^{\mathrm{ord}}(\mathbb{C}),\,
\mathcal{F}_n^{\mathrm{ord}}\bigr)^{R\text{-}\Sigma_n}.
\]
The left side is $\Barch(A)_n$ and the right side is
$(\Barch^{\mathrm{ord}}(A)_n)^{R\text{-}\Sigma_n}$,
completing the proof.
\end{proof}

\begin{corollary}[Pole-free descent is naive]
\ClaimStatusProvedHere
\label{cor:pole-free-descent}
\index{descent!pole-free algebras}
For a pole-free $E_\infty$-chiral algebra \textup{(}a commutative chiral
algebra whose chiral product extends across the diagonal without
singularities\textup{)}, the collision residue vanishes:
$r(z)=0$.  Hence $\nabla=d$ is the trivial flat connection,
$R(z)=\id$, and~\eqref{eq:descent-identification-vol1} reduces to
the ordinary $\Sigma_n$-coinvariant:
\[
\Barch(A)_n
\;\simeq\;
\bigl(\Barch^{\mathrm{ord}}(A)_n\bigr)\big/\Sigma_n.
\]
For $E_\infty$-chiral algebras with OPE poles
\textup{(}all vertex algebras with nontrivial singular OPE:
$\hat\mathfrak{g}_k$, $\mathrm{Vir}_c$, $\mathcal{H}_k$,
lattice algebras, $W$-algebras --- these are all $E_\infty$-chiral,
since every vertex algebra is $E_\infty$\textup{)}, the $R$-matrix
$R(z)\neq\id$ carries nontrivial spectral dependence derived from
the local OPE, and the descent is genuinely $R$-twisted.
\end{corollary}

\begin{proof}
If the OPE has no poles, every collision residue vanishes identically,
so the connection form in~\eqref{eq:kohno-connection} is zero and the
monodromy~\eqref{eq:r-matrix-monodromy-vol1} is the identity.
The $R$-twisted $\Sigma_n$-action~\eqref{eq:r-twisted-action}
reduces to the plain permutation $\tau_{i,i+1}$, and the
$R$-twisted invariants coincide with the ordinary
$\Sigma_n$-coinvariants.
\end{proof}

\begin{remark}[The $R$-matrix as Maurer--Cartan element]
\label{rem:r-matrix-mc-vol1}
\index{R-matrix!Maurer--Cartan interpretation}
The classical $r$-matrix
$r(z)\in\operatorname{End}(\bar A\otimes\bar A)
\otimes\mathcal{O}(*\Delta)$
is a Maurer--Cartan element in the arity-$2$ component of the
$E_1$ modular convolution algebra
(cf.\ the $E_1$ shadow Postnikov tower in the body of the text).
The CYBE~\eqref{eq:cybe-vol1} is the linearised
Maurer--Cartan equation; the YBE on $R(z)$ is its
exponentiated form.  The $R$-matrix is the descent datum because the
convolution dg~Lie algebra governing $\Sigma_n$-equivariant structures
on local systems over $\mathrm{Conf}_n$ is precisely the
Gerstenhaber--Schack complex of bilinear operations, and
Maurer--Cartan elements in it parametrise the nontrivial descents.
Thus the entire apparatus of twisted descent is controlled by a single
MC element in the arity-$2$ sector.
\end{remark}

\begin{remark}[Motivic interpretation of ordered-to-symmetric descent]
\label{rem:motivic-interpretation}
\index{motivic interpretation!ordered bar complex}
\index{descent!motivic content}
\index{mixed Tate motives!configuration space periods}
The averaging map
$\mathrm{av}\colon\Barch^{\mathrm{ord}}(A)\to\Barch(A)$
is a motivic projection: it kills precisely the transcendental
arithmetic content of the configuration-space integrals.

The ordered configuration space
$\mathrm{Conf}_k^{\mathrm{ord}}(\mathbb{C})$
carries a canonical mixed Hodge structure whose periods are
the iterated integrals of the Knizhnik--Zamolodchikov connection.
The $R$-matrix monodromy of
Proposition~\ref{prop:r-matrix-descent-vol1} is computed by these
periods, and the resulting transcendental data is stratified by
arity:
\begin{itemize}
\item At arity~$2$: the monodromy $R(z)$ of the flat connection
$\nabla = d - r(z)\,d\log z$ involves $2\pi i$, a period of
the Lefschetz motive~$\mathbb{L}$ (weight~$1$).

\item At arity~$3$: the Drinfeld associator
$\Phi_{\mathrm{KZ}}$
--- the monodromy of the KZ connection on
$\mathrm{Conf}_3^{\mathrm{ord}}(\mathbb{C})$ ---
involves multiple zeta values
$\zeta(2),\,\zeta(3),\,\zeta(2,1),\ldots$\
(periods of mixed Tate motives over~$\mathbb{Z}$).

\item At arity~$k$: the iterated monodromy involves
periods of mixed Tate motives of weight~$\leq k-1$,
by the theorem of Brown~\cite{brown-mzv}.
\end{itemize}

\noindent
The symmetric bar complex $\Barch(A)$ retains none of this.
Its invariants --- the curvature~$\kappa$, the complementarity
class~$C$, the shadow tower~$Q$ --- are rational numbers, periods
of the trivial motive~$\mathbb{Q}(0)$.  The kernel
$\ker(\mathrm{av})$ is the motivic surplus: it carries the full
weight-graded mixed Tate content that the symmetric projection
annihilates.

This is the arithmetic shadow of a structural phenomenon.
Drinfeld's Grothendieck--Teichm\"uller group~$\widehat{GT}$
acts on the tower of KZ associators --- equivalently, on the
pro-nilpotent completion of the ordered bar complex at each
arity.  The Deligne--Ihara conjecture identifies
$\mathrm{Lie}(\widehat{GT})$ with the free Lie algebra on
generators in odd degrees~$\geq 3$, one for each odd
zeta value~$\zeta(2n+1)$.  The ordered bar complex of a
chiral algebra is a \emph{representation} of this arithmetic
symmetry group: the $\widehat{GT}$-action permutes the
associators that parametrise equivalent ordered-to-symmetric
descents.  Different choices of Drinfeld associator give
different $R$-twisted $\Sigma_n$-actions with isomorphic
but not identical descent data.

The ordered bar complex retains the transcendental arithmetic of
the configuration-space integrals.  The symmetric bar complex
projects to the rational shadow.  Every motivic weight is
visible in the ordered complex; every motivic weight is killed
by the averaging map.
\end{remark}

\section{Opposite-duality, anti-involutions, and enveloping duality}
\label{sec:opposite}

The second foundational phenomenon is order reversal.  It transforms opposite multiplication on the
algebra side into opposite deconcatenation on the bar side.

\begin{theorem}[Opposite-duality for ordered bar coalgebras]
\ClaimStatusProvedHere
\label{thm:opposite}
Let $A$ be a strongly admissible associative chiral algebra.  There is a canonical isomorphism of
coaugmented conilpotent dg coalgebras
\[
\mathsf{R}_A:\Barch(A^{\op})\xrightarrow{\sim}\Barch(A)^{\cop}.
\]
If $A$ is finite-type Koszul so that linear duality returns an associative chiral algebra, then
\[
(A^{\op})^!\simeq (A^!)^{\op}.
\]
\end{theorem}

\begin{proof}
On suspended bar words define
\[
\mathsf{R}_A[\susp a_1|\cdots|\susp a_n]
=
(-1)^{\sum_{i<j}(|a_i|+1)(|a_j|+1)}
[\susp a_n|\cdots|\susp a_1].
\]
Geometrically, this is induced by the order-reversal involution on the ordered configuration space.

The bar differential is the alternating sum of adjacent collision maps.  The $i$-th face for
$A^{\op}$ multiplies the $i$-th and $(i+1)$-st slots in the opposite order.  Reversal carries that
face to the $(n-i)$-th face for $A$ with exactly the sign written above.  Hence $\mathsf{R}_A$ is a
chain map.

The coproduct on $\Barch(A)$ is deconcatenation.  Reversal sends a cut after the $p$-th letter to a
cut before the $(n-p)$-th letter.  Therefore
\[
\Delta\circ\mathsf{R}_A
=
\tau\circ(\mathsf{R}_A\otimes\mathsf{R}_A)\circ\Delta,
\]
where $\tau$ flips the tensor factors, which is exactly the statement that the target is the
opposite coalgebra.  The claim about linear duals follows because linear duality exchanges opposite
coproduct with opposite product.
\end{proof}

\begin{corollary}[Anti-involutions survive duality]
\ClaimStatusProvedHere
\label{cor:anti}
Any anti-involution $\iota:A\to A^{\op}$ induces an anti-involution
\[
\iota^!:A^!\to (A^!)^{\op}.
\]
If $\iota^2=\id$, then $(\iota^!)^2=\id$.
\end{corollary}

\begin{proof}
Apply Theorem~\ref{thm:opposite} functorially.
\end{proof}

\begin{lemma}[Closure of admissibility under opposite and enveloping constructions]
\ClaimStatusProvedHere
\label{lem:closure}
If $A$ is strongly admissible, then so are $A^{\op}$ and $A^e=A\chotimes A^{\op}$.
\end{lemma}

\begin{proof}
All hypotheses in the definition of strong admissibility are stable under replacing multiplication by
its opposite and under taking completed tensor products of filtered objects.
\end{proof}

\begin{corollary}[Enveloping duality]
\ClaimStatusProvedHere
\label{cor:enveloping}
Let $A$ be strongly admissible, and assume $(A,A^{\op})$ is EZ-admissible.  Set $C=\Barch(A)$.
Then there is a canonical quasi-isomorphism of conilpotent dg coalgebras
\[
\Barch(A^e)\xrightarrow{\sim} C^e=C\widehat\otimes C^{\cop}.
\]
If $A$ is finite-type Koszul, then
\[
(A^e)^!\simeq (A^!)^e.
\]
\end{corollary}

\begin{proof}
By the ordered shuffle theorem,
\[
\Barch(A^e)=\Barch(A\chotimes A^{\op})\xrightarrow{\sim}\Barch(A)\widehat\otimes \Barch(A^{\op}).
\]
Now apply Theorem~\ref{thm:opposite} to the second factor:
\[
\Barch(A^{\op})\xrightarrow{\sim}\Barch(A)^{\cop}.
\]
This proves the coalgebra statement.  The algebra statement follows after finite-type linear duality.
\end{proof}

\section{The two-sided twisted bicomodule and the tangent bar construction}
\label{sec:tangent}

The correct coalgebra-side avatar of an $A$-bimodule admits two a priori
different constructions:
the first is obtained by differentiating the ordered bar construction at a square-zero extension;
the second is obtained by twisting a two-sided tensor product with the universal twisting cochain.
These two constructions are the same.

\begin{definition}[Two-sided twisted bicomodule]
\label{def:Kbi}
Let $A$ be strongly admissible, let $C=\Barch(A)$, and let $M$ be an $A$-bimodule.
Define the \emph{two-sided twisted bicomodule}
\[
\KK_A^{\mathrm{bi}}(M):=C\widehat\otimes M\widehat\otimes C
\]
with differential
\[
d_{\KK}
=
d_C\otimes 1\otimes 1 + 1\otimes d_M\otimes 1 + 1\otimes 1\otimes d_C + d_L+d_R,
\]
where
\[
d_L(c\otimes m\otimes d)
=
\sum (-1)^{|c_{(1)}|}\, c_{(1)}\otimes \tau(c_{(2)})\cdot m\otimes d,
\]
\[
d_R(c\otimes m\otimes d)
=
\sum (-1)^{|c|+|m|}\, c\otimes m\cdot \tau(d_{(1)})\otimes d_{(2)},
\]
with the usual Koszul signs determined by the grading conventions.
The left and right $C$-coactions are induced by comultiplication on the two outer factors.
\end{definition}

\begin{lemma}
\ClaimStatusProvedHere
\label{lem:Kbi-dg}
$\KK_A^{\mathrm{bi}}(M)$ is a dg $C$-bicomodule.
\end{lemma}

\begin{proof}
The untwisted part of the differential squares to zero.  The quadratic terms in $d_L$ and in $d_R$
vanish by the Maurer--Cartan equation
\[
d\tau+\tau\star\tau=0.
\]
The mixed commutator $d_Ld_R+d_Rd_L$ vanishes because the left and right $A$-actions on the
bimodule $M$ commute.  Compatibility with the bicomodule structure is immediate from coassociativity
of the coproduct on the outer $C$-factors.
\end{proof}

\begin{definition}[Square-zero extension and one-defect bar complex]
Let $M$ be an $A$-bimodule and form the square-zero extension
\[
A_\eps:=A\oplus \eps M,\qquad \eps^2=0.
\]
Inside the reduced ordered bar complex $\Barch(A_\eps)$, let $T_A^{\mathrm{bar}}(M)$ be the
subcomplex spanned by bar words containing exactly one $\eps M$-slot.
We call this the \emph{one-defect ordered bar complex}.
\end{definition}

\begin{proposition}
\ClaimStatusProvedHere
\label{prop:one-defect}
Modulo $\eps^2$, there is a canonical decomposition
\[
\Barch(A_\eps)\cong \Barch(A)\oplus \eps\, T_A^{\mathrm{bar}}(M).
\]
Moreover, $T_A^{\mathrm{bar}}(M)$ is naturally a $C$-bicomodule.
\end{proposition}

\begin{proof}
A reduced bar word in $\Barch(A_\eps)$ contains some number of $\eps M$-slots.
Since $\eps^2=0$, any word with at least two such slots vanishes modulo $\eps^2$.
Therefore only the zero-defect and one-defect pieces survive.

The bar differential merges adjacent slots.  Multiplying two ordinary $A$-slots remains in $A$.
Multiplying an $A$-slot with the unique $M$-slot remains in $M$ by the bimodule structure.
There is no nonzero product of two defect slots.  Hence the one-defect sector is a subcomplex.
Deconcatenation cuts a one-defect word into a defect-free part and a one-defect part, or vice versa,
so the one-defect sector is naturally a bicomodule over the ordinary bar coalgebra $C=\Barch(A)$.
\end{proof}

\begin{theorem}[Tangent identification]
\ClaimStatusProvedHere
\label{thm:tangent=K}
There is a canonical isomorphism of dg $C$-bicomodules
\[
T_A^{\mathrm{bar}}(M)\xrightarrow{\sim}\KK_A^{\mathrm{bi}}(M).
\]
Equivalently, the one-defect ordered bar complex of a square-zero extension is exactly the two-sided
twisted bicomodule of the defect bimodule.
\end{theorem}

\begin{proof}
As a graded object,
\[
C=\bigoplus_{p\geq 0} (\susp\bar A)^{\otimes p},
\]
hence
\[
C\widehat\otimes M\widehat\otimes C
\cong
\bigoplus_{p,q\geq 0}
(\susp\bar A)^{\otimes p}\otimes M\otimes(\susp\bar A)^{\otimes q}.
\]
This is in tautological bijection with one-defect bar words
\[
[a_1|\cdots|a_p|m|a_{p+1}|\cdots|a_{p+q}]
\]
in which the unique defect sits between a left bar segment of length $p$ and a right bar segment of
length $q$.

Under this identification, the ordinary bar differential on the left segment and the right segment
matches the internal bar differential coming from the two outer copies of $C$, while the two bar
faces adjacent to the defect slot match exactly the two twisting terms $d_L$ and $d_R$.
Therefore the differentials agree.  Deconcatenation also agrees on the nose, because cutting a
one-defect word is the same as cutting the left and right bar segments around the unique defect.
\end{proof}

\begin{corollary}[Infinitesimal dual coalgebra]
\ClaimStatusProvedHere
\label{cor:infdual}
Modulo $\eps^2$, the bar coalgebra of the square-zero extension is
\[
\Barch(A_\eps)\cong C\oplus \eps\,\KK_A^{\mathrm{bi}}(M).
\]
If $A$ is finite-type Koszul, then after linear duality
\[
(A_\eps)^!\cong A^!\oplus \eps\, M^!
\qquad (\mathrm{mod}\ \eps^2),
\]
where
\[
M^!:=\Omega_C^{\mathrm{bi}}\bigl(\KK_A^{\mathrm{bi}}(M)\bigr).
\]
\end{corollary}

\begin{proof}
Combine Proposition~\ref{prop:one-defect} and Theorem~\ref{thm:tangent=K}.  The statement after
dualization is the corresponding first-order cobar reconstruction.
\end{proof}

\begin{proposition}[Infinitesimal annular variation]
\ClaimStatusProvedHere
\label{prop:infann}
Modulo $\eps^2$ one has
\[
\HH^{\mathrm{ch}}_\bullet(A_\eps)
\cong
\HH^{\mathrm{ch}}_\bullet(A)\oplus \eps\, \HH^{\mathrm{ch}}_\bullet(A,M).
\]
On the coalgebra side this becomes
\[
\coHoch^{\mathrm{ch}}_\bullet(\Barch(A_\eps))
\cong
\coHoch^{\mathrm{ch}}_\bullet(C)\oplus \eps\,
\coHoch^{\mathrm{ch}}_\bullet(C,\KK_A^{\mathrm{bi}}(M)).
\]
\end{proposition}

\begin{proof}
The Hochschild complex of the square-zero extension decomposes, modulo $\eps^2$, into the sector with
no defect and the sector with exactly one defect; the second computes Hochschild homology with
coefficients in $M$.  The coalgebra statement follows from
Corollary~\ref{cor:infdual} and the coHochschild identification proved later in
Theorem~\ref{thm:HH-coHH-homology}.
\end{proof}

\section{PBW-complete two-sided Koszul duality}
\label{sec:bimod-bicomod}

One-sided module/comodule duality extends to the
two-sided theory by passing from $A$ to its enveloping algebra $A^e$ and applying
Corollary~\ref{cor:enveloping}.

\begin{definition}[PBW-complete and coderived bicomodule categories]
Let $A$ be strongly admissible and set $C=\Barch(A)$.
\begin{enumerate}[label=(\alph*)]
\item $\Dpbw_{\mathrm{bi}}(A)$ denotes the derived category of PBW-complete $A$-bimodules obtained
from complete semifree $A^e$-resolutions.

\item $\Dco_{\mathrm{bi}}(C)$ denotes the coderived category of conilpotent $C$-bicomodules.
Equivalently, by Lemma~\ref{lem:bicom-e}, it is the coderived category of left $C^e$-comodules.
\end{enumerate}
\end{definition}

\begin{definition}[Two-sided twisted cobar]
Let $N$ be a $C$-bicomodule.  Define
\[
\Omega_C^{\mathrm{bi}}(N):=A\widehat\otimes N\widehat\otimes A
\]
with the differential dual to the one in Definition~\ref{def:Kbi}.
\end{definition}

\begin{theorem}[PBW-complete bimodule/bicomodule equivalence]
\ClaimStatusProvedHere
\label{thm:bimod-bicomod}
Let $A$ be strongly admissible, and assume $(A,A^{\op})$ is EZ-admissible.
Set $C=\Barch(A)$.
Then there is an exact equivalence
\[
\KK_A^{\mathrm{bi}}:\Dpbw_{\mathrm{bi}}(A)\xrightarrow{\sim}\Dco_{\mathrm{bi}}(C)
\]
with inverse $\Omega_C^{\mathrm{bi}}$.  In the curved regime the same statement holds after replacing
the ordinary derived category on the algebra side by the corresponding completed
contra\-derived/coderived version.
\end{theorem}

\begin{proof}
By the enveloping algebra formalism, $A$-bimodules are left $A^e$-modules.
By Corollary~\ref{cor:enveloping},
\[
\Barch(A^e)\xrightarrow{\sim} C^e.
\]
Apply the one-sided module/comodule Koszul duality from Hypothesis~\ref{hyp:inputs}(II) to the
associative chiral algebra $A^e$.  This gives an equivalence between complete $A^e$-modules and
conilpotent $C^e$-comodules.  Since left $C^e$-comodules are the same as $C$-bicomodules by
Lemma~\ref{lem:bicom-e}, we obtain the asserted equivalence.

The functor is precisely the two-sided twisted bar construction
\[
M\longmapsto C\widehat\otimes M\widehat\otimes C,
\]
because that is the one-sided bar functor for the algebra $A^e$ rewritten using the enveloping
identification.  The inverse is the corresponding two-sided cobar functor.
\end{proof}

\begin{theorem}[Diagonal correspondence]
\ClaimStatusProvedHere
\label{thm:diagonal}
Under the equivalence of Theorem~\ref{thm:bimod-bicomod}, the diagonal bimodule ${}_A A_A$ is sent
to the diagonal bicomodule $C_\Delta$.  Equivalently,
\[
\KK_A^{\mathrm{bi}}(A)\simeq C_\Delta
\]
in $\Dco_{\mathrm{bi}}(C)$.
\end{theorem}

\begin{proof}
By definition,
\[
\KK_A^{\mathrm{bi}}(A)=C\widehat\otimes A\widehat\otimes C
\]
with the two-sided twisting differential.

First contract the left half:
the standard bar--cobar counit resolution yields a quasi-isomorphism
\[
C\widehat\otimes_\tau A\xrightarrow{\sim} \mathbf 1
\]
compatible with the right $C$-coaction.  Hence
\[
\KK_A^{\mathrm{bi}}(A)\simeq C
\]
as a right $C$-comodule.

Similarly, contracting the right half using
\[
A\widehat\otimes_\tau C\xrightarrow{\sim}\mathbf 1
\]
gives
\[
\KK_A^{\mathrm{bi}}(A)\simeq C
\]
as a left $C$-comodule.

To identify the opposite coactions, filter $\KK_A^{\mathrm{bi}}(A)$ by total coradical degree in the
two outer copies of $C$.  On the associated graded the twisting terms disappear, and the bicomodule
structure is literally the left and right coproducts on the two outer factors.  Under either of the
two contractions above, the associated graded collapses to a single copy of $C$, and the surviving
left and right coactions are exactly the two legs of $\Delta_C$.  Because the filtration is complete
and bounded below in the strongly admissible regime, this associated-graded identification lifts to
the filtered complex.  Therefore the resulting bicomodule is $C_\Delta$.
\end{proof}

\begin{corollary}[The diagonal is the unit for composition]
\ClaimStatusProvedHere
\label{cor:unit}
For every $X\in \Dco_{\mathrm{bi}}(C)$ there are canonical isomorphisms
\[
C_\Delta \square_C^{\mathbf R} X \simeq X \simeq X \square_C^{\mathbf R} C_\Delta.
\]
\end{corollary}

\begin{proof}
Transport the tautologies
\[
A\otimes_A^{\mathbf L} M\simeq M\simeq M\otimes_A^{\mathbf L} A
\]
for $A$-bimodules along Theorem~\ref{thm:bimod-bicomod} and use
Theorem~\ref{thm:diagonal}.
\end{proof}

\begin{corollary}[Tensor--cotensor gluing]
\ClaimStatusProvedHere
\label{cor:tensor-cotensor}
For all $M,N\in \Dpbw_{\mathrm{bi}}(A)$ there is a canonical isomorphism
\[
\KK_A^{\mathrm{bi}}(M\otimes_A^{\mathbf L} N)
\simeq
\KK_A^{\mathrm{bi}}(M)\square_C^{\mathbf R}\KK_A^{\mathrm{bi}}(N)
\]
in $\Dco_{\mathrm{bi}}(C)$.
\end{corollary}

\begin{proof}
Transport the derived tensor product across the equivalence of
Theorem~\ref{thm:bimod-bicomod}.
\end{proof}

\begin{remark}
Corollary~\ref{cor:tensor-cotensor} is the coalgebraic gluing law for intervals.
The diagonal bicomodule $C_\Delta$ is therefore the universal boundary condition for interval
composition.
\end{remark}

\section{Hochschild--coHochschild duality}
\label{sec:HH-coHH}

Annular traces and centers on the algebra side are identified with coHochschild theories of the bar
coalgebra.

\begin{theorem}[Associative chiral Hochschild/coHochschild homology]
\ClaimStatusProvedHere
\label{thm:HH-coHH-homology}
Let $A$ be strongly admissible, let $(A,A^{\op})$ be EZ-admissible, and set $C=\Barch(A)$.
Then for every complete $A$-bimodule $M$ there is a canonical isomorphism
\[
\HH^{\mathrm{ch}}_\bullet(A,M)\simeq
\coHoch^{\mathrm{ch}}_\bullet(C,\KK_A^{\mathrm{bi}}(M)).
\]
In particular,
\[
\HH^{\mathrm{ch}}_\bullet(A)\simeq \coHoch^{\mathrm{ch}}_\bullet(C)
\simeq
C_\Delta \square_{C^e}^{\mathbf R} C_\Delta.
\]
\end{theorem}

\begin{proof}
By definition,
\[
\HH^{\mathrm{ch}}_\bullet(A,M)=\Tor_\bullet^{A^e}(A,M).
\]
Apply the equivalence of Theorem~\ref{thm:bimod-bicomod} to the pair $(A,M)$.
The diagonal bimodule $A$ goes to $C_\Delta$ by Theorem~\ref{thm:diagonal}, and $M$ goes to
$\KK_A^{\mathrm{bi}}(M)$ by definition.  Derived tensor over $A^e$ is transported to derived cotensor
over $C^e$.  Therefore
\[
\Tor_\bullet^{A^e}(A,M)\simeq \Cotor_\bullet^{C^e}(C_\Delta,\KK_A^{\mathrm{bi}}(M)),
\]
which is the stated formula.  Taking $M=A$ yields the second assertion.
\end{proof}

\begin{theorem}[Associative chiral Hochschild/coHochschild cohomology]
\ClaimStatusProvedHere
\label{thm:HH-coHH-cohomology}
Under the same hypotheses, for every complete $A$-bimodule $M$ there is a canonical isomorphism
\[
\HH_{\mathrm{ch}}^\bullet(A,M)\simeq
\coHoch_{\mathrm{ch}}^\bullet(C,\KK_A^{\mathrm{bi}}(M)).
\]
In particular,
\[
\HH_{\mathrm{ch}}^\bullet(A)\simeq
\coHoch_{\mathrm{ch}}^\bullet(C)
=
\Coext_{C^e}^\bullet(C_\Delta,C_\Delta).
\]
\end{theorem}

\begin{proof}
Derived Hom is preserved by the equivalence of Theorem~\ref{thm:bimod-bicomod}.  Therefore
\[
\RHom_{A^e}(A,M)\simeq \RHom_{C^e}(C_\Delta,\KK_A^{\mathrm{bi}}(M)),
\]
which is the desired coHochschild cohomology complex.  The specialization to $M=A$ follows from the
diagonal correspondence.
\end{proof}

\begin{corollary}[The annulus as self-cotrace]
\ClaimStatusProvedHere
\label{cor:annulus}
The annular closed object attached to the universal boundary condition $C_\Delta$ is the coHochschild
object
\[
\mathcal H_C^{\mathrm{cl}}
:=
C_\Delta \square_{C^e}^{\mathbf R} C_\Delta
\simeq
\coHoch^{\mathrm{ch}}_\bullet(C)
\simeq
\HH^{\mathrm{ch}}_\bullet(A).
\]
\end{corollary}

\begin{proof}
This is the special case $M=A$ of Theorem~\ref{thm:HH-coHH-homology}.
\end{proof}

\begin{corollary}[Cap action]
\ClaimStatusProvedHere
\label{cor:cap}
There is a canonical action
\[
\coHoch^{\mathrm{ch}}_\bullet(C)\square_C^{\mathbf R} C_\Delta \longrightarrow C_\Delta
\]
obtained by transporting the Hochschild cap action
\[
\HH^{\mathrm{ch}}_\bullet(A)\otimes A\longrightarrow A
\]
across the equivalence of Theorem~\ref{thm:bimod-bicomod}.
\end{corollary}

\begin{proof}
Apply Theorem~\ref{thm:bimod-bicomod} to the classical cap action viewed as a map of bimodules.
\end{proof}

\begin{remark}
The cohomology theorem is the center theorem.  The homology theorem is the annulus theorem.  Together
they say that the algebraic center and algebraic trace of $A$ can be read directly from the
coalgebraic diagonal $C_\Delta$.
\end{remark}

\section{The diagonal bicomodule as universal boundary condition}
\label{sec:open-trace}

There are two operations one wants on the boundary side.

First, interval composition, already realized by derived cotensor and controlled by
Corollary~\ref{cor:tensor-cotensor}.  Second, the ordered pair-of-pants operation, which is not
cotensor but rather the transport of the external tensor product of bimodules.  This second
operation is inherently ordered; in general we do \emph{not} claim it is symmetric without extra
structure.

\begin{definition}[Ordered fusion product]
Let $X,Y\in \Dco_{\mathrm{bi}}(C)$.  Define
\[
X\star Y
:=
\KK_A^{\mathrm{bi}}\Bigl(
\Omega_C^{\mathrm{bi}}(X)\widehat\otimes \Omega_C^{\mathrm{bi}}(Y)
\Bigr),
\]
where $\widehat\otimes$ is the completed tensor product of bimodules with outer $A$-actions:
the left $A$-action is on the first factor and the right $A$-action is on the second factor.
We call $\star$ the \emph{ordered fusion product}.
\end{definition}

\begin{remark}
The adjective ``ordered'' is essential.  Unless further symmetry data are supplied, $\star$ is only
a monoidal product reflecting the ordering of incoming open boundaries.  It is the algebraic shadow
of an ordered pair-of-pants, not yet of the fully symmetric open cobordism category.
\end{remark}

\begin{theorem}[Ordered pair-of-pants algebra]
\ClaimStatusProvedHere
\label{thm:pair-of-pants}
The diagonal bicomodule $C_\Delta$ carries a canonical associative algebra structure with respect to
the ordered fusion product:
\[
\mu_P:C_\Delta\star C_\Delta\longrightarrow C_\Delta,
\qquad
\eta_P:\mathbf 1_\star\longrightarrow C_\Delta,
\]
where $\mathbf 1_\star:=\KK_A^{\mathrm{bi}}(k)$ is the image of the trivial bimodule.
\end{theorem}

\begin{proof}
By Theorem~\ref{thm:diagonal}, $C_\Delta\simeq \KK_A^{\mathrm{bi}}(A)$.  The multiplication and unit
of the algebra $A$,
\[
\mu_A:A\widehat\otimes A\to A,\qquad \eta_A:k\to A,
\]
transport across $\KK_A^{\mathrm{bi}}$ to maps
\[
\mu_P:C_\Delta\star C_\Delta\to C_\Delta,\qquad \eta_P:\mathbf 1_\star\to C_\Delta.
\]
Associativity and unitality follow immediately from associativity and unitality of $A$.
\end{proof}

\begin{definition}[Ordered open trace formalism]
An \emph{ordered genus-zero open trace formalism} consists of:
\begin{enumerate}[label=(\roman*)]
\item a composition product $\square_C^{\mathbf R}$ with unit object $C_\Delta$,
\item an ordered fusion product $\star$,
\item an associative algebra structure $(C_\Delta,\mu_P,\eta_P)$ for $\star$,
\item a closed object $\mathcal H_C^{\mathrm{cl}}$ obtained by annular closure,
\item an annulus action of $\mathcal H_C^{\mathrm{cl}}$ on $C_\Delta$,
\end{enumerate}
subject to the sewing relations induced from the algebra side.
\end{definition}

\begin{theorem}[Ordered genus-zero open trace formalism]
\ClaimStatusProvedHere
\label{thm:ordered-open}
The data
\[
\bigl(\Dco_{\mathrm{bi}}(C),\square_C^{\mathbf R},C_\Delta,\star,\mu_P,\eta_P,\mathcal H_C^{\mathrm{cl}}\bigr)
\]
constitute an ordered genus-zero open trace formalism.  More precisely:
\begin{enumerate}[label=(\alph*)]
\item $C_\Delta$ is the unit for interval composition $\square_C^{\mathbf R}$.
\item $(C_\Delta,\mu_P,\eta_P)$ is an associative algebra object for ordered fusion.
\item Annular closure is
\[
\mathcal H_C^{\mathrm{cl}}\simeq \coHoch^{\mathrm{ch}}_\bullet(C).
\]
\item The annulus acts on the boundary object by the action of
Corollary~\ref{cor:cap}.
\item All ordered genus-zero sewing identities are transported from the corresponding identities for
$A$-bimodules.
\end{enumerate}
\end{theorem}

\begin{proof}
Statement (a) is Corollary~\ref{cor:unit}.  Statement (b) is
Theorem~\ref{thm:pair-of-pants}.  Statement (c) is
Corollary~\ref{cor:annulus}.  Statement (d) is Corollary~\ref{cor:cap}.
Finally, every ordered genus-zero sewing relation on the coalgebra side is the image under the
equivalence $\KK_A^{\mathrm{bi}}$ of the corresponding relation on the algebra side.
\end{proof}

\begin{remark}[What is proved and what is not]
The theorem gives the algebraic ordered open package.  To upgrade it to a fully symmetric
one-brane open TCFT one needs additional structure that symmetrizes ordered fusion (for example, a
boundary symmetry datum or a direct geometric realization by bordered configuration spaces).
This is formulated below as part of the conjectural frontier.
\end{remark}

\subsection{Calabi--Yau transport}

The ordered open trace formalism acquires a Frobenius structure under a Calabi--Yau hypothesis on
$A$.

\begin{definition}[Complete chiral Calabi--Yau structure]
A strongly admissible associative chiral algebra $A$ is called \emph{$d$-Calabi--Yau} if:
\begin{enumerate}[label=(\alph*)]
\item $A$ is perfect as an $A^e$-module, and
\item there is an isomorphism in $\Dpbw_{\mathrm{bi}}(A)$
\[
\varphi_A:A\xrightarrow{\sim}\RHom_{A^e}(A,A^e)[d].
\]
\end{enumerate}
Equivalently, $A$ carries a nondegenerate $d$-shifted derived trace pairing.
\end{definition}

\begin{theorem}[Shifted ordered Frobenius structure]
\ClaimStatusProvedHere
\label{thm:CY}
Assume $A$ is $d$-Calabi--Yau.  Then the underlying complex of $C_\Delta$ carries a canonical
$d$-shifted ordered Frobenius structure.  More precisely:
\begin{enumerate}[label=(\roman*)]
\item the ordered pair-of-pants multiplication $\mu_P$ and unit $\eta_P$ are as above;
\item there is a trace
\[
\eps_P:C_\Delta\to k[-d]
\]
transported from the Calabi--Yau trace on $A$;
\item there is a coproduct
\[
\Delta_P:C_\Delta\to C_\Delta\star C_\Delta[d]
\]
adjoint to $\mu_P$ under $\eps_P$;
\item the ordered Frobenius identities
\[
(\mu_P\star \id)\circ(\id\star \Delta_P)
=
\Delta_P\circ \mu_P
=
(\id\star\mu_P)\circ(\Delta_P\star \id)
\]
hold.
\end{enumerate}
\end{theorem}

\begin{proof}
On the algebra side, the Calabi--Yau structure identifies $A$ with its bimodule dual shifted by
$d$.  Therefore multiplication
\[
\mu_A:A\widehat\otimes A\to A
\]
admits a unique adjoint coproduct
\[
\Delta_A:A\to A\widehat\otimes A[d]
\]
with respect to the trace pairing $\eps_A:A\to k[-d]$; explicitly one may define $\Delta_A$ as the
adjoint of $\mu_A$ under $\varphi_A$.  The ordered Frobenius identities are formal consequences of
associativity of $\mu_A$ and the nondegeneracy of the pairing.

Transport the entire package $(\mu_A,\eta_A,\eps_A,\Delta_A)$ across the equivalence
$\KK_A^{\mathrm{bi}}$ and use Theorem~\ref{thm:diagonal}.  The resulting structure maps are
$\mu_P,\eta_P,\eps_P,\Delta_P$, and all identities are preserved because the equivalence preserves
composition.
\end{proof}

\begin{corollary}[Cardy operator on the coalgebra side]
\ClaimStatusProvedHere
\label{cor:cardy}
Under the hypotheses of Theorem~\ref{thm:CY}, the ordered annulus endomorphism of the universal
boundary object is
\[
\mathsf{Ann}_{C_\Delta}:=\mu_P\circ \Delta_P\in \operatorname{End}(C_\Delta)[d].
\]
Its annular closure factors through the closed object
\[
\mathcal H_C^{\mathrm{cl}}\simeq \coHoch^{\mathrm{ch}}_\bullet(C).
\]
\end{corollary}

\begin{proof}
Transport the usual annulus endomorphism $\mu_A\circ \Delta_A$ of the Calabi--Yau algebra $A$ across
$\KK_A^{\mathrm{bi}}$ and use Corollary~\ref{cor:annulus}.
\end{proof}

\section{A master theorem and its conceptual form}
\label{sec:master}

The preceding results assemble into a single statement.

\begin{theorem}[Master theorem]
\ClaimStatusProvedHere
\label{thm:master}
Let $A$ be a strongly admissible augmented associative chiral algebra on $X$, assume
$(A,A^{\op})$ is EZ-admissible, and set
\[
C:=\Barch(A).
\]
Then:
\begin{enumerate}[label=(\arabic*)]
\item \textbf{Ordered shuffle.}
There is a quasi-isomorphism of conilpotent dg coalgebras
\[
\Barch(A\chotimes B)\simeq \Barch(A)\widehat\otimes \Barch(B)
\]
for every EZ-admissible partner $B$.

\item \textbf{Opposite-duality.}
There is a canonical isomorphism
\[
\Barch(A^{\op})\simeq \Barch(A)^{\cop}.
\]

\item \textbf{Enveloping duality.}
There is a canonical quasi-isomorphism
\[
\Barch(A^e)\simeq C^e.
\]

\item \textbf{Two-sided bimodule/bicomodule duality.}
There is an exact equivalence
\[
\Dpbw_{\mathrm{bi}}(A)\simeq \Dco_{\mathrm{bi}}(C).
\]

\item \textbf{Diagonal correspondence.}
Under this equivalence one has
\[
{}_A A_A \longleftrightarrow C_\Delta.
\]

\item \textbf{Composition.}
The diagonal bicomodule is the unit for derived cotensor:
\[
C_\Delta\square_C^{\mathbf R} X\simeq X\simeq X\square_C^{\mathbf R}C_\Delta.
\]

\item \textbf{Hochschild/coHochschild homology.}
For every complete $A$-bimodule $M$,
\[
\HH^{\mathrm{ch}}_\bullet(A,M)
\simeq
\coHoch^{\mathrm{ch}}_\bullet(C,\KK_A^{\mathrm{bi}}(M)).
\]

\item \textbf{Hochschild/coHochschild cohomology.}
For every complete $A$-bimodule $M$,
\[
\HH_{\mathrm{ch}}^\bullet(A,M)
\simeq
\coHoch_{\mathrm{ch}}^\bullet(C,\KK_A^{\mathrm{bi}}(M)).
\]

\item \textbf{Tangent theory.}
For a square-zero extension $A_\eps=A\oplus \eps M$,
\[
\Barch(A_\eps)\cong C\oplus \eps\,\KK_A^{\mathrm{bi}}(M)\qquad (\mathrm{mod}\ \eps^2).
\]

\item \textbf{Ordered open trace formalism.}
The coalgebraic boundary object $C_\Delta$ carries ordered pair-of-pants, annulus, and composition
operations transported from the algebra side, and if $A$ is $d$-Calabi--Yau then the underlying
complex of $C_\Delta$ becomes a $d$-shifted ordered Frobenius object.
\end{enumerate}
\end{theorem}

\begin{proof}
This is a summary of Theorems~\ref{thm:shuffle}, \ref{thm:opposite},
\ref{thm:bimod-bicomod}, \ref{thm:diagonal}, \ref{thm:HH-coHH-homology},
\ref{thm:HH-coHH-cohomology}, \ref{thm:tangent=K}, and \ref{thm:ordered-open}, together with
Corollaries~\ref{cor:enveloping}, \ref{cor:unit}, and \ref{cor:infdual}, and
Theorem~\ref{thm:CY}.
\end{proof}


\section{Configuration space geometry of the ordered bar complex}
\label{sec:config-space-geometry}

The algebraic theory of the preceding sections rests on
configuration spaces whose geometry we now make explicit.
The ordered bar complex lives on the product of holomorphic
and topological configuration spaces; the bar differential
comes from the holomorphic direction; the coproduct comes
from the topological direction.

\subsection{Ordered configuration spaces}
\label{subsec:ordered-config-spaces-geom}

\begin{definition}[Ordered real configuration space]
\ClaimStatusProvedHere
\label{def:ordered-real-config}
For $k\ge 1$, the \emph{ordered real configuration space} is
\[
\mathrm{Conf}_k^{<}(\mathbb R)
\;:=\;
\bigl\{(t_1,\dots,t_k)\in\mathbb R^k
\;\big|\;
t_1<t_2<\cdots<t_k\bigr\}.
\]
This is a single open convex cell in~$\mathbb R^k$,
hence contractible.
\end{definition}

\begin{proposition}[Topology of ordered real configurations;
\ClaimStatusProvedHere]
\label{prop:ordered-real-config-topology}
\begin{enumerate}[label=\textup{(\roman*)}]
\item $\mathrm{Conf}_k^{<}(\mathbb R)$ is contractible
  (it is a connected component of
  $\mathrm{Conf}_k^{\mathrm{ord}}(\mathbb R)$, which
  consists of $k!$~contractible chambers).
\item The natural projection
  $\pi_k\colon \mathrm{Conf}_k^{\mathrm{ord}}(\mathbb R)
  \to \mathrm{Conf}_k(\mathbb R)/\Sigma_k$
  is a $k!$-sheeted covering.  The fibre over a point
  $\{t_1,\dots,t_k\}\in\mathrm{Conf}_k(\mathbb R)/\Sigma_k$
  consists of the $k!$~orderings of the~$t_i$.
\item $\mathrm{Conf}_k^{<}(\mathbb R)$ is a fundamental domain
  for the $\Sigma_k$-action: every unordered configuration
  can be uniquely ordered.  This discreteness --- each fibre
  is a set, not a path-connected space --- is the geometric
  reason that the $E_1$-structure (associativity) carries no
  homotopy-coherence data beyond the sequence of
  associahedra.
\end{enumerate}
\end{proposition}

\begin{definition}[Ordered holomorphic configuration space]
\ClaimStatusProvedHere
\label{def:ordered-hol-config}
For a smooth curve~$X$ and $k\ge 1$, the \emph{ordered
configuration space} is
\[
\mathrm{Conf}_k^{\mathrm{ord}}(X)
\;:=\;
\bigl\{(z_1,\dots,z_k)\in X^k
\;\big|\;
z_i\neq z_j\;\text{for all}\;i\neq j\bigr\}.
\]
The symmetric group~$\Sigma_k$ acts freely by permuting
coordinates.  The unordered configuration space is the
quotient:
$\mathrm{Conf}_k(X)
=\mathrm{Conf}_k^{\mathrm{ord}}(X)/\Sigma_k$.
For $X=\mathbb C$:
$\pi_1(\mathrm{Conf}_k^{\mathrm{ord}}(\mathbb C))=P_k$,
the pure braid group;
$\pi_1(\mathrm{Conf}_k(\mathbb C))=B_k$,
the full braid group.
\end{definition}

\subsection{The product space and its operadic role}
\label{subsec:product-space}

\begin{construction}[The SC$^{\mathrm{ch,top}}$ operation space;
\ClaimStatusProvedHere]
\label{constr:sc-operation-space}
The two-coloured Swiss-cheese operad
$\mathrm{SC}^{\mathrm{ch,top}}$ has operation spaces
\[
\mathrm{SC}^{\mathrm{ch,top}}(k,m)
\;=\;
\FM_k(\mathbb C)\times E_1(m),
\]
where $\FM_k(\mathbb C)$ is the Fulton--MacPherson
compactification of $\mathrm{Conf}_k(\mathbb C)$
(the closed/holomorphic colour) and $E_1(m)$ is
the little-intervals operad (the open/topological colour).
Concretely, $E_1(m)\simeq \mathrm{Conf}_m^{<}(\mathbb R)$
up to translation and scaling, which is contractible.

For the ordered bar complex of an $E_1$-chiral algebra,
the relevant spaces at tensor degree~$k$ are
\[
\FM_k(\mathbb C)\times \mathrm{Conf}_k^{<}(\mathbb R).
\]
A bar element of degree~$k$ is parametrised by this product:
the coordinates $z_1,\dots,z_k\in\mathbb C$ live on the
holomorphic curve, while the ordering
$t_1<t_2<\cdots<t_k$ records the position in the
topological direction~$\mathbb R_t$.  The bar differential
acts in the $\mathbb C$-direction (extracting OPE residues);
the deconcatenation coproduct acts in the $\mathbb R$-direction
(cutting the ordered sequence).
\end{construction}

\begin{remark}[The two colours]
\label{rem:two-colours-geom}
The product $\FM_k(\mathbb C)\times\mathrm{Conf}_k^{<}(\mathbb R)$
is the geometric fingerprint of a 3d holomorphic-topological
field theory on $\mathbb C_z\times\mathbb R_t$:
observables factorise holomorphically in~$z$ and
associatively in~$t$.  The bar differential is the
closed (holomorphic) colour.  The bar coproduct is the
open (topological) colour.
\end{remark}

\subsection{The ordered FM compactification}
\label{subsec:ordered-fm-compact}

\begin{construction}[Ordered Fulton--MacPherson compactification;
\ClaimStatusProvedHere]
\label{constr:ordered-fm-compact}
The ordered FM compactification
$\overline{\FM}_k^{\mathrm{ord}}(\mathbb C)$ is the
iterated real oriented blowup of
$\mathrm{Conf}_k^{\mathrm{ord}}(\mathbb C)$ along the
\emph{ordered} diagonals
$D_{i,i+1}=\{z_i=z_{i+1}\}$ for $i=1,\dots,k-1$.
(In the unordered FM compactification, one blows up along
all diagonals $D_S=\{z_i=z_j:i,j\in S\}$ for all subsets
$S\subseteq\{1,\dots,k\}$ with $|S|\ge 2$.  In the ordered
case, only \emph{consecutive} diagonals appear.)

The codimension-one boundary faces correspond to
\emph{consecutive} collisions $z_i\to z_{i+1}$, not to
arbitrary subsets.  A non-consecutive collision such as
$z_1\to z_3$ (with $z_2$ remaining separated) is not a
codimension-one face of the ordered compactification.

\medskip
\noindent\emph{Low-degree examples.}

For $k=2$:
$\overline{\FM}_2^{\mathrm{ord}}(\mathbb C)
\cong\overline{\mathbb C}$ (the Riemann sphere), with a
single boundary divisor $D_{1,2}=\{z_1=z_2\}$.

For $k=3$: two codimension-one faces: $D_{1,2}$ (first pair
collides) and $D_{2,3}$ (last pair collides).  The face
$D_{1,3}$ (first and third collide without the middle) is
codimension~$2$: it requires both $D_{1,2}$ and $D_{2,3}$ to
degenerate simultaneously.  This is the associahedron~$K_3$:
a line segment with two endpoints.
\end{construction}

\subsection{Boundary strata: planted forests}
\label{subsec:boundary-strata-planted}

\begin{construction}[Boundary stratification by planted forests;
\ClaimStatusProvedHere]
\label{constr:planted-forests}
The boundary strata of $\overline{\FM}_k^{\mathrm{ord}}(\mathbb C)$
are indexed by \emph{planted forests}: ordered sequences of
planar rooted trees.

A \emph{planted forest} $F=(T_1,T_2,\dots,T_r)$ is an
\emph{ordered} sequence of planar rooted trees whose leaves,
read from left to right across all trees, are labelled
$1,2,\dots,k$ in order.  The trees $T_1,\dots,T_r$ correspond
to the $r$~clusters of points that remain separated at the
given codimension; within each tree, the internal vertices
record the collision history of points within that cluster.

The stratum indexed by~$F$ is the product
\[
\partial_F\overline{\FM}_k^{\mathrm{ord}}(\mathbb C)
\;\cong\;
\prod_{j=1}^r
\overline{\FM}_{|T_j|}^{\mathrm{ord}}(\mathbb C),
\]
where $|T_j|$ is the number of leaves of~$T_j$.

For the codimension-one faces: each is indexed by a forest
with $k{-}1$~trees, $k{-}2$ of which are single leaves and one
of which is a single binary tree with two leaves.
Concretely, the face $D_{i,i+1}$ corresponds to the forest in
which points $i$ and $i{+}1$ form a binary tree (they collide)
and all other points are isolated.

\emph{Contrast with the unordered case.}
In the unordered FM compactification, boundary strata are
indexed by \emph{trees} (rooted trees with unordered children).
In the ordered case, the strata are indexed by
\emph{planted forests} (ordered sequences of planar trees) ---
the ordering of trees within the forest records the linear
order on the topological direction.
\end{construction}

\subsection{The bar differential from collision residues}
\label{subsec:bar-diff-collision}

\begin{construction}[Bar differential; \ClaimStatusProvedHere]
\label{constr:bar-diff-collision}
Let $\cA$ be a strongly admissible $E_1$-chiral algebra with
structure constants $e_I{}_{(n)}\,e_J=\sum_K c^K_{IJ;n}\,e_K$
relative to a homogeneous basis $\{e_I\}$ of the
augmentation ideal~$\bar\cA$.  The bar differential on
$\Barch^{\mathrm{ord}}(\cA)$ decomposes as
\[
d_{\Barch^{\mathrm{ord}}}
\;=\;
d_{\mathrm{int}}+\sum_{i=1}^{k-1} d_{i,i+1},
\]
where:
\begin{itemize}
\item $d_{\mathrm{int}}$ applies the internal differential
  of~$\cA$ to each tensor factor.
\item $d_{i,i+1}$ extracts the collision residue at the
  codimension-one stratum $D_{i,i+1}=\{z_i=z_{i+1}\}$,
  using the ordered logarithmic propagator
  $\eta_{i,i+1}=d\log(z_i-z_{i+1})$.
\end{itemize}

Explicitly, the component $d_{i,i+1}$ acts on a bar word
$[s^{-1}e_{I_1}|\cdots|s^{-1}e_{I_k}]$ by contracting
the $i$-th and $(i{+}1)$-st slots via all OPE products:
\[
d_{i,i+1}[s^{-1}e_{I_1}|\cdots|s^{-1}e_{I_k}]
\;=\;
\sum_{K,n}\pm\,c^K_{I_i I_{i+1};n}\;
[s^{-1}e_{I_1}|\cdots|s^{-1}e_K|\cdots|s^{-1}e_{I_k}]
\;\otimes\;[\delta_{D_{i,i+1}}^{(n)}],
\]
where the contracted pair is replaced by a single slot and
$[\delta_{D_{i,i+1}}^{(n)}]$ is the $n$-th residue class at the
collision divisor.

The sum runs only over \emph{consecutive} collisions
$i,i{+}1$: this is the ordered constraint.  In the unordered
bar complex (Vol.~I, Theorem~A), one sums over all pairs
$\{i,j\}$ and then takes $\Sigma_k$-coinvariants.

The identity $d^2=0$ follows from the Stokes theorem applied
to the boundary strata of
$\overline{\FM}_k^{\mathrm{ord}}(\mathbb C)$: the
codimension-two strata receive two contributions from the
two orderings of consecutive collisions, and the algebraic
content of their cancellation is precisely the $E_1$ axiom
(associativity of the zeroth product).
\end{construction}

\subsection{The deconcatenation coproduct}
\label{subsec:deconcatenation}

\begin{construction}[Deconcatenation coproduct;
\ClaimStatusProvedHere]
\label{constr:deconcatenation}
The ordered bar coalgebra carries a coassociative coproduct
$\Delta$ defined by \emph{deconcatenation}: cutting the
ordered sequence at every possible position.  For a bar word
of length~$k$:
\[
\Delta[s^{-1}e_{I_1}|\cdots|s^{-1}e_{I_k}]
\;=\;
\sum_{p=0}^{k}
[s^{-1}e_{I_1}|\cdots|s^{-1}e_{I_p}]
\otimes
[s^{-1}e_{I_{p+1}}|\cdots|s^{-1}e_{I_k}],
\]
where the $p=0$ and $p=k$ terms produce the coaugmentation.

This coproduct is the topological structure: it comes from
the ordered configuration space
$\mathrm{Conf}_k^{<}(\mathbb R)$ by cutting the
ordered real line at a point between $t_p$ and $t_{p+1}$.
The $p$ points to the left form one ordered subconfiguration;
the $k{-}p$ points to the right form another.  The coassociativity
of~$\Delta$ is the associativity of interval composition.

\emph{Contrast with the holomorphic differential.}
The bar differential~$d$ extracts residues in the holomorphic
direction (collisions of points on~$\mathbb C$).  The coproduct
$\Delta$ cuts in the topological direction (splitting the
ordered sequence on~$\mathbb R$).  Together, the pair
$(d,\Delta)$ encodes both colours of the Swiss-cheese
algebra:
\begin{center}
\begin{tabular}{lll}
\textbf{Structure} & \textbf{Direction} &
  \textbf{Geometric source} \\
\hline
Bar differential~$d$ & holomorphic ($\mathbb C_z$) &
  OPE residues on $\FM_k(\mathbb C)$ \\
Coproduct~$\Delta$ & topological ($\mathbb R_t$) &
  Deconcatenation of $\mathrm{Conf}_k^{<}(\mathbb R)$
\end{tabular}
\end{center}
\end{construction}

\subsection{The covering space frame}
\label{subsec:covering-space-frame}

\begin{construction}[Ordered-to-unordered covering;
\ClaimStatusProvedHere]
\label{constr:covering-space}
The projection
\[
\pi_k\colon
\mathrm{Conf}_k^{\mathrm{ord}}(\mathbb C)
\;\longrightarrow\;
\mathrm{Conf}_k(\mathbb C)
\;=\;
\mathrm{Conf}_k^{\mathrm{ord}}(\mathbb C)/\Sigma_k
\]
is a principal $\Sigma_k$-bundle, \'etale with fibre~$\Sigma_k$.
The fibre is a $k!$-element set.

In the topological direction, the analogous covering
$\mathrm{Conf}_k^{\mathrm{ord}}(\mathbb R)\to
\mathrm{Conf}_k(\mathbb R)/\Sigma_k$
is also a $k!$-sheeted covering, but here it is trivialised
by the total order on~$\mathbb R$: the component
$\mathrm{Conf}_k^{<}(\mathbb R)$ is a fundamental domain.
This triviality is why the topological direction carries
an $E_1$-structure (associative, no higher coherences from
the covering) while the holomorphic direction carries
monodromy (the $R$-matrix).

The bar complex at tensor degree~$k$ defines a local system
$\mathcal{F}_k^{\mathrm{ord}}$ on
$\mathrm{Conf}_k^{\mathrm{ord}}(\mathbb C)$
whose stalk at $(z_1,\dots,z_k)$ is
$(s^{-1}\bar\cA)^{\otimes k}$, with flat connection induced
by the residue differential $d_{\mathrm{res}}$.  The
$R$-matrix is the monodromy of this flat connection.

To descend $\mathcal{F}_k^{\mathrm{ord}}$ from ordered to
unordered configurations requires a $\Sigma_k$-equivariant
structure: isomorphisms
$R_\sigma\colon
\mathcal{F}_k^{\mathrm{ord}}\big|_{(z_1,\dots,z_k)}
\xrightarrow{\sim}
\mathcal{F}_k^{\mathrm{ord}}\big|_{(z_{\sigma(1)},\dots,z_{\sigma(k)})}$
for each $\sigma\in\Sigma_k$, satisfying the cocycle condition
and compatible with the flat connection.

For an $E_\infty$-chiral algebra (which is any vertex algebra,
including those with OPE poles), the equivariant structure
is canonical: $R_\sigma=\epsilon(\sigma)\cdot\tau_\sigma$,
where $\tau_\sigma$ permutes tensor factors and
$\epsilon(\sigma)$ is the Koszul sign.  The descent then
recovers the ordinary $\Sigma_k$-coinvariant bar complex
(Vol.~I, Theorem~A).

For a genuinely $E_1$-chiral algebra (where the ordering is
primitive, not derived from a local OPE), the equivariant
structure must be \emph{constructed} from the OPE data.
This construction is the $R$-matrix, and its existence
imposes the Yang--Baxter equation (the content of $d^2=0$
at tensor degree~$3$).
\end{construction}

\begin{remark}[Why ordered spaces]
\label{rem:why-ordered-spaces}
The use of ordered versus unordered configuration spaces
reflects the operadic symmetry of the algebra:
\begin{itemize}
\item An $E_\infty$-algebra (any vertex algebra, any BD chiral
  algebra) has operations that are $\Sigma_n$-equivariant:
  the factorisation structure is defined on the \emph{unordered}
  configuration space $\mathrm{Conf}_k(X)/\Sigma_k$.  The
  ordering on the bar complex is a computational convenience
  that can be quotiented out.
\item An $E_1$-algebra (an associative chiral algebra, a
  quantum vertex algebra in the sense of Etingof--Kazhdan)
  has operations indexed by \emph{ordered} compositions:
  its bar complex uses
  $\mathrm{Conf}_k^{\mathrm{ord}}(X)$.  The ordering
  records the topological direction: in
  $\mathbb C_z\times\mathbb R_t$, points along
  $\mathbb R_t$ are ordered by time, and the
  $E_1$-structure is the time-ordered product.
\item The Swiss-cheese operad $\mathrm{SC}^{\mathrm{ch,top}}$
  uses \emph{both}: unordered configurations in the
  holomorphic direction (closed colour) and ordered
  configurations in the topological direction (open colour).
  The ordered bar complex of a vertex algebra, which is
  $E_\infty$-chiral, uses the ordered configuration space
  for the open/$E_1$ colour of $\mathrm{SC}^{\mathrm{ch,top}}$
  --- this is the topological direction, not a claim that the
  algebra is nonlocal.  The holomorphic colour remains
  $E_\infty$.
\end{itemize}
The $R$-matrix $R(z)\in\operatorname{End}(V\otimes V)(\!(z)\!)$
is the data needed to pass from ordered to unordered
configurations in the holomorphic direction: it is the
monodromy of the local system on
$\mathrm{Conf}_2^{\mathrm{ord}}(\mathbb C)$ around the
generator of $\pi_1=\mathbb Z$.
\end{remark}


\section{Explicit computations for the standard families}
\label{sec:explicit-computations}

The master theorem applies to every strongly admissible chiral algebra.
The following computations treat the standard families in order of
complexity, providing closed-form answers against which the abstract
theory can be checked.

\subsection{The Heisenberg ordered bar complex}
\label{subsec:heisenberg-ordered-bar}

Let $V = \bC\cdot J$ be the one-dimensional vector space spanned by
the Heisenberg current, let $k\in\bC$ be the level, and let $\cH_k$
be the Heisenberg chiral algebra with OPE
\begin{equation}\label{eq:heisenberg-ope-ordered}
J(z)\,J(w) \;\sim\; \frac{k}{(z-w)^2}.
\end{equation}
The Heisenberg algebra is $E_\infty$-chiral: it is a local vertex
algebra with $\Sigma_n$-equivariant factorization structure.  Its OPE
has a double pole, but poles do not break $E_\infty$-locality.

The ordered bar complex $\Barch(\cH_k)$ is the cofree conilpotent
coalgebra on the desuspension $\susp^{-1}\overline{\cH}_k$, with
the ordering of tensor factors recording the linear order
$t_1 < \cdots < t_n$ on $\bR$.  We compute the complex degree by
degree.

\begin{theorem}[The Heisenberg ordered bar complex;
\ClaimStatusProvedHere]
\label{thm:heisenberg-ordered-bar}
The ordered bar complex $\Barch(\cH_k)$ is concentrated in the
following data:
\begin{enumerate}[label=\textup{(\arabic*)}]
\item \textbf{Degree~$1$.}\enspace
$\Barch_1(\cH_k) = \bC\cdot[\susp^{-1}J]$, with
$d_{\Barch}[\susp^{-1}J] = 0$
\textup{(}the bar differential is binary;
a single element has no adjacent pair to collide\textup{)}.

\item \textbf{Degree~$2$.}\enspace
$\Barch_2(\cH_k) = \bC\cdot[\susp^{-1}J \,|\, \susp^{-1}J]$,
a \emph{single} generator
\textup{(}the ordered complex retains the linear order;
there is no $\Sigma_2$-identification\textup{)}.
The bar differential is
\begin{equation}\label{eq:heisenberg-d2-ordered}
d_{\Barch}[\susp^{-1}J \,|\, \susp^{-1}J]
\;=\; J_{(0)}J \;=\; 0.
\end{equation}
The mode product $J_{(0)}J$ vanishes for the Heisenberg algebra:
the only nonvanishing OPE mode is $J_{(1)}J = k$, but the
ordered bar differential at degree~$2$ uses the zeroth mode
product $a_{(0)}b$ \textup{(}the $d\log$~kernel in the
collision residue absorbs one power of the pole, so the double
pole $k/(z-w)^2$ contributes through $\Res_{z=w}\frac{k}{(z-w)^2}
\cdot(z-w)\,d\log(z-w) = k\cdot d\log(z-w)|_{z=w}$, which
extracts $J_{(0)}J$, not $J_{(1)}J$\textup{)}.

Since $J_{(0)}J = 0$, the bar differential vanishes on
$\Barch_2$.

\item \textbf{Degree~$n$.}\enspace
$\Barch_n(\cH_k) = \bC\cdot[\susp^{-1}J \,|\,
\cdots \,|\, \susp^{-1}J]$ \textup{(}$n$~copies\textup{)},
with $d_{\Barch} = 0$ on every element.  The argument is
uniform: every summand in the bar differential involves an
adjacent collision $J_{(0)}J = 0$.

\item \textbf{Coproduct.}\enspace The deconcatenation coproduct
$\Delta$ on $\Barch(\cH_k)$ is
\begin{equation}\label{eq:heisenberg-delta-ordered}
\Delta[\susp^{-1}J \,|\, \cdots \,|\, \susp^{-1}J]
\;=\;
\sum_{i=0}^{n}
[\susp^{-1}J \,|\, \cdots \,|\, \susp^{-1}J]_{i}
\;\otimes\;
[\susp^{-1}J \,|\, \cdots \,|\, \susp^{-1}J]_{n-i},
\end{equation}
where the subscript denotes the number of tensor factors.
This is the tensor coalgebra coproduct, encoding interval-cutting
in $\bR$.

\item \textbf{Swiss-cheese structure.}\enspace
Since $d_{\Barch} = 0$, the entire Swiss-cheese algebra
$(\Barch(\cH_k), d_{\Barch}, \Delta)$ reduces to the cofree
conilpotent coalgebra $T^c(\susp^{-1}V)$ with trivial
differential and deconcatenation coproduct.  This is a
$\mathrm{SC}^{\mathrm{ch,top}}$-algebra in which every
closed-colour operation vanishes.
\end{enumerate}
\end{theorem}

\begin{proof}
Since $\cH_k$ is generated by a single field $J$ with
$J_{(0)}J = 0$, the bar differential (which extracts the
collision residue via the zeroth mode) vanishes on every bar
word.  More precisely: the OPE $J(z)J(w) \sim k/(z-w)^2$ has
$J_{(n)}J = 0$ for all $n \neq 1$, and $J_{(1)}J = k$.
The ordered bar differential extracts $J_{(0)}J$, which
vanishes.  All higher-arity components of $d_{\Barch}$
involve iterated zeroth-mode products, which vanish
\emph{a fortiori}.

The deconcatenation coproduct is manifestly coassociative
(it is the tensor coalgebra structure), and
$\Delta\circ d_{\Barch} = (d_{\Barch}\otimes\id +
\id\otimes d_{\Barch})\circ\Delta$ is automatic since
both sides vanish.
\end{proof}

\begin{remark}[The role of $J_{(1)}J = k$]
\label{rem:heisenberg-j1j}
The mode $J_{(1)}J = k$ is not lost: it appears in the
\emph{symmetric} bar complex $\barB^{\Sigma}(\cH_k)$
(Vol~I, Theorem~A) as the bar differential
$d_{\barB}[\susp^{-1}J \,|\, \susp^{-1}J] = k$.
In the ordered complex, the same information is carried
by the \emph{$R$-matrix} rather than the differential
(see Theorem~\ref{thm:heisenberg-rmatrix} below).
This is the ordered/symmetric distinction at work: the
symmetric bar complex quotients by $\Sigma_n$ using the
$R$-matrix as twisting datum, and the $J_{(1)}J = k$
residue reappears in that quotient.
\end{remark}

\begin{theorem}[Collision residue and $R$-matrix;
\ClaimStatusProvedHere]
\label{thm:heisenberg-rmatrix}
The collision residue of the Heisenberg OPE on
$\Conf_2^{<}(\bC)$ is
\begin{equation}\label{eq:heisenberg-collision-residue}
r^{\mathrm{coll}}(z)
\;=\;
\frac{k}{z},
\end{equation}
a simple pole in the collision residue arising from the
double pole in the OPE \textup{(}the $d\log$~kernel absorbs
one power: $\Res_{w=0}\frac{k}{w^2}\cdot w\,d\log w = k/z$;
cf.\ AP19\textup{)}.

The flat connection on the trivial bundle over
$\Conf_2^{<}(\bC)$ determined by the bar propagator is
\begin{equation}\label{eq:heisenberg-connection}
\nabla \;=\; d \;-\; k\cdot d\log(z),
\end{equation}
and its monodromy around the origin is
\begin{equation}\label{eq:heisenberg-monodromy}
\operatorname{Mon}_0(\nabla) \;=\; e^{-2\pi ik},
\end{equation}
a scalar phase.

The spectral $R$-matrix on modules is
\begin{equation}\label{eq:heisenberg-R-matrix}
R(z) \;=\; e^{k\hbar/z}.
\end{equation}
Since $\cH_k$ is $E_\infty$-chiral, this $R$-matrix is
\emph{derived} from the local OPE via analytic continuation:
it is the monodromy of the flat connection
$\nabla = d - r^{\mathrm{coll}}(z)\,dz/z$, not independent
input.  The Yang--Baxter equation is satisfied trivially
because $R(z)$ is scalar \textup{(}$\cH_k$ has rank~$1$\textup{)}.
\end{theorem}

\begin{proof}
The OPE $J(z)J(w) \sim k/(z-w)^2$ has a double pole.
The bar propagator is $\omega_{12} = d\log(z_1 - z_2)$;
integrating the OPE coefficient against this propagator
over the exceptional divisor $E_{12} \cong S^1$ in
$\FM_2(\bC)$ extracts
\[
r^{\mathrm{coll}}(z)
\;=\;
\Res_{w=z}\frac{k}{(z-w)^2}\cdot(z-w)
\;=\; \frac{k}{z}
\]
(the factor $(z-w)$ comes from $d\log(z-w) = dw/(z-w)$;
taking the residue at $w = z$ gives $k/z$).

The connection $\nabla = d - (k/z)\,dz$ on the trivial
line bundle over $\bC^\times$ has flat sections
$f(z) = z^k$; the monodromy of this multivalued
function around the origin is $e^{2\pi ik}$ on the
solution space, which acts as $e^{-2\pi ik}$ on the
parallel transport.

The spectral $R$-matrix is obtained by regularizing the
connection on modules.  For the rank-$1$ Heisenberg at
level~$k$, the $R$-matrix acts on $\cF_\mu\otimes\cF_\nu$
(Fock modules of charges $\mu,\nu$) as
$R(z) = e^{k\hbar/z}$ where $\hbar = \mu\nu$; since all
computations are scalar, this reduces to the exponential
of the collision residue.
\end{proof}

\begin{theorem}[Open-colour Koszul dual: the abelian Yangian;
\ClaimStatusProvedHere]
\label{thm:heisenberg-yangian}
The open-colour \textup{(}ordered, associative\textup{)} Koszul
dual of $\cH_k$ is the abelian Yangian
\begin{equation}\label{eq:heisenberg-yangian}
\cH_k^{!,\mathrm{ord}}
\;\simeq\;
Y\bigl(\mathfrak{u}(1)\bigr)
\;\simeq\;
\bC[h_0, h_1, h_2, \ldots],
\end{equation}
the polynomial algebra on generators $h_n$ of degree~$n+1$,
with product inherited from the linear dual of the
deconcatenation coproduct.  The spectral-parameter
presentation identifies $h_n = \partial_u^n h(u)|_{u=0}$,
where $h(u) = \sum_{n\ge 0} h_n u^{-n-1}$ is the generating
current of the Yangian.
\end{theorem}

\begin{proof}
Since the bar differential vanishes on $\Barch(\cH_k)$,
the Koszul dual is the linear dual of the bar coalgebra:
\[
\cH_k^{!,\mathrm{ord}}
\;=\;
\bigl(\Barch(\cH_k)\bigr)^\vee
\;=\;
\bigl(T^c(\susp^{-1}V)\bigr)^\vee
\;\simeq\;
T(\susp V^\vee)
\;\simeq\;
\bC[h],
\]
where $h = (\susp^{-1}J)^\vee$ is the dual generator.
The coalgebra $T^c(\susp^{-1}V)$ has deconcatenation
coproduct; dualizing produces the concatenation
(polynomial) product on $T(\susp V^\vee) = \bC[h]$.

Since $V$ is one-dimensional and $d_{\Barch} = 0$,
the Koszul dual algebra is freely generated by $h$ with
no relations.  Restoring the spectral grading from the
FM configuration-space stratification, the degree-$n$
component corresponds to $\FM_n(\bC)$, giving the
generators $h_0, h_1, h_2, \ldots$ of the abelian
Yangian $Y(\mathfrak{u}(1))$.

Equivalently: the Koszul dual is the endomorphism algebra
of the line operator $b$ in the open-colour category
$\cC_{\mathrm{line}}$, which for the abelian theory is
$Y(\mathfrak{u}(1)) \simeq \cH_{-k}$ as an associative
algebra (the level-negation is the complementarity
involution $\kappa\mapsto -\kappa$ for the Kac--Moody
family).
\end{proof}

\begin{theorem}[Formality: class~G, shadow depth~$2$;
\ClaimStatusProvedHere]
\label{thm:heisenberg-formality}
The Heisenberg ordered bar complex is formal:
the $\Ainf$-structure on $H^\bullet(\Barch(\cH_k))$
has $m_k = 0$ for all $k\ge 3$.  The Heisenberg
belongs to class~G \textup{(}shadow depth $r_{\max} = 2$,
the minimum for a nontrivial chiral algebra\textup{)}.
\end{theorem}

\begin{proof}
Since $d_{\Barch} = 0$, the bar complex \emph{is} its own
cohomology.  There is nothing to transfer: the $\Ainf$-structure
on cohomology inherits the trivial differential and the
deconcatenation coproduct, with all higher operations vanishing.
The vanishing of $m_{k\ge 3}$ is equivalent to saying that the
OPE has a single pole order (the double pole),
and the collision residue $r^{\mathrm{coll}} = k/z$ has a
simple pole, which falls in class~G.
\end{proof}

\begin{remark}[Contrast with the chiral \textup{(}symmetric\textup{)}
Koszul dual]
\label{rem:heisenberg-chiral-vs-ordered}
The chiral Koszul dual of $\cH_k$ (Vol~I) is the
commutative chiral algebra
\[
\cH_k^{!,\mathrm{ch}}
\;=\;
\operatorname{Sym}^{\mathrm{ch}}(V^*),
\]
a free commutative vertex algebra on one generator.
The ordered Koszul dual $Y(\mathfrak{u}(1))$ carries
strictly more structure: the spectral parameter~$u$
(arising from the FM stratification in the ordered
complex) is invisible in the symmetric Koszul dual.
The passage from ordered to symmetric is
$R$-matrix-twisted $\Sigma_n$-coinvariants: since
$R(z) = e^{k\hbar/z} \neq \tau$ (the Koszul sign),
the descent $\Barch(\cH_k)\to\barB^{\Sigma}(\cH_k)$
is genuinely twisted, and the spectral parameter
information is lost in the quotient (absorbed into
the bar differential of $\barB^{\Sigma}$, where
$d_{\barB}[\susp^{-1}J\,|\,\susp^{-1}J] = k \neq 0$).

This is the three-tier picture in its simplest instance:
the Heisenberg lives in tier~(ii)
($E_\infty$-chiral with $R(z)\neq\tau$, derived from
the local OPE).  Tier~(i) would require $R(z) = \tau$
(no OPE poles); tier~(iii) would require $R(z)$ to be
independent input (genuinely $E_1$).
\end{remark}

\subsection{The \texorpdfstring{$\beta\gamma$}{beta-gamma} ordered bar complex}
\label{subsec:betagamma-ordered-bar}

Let $A = \mathcal{F}_{\beta\gamma}$ denote the $\beta\gamma$ system:
a rank-one symplectic boson with generators $\beta, \gamma$ of conformal
weights $(1,0)$ and the single nontrivial OPE
\begin{equation}\label{eq:bg-ope}
\beta(z)\,\gamma(w) \;\sim\; \frac{1}{z-w}, \qquad
\gamma(z)\,\beta(w) \;\sim\; \frac{-1}{z-w}.
\end{equation}
The $\beta\gamma$ system is $\Einf$-chiral: it is a local vertex algebra
with $\Sigma_n$-equivariant factorization structure.  The OPE has a
simple pole, so this is a vertex algebra with poles (tier~(ii) in the
three-tier picture), not a pole-free commutative algebra.

\subsubsection{Collision residue and the nilpotent kernel}

The ordered bar complex uses the $d\log$ propagator.  By the
pole-absorption principle (AP19), the $d\log(z{-}w)$ kernel absorbs one
power of the OPE pole.  Since the $\beta\gamma$ OPE has a simple pole
$1/(z{-}w)$, the collision residue is the constant:
\begin{equation}\label{eq:bg-collision-residue}
r^{\mathrm{coll}} \;=\; \operatorname{Res}_{z=w}
  \frac{1}{z-w}\,d\log(z-w)
\;=\; 1.
\end{equation}
The collision residue kernel, as an element of $\operatorname{End}(A \otimes A)$,
is the \emph{nilpotent kernel}
\begin{equation}\label{eq:bg-theta}
\Theta \;:=\; \beta \otimes \gamma \;-\; \gamma \otimes \beta
\;\in\; A \otimes A,
\end{equation}
encoding the antisymmetric channel of the $\beta\gamma$ OPE.  The key
structural property is nilpotency:
\[
\Theta^2 \;=\;
  (\beta \otimes \gamma - \gamma \otimes \beta)^2
\;=\; \beta\gamma \otimes \gamma\beta
  - \beta\gamma \otimes \beta\gamma
  - \gamma\beta \otimes \gamma\beta
  + \gamma\beta \otimes \beta\gamma
\;=\; 0,
\]
where the vanishing follows from $\beta_{(0)}\gamma = 1$ (a scalar,
hence central) and $\gamma_{(0)}\beta = -1$.  Thus $\Theta^2 = 0$:
the collision residue is \emph{square-zero}.

\subsubsection{Generators and differential}

\begin{theorem}[$\beta\gamma$ ordered bar complex;
\ClaimStatusProvedHere]
\label{thm:bg-ordered-bar}
The ordered bar complex $\Barch(\mathcal{F}_{\beta\gamma})$ has the
following structure.

\emph{Degree~$1$.}  Two generators:
\[
\susp^{-1}\beta, \qquad \susp^{-1}\gamma.
\]
The bar differential vanishes: $d(\susp^{-1}\beta) = 0$,
$d(\susp^{-1}\gamma) = 0$.

\emph{Degree~$2$.}  Four generators:
\[
[\beta|\beta], \quad [\beta|\gamma], \quad
[\gamma|\beta], \quad [\gamma|\gamma].
\]
The bar differential acts by the zeroth product
$\beta_{(0)}\gamma = 1$:
\begin{align}
d[\beta|\gamma] &= \susp^{-1}(\beta_{(0)}\gamma)
  = \susp^{-1}(1) = 0, \label{eq:bg-d-bg} \\
d[\gamma|\beta] &= \susp^{-1}(\gamma_{(0)}\beta)
  = \susp^{-1}(-1) = 0, \label{eq:bg-d-gb}
\end{align}
where the output is zero because the unit $1 \in A$ is killed by the
augmentation (the bar complex is reduced).  The remaining generators
satisfy $d[\beta|\beta] = 0$ and $d[\gamma|\gamma] = 0$ since
$\beta_{(0)}\beta = 0 = \gamma_{(0)}\gamma$ (no self-OPE).

\emph{Higher degrees.}  In degree $p \geq 3$, the bar differential
sums over adjacent collisions.  The only nontrivial residues come from
adjacent $\beta\gamma$ or $\gamma\beta$ pairs, each contributing the
scalar $\pm 1$.  Since these scalars are killed by the augmentation in
the reduced complex, the bar complex is concentrated in low degrees
with acyclic higher components.
\end{theorem}

\begin{proof}
The OPE is $\beta(z)\gamma(w) \sim 1/(z-w)$, a simple pole.
After $d\log$ absorption, the collision residue is the constant $1$.
In the reduced bar complex, insertion of the unit $1 \in A$ is zero
(the augmentation ideal does not contain the vacuum).  Hence the bar
differential on $[\beta|\gamma]$ produces $\susp^{-1}(1) = 0$ in the
reduced complex.  The same argument applies at all higher degrees: every
nontrivial adjacent collision produces a scalar, which is annihilated
by the reduction.  The only surviving cohomology classes are the
degree-$1$ generators $\susp^{-1}\beta$ and $\susp^{-1}\gamma$.
\end{proof}

\subsubsection{Parallel transport and R-matrix}

\begin{theorem}[$\beta\gamma$ parallel transport and R-matrix;
\ClaimStatusProvedHere]
\label{thm:bg-rmatrix}
The parallel transport operator on $\operatorname{Conf}_2^{\mathrm{ord}}(\bC)$
for the $\beta\gamma$ system is:
\begin{equation}\label{eq:bg-transport}
T(u, u_0) \;=\; \id + \Theta \cdot \log\!\Bigl(
  \frac{u + x - y}{u_0 + x - y}\Bigr),
\end{equation}
where $\Theta = \beta \otimes \gamma - \gamma \otimes \beta$ is the
nilpotent kernel~\eqref{eq:bg-theta}, $u$ is the transport parameter,
and $x, y$ are the positions of the two points.  The series terminates
at linear order because $\Theta^2 = 0$.

The monodromy around a loop encircling $z_1 = z_2$ is
\begin{equation}\label{eq:bg-monodromy}
M \;=\; \id + 2\pi i\,\Theta,
\end{equation}
and the spectral R-matrix of the $\beta\gamma$ system is
\begin{equation}\label{eq:bg-rmatrix}
R \;=\; \id + 2\pi i\,\Theta.
\end{equation}
This is \emph{nilpotent}: $R = \id + 2\pi i\,\Theta$ with $\Theta^2=0$,
so $(R - \id)^2 = 0$.  The exponential series terminates at linear order.
\end{theorem}

\begin{proof}
The connection on $\operatorname{Conf}_2^{\mathrm{ord}}(\bC)$ is
$\nabla = d - \Theta \cdot d\log(z_1 - z_2)$.  Integrating:
\[
T(u, u_0)
= \mathcal{P}\exp\Bigl(-\int_{u_0}^{u}\Theta\,
  \frac{ds}{s + x - y}\Bigr)
= \exp\Bigl(\Theta\,\log\frac{u+x-y}{u_0+x-y}\Bigr).
\]
Since $\Theta^2 = 0$, the exponential terminates at the linear term:
$\exp(\Theta \cdot t) = \id + \Theta \cdot t$.  The monodromy is
$M = T$ evaluated around a loop encircling the collision locus;
the $d\log$ integral around a small loop contributes $2\pi i$:
\[
M = \exp(2\pi i\,\Theta) = \id + 2\pi i\,\Theta.
\qedhere
\]
\end{proof}

\begin{remark}[Contrast with Heisenberg: exponential vs.\ nilpotent]
\label{rem:bg-vs-heisenberg}
The Heisenberg R-matrix $R(z) = \exp(k\hbar/z)$
(Theorem~\ref{thm:heisenberg-rmatrix}) is an infinite power series.
The $\beta\gamma$ R-matrix $R = \id + 2\pi i\,\Theta$ terminates at
linear order because the collision residue is nilpotent.  This
dichotomy --- exponential vs.\ polynomial --- reflects the pole-order
classification: the Heisenberg OPE has a double pole with scalar
residue (class~G, shadow depth $r_{\max} = 2$), while the
$\beta\gamma$ OPE has a simple pole with nilpotent residue (class~C,
shadow depth $r_{\max} = 4$; the stratum-separation mechanism of
Volume~I, Theorem~\ref*{V1-thm:betagamma-global-depth}).
The two rank-one free theories thus occupy different positions in the
pole-order landscape despite both being free.
\end{remark}

\subsubsection{Koszul dual: the finite-dimensional exterior algebra}

\begin{theorem}[$\beta\gamma$ ordered Koszul dual;
\ClaimStatusProvedHere]
\label{thm:bg-koszul-dual}
The ordered Koszul dual of the $\beta\gamma$ system is a
\emph{finite-dimensional} exterior algebra:
\begin{equation}\label{eq:bg-koszul-dual}
\mathcal{F}_{\beta\gamma}^{!,\mathrm{ord}}
\;\simeq\; \bigwedge\nolimits^{\bullet}(\susp\,\beta^* \oplus \susp\,\gamma^*),
\end{equation}
a $4$-dimensional algebra \textup{(}in degrees $0,1,1,2$\textup{)} with
$b^* \wedge c^* = -c^* \wedge b^*$ and $(b^*)^2 = 0 = (c^*)^2$.  In
particular, $\mathcal{F}_{\beta\gamma}$ is of class~C with shadow depth
$r_{\max} = 4$ \textup{(}Volume~I,
Theorem~\textup{\ref*{V1-thm:betagamma-global-depth}}\textup{)}.
\end{theorem}

\begin{proof}
The bar complex has generators in degrees~$1$ and~$2$ with trivial
differential in the reduced complex (Theorem~\ref{thm:bg-ordered-bar}).
The linear dual of the bar coalgebra is the Koszul dual algebra.
Since the collision residue is nilpotent ($\Theta^2 = 0$), the dual
algebra is generated by $\susp\,\beta^*$ and $\susp\,\gamma^*$ with
the relation $(b^*)^2 = 0 = (c^*)^2$ and $b^* \wedge c^* = -c^* \wedge b^*$
(the antisymmetry of $\Theta$).  This is the exterior algebra
$\bigwedge^{\bullet}(\susp\,\beta^* \oplus \susp\,\gamma^*)$, which
is $4$-dimensional.

Shadow depth: the shadow depth $r_{\max} = 4$ is proved
in Volume~I (Theorem~\ref*{V1-thm:betagamma-global-depth}) via
the stratum-separation mechanism.  The T-line carries
$S_3 = 2 \neq 0$ and $S_4 \neq 0$, but the quintic
obstruction vanishes because the cubic and quartic shadows
live on different strata of the deformation complex.
\end{proof}

\begin{remark}[The $bc$ ghost system as Koszul dual]
\label{rem:bc-ghost}
The exterior algebra
$\bigwedge^{\bullet}(\susp\,\beta^* \oplus \susp\,\gamma^*)$
is the finite-dimensional shadow of the $bc$ ghost system.  This matches
the Koszul duality $(\beta\gamma)^! \cong \mathcal{F}_{bc}$ established
in the symmetric bar theory: the symplectic (bosonic) pairing of
$\beta\gamma$ dualizes to the antisymmetric (fermionic) pairing of $bc$
under Verdier duality on the bar complex.  The ordered bar complex makes
this duality completely explicit: the nilpotent kernel
$\Theta = \beta \otimes \gamma - \gamma \otimes \beta$ is the
\emph{antisymmetric} bilinear form that generates the exterior algebra.

By contrast, the Heisenberg ordered Koszul dual
$\cH_k^{!,\mathrm{ord}} \simeq Y(\mathfrak{u}(1))$
(Theorem~\ref{thm:heisenberg-yangian}) is infinite-dimensional:
the polynomial algebra on infinitely many generators.  The
$\beta\gamma$ system thus provides the simplest example of a
chiral algebra whose ordered Koszul dual is \emph{finite-dimensional},
a direct consequence of the nilpotency of $\Theta$.
\end{remark}

\subsubsection{The \texorpdfstring{$bc$}{bc} ghost ordered bar complex}

The $bc$ ghost system of conformal weights $(\lambda, 1{-}\lambda)$
is the \emph{fermionic} free-field system with OPE
\begin{equation}\label{eq:bc-ope-ordered}
b(z)\,c(w) \;\sim\; \frac{1}{z-w}\,.
\end{equation}
It is $\Einf$-chiral and Koszul dual to $\beta\gamma$:
\begin{equation}\label{eq:bc-betagamma-koszul-ordered}
bc_\lambda^{!,\mathrm{ch}} = \beta\gamma_\lambda,
\qquad
\beta\gamma_\lambda^{!,\mathrm{ch}} = bc_\lambda.
\end{equation}
This is the chiral incarnation of classical
$\bigwedge/\mathrm{Sym}$ duality: $bc$ is the chiral exterior
algebra $\bigwedge^{\mathrm{ch}}(W)$ on
$W = \bC b \oplus \bC c$ (fermionic), and its (chiral,
symmetric) Koszul dual is the chiral symmetric algebra
$\mathrm{Sym}^{\mathrm{ch}}(W^*)$
(bosonic), which is precisely the $\beta\gamma$ system.

\begin{theorem}[The $bc$ ordered bar complex;
\ClaimStatusProvedHere]
\label{thm:bc-ordered-bar}
The ordered bar complex of the $bc$ ghost system has the same
structural features as the $\beta\gamma$ complex:
\begin{enumerate}[label=\textup{(\arabic*)}]
\item The bar differential is
nontrivial: $b_{(0)}c = 1 \neq 0$, $c_{(0)}b = -1$, and
$b_{(0)}b = c_{(0)}c = 0$.

\item The collision residue kernel is the nilpotent element
\begin{equation}\label{eq:bc-theta}
\Theta_{bc}
\;:=\;
b\otimes c \;-\; c\otimes b
\;\in\; W\otimes W,
\end{equation}
with $\Theta_{bc}^2 = 0$ by the same argument as for
$\Theta_{\beta\gamma}$.

\item The Steinberg connection is
$\nabla_{\mathfrak{S}} = d + (\Theta_{bc}/z)\,dz/z$,
regular singular with monodromy
\begin{equation}\label{eq:bc-monodromy}
\operatorname{Mon}_0(\nabla_{\mathfrak{S}})
\;=\;
\id + 2\pi i\,\Theta_{bc}.
\end{equation}

\item The open-colour Koszul dual exchanges statistics:
\begin{equation}\label{eq:bc-line-dual}
\mathcal{F}_{bc}^{!,\mathrm{ord}}
\;\simeq\;
\mathrm{Sym}^{\bullet}\!\bigl(\susp\,b^* \oplus \susp\,c^*\bigr),
\end{equation}
a polynomial algebra.  The fermionic $bc$ has a bosonic
\textup{(}symmetric\textup{)} ordered dual, just as the
chiral duality $bc^! = \beta\gamma$ exchanges
$\bigwedge \leftrightarrow \mathrm{Sym}$.

\item Shadow depth: $r_{\max} = 4$ \textup{(}class~C\textup{)},
the same as $\beta\gamma$.  The complementarity constant
vanishes: $K(bc_\lambda) = c_{bc}(\lambda) + c_{\beta\gamma}(\lambda)
= 0$ for all $\lambda$.
\end{enumerate}
\end{theorem}

\begin{proof}
The computation is parallel to the $\beta\gamma$ case.
The OPE $b(z)c(w) \sim 1/(z{-}w)$ gives
$b_{(0)}c = 1$, $c_{(0)}b = -1$ (skew-symmetry),
$b_{(0)}b = c_{(0)}c = 0$.  The collision residue
$\Theta_{bc} = b\otimes c - c\otimes b$ and its nilpotency
follow by the same argument as
$\Theta_{\beta\gamma}$.  The Steinberg
connection and monodromy are identical in structure.
For the statistics exchange: the fermionic system has bar
cohomology yielding symmetric (polynomial) generators rather
than exterior ones, because the Koszul complex of the
zeroth-mode relations in the fermionic sector produces
$\mathrm{Sym}$ rather than $\bigwedge$.  Shadow depth:
Volume~I, by the same stratum-separation mechanism as
$\beta\gamma$.  Complementarity: $K = 0$ from
$c_{bc}(\lambda) = -2(6\lambda^2{-}6\lambda{+}1)$ and
$c_{\beta\gamma}(\lambda) = 2(6\lambda^2{-}6\lambda{+}1)$.
\end{proof}

\subsubsection{The Wakimoto connection: free-field realisation of
  interacting Yangians}

The Wakimoto embedding realises affine Kac--Moody algebras
inside free-field tensor products.  At the level of ordered
bar complexes, this provides a free-field description of
the dg Yangian.

\begin{theorem}[Wakimoto bar complex descent;
\ClaimStatusProvedHere]
\label{thm:wakimoto-ordered-bar}
Let $\fg$ be a finite-dimensional simple Lie algebra of
rank~$r$ and let $\widehat{\fg}_k$ be the associated affine
Kac--Moody algebra at non-critical level
$k \neq -h^\vee$.  The Wakimoto module decomposes as
\[
\mathcal{M}_{\mathrm{Wak}}
\;=\;
\bigotimes_{\alpha\in\Delta_+}
(\beta\gamma)_\alpha
\;\otimes\;
\cH_{\mathfrak{h}},
\]
a tensor product of $|\Delta_+|$ copies of $\beta\gamma$
systems \textup{(}one per positive root\textup{)} and a
rank-$r$ Heisenberg system on the Cartan.  The ordered bar
complex of this tensor product carries a BRST differential
$Q_{\mathrm{DS}}$ \textup{(}the Drinfeld--Sokolov
reduction\textup{)}, and
\begin{equation}\label{eq:wakimoto-ordered-bar-descent}
H^\bullet\bigl(
  \Barch(\mathcal{M}_{\mathrm{Wak}}),\;
  d_{\Barch} + Q_{\mathrm{DS}}
\bigr)
\;\simeq\;
\Barch(\widehat{\fg}_k),
\end{equation}
recovering the interacting ordered bar complex of
$\widehat{\fg}_k$ as the BRST cohomology of the free-field
ordered bar complex.
At the level of Koszul duals, the dg Yangian
$Y^{\mathrm{dg}}(\fg)$ is the BRST reduction of the
free-field Yangian:
\begin{equation}\label{eq:yangian-wakimoto}
Y^{\mathrm{dg}}_\hbar(\fg)
\;\simeq\;
H^\bullet\!\left(
  \bigotimes_{\alpha\in\Delta_+}
  \mathcal{F}_{\beta\gamma,\alpha}^{!,\mathrm{ord}}
  \;\otimes\;
  Y(\mathfrak{u}(1))^{\otimes r},\;
  Q_{\mathrm{DS}}^!
\right)\!.
\end{equation}
The DS differential $Q_{\mathrm{DS}}$ on the free-field
Yangian produces the interacting Yangian.
\end{theorem}

\begin{proof}
The Wakimoto decomposition is standard
(Wakimoto~1986, Feigin--Frenkel~1990).
The BRST charge $Q_{\mathrm{DS}}$ commutes with the bar
differential ($Q_{\mathrm{DS}}^2 = 0$ and
$[d_{\Barch}, Q_{\mathrm{DS}}] = 0$), so a spectral
sequence with $E_1 = H^\bullet(\Barch(\mathcal{M}_{\mathrm{Wak}}),
d_{\Barch})$ and $d_1 = Q_{\mathrm{DS}}$ converges to
$\Barch(\widehat{\fg}_k)$ (Volume~I,
Theorem~\ref*{V1-thm:wakimoto-bar}).  Applying linear
duality and the Koszul dual functor to both sides gives
the Yangian-level
statement~\eqref{eq:yangian-wakimoto}.
Theorem~\ref{thm:ordered-associative-ds-principal}
(principal DS commutes with ordered chiral duality) ensures
the Koszul dual functor and DS reduction commute in the
principal case, which covers the Wakimoto
embedding.\qedhere
\end{proof}

\subsubsection{Contrast table: shadow depth and $R$-matrix across
  the standard families}

\begin{table}[ht]
\centering
\renewcommand{\arraystretch}{1.4}
\begin{tabular}{llccll}
\toprule
\textbf{System} & \textbf{$R$-matrix}
  & \textbf{Depth} & \textbf{Class}
  & \textbf{$A^{!,\mathrm{ord}}_{\mathrm{line}}$}
  & \textbf{$\Einf$}\\
\midrule
$\cH_k$ & $e^{k\hbar/z}$ \textup{(scalar)}
  & $2$ & $\mathbf{G}$ & $\bC[h_0,h_1,\ldots]$
  & $\checkmark$ \\[2pt]
$\beta\gamma$ & $\id + 2\pi i\,\Theta$
  \textup{(nilpotent)}
  & $4$ & $\mathbf{C}$
  & $\bigwedge^*(\beta^*,\gamma^*)$
  & $\checkmark$ \\[2pt]
$bc$ & $\id + 2\pi i\,\Theta_{bc}$
  \textup{(nilpotent)}
  & $4$ & $\mathbf{C}$
  & $\mathrm{Sym}^*(b^*,c^*)$
  & $\checkmark$ \\[2pt]
$\widehat{\fg}_k$ & $\exp(\hbar\,\Omega/z)$
  & $3$ & $\mathbf{L}$ & $Y_\hbar(\fg)$
  & $\checkmark$ \\[2pt]
$\mathrm{Vir}_c$ & \textup{[triple pole]}
  & $\infty$ & $\mathbf{M}$ & $Y^{\mathrm{dg}}(\mathrm{Vir})$
  & $\checkmark$ \\
\bottomrule
\end{tabular}
\caption{Ordered bar complex data for the standard families.
All entries are $\Einf$-chiral.  The $R$-matrix is derived
from the local OPE in every case.  Shadow depth measures
complexity of the $\Ainf$ structure within the Koszul world.
The free-field systems $\beta\gamma$ and $bc$ have nilpotent
$R$-matrices \textup{(}the exponential terminates at linear
order\textup{)}, while the Heisenberg and affine families have
genuine exponentials.}
\label{tab:free-field-contrast}
\end{table}

\begin{remark}[The Wakimoto escalation]
\label{rem:wakimoto-escalation}
The contrast table exhibits the Wakimoto escalation pattern.
The free-field atoms $\cH_k$ (class~G, depth~$2$) and
$\beta\gamma$ (class~C, depth~$4$) are the elementary
building blocks.  Composing them via the Wakimoto embedding
and applying DS reduction produces the interacting families:
$\widehat{\fg}_k$ (class~L, depth~$3$) and, after a second
DS reduction, $\mathrm{Vir}_c$ (class~M, depth~$\infty$).
The shadow depth escalates through the chain
\[
\underbrace{\cH_k \;\otimes\; \beta\gamma^{\otimes |\Delta_+|}}%
_{\text{free-field: depths $2$ and $4$}}
\;\xrightarrow{\;Q_{\mathrm{DS}}\;}
\underbrace{\widehat{\fg}_k}_{\text{depth $3$}}
\;\xrightarrow{\;Q_{\mathrm{DS}}\;}
\underbrace{\mathrm{Vir}_c}_{\text{depth $\infty$}}.
\]
The second reduction (affine $\to$ Virasoro via Sugawara)
escalates from finite to infinite depth: the nonlinearity
of the Sugawara construction manufactures the quartic pole
in the Virasoro OPE from the double poles of the affine
OPE, and this quartic pole generates the infinite shadow
tower.  This is the
pole-order dichotomy of Volume~II at the level of ordered
bar complexes: DS reduction transports class~L (finite
depth, formal Swiss-cheese) to class~M (infinite depth,
non-formal Swiss-cheese).
\end{remark}

\subsection{Lattice vertex algebras}
\label{subsec:lattice-ordered-bar}

Let $\Lambda$ be an even lattice of rank $d$ with bilinear form
$\langle\cdot,\cdot\rangle$ and let $\varepsilon \colon \Lambda \times
\Lambda \to \bC^\times$ be a normalized bimultiplicative $2$-cocycle.
The lattice vertex algebra $\Vlat_\Lambda^\varepsilon$ has vertex operators
$Y(e^\alpha, z)$ for $\alpha \in \Lambda$ with OPE
\begin{equation}\label{eq:lattice-ope-ordered}
Y(e^\alpha, z)\,Y(e^\beta, w) \;\sim\;
  \varepsilon(\alpha,\beta)\,(z - w)^{\langle\alpha,\beta\rangle}\,
  Y(e^{\alpha+\beta}, w) + \cdots
\end{equation}
The $\Eone$/$\Einf$ character of $\Vlat_\Lambda^\varepsilon$ depends on
the symmetry of the cocycle, as established in
\S\ref{sec:lattice:e1-families} of the lattice foundations chapter.

\subsubsection{Symmetric cocycle: \texorpdfstring{$\Einf$}{E-infinity}-chiral}

\begin{theorem}[Ordered bar complex with symmetric cocycle;
\ClaimStatusProvedHere]
\label{thm:lattice-symmetric-ordered-bar}
Let $\varepsilon$ be a symmetric cocycle satisfying the Borcherds
symmetry $\varepsilon(\alpha,\beta) = (-1)^{\langle\alpha,\beta\rangle}
\varepsilon(\beta,\alpha)$.  Then $\Vlat_\Lambda^\varepsilon$ is
$\Einf$-chiral, and the ordered bar complex
$\Barch(\Vlat_\Lambda^\varepsilon)$ has:

\begin{enumerate}[label=\textup{(\arabic*)}]
\item \textbf{Collision residue.}
The collision residue kernel on the diagonal $\alpha$-sector is:
\begin{equation}\label{eq:lattice-symmetric-residue}
r(z) \;=\; \frac{1}{z}\sum_{\alpha \in \Lambda}\,
  \langle\alpha,\alpha\rangle\,
  (e^\alpha \otimes e^{-\alpha}).
\end{equation}
This is the \emph{diagonal} residue: only the Heisenberg-type
contributions $\alpha \otimes (-\alpha)$ survive, weighted by the norm
$\langle\alpha,\alpha\rangle$.

\item \textbf{R-matrix.}
Since $\varepsilon$ is symmetric, the R-matrix is
determined by the local OPE \textup{(}tier~(ii) of the three-tier
picture\textup{)}.
The ordered-to-unordered descent
$\Barch \to \barBgeom$
is the R-matrix-twisted $\Sigma_n$-quotient, with twist
$R(z)$ derived from the OPE monodromy:
\begin{equation}\label{eq:lattice-symmetric-rmatrix}
R_{\alpha,\beta}(z) \;=\;
  (-1)^{\langle\alpha,\beta\rangle}\,
  e^{\pi i\langle\alpha,\beta\rangle},
\end{equation}
the standard lattice braiding phase.  This is \emph{not} independent
input --- it is the monodromy of the flat connection on
$\operatorname{Conf}_2(\bC)$ determined by the OPE poles,
confirming the $\Einf$ character.
\end{enumerate}
\end{theorem}

\begin{proof}
The OPE $Y(e^\alpha,z)\,Y(e^\beta,w) \sim
\varepsilon(\alpha,\beta)(z-w)^{\langle\alpha,\beta\rangle}
Y(e^{\alpha+\beta},w) + \cdots$ has a pole of order
$-\langle\alpha,\beta\rangle$ at $z = w$.  After $d\log$ absorption,
the collision residue extracts the order-$1$ coefficient:
$\operatorname{Res}_{z=w}[(z-w)^{-m}\,d\log(z-w)] = \delta_{m,1}$.
For the diagonal sector $\beta = -\alpha$, the pole order is
$-\langle\alpha,-\alpha\rangle = \langle\alpha,\alpha\rangle$, which
is always even for an even lattice.  The non-diagonal residue at
order~$1$ gives~\eqref{eq:lattice-symmetric-residue}.

For the R-matrix: the connection is
$\nabla = d - r(z)\,d\log(z_1 - z_2)$, and the monodromy around
$z_1 = z_2$ is $\exp(2\pi i \cdot r)$.  On the lattice sector
$(e^\alpha, e^\beta)$, this is
$(-1)^{\langle\alpha,\beta\rangle}\,
e^{\pi i \langle\alpha,\beta\rangle}$, matching the standard
lattice braiding.  The Borcherds symmetry of $\varepsilon$ ensures
this monodromy is \emph{derived from} the local OPE, confirming
$\Einf$.
\end{proof}

\subsubsection{Non-symmetric cocycle: genuinely \texorpdfstring{$\Eone$}{E1}-chiral}

\begin{theorem}[Ordered bar complex with non-symmetric cocycle;
\ClaimStatusProvedHere]
\label{thm:lattice-nonsymmetric-ordered-bar}
Let $\varepsilon$ be a non-symmetric cocycle: there exist
$\alpha, \beta \in \Lambda$ with
$\varepsilon(\alpha,\beta)/\varepsilon(\beta,\alpha) \neq
(-1)^{\langle\alpha,\beta\rangle}$.  Then
$\Vlat_\Lambda^\varepsilon$ is genuinely $\Eone$-chiral
\textup{(}nonlocal\textup{)},
and the ordered bar complex $\Barch(\Vlat_\Lambda^\varepsilon)$ has:

\begin{enumerate}[label=\textup{(\arabic*)}]
\item \textbf{Differential.}  In bar degree~$2$:
\begin{equation}\label{eq:lattice-nonsym-diff}
d([e^{\alpha_1} \mid e^{\alpha_2}])
\;=\; \varepsilon(\alpha_1,\alpha_2)\,
  \delta_{\langle\alpha_1,\alpha_2\rangle,\,-1}\,
  [e^{\alpha_1 + \alpha_2}].
\end{equation}
The cocycle $\varepsilon(\alpha_1,\alpha_2)$ appears with its value at
the \emph{specific ordered pair} $(\alpha_1, \alpha_2)$.  For a
non-symmetric cocycle, the coefficients
$\varepsilon(\alpha,\beta) \neq
(-1)^{\langle\alpha,\beta\rangle}\varepsilon(\beta,\alpha)$, so
permuting the pair changes the differential.  No $\Sigma_2$-quotient is
possible.

\item \textbf{R-matrix as independent input.}
The R-matrix $R_{\alpha,\beta}$ is part of the defining data of the
$\Eone$-chiral algebra, not derivable from a local OPE.  The ratio
\begin{equation}\label{eq:lattice-nonsym-rmatrix-ratio}
\frac{\varepsilon(\alpha,\beta)}{\varepsilon(\beta,\alpha)}
\end{equation}
measures the asymmetry of the cocycle and enters the R-matrix directly.
When this ratio is not $(-1)^{\langle\alpha,\beta\rangle}$, the
braiding is genuinely nonlocal: it cannot be obtained as the monodromy
of any connection on $\operatorname{Conf}_2(\bC)$ determined by a local
OPE.  This is tier~(iii) of the three-tier picture.
\end{enumerate}
\end{theorem}

\begin{proof}
The differential follows from the
lattice OPE~\eqref{eq:lattice-ope-ordered} and the ordered bar
construction: adjacent collisions in the ordered bar complex use the
$d\log$ kernel, selecting simple poles
($\langle\alpha_i,\alpha_{i+1}\rangle = -1$) with coefficient
$\varepsilon(\alpha_i, \alpha_{i+1})$ at the specific ordered pair.

For the R-matrix claim: the Borcherds symmetry
$\varepsilon(\alpha,\beta) =
(-1)^{\langle\alpha,\beta\rangle}\varepsilon(\beta,\alpha)$
is precisely the condition for the $\Sigma_2$-action on the ordered bar
complex to be compatible with the differential.  When this fails, the
descent $\Barch \to \barBgeom$ requires independent braiding data, and
the algebra is genuinely $\Eone$.
\end{proof}

\subsubsection{Koszul dual of lattice algebras}

\begin{theorem}[Ordered Koszul dual of lattice algebras;
\ClaimStatusProvedHere]
\label{thm:lattice-ordered-koszul-dual}
The ordered Koszul dual of $\Vlat_\Lambda^\varepsilon$ depends on
the lattice:
\begin{enumerate}[label=\textup{(\roman*)}]
\item For a self-dual even lattice $\Lambda = \Lambda^*$
  with symmetric cocycle, the Koszul dual
  $(\Vlat_\Lambda^\varepsilon)^!$ is again a lattice vertex algebra
  $(\Vlat_\Lambda^{\varepsilon^{-1}})^c$ with inverse cocycle
  \textup{(}Theorem~\textup{\ref{thm:lattice:koszul-dual}}\textup{)}.
\item For a non-self-dual lattice or non-symmetric cocycle, the
  ordered Koszul dual carries strictly more structure: the
  ordered bar complex has independent generators in each ordering,
  and the dual algebra reflects this asymmetry.
\item In all cases, the lattice grading is preserved: the Koszul dual
  decomposes as $\bigoplus_{\gamma \in \Lambda}
  (\Vlat_\Lambda^\varepsilon)^!_\gamma$ with the differential
  preserving total lattice degree.
\end{enumerate}
\end{theorem}

\begin{proof}
Part~(i) follows from Theorem~\ref{thm:lattice:koszul-dual} and its
proof: the bar complex coproduct has structure constants
$\varepsilon(\alpha,\beta)^{-1}$, and linear duality recovers the
lattice algebra with inverse cocycle.  Self-duality $\Lambda = \Lambda^*$
ensures the dual lives on the same lattice.

Part~(ii): for a non-symmetric cocycle, the ordered bar complex in
degree~$2$ has generators $[e^\alpha \mid e^\beta]$ and
$[e^\beta \mid e^\alpha]$ that are \emph{independent} (no
$\Sigma_2$-identification).  Their duals generate the Koszul dual
algebra with product reflecting the asymmetric cocycle values.

Part~(iii): the bar differential preserves the total lattice degree
$\gamma = \sum_i \alpha_i \in \Lambda$
(Theorem~\ref{thm:lattice:bar-structure}(i)), so the Koszul dual
inherits the same decomposition.
\end{proof}

\begin{remark}[The three-tier picture for lattice algebras]
\label{rem:lattice-three-tier}
Lattice vertex algebras provide the cleanest illustration of the
three-tier picture:
\begin{enumerate}[label=(\roman*)]
\item \emph{Pole-free commutative} ($R = \tau$): degenerate case where
  $\langle\alpha,\beta\rangle = 0$ for all $\alpha, \beta$ (the lattice
  is trivial or the restriction to a sublattice with vanishing pairing).
\item \emph{$\Einf$ with poles} ($R \neq \tau$, derived from locality):
  symmetric cocycle, OPE with poles from
  $\langle\alpha,\beta\rangle < 0$.  The R-matrix is the standard lattice
  braiding, determined by the OPE monodromy.
\item \emph{Genuinely $\Eone$} ($R \neq \tau$, independent input):
  non-symmetric cocycle.  The R-matrix is part of the defining data.
  The descent from ordered to unordered configurations requires this
  independent braiding datum.
\end{enumerate}
The lattice family thus realizes all three tiers explicitly, with the
cocycle symmetry as the single discrete parameter controlling the
$\Einf$/$\Eone$ transition.
\end{remark}


\subsection{The Drinfeld presentation of $Y_\hbar(\mathfrak{sl}_2)$}
\label{subsec:drinfeld-yangian-sl2}

The ordered bar complex $\Barch(V_k(\mathfrak{sl}_2))$ produces the
Yang $R$-matrix, the RTT presentation, and the Gauss decomposition.
The \emph{Drinfeld presentation} --- the generating-current
form that reveals the Yangian as a deformation of the current algebra
$U(\mathfrak{sl}_2[t])$ --- emerges from the ordered bar data.  The affine vertex algebra
$V_k(\mathfrak{sl}_2)$ is $\Einf$-chiral (it is a local vertex algebra
with $\Sigma_n$-equivariant factorisation structure); the Yangian
$Y_\hbar(\mathfrak{sl}_2)$ is the ordered Koszul dual living on the
open boundary $\bR_t$.

\subsubsection*{Generators}

The Drinfeld generators are obtained from the Gauss decomposition of the
transfer matrix $T(u) = F(u)\,K(u)\,E(u)$, where $F(u)$ is
lower-triangular unipotent, $K(u) = \operatorname{diag}(k_1(u),k_2(u))$
is diagonal, and $E(u)$ is upper-triangular unipotent.  Define the
Drinfeld currents:
\begin{align}
\label{eq:drinfeld-E-current}
E(u) &= \sum_{r \ge 0} E_r\,u^{-r-1}, \\
\label{eq:drinfeld-F-current}
F(u) &= \sum_{r \ge 0} F_r\,u^{-r-1}, \\
\label{eq:drinfeld-H-current}
H(u) &= 1 + \hbar\sum_{r \ge 0} H_r\,u^{-r-1}
\;=\; \frac{k_1(u)}{k_2(u)}.
\end{align}
The modes $\{E_r,\,F_r,\,H_r\}_{r \ge 0}$ generate $Y_\hbar(\mathfrak{sl}_2)$
as an associative algebra.  At $\hbar = 0$, the identification
$E_r \mapsto e \otimes t^r$, $F_r \mapsto f \otimes t^r$,
$H_r \mapsto h \otimes t^r$ recovers $U(\mathfrak{sl}_2[t])$.

\subsubsection*{The Drinfeld relations from the bar complex}

\begin{theorem}[Drinfeld presentation of $Y_\hbar(\mathfrak{sl}_2)$]
\ClaimStatusProvedHere
\label{thm:drinfeld-yangian-sl2}
The Yangian $Y_\hbar(\mathfrak{sl}_2)$, with $\hbar = 1/(k+2)$, is the
associative algebra generated by $\{E_r,\,F_r,\,H_r\}_{r \ge 0}$ subject
to the relations:
\begin{enumerate}[label=\textup{(D\arabic*)},leftmargin=3.5em]
\item\label{D1} $[H_r, H_s] = 0$ for all $r,s \ge 0$.

\item\label{D2}
$[H(u), E(v)] \;=\; \dfrac{2\hbar}{u-v}\bigl(E(u) - E(v)\bigr)$.

\item\label{D3} $[H(u), F(v)] \;=\; -\,\dfrac{2\hbar}{u-v}\bigl(F(u) - F(v)\bigr)$.

\item\label{D4} $[E(u), F(v)] \;=\; \dfrac{1}{u-v}\bigl(H(u) - H(v)\bigr)$.

\item\label{D5} $(u{-}v)\,[E(u), E(v)]
\;=\; \hbar\bigl(E(u)\,E(v) + E(v)\,E(u)\bigr)$
\textup{(}deformed Serre relation\textup{)}.

\item\label{D6} $(u{-}v)\,[F(u), F(v)]
\;=\; -\,\hbar\bigl(F(u)\,F(v) + F(v)\,F(u)\bigr)$
\textup{(}deformed Serre relation\textup{)}.
\end{enumerate}
\end{theorem}

\begin{proof}
Relations~\ref{D1}--\ref{D4} are extracted from the Gauss
decomposition of the RTT equation
$R_{12}(u{-}v)\,T_1(u)\,T_2(v) = T_2(v)\,T_1(u)\,R_{12}(u{-}v)$.
The key step is the factorisation $T(u) = F(u)\,K(u)\,E(u)$: the
diagonal entries produce~\ref{D1}, the off-diagonal--diagonal
cross-terms produce~\ref{D2}--\ref{D3}, and the purely off-diagonal
cross-terms produce~\ref{D4}.

For \ref{D1}: the RTT relations $[A(u),A(v)] = 0 = [D(u),D(v)]$
imply that $k_1(u),k_2(u)$ (and hence $H(u) = k_1(u)/k_2(u)$)
generate a commutative subalgebra.

For \ref{D2}: the RTT relation
$(u{-}v)[A(u),B(v)] = A(v)B(u) - A(u)B(v)$, rewritten via Gauss
decomposition in terms of $H(u)$ and $E(u)$, yields the stated
current-algebra relation.  Expanding
$A(u) = k_1(u)(1 + F(u)\,k_2(u)\,k_1(u)^{-1}\,E(u))$ and
$B(u) = k_1(u)\,E(u)$,
the diagonal factor cancels between left and right sides, and the
residual identity in the off-diagonal is precisely~\ref{D2}.
Relation~\ref{D3} follows by the same argument applied to
$(u{-}v)[A(u),C(v)] = C(v)A(u) - C(u)A(v)$.

For \ref{D4}: the RTT relation
$(u{-}v)[B(u),C(v)] = D(v)A(u) - D(u)A(v)$
extracts, via Gauss decomposition, to
$(u{-}v)[E(u),F(v)] = H(u) - H(v)$.

For \ref{D5} and \ref{D6}: the RTT relation $[B(u),B(v)] = 0$
encodes $[E(u),E(v)]$ modulo diagonal factors.  The
precise computation: $[B(u),B(v)] = 0$ with $B(u) = k_1(u)\,E(u)$
gives, after extracting the non-commuting diagonal,
$(u{-}v)[E(u),E(v)] = \hbar(E(u)E(v) + E(v)E(u))$.
The sign in~\ref{D6} follows from $[C(u),C(v)] = 0$ by the same
argument with $F$ in place of $E$.

Every relation is a direct algebraic consequence of the RTT
presentation, itself extracted from the degree-$2$ bar differential
and the Yang $R$-matrix $R(u) = u\,I + P$.
\end{proof}

\begin{remark}[Provenance from the bar complex]
\label{rem:drinfeld-bar-provenance}
The chain of identifications is:
\[
\text{OPE of } \widehat{\mathfrak{sl}}_2
\;\xrightarrow{\;\text{bar differential}\;}\;
r(z) = \frac{\hbar\,\Omega}{z}
\;\xrightarrow{\;\text{exponentiation}\;}\;
R(u) = u\,I + P
\;\xrightarrow{\;\text{RTT}\;}\;
\text{(D1)--(D6)}.
\]
The first arrow is AP19 (the $d\log$ kernel absorbs one pole from the
affine double-pole OPE); the second is the monodromy of the Kohno
connection on $\mathrm{Conf}_2^{\mathrm{ord}}(\bC)$ specialised to the
fundamental representation; the third is the Gauss decomposition.  The
entire Drinfeld presentation is a derived algebraic consequence of the
OPE of $V_k(\mathfrak{sl}_2)$.
\end{remark}

\subsubsection*{Gauss decomposition}

\begin{proposition}[Gauss decomposition]
\ClaimStatusProvedHere
\label{prop:gauss-decomposition-sl2}
The transfer matrix $T(u) = \bigl(\begin{smallmatrix}
A(u) & B(u) \\ C(u) & D(u) \end{smallmatrix}\bigr)$
factors uniquely as
\[
T(u) \;=\;
\underbrace{\begin{pmatrix} 1 & 0 \\ F(u) & 1 \end{pmatrix}}_{F(u)}
\;\cdot\;
\underbrace{\begin{pmatrix} k_1(u) & 0 \\ 0 & k_2(u) \end{pmatrix}}_{K(u)}
\;\cdot\;
\underbrace{\begin{pmatrix} 1 & E(u) \\ 0 & 1 \end{pmatrix}}_{E(u)},
\]
with
$A(u) = k_1(u) + F(u)\,k_2(u)\,E(u)$,
$B(u) = k_1(u)\,E(u)$,
$C(u) = F(u)\,k_2(u)$,
$D(u) = k_2(u)(1 + F(u)\,E(u))$.
The Drinfeld Cartan current is the ratio
$H(u) = k_1(u)/k_2(u)$.
\end{proposition}

\begin{proof}
The upper-triangular unipotent factor is determined by
$E(u) = k_1(u)^{-1}\,B(u)$; the lower-triangular unipotent factor by
$F(u) = C(u)\,k_2(u)^{-1}$; the diagonal entries by
$k_2(u) = D(u) - C(u)\,A(u)^{-1}\,B(u)$ (the Schur complement) and
$k_1(u) = A(u) - B(u)\,D(u)^{-1}\,C(u)$.  The quantum determinant
$\operatorname{qdet} T(u) = k_1(u)\,k_2(u{-}\hbar)$ is central in
$Y_\hbar(\mathfrak{sl}_2)$, confirming the factorisation is well-defined
over the Yangian.
\end{proof}

\subsubsection*{Twisted coproduct on Drinfeld generators}

\begin{theorem}[Twisted coproduct]
\ClaimStatusProvedHere
\label{thm:twisted-coproduct-sl2}
The twisted coproduct $\Delta_z \colon Y_\hbar(\mathfrak{sl}_2)
\to Y_\hbar(\mathfrak{sl}_2) \otimes Y_\hbar(\mathfrak{sl}_2)$
acts on the Drinfeld generators as:
\begin{align}
\label{eq:twisted-coprod-H}
\Delta_z(H(u)) &= H(u) \otimes H(u{-}z), \\
\label{eq:twisted-coprod-E}
\Delta_z(E(u)) &= E(u) \otimes 1 + H(u{-}z) \otimes E(u{-}z), \\
\label{eq:twisted-coprod-F}
\Delta_z(F(u)) &= 1 \otimes F(u{-}z) + F(u) \otimes H(u{-}z)^{-1}.
\end{align}
At $z = 0$ this recovers the standard Hopf coproduct of $Y_\hbar(\mathfrak{sl}_2)$.
\end{theorem}

\begin{proof}
The RTT coproduct is $\Delta_z(T(u)) = T(u) \dot\otimes T(u{-}z)$
(matrix multiplication in the auxiliary space, with spectral parameter
shifted by $z$ in the second tensor factor).  Applying Gauss
decomposition to both sides:
\[
\Delta_z(T(u))
= F(u)\,K(u)\,E(u) \;\dot\otimes\; F(u{-}z)\,K(u{-}z)\,E(u{-}z).
\]
The diagonal entry gives
$\Delta_z(K(u)) = K(u) \otimes K(u{-}z)$,
whence~\eqref{eq:twisted-coprod-H} follows by taking the ratio
$k_1/k_2$.  The off-diagonal entry
$\Delta_z(B(u)) = A(u) \otimes B(u{-}z) + B(u) \otimes D(u{-}z)$,
rewritten via $B = k_1\,E$ and $A = k_1 + F\,k_2\,E$,
yields~\eqref{eq:twisted-coprod-E} after extracting the $k_1$
factor.  Equation~\eqref{eq:twisted-coprod-F} follows by the analogous
computation on $C(u) = F(u)\,k_2(u)$.
\end{proof}

\begin{remark}[Factorisation-algebraic origin of $\Delta_z$]
\label{rem:twisted-coproduct-origin}
The spectral parameter $z$ is the relative position along $\bR_t$:
it is the topological interval coordinate of the
$\mathrm{SC}^{\mathrm{ch,top}}$-structure.  The twisted coproduct
$\Delta_z$ is the factorisation product of line operators at positions
$t_1$ and $t_2 = t_1 + z$ on the topological interval.  The formula
$\Delta_z(T(u)) = T(u) \dot\otimes T(u{-}z)$ is the bar-coalgebra
deconcatenation specialised to degree~$2$ in the ordered bar
complex.  At $z = 0$ the two insertion points collide and one recovers
the ordinary (untwisted) Hopf coproduct.
\end{remark}

\subsubsection*{PBW basis and Hilbert series}

\begin{theorem}[PBW theorem for $Y_\hbar(\mathfrak{sl}_2)$]
\ClaimStatusProvedHere
\label{thm:PBW-yangian-sl2}
The ordered monomials
\[
\bigl\{\,
f_{r_1}\cdots f_{r_a}\;\;
h_{s_1}\cdots h_{s_b}\;\;
e_{t_1}\cdots e_{t_c}
\;\bigm|\;
0 \le r_1 \le \cdots \le r_a,\;
0 \le s_1 \le \cdots \le s_b,\;
0 \le t_1 \le \cdots \le t_c
\,\bigr\}
\]
form a basis of $Y_\hbar(\mathfrak{sl}_2)$ over~$\bC$.  The Hilbert series
with respect to the grading $\deg(X_r) = r$ is
\[
\operatorname{HS}_{Y_\hbar(\mathfrak{sl}_2)}(q)
\;=\; \prod_{d \ge 0} \frac{1}{(1-q^d)^3}.
\]
\end{theorem}

\begin{proof}
The PBW property follows from the Drinfeld relations~\ref{D1}--\ref{D6}
by a standard diamond lemma argument: the leading terms of
\ref{D2}--\ref{D6} (in the degree-lexicographic order
$f_{r_1}\cdots h_{s_1}\cdots e_{t_1}\cdots$) show that every
out-of-order monomial reduces to a sum of ordered monomials of
equal or lower total degree.  The three families of generators
$\{e_r\}_{r \ge 0}$, $\{f_r\}_{r \ge 0}$, $\{h_r\}_{r \ge 0}$
each contribute a factor of $\prod_{d \ge 0}(1-q^d)^{-1}$ to the
Hilbert series, giving the stated product formula.  The PBW
basis is compatible with the filtration
$F_{\le n}\,Y_\hbar(\mathfrak{sl}_2) = \langle X_{r_1}\cdots X_{r_k}
\mid r_1 + \cdots + r_k \le n \rangle$, and the associated graded
is $\gr\,Y_\hbar(\mathfrak{sl}_2) \simeq U(\mathfrak{sl}_2[t])$.
\end{proof}

\subsubsection*{Classical limit: comparison with $U(\mathfrak{sl}_2[t])$}

\begin{theorem}[Classical limit]
\ClaimStatusProvedHere
\label{thm:classical-limit-sl2}
At $\hbar = 0$, the Yangian $Y_\hbar(\mathfrak{sl}_2)$ degenerates to the
universal enveloping algebra of the current Lie algebra
$\mathfrak{sl}_2[t] = \mathfrak{sl}_2 \otimes \bC[t]$.  Explicitly:
\begin{enumerate}[label=\textup{(\roman*)}]
\item Relations \ref{D2}--\ref{D3} at $\hbar = 0$ become
$[H(u), E(v)] = 0 = [H(u), F(v)]$ after clearing denominators
and passing to modes, yielding
$[h_r, e_s] = 2\,e_{r+s}$, $[h_r, f_s] = -2\,f_{r+s}$
--- the current-algebra Lie brackets
$[h \otimes t^r,\, e \otimes t^s] = 2\,(e \otimes t^{r+s})$.
\item Relation \ref{D4} at $\hbar = 0$ becomes
$[E_r, F_s] = H_{r+s}$, i.e.\
$[e \otimes t^r,\, f \otimes t^s] = h \otimes t^{r+s}$.
\item Relations \ref{D5}--\ref{D6} at $\hbar = 0$ become
$[E(u),E(v)] = 0 = [F(u),F(v)]$ after clearing $(u{-}v)$
--- the Serre relation for
$\mathfrak{sl}_2[t]$ (abelian root spaces).
\item The twisted coproduct $\Delta_z$ at $\hbar = 0$
degenerates to the cocommutative coproduct
$\Delta(X_r) = X_r \otimes 1 + 1 \otimes X_r$
of $U(\mathfrak{sl}_2[t])$.
\end{enumerate}
The PBW filtration on $Y_\hbar(\mathfrak{sl}_2)$ has associated graded
$\gr\,Y_\hbar(\mathfrak{sl}_2) \cong U(\mathfrak{sl}_2[t])$,
confirming that the Yangian is a filtered deformation of the current
algebra, with deformation parameter $\hbar = 1/(k+2)$ sourced by the
level of the affine input $V_k(\mathfrak{sl}_2)$.
\end{theorem}

\begin{proof}
Setting $\hbar = 0$ in relations~\ref{D1}--\ref{D6} is immediate.
For (i): \ref{D2} becomes $[H(u),E(v)] = 0$ at $\hbar = 0$;
expanding $H(u) = 1 + \hbar\sum H_r u^{-r-1}$ and taking
the leading $\hbar$-linear term recovers
$[h_r, e_s] = 2\,e_{r+s}$.
For (ii): \ref{D4} is $\hbar$-independent.
For (iii): the right-hand side of \ref{D5} vanishes at $\hbar = 0$,
giving $[E(u),E(v)] = 0$.
For (iv): the shift $u \mapsto u - z$ in the second tensor factor of
$\Delta_z$ is the spectral-parameter manifestation of the coproduct;
at $\hbar = 0$ the modes become primitive.
The associated-graded statement follows from the PBW theorem
(Theorem~\ref{thm:PBW-yangian-sl2}).
\end{proof}

\begin{remark}[The deformation parameter and the level]
\label{rem:hbar-level}
The identification $\hbar = 1/(k+2)$ has a precise bar-complex origin:
the collision residue $r(z) = \hbar\,\Omega/z$ inherits $\hbar$ from
the normalisation of the affine OPE $J^a(z)\,J^b(w) \sim
k\,\kappa^{ab}/(z{-}w)^2 + f^{ab}{}_c\,J^c/(z{-}w)$, where the
$d\log$ kernel absorbs one power (AP19) and the remaining pole is
rescaled by $1/(k + h^\vee)$ with $h^\vee = 2$ for $\mathfrak{sl}_2$.
At the critical level $k = -h^\vee = -2$, the parameter $\hbar$
diverges and the Yangian construction degenerates --- this is the
ordered-bar reflection of the Sugawara obstruction at critical level.
\end{remark}

\subsection{The Virasoro ordered bar complex and the gravitational Yangian}
\label{subsec:virasoro-ordered-bar}

The ordered bar complex $\Barch(\mathrm{Vir}_c)$ of the
Virasoro algebra provides the most striking contrast with the
Heisenberg case.  The Virasoro algebra $\mathrm{Vir}_c$ is an
$E_\infty$-chiral algebra (it is local, with $\Sigma_n$-equivariant
factorization structure), and the computation below applies the
$E_1$-chiral Koszul duality \emph{operation} to this $E_\infty$ input.
The Virasoro OPE is
\begin{equation}\label{eq:virasoro-ope-ordered}
T(z)\,T(w) \;\sim\;
\frac{c/2}{(z-w)^4}
+ \frac{2T(w)}{(z-w)^2}
+ \frac{\partial T(w)}{z-w}.
\end{equation}

\subsubsection*{Degree~$1$: the desuspended generator}

The algebra $\mathrm{Vir}_c$ has a single generator~$T$ of conformal
weight~$2$.  In the ordered bar complex, the degree-$1$ component is
\[
\Barch_1(\mathrm{Vir}_c)
\;=\;
\bC \cdot \susp^{-1}T,
\]
the one-dimensional desuspended generator.  Since there are no bar words
of length~$0$ for the bar differential to target, the differential
vanishes on degree~$1$:
\[
d_{\Barch}(\susp^{-1}T) \;=\; 0.
\]

\subsubsection*{Degree~$2$: the binary bar differential and the OPE modes}

The degree-$2$ component is spanned by the single bar word
$[\susp^{-1}T \,|\, \susp^{-1}T]$.  The bar differential extracts the
OPE modes of the Virasoro self-coupling:
\begin{align}
T_{(0)}T &= \partial T,
\label{eq:vir-ordered-mode-0}\\
T_{(1)}T &= 2T,
\label{eq:vir-ordered-mode-1}\\
T_{(3)}T &= \tfrac{c}{2}.
\label{eq:vir-ordered-mode-3}
\end{align}
The bar differential on the degree-$2$ generator is
\[
d_{\Barch}[\susp^{-1}T \,|\, \susp^{-1}T]
\;=\;
\susp^{-1}\bigl(T_{(0)}T\bigr)
\;=\;
\susp^{-1}(\partial T).
\]
The FM boundary integral over the codimension-$1$ stratum $E_{12}$ of
$\overline{\mathrm{FM}}_2(\mathbb{C})$ extracts precisely the
simple-pole mode $T_{(0)}T = \partial T$.  The higher modes $T_{(1)}T$,
$T_{(3)}T$ enter through the higher $\Ainf$ operations $m_k$
($k \ge 3$), not through the bar differential.

\subsubsection*{The collision residue $r$-matrix: a triple-pole structure}

\begin{proposition}[Virasoro collision residue; \ClaimStatusProvedHere]
\label{prop:vir-collision-residue}
The collision residue $r$-matrix of $\mathrm{Vir}_c$, obtained from the
OPE by the $d\log$ kernel absorption (the bar complex uses
$d\log(z-w)$ propagators, so pole orders drop by one relative to the
OPE), is
\begin{equation}\label{eq:vir-r-collision}
r^{\mathrm{coll}}(z)
\;=\;
\frac{c/2}{z^3}\,\mathbf{1}\otimes\mathbf{1}
\;+\;
\frac{2T}{z}\otimes\mathbf{1}.
\end{equation}
This is a \textbf{triple-pole} $r$-matrix: the leading singularity is
$z^{-3}$, corresponding to a quartic OPE pole reduced by one power
through the $d\log$ absorption.
\end{proposition}

\begin{proof}
The Virasoro OPE has poles of orders $4$, $2$, and $1$.  The $d\log$
kernel $d\log(z-w) = dz/(z-w)$ absorbs one power of $(z-w)^{-1}$ in
each term, yielding collision residue poles of orders $3$, $1$, and
$0$.  The simple-pole OPE term $\partial T(w)/(z-w)$ becomes a regular
term after absorption, contributing to the bar differential but not to
the $r$-matrix singularity.  The result is~\eqref{eq:vir-r-collision}.

Compare with the Heisenberg collision residue
$r^{\mathrm{coll}}_{\cH}(z) = k/z$ (simple pole,
Theorem~\ref{thm:heisenberg-ordered-bar}).  The Virasoro triple pole is
the source of the infinite $\Ainf$ tower (class~$\mathbf{M}$,
infinite shadow depth), while the Heisenberg simple pole gives
class~$\mathbf{G}$ (formal, $m_k = 0$ for $k \ge 3$).
\end{proof}

\subsubsection*{The classical Yang--Baxter equation}

\begin{proposition}[Virasoro CYBE; \ClaimStatusProvedHere]
\label{prop:vir-CYBE-ordered}
The collision residue $r$-matrix~\eqref{eq:vir-r-collision} satisfies
the classical Yang--Baxter equation
\begin{equation}\label{eq:vir-CYBE-ordered}
[r^{12}(u),\, r^{13}(u+v)]
+ [r^{12}(u),\, r^{23}(v)]
+ [r^{13}(u+v),\, r^{23}(v)]
\;=\; 0.
\end{equation}
\end{proposition}

\begin{proof}
The CYBE is equivalent to the PVA Jacobi identity for the Virasoro
$\lambda$-bracket
$\{T{}_\lambda T\} = \partial T + 2T\lambda + (c/12)\lambda^3$
under the Laplace correspondence
$r(z) = \int_0^\infty e^{-\lambda z}\{T{}_\lambda T\}\,d\lambda$.
The central term $(c/2)/z^3 \cdot \mathbf{1}\otimes\mathbf{1}$
commutes with all tensor factors and drops out of every commutator
in~\eqref{eq:vir-CYBE-ordered}.  The remaining $T$-dependent term
$2T/z$ satisfies the CYBE by the Jacobi identity of the underlying
Lie bracket $T_{(0)}T = \partial T$, which is the Witt algebra relation.
Explicitly, the PVA Jacobi identity
$\{T{}_\lambda\{T{}_\mu T\}\} - \{T{}_\mu\{T{}_\lambda T\}\}
= \{\{T{}_\lambda T\}{}_{\lambda+\mu}\, T\}$
becomes, after Laplace transform in both $\lambda$ and $\mu$, the
three-term CYBE relation for $r(z)$.
\end{proof}

\subsubsection*{Gravitational Yangian collapse: all levels are linearly dependent}

\begin{theorem}[Gravitational Yangian collapse; \ClaimStatusProvedHere]
\label{thm:grav-yangian-collapse}
Let $c' = 26 - c$ denote the Koszul dual central charge.  On the
Koszul dual $\mathrm{Vir}_{c'}$, define the level-$r$ generators
\begin{equation}\label{eq:grav-yangian-generators}
\mathsf{T}_n^{(r)}
\;:=\;
\binom{n+r-1}{r}\,L_{-n-r},
\qquad
\Omega_r
\;=\;
\sum_{n \in \mathbb{Z}}
\mathsf{T}_n^{(r)} \otimes L_n,
\end{equation}
so that $r(z) = \sum_{r \ge 0}\Omega_r\,z^{-r-1}$ is the
collision residue $r$-matrix.  Then:
\begin{enumerate}[label=\textup{(\roman*)}]
\item \emph{Linear dependence.}\enspace
For every $r \ge 1$ and $n \in \mathbb{Z}$, the level-$r$ generator
$\mathsf{T}_n^{(r)} = \binom{n+r-1}{r}\,L_{-n-r}$ is a
\emph{scalar multiple} of a single Virasoro mode.  There is no
independent Yangian partner generator at any level.

\item \emph{Lie-algebraic isomorphism.}\enspace
As a Lie algebra,
$Y^{\mathrm{dg}}(\mathrm{Vir}_c) \cong \mathrm{Vir}_{c'}$.
The entire Yangian Lie structure collapses onto the Virasoro algebra
at the dual central charge.

\item \emph{All-level bracket.}\enspace
The bracket at arbitrary levels $r,s \ge 0$ is
\begin{equation}\label{eq:grav-yangian-bracket}
[\mathsf{T}_m^{(r)},\,\mathsf{T}_n^{(s)}]
\;=\;
\binom{m+r-1}{r}\binom{n+s-1}{s}\,\bigl[
(n{+}s{-}m{-}r)\,L_{-m-n-r-s}
- \tfrac{c'}{12}((m{+}r)^3 - (m{+}r))\,\delta_{m+n+r+s,0}
\bigr],
\end{equation}
which reduces to a single Virasoro commutator on the right-hand side.
\end{enumerate}
\end{theorem}

\begin{proof}
Part~(i) is immediate from the definition: $\mathsf{T}_n^{(r)}
= \binom{n+r-1}{r}\,L_{-n-r}$, so the level-$r$ generator is
proportional to the single mode $L_{-n-r}$ with a binomial
coefficient.  Unlike the affine Yangian $Y(\mathfrak{sl}_2)$, where
the level-$1$ generators $(e_1, f_1, h_1)$ are linearly
independent of $(e_0, f_0, h_0)$, here all levels collapse
because the Virasoro algebra has a single generator~$T$ of spin~$2$:
the mode $L_{-n-r}$ depends only on the sum $n + r$, and the
binomial coefficient encodes only the ``splitting'' of $n+r$ into
$(n,r)$.

Part~(ii) follows from part~(i): the Lie subalgebra generated by
$\{\mathsf{T}_n^{(r)}\}_{n,r}$ is spanned by $\{L_m\}_{m \in
\mathbb{Z}}$, which is $\mathrm{Vir}_{c'}$.

Part~(iii): compute
$[\mathsf{T}_m^{(r)}, \mathsf{T}_n^{(s)}]
= \binom{m+r-1}{r}\binom{n+s-1}{s}\,[L_{-m-r}, L_{-n-s}]$.
The Virasoro commutation relation gives
$[L_{-m-r}, L_{-n-s}] = (n{+}s{-}m{-}r)\,L_{-m-n-r-s}
+ \frac{c'}{12}((m{+}r)^3 - (m{+}r))\,\delta_{m+n+r+s,0}$.
\end{proof}

\begin{remark}[Contrast with the affine Yangian]
\label{rem:virasoro-vs-affine-yangian}
For $\widehat{\mathfrak{g}}_k$ with $\mathfrak{g}$ simple, the
collision residue is $r^{\mathrm{coll}}_{\mathrm{KM}}(z) =
\Omega_{\mathfrak{g}}/z$ (simple pole), and the level-$1$ Yangian
generators are genuinely independent of the level-$0$ generators.
The spectral factorization quantum group is the full dg-shifted
Yangian $Y^{\mathrm{dg}}(\widehat{\mathfrak{g}}_k)$ with
non-primitive coproduct and independent higher-level data.  The
gravitational collapse is special to algebras with a single
generator: since $T$ has spin~$2$, all level-$r$ generators
$T^{(r)} = L_{-1}^r \cdot T_0 / r!$ (schematically) are
linearly dependent on the basic modes, and no new algebraic degrees
of freedom appear at higher levels.
\end{remark}

\subsubsection*{Non-formality: the infinite $\Ainf$ tower (class~$\mathbf{M}$)}

\begin{theorem}[Virasoro non-formality; \ClaimStatusProvedHere]
\label{thm:vir-non-formality}
For every $k \ge 3$ and generic central charge~$c$, the higher
$\Ainf$ operation $m_k$ of the ordered bar complex
$\Barch(\mathrm{Vir}_c)$ is nonzero:
\[
m_k \;\neq\; 0
\qquad\text{for all } k \ge 3.
\]
The Virasoro algebra is class~$\mathbf{M}$ (infinite shadow depth):
the $\Ainf$ tower does not truncate at any finite arity.
\end{theorem}

\begin{proof}
The quartic OPE pole forces a nonzero associator at arity~$3$:
the PVA Jacobi identity gives
$A_3(T,T,T;\,\lambda_{12},\lambda_{23})
= -\{T{}_{\lambda_{23}}\{T{}_{\lambda_{12}} T\}\} \neq 0$,
and the scalar contact term
$\frac{c}{12}\lambda_{23}^3(2\lambda_{12} + \lambda_{23})$
is proportional to~$c$.  The correction $m_3 = -A_3$
forces $m_4$ via the arity-$4$ Stasheff identity, $m_4$ forces
$m_5$, and so on in perpetuity.  The mechanism is the obstruction
cascade: each Stasheff identity $\sum m_i \circ m_j = 0$ at arity
$k+1$ produces a nonzero obstruction from the compositions of
lower-arity operations, and the unique solution is $m_{k+1} \neq 0$.

At $c = 0$ the quartic pole vanishes, but $m_k \neq 0$ persists for
all $k \ge 3$: the non-central double-pole coefficient $2T$ (a
field, not a scalar) generates its own non-vanishing associator.
The $\Ainf$ tower is $c$-independent in its $T$-dependent terms
and has an additional $c$-dependent scalar contact sector.
\end{proof}

\begin{remark}[What non-formality is not]
\label{rem:non-formality-not-E1}
Non-formality ($m_k \neq 0$ for all $k$) is a property of the
\emph{bar complex}, not of the algebra's operadic type.
$\mathrm{Vir}_c$ is and remains an $E_\infty$-chiral algebra: it is
local, with operations defined on unordered configuration spaces and
$\Sigma_n$-equivariant factorization structure.  The ordered bar
complex is obtained by applying the $E_1$-chiral Koszul duality
\emph{functor} to the $E_\infty$ input; non-formality of the output
does not make the input non-local.
\end{remark}

\subsubsection*{Curvature and the critical/self-dual dichotomy}

\begin{proposition}[Gravitational Yangian curvature; \ClaimStatusProvedHere]
\label{prop:grav-yangian-curvature}
The curvature of the gravitational dg-shifted Yangian is
\begin{equation}\label{eq:grav-yangian-kappa}
\kappa\bigl(Y^{\mathrm{dg}}(\mathrm{Vir}_c)\bigr)
\;=\;
\frac{26 - c}{2}.
\end{equation}
Two central charges are distinguished:
\begin{enumerate}[label=\textup{(\roman*)}]
\item \emph{Critical string ($c = 26$).}\enspace
$\kappa = 0$: the gravitational Yangian is flat.  The genus tower
$F_g = 0$ for all $g \ge 1$; no genus obstructions.  This is anomaly
cancellation for the gravitational sector.

\item \emph{Koszul self-dual point ($c = 13$).}\enspace
$\kappa = 13/2$.  Since $\mathrm{Vir}_{13}^! = \mathrm{Vir}_{13}$
(the Koszul involution $c \mapsto 26 - c$ has the unique fixed
point $c = 13$), the complementarity relation gives
$2F_g(\mathrm{Vir}_{13}) = 13\lambda_g^{\mathrm{FP}}$,
so $F_g(\mathrm{Vir}_{13}) = \tfrac{13}{2}\lambda_g^{\mathrm{FP}}
\neq 0$.  The genus tower does \emph{not} vanish at the
self-dual point; rather, it is \emph{invariant} under the
Koszul involution $c \mapsto 26-c$ (gravitational S-duality).
\end{enumerate}
\end{proposition}

\begin{proof}
By Vol~I Theorem~D, $\kappa(\mathrm{Vir}_c) = c/2$, and the
Koszul dual is $\mathrm{Vir}_{c'}$ with $c' = 26 - c$, so
$\kappa(Y^{\mathrm{dg}}(\mathrm{Vir}_c)) = \kappa(\mathrm{Vir}_{c'})
= (26 - c)/2$.  At $c = 26$: $\kappa = 0$ and $F_g = 0$ for all
$g \ge 1$.  At $c = 13$: self-duality
$\mathrm{Vir}_{13}^! = \mathrm{Vir}_{13}$ and the complementarity
$\kappa + \kappa^! = 13 \neq 0$ give
$2F_g = 13\lambda_g^{\mathrm{FP}}$, so
$F_g = \tfrac{13}{2}\lambda_g^{\mathrm{FP}} \neq 0$.
\end{proof}

\subsubsection*{The gauge-gravity dichotomy}

\begin{corollary}[Gauge-gravity complexity dichotomy; \ClaimStatusProvedHere]
\label{cor:gauge-gravity-dichotomy-ordered}
The two standard HT families occupy opposite corners of the
$(\Ainf,\,\Delta)$-complexity plane:
\[
\begin{array}{c|cc}
& \Ainf\textup{ products } m_k\;(k \ge 3) &
  \textup{coproduct } \Delta_z \\
\hline
\textup{Gauge \textup{(}CS, class~$\mathbf{L}$\textup{)}}
  & m_k = 0\;\textup{(}formal\textup{)}
  & \textup{nontrivial Yangian } \Delta_z \\[3pt]
\textup{Gravity \textup{(}Vir, class~$\mathbf{M}$\textup{)}}
  & m_k \neq 0\;\forall\, k\;\textup{(}non-formal\textup{)}
  & \textup{strictly primitive } \Delta_z
\end{array}
\]

\medskip
\noindent
The structural dichotomy has three aspects.
\begin{enumerate}[label=\textup{(\arabic*)}]
\item \emph{Pole order.}\enspace
Kac--Moody has a double OPE pole $\Rightarrow$ simple collision
residue $\Rightarrow$ class~$\mathbf{L}$ $\Rightarrow$ formal.
Virasoro has a quartic OPE pole $\Rightarrow$ triple collision
residue $\Rightarrow$ class~$\mathbf{M}$ $\Rightarrow$ non-formal.

\item \emph{Generator count.}\enspace
Kac--Moody has $\dim\mathfrak{g}$ generators $\Rightarrow$
independent level-$1$ Yangian partners $\Rightarrow$ nontrivial
coproduct.  Virasoro has one generator~$T$ $\Rightarrow$ all
levels collapse (Theorem~\ref{thm:grav-yangian-collapse})
$\Rightarrow$ strictly primitive coproduct.

\item \emph{DS transport.}\enspace
Drinfeld--Sokolov reduction transports class~$\mathbf{L}$ to
class~$\mathbf{M}$: the Sugawara double Wick contraction
manufactures the quartic pole from the affine double pole.
The resolvent tree formula
$m_k^{\mathcal{W}} = \sum_{\mathrm{Catalan}}(m_2^{\mathrm{aff}}
\circ h_{\mathrm{DS}})$
generates the infinite $\Ainf$ tower from the single quadratic
product of the affine input.
\end{enumerate}
\end{corollary}

\begin{proof}
The $\Ainf$ column: class~$\mathbf{L}$ formal by the pole-order
classification; class~$\mathbf{M}$ non-formal by
Theorem~\ref{thm:vir-non-formality}.  The coproduct column: for
gauge theory $\Delta_z$ is the standard Yangian coproduct
(non-primitive, encoding the spectral braiding); for gravity
$\Delta_z(x) = \tau_z(x) \otimes 1 + 1 \otimes x$ is strictly
primitive (the ghost-number obstruction kills all non-primitive
terms in the HPL transfer, since $|h_{\mathrm{DS}}| = -1$ in ghost
degree and projecting to ghost-number-$0$ $\mathcal{W}$-algebra
cohomology annihilates all trees with net nonzero ghost number).
The DS transport is the resolvent tree formula.
\end{proof}

\begin{remark}[Summary of the gravitational dg-shifted Yangian]
\label{rem:grav-yangian-summary}
The gravitational dg-shifted Yangian
$Y^{\mathrm{dg}}(\mathrm{Vir}_c)$ is the Virasoro algebra
$\mathrm{Vir}_{26-c}$ equipped with:
\begin{enumerate}[label=(\alph*)]
\item the infinite $\Ainf$ product tower
$\{m_k\}_{k \ge 3}$ (non-formality, class~$\mathbf{M}$);
\item the strictly primitive coproduct
$\Delta_z(x) = \tau_z(x) \otimes 1 + 1 \otimes x$;
\item the triple-pole collision residue
$r(z) = (c'/2)/z^3 + 2T/z$;
\item curvature $\kappa = (26-c)/2$.
\end{enumerate}
The Lie-algebraic Yangian structure is trivial (all levels collapse);
the entire nontrivial content resides in the $\Ainf$ tower.  The
gravitational Yangian is thus a \emph{curved $\Ainf$-deformation}
of $U(\mathrm{Vir}_{26-c})$, not a filtered deformation of
$U(\mathrm{Witt}[t])$.
\end{remark}


\section{The double bar and what it recovers}
\label{sec:double-bar}

The ordered Koszul dual of $\hat\mathfrak{g}_k$ is
$Y_\hbar(\mathfrak{g})$.  What is the ordered Koszul dual of
$Y_\hbar(\mathfrak{g})$?

If the ordered Koszul duality were an involution --- as the full
$\mathrm{SC}^{\mathrm{ch,top}}$ duality is
(Theorem~\ref{thm:two-colour-double-kd} below) --- then
$\Barch^{\mathrm{ord}}(Y_\hbar(\mathfrak{g}))$ would recover
$\hat\mathfrak{g}_k$.  It does not.  The open-colour double bar
recovers the current algebra $U(\mathfrak{g}[t])$ \emph{without}
the central extension~$k$.  The level is invisible to the ordered bar
because it enters the Yangian as a deformation parameter
$\hbar = 1/(k{+}h^\vee)$, not as a structural feature of the spectral
OPE.

\begin{theorem}[Central extension is invisible to the open-colour
double bar]
\ClaimStatusProvedHere
\label{thm:central-extension-invisible}
\label{prop:double-bar-sl2}
Let $\mathfrak{g}$ be a simple Lie algebra and let
$A = \hat\mathfrak{g}_k$ at non-critical level.  The open-colour
double Koszul dual satisfies
\[
(A^{!}_{\mathrm{line}})^{!}_{\mathrm{line}}
\;\simeq\;
U(\mathfrak{g}[t]),
\qquad
\textup{not}\;\;
\hat\mathfrak{g}_k.
\]
\end{theorem}

\begin{proof}
The level $k$ of $\hat\mathfrak{g}_k$ enters the ordered Koszul dual
$Y_\hbar(\mathfrak{g})$ as the scalar $\hbar = 1/(k + h^\vee)$.
The holomorphic double pole $k\,\delta^{ab}/(z-w)^2$ --- the source
of the central term $k\,n\,\delta_{n+m,0}$ in the mode algebra ---
does not survive as a double pole in the spectral variable: the
Yangian $R$-matrix $R(u) = 1 + \hbar\,r(u) + O(\hbar^2)$ has its
singularity at $u = 0$ with a \emph{simple} pole $r(u) = \Omega/u$,
where $\Omega = \sum_a T^a \otimes T^a$ is the Casimir tensor.

The ordered bar differential of $Y_\hbar(\mathfrak{g})$ uses the
propagator $d\log(u_i - u_j)$ on
$\mathrm{Conf}_n^{<}(\mathbb{A}^1_u)$.  By the $d\log$ absorption
principle, a simple spectral pole yields the zeroth product
$a_{(0)}b$; a double pole would yield a derivative term
$a_{(1)}b$ carrying the central extension.  But the spectral OPE has
no double pole.  Degree by degree:

\smallskip\noindent\emph{Degree $1$.}
Yangian generators $E_n, F_n, H_n$ map to current modes $J^a_n$.

\smallskip\noindent\emph{Degree $2$.}
The spectral OPE
$[E(u), F(v)] \sim \hbar/(u-v)\cdot(\ldots)$ has only simple poles, so
\[
d_{\mathrm{bar}}(s^{-1}E_n \otimes s^{-1}F_m)
\;=\;
s^{-1}[E_n, F_m]
\;=\;
s^{-1}H_{n+m},
\]
with no central term.  For the $E$--$E$ component, the pole at
$u - v = \hbar$ is shifted off the collision diagonal, giving
$d_{\mathrm{res}}(s^{-1}E \otimes s^{-1}E) = 0$, i.e.\
$[e,e] = 0$ in $\mathfrak{g}[t]$.

\smallskip\noindent\emph{Degree $3$.}
$d^2 = 0$ reduces to the classical Yang--Baxter equation for
$r(u) = \Omega/u$, which holds by the Jacobi identity for
$\mathfrak{g}$.

\smallskip
Hence the bar differential recovers only the Lie bracket of
$\mathfrak{g}[t]$, not the central extension.
\end{proof}

The loss is not a failure of the machinery --- it is a precise
diagnosis.  The double pole $k/(z-w)^2$ lives in the
\emph{holomorphic} OPE on $\mathrm{FM}_2(\mathbb{C})$.  Passage to
the Yangian transports this pole into the scalar $\hbar$, which
parametrises the Yangian but is invisible to its spectral OPE on
$\mathrm{Conf}^{<}(\mathbb{A}^1)$.  The open colour sees the Lie
bracket; the level is carried by the closed colour alone.

\begin{theorem}[Two-colour double Koszul duality is involutive]
\ClaimStatusProvedHere
\label{thm:two-colour-double-kd}
Let $A$ be a strongly admissible $\Einf$-chiral algebra on a curve
$X$.  The full $\mathrm{SC}^{\mathrm{ch,top}}$ double Koszul duality
satisfies $(A^!)^! \simeq A$, recovering all central extensions and
higher-pole data.
\end{theorem}

\begin{proof}
The full duality operates on both colours: the \emph{closed colour}
(symmetric bar $\barB^{\Sigma}(A)$ on $\mathrm{Ran}(X)/\Sigma_n$,
dualised by $D_{\mathrm{Ran}}$) carries the holomorphic OPE including
the double pole; the \emph{open colour} (ordered bar
$\Barch^{\mathrm{ord}}(A)$ on $\mathrm{Conf}_n^{<}(\mathbb{R})$,
dualised by linear duality) carries the spectral data.  Verdier/Ran
duality preserves the full pole structure, so the double pole
survives.  The bar--cobar involution
$\Omega(\Barch(\Omega(\Barch(A)))) \simeq A$
(Hypothesis~\ref{hyp:inputs}(I), applied twice) then recovers $A$ in
its entirety.  For $\hat\mathfrak{g}_k$: the closed-colour Verdier
dual is $\hat\mathfrak{g}_{-k-2h^\vee}$ (Feigin--Frenkel dual level),
and the second application returns to $\hat\mathfrak{g}_k$ with the
central extension intact.
\end{proof}

\begin{corollary}[Non-redundancy of the two colours]
\ClaimStatusProvedHere
\label{cor:two-colours-non-redundant}
\index{double bar!open-colour}
\index{central extension!invisible to open colour}
For any $\Einf$-chiral algebra $A$ with non-trivial OPE poles, the
open and closed bar complexes carry genuinely non-redundant
information.  Neither colour alone recovers $A$ by iterating Koszul
duality: the open-colour double bar loses the central extension
\textup{(Theorem~\ref{thm:central-extension-invisible})}, while the
closed-colour symmetric bar does not determine the $R$-matrix
monodromy for the ordered-to-unordered descent
\textup{(Proposition~\ref{prop:r-matrix-descent})}.
\end{corollary}

\begin{remark}[The two-colour architecture at a glance]
\label{rem:double-bar-architecture}
The double-bar computation for $\hat\mathfrak{sl}_{2,k}$ distils the
architecture:
\[
\begin{array}{c|c|c}
& \text{Closed (holomorphic)} & \text{Open (topological)} \\
\hline
\text{Pole order} & \text{double: } k/(z-w)^2
  & \text{simple: } \hbar/(u-v) \\
\text{Extracted datum} & \text{central extension}
  & \text{Lie bracket of } \mathfrak{g}[t] \\
\text{Double KD output} & \hat\mathfrak{g}_k \text{ (full)}
  & U(\mathfrak{g}[t]) \text{ (no central ext.)}
\end{array}
\]
The two colours of Swiss-cheese are not decorative.  The open colour
sees the Yangian; the closed colour sees the level.  Neither alone
recovers the algebra.
\end{remark}

\section{The universal Kac--Moody Yangian theorem}
\label{sec:km-yangian}

The full nonabelian Kac--Moody family yields the definitive result.  For every simple Lie algebra~$\mathfrak{g}$
and non-critical level $k\neq -h^\vee$, the open-colour Koszul
dual is the Yangian $Y_\hbar(\mathfrak{g})$.  This is the
ordered analogue of Theorem~\ref{thm:universal-kac-moody-koszul},
which gives the symmetric Koszul duality
$\widehat{\mathfrak{g}}_k^!\simeq
\widehat{\mathfrak{g}}_{-k-2h^\vee}$.

All affine Kac--Moody algebras $\widehat{\mathfrak{g}}_k$ are
$E_\infty$-chiral: they are local vertex algebras with
$\Sigma_n$-equivariant factorization structure.  The OPE poles
do not break $E_\infty$-locality.

\subsection{The ordered bar complex of
$\widehat{\mathfrak{g}}_k$}
\label{subsec:km-ordered-bar}

Fix a simple Lie algebra~$\mathfrak{g}$ with invariant form
$(\cdot,\cdot)$ normalised so that $(\theta|\theta)=2$
for the highest root~$\theta$, dual Coxeter number~$h^\vee$,
and structure constants $f^{ab}{}_{c}$ in an orthonormal basis
$\{T^a\}_{a=1}^{\dim\mathfrak{g}}$.
Let $\widehat{\mathfrak{g}}_k$ be the affine vertex algebra at
level $k\neq -h^\vee$, with OPE
\begin{equation}\label{eq:km-ope-ordered}
J^a(z)\,J^b(w)
\;\sim\;
\frac{k\,(T^a,T^b)}{(z-w)^2}
\;+\;
\frac{f^{ab}{}_{c}\,J^c(w)}{z-w}.
\end{equation}
Write $\hbar = 1/(k+h^\vee)$ for the deformation parameter
(well-defined away from critical level) and
\[
\Omega
\;=\;
\sum_{a=1}^{\dim\mathfrak{g}} T^a\otimes T^a
\;\in\;
\mathfrak{g}\otimes\mathfrak{g}
\]
for the quadratic Casimir tensor.

\begin{theorem}[Universal Kac--Moody Yangian theorem;
\ClaimStatusProvedHere]
\label{thm:km-yangian}
For every simple Lie algebra~$\mathfrak{g}$ and non-critical
level $k\neq -h^\vee$, the open-colour Koszul dual of the
affine vertex algebra $\widehat{\mathfrak{g}}_k$ is the
Yangian:
\begin{equation}\label{eq:km-yangian-duality}
A^{!,\mathrm{ord}}_{\mathrm{line}}
(\widehat{\mathfrak{g}}_k)
\;\simeq\;
Y_\hbar(\mathfrak{g}),
\qquad
\hbar = \frac{1}{k+h^\vee}.
\end{equation}
The duality is:
\begin{enumerate}[label=\textup{(\arabic*)}]
\item functorial in~$\mathfrak{g}$
\textup{(}respects Lie algebra homomorphisms\textup{)};

\item compatible with the symmetric Koszul duality
$\widehat{\mathfrak{g}}_k^!\simeq
\widehat{\mathfrak{g}}_{-k-2h^\vee}$
via $R$-twisted $\Sigma_n$-descent
\textup{(}Proposition~\textup{\ref{prop:r-matrix-descent})};

\item valid for all nine infinite families and
exceptional types \textup{(}Table~\textup{\ref{tab:km-yangian-data})}.
\end{enumerate}
\end{theorem}

\begin{proof}
The proof has five steps.

\emph{Step~1: Collision residue.}
The OPE~\eqref{eq:km-ope-ordered} has a double pole
$k\,(T^a,T^b)/(z-w)^2$ and a simple pole
$f^{ab}{}_{c}\,J^c(w)/(z-w)$.  The bar kernel
$d\log(z-w)$ absorbs one power (AP19), so the collision
residue on $\mathrm{Conf}_2^{\mathrm{ord}}(\bC)$ is
\begin{equation}\label{eq:km-collision-residue}
r_{12}(z)
\;=\;
\frac{\hbar\,\Omega_{12}}{z}
\;+\;
(r_0)_{12},
\end{equation}
where $r_0 = \sum_{a,b,c} f^{ab}{}_{c}\,T^a\wedge T^b\otimes T_c$
encodes the Lie bracket and $\hbar\,\Omega/z$ is the spectral
term from the double pole.  Here
$(r_0)_{12}$ acts in the $(1,2)$ tensor slots.

\emph{Step~2: Bar differential at degree~$2$.}
The ordered bar differential on
$\Barch_2^{\mathrm{ord}}(\widehat{\mathfrak{g}}_k)$ extracts
the Lie bracket from the simple-pole collision:
\begin{equation}\label{eq:km-bar-d2}
d_{\Barch}[\susp^{-1}J^a \,|\, \susp^{-1}J^b]
\;=\;
f^{ab}{}_{c}\,[\susp^{-1}J^c].
\end{equation}
This is precisely the Lie bracket $[T^a, T^b] = f^{ab}{}_{c}\,T^c$
of~$\mathfrak{g}$.  At degree~$2$, the ordered bar complex
reproduces the Chevalley--Eilenberg complex
of~$\mathfrak{g}$.  The double-pole term
$k\,(T^a,T^b)/(z-w)^2$ contributes through the collision
residue $\hbar\,\Omega/z$ to the \emph{$R$-matrix}, not
to the bar differential (this is the ordered/symmetric
distinction: in the symmetric bar complex, the same data
appears in the bar differential as the curvature
$m_0 = \frac{k+h^\vee}{2h^\vee}\kappa$).

\emph{Step~3: $d^2=0$ at degree~$3$ is the CYBE.}
The bar differential satisfies $d^2 = 0$ on the ordered bar
complex.  At degree~$3$, this condition applied to
$[\susp^{-1}J^a \,|\, \susp^{-1}J^b \,|\, \susp^{-1}J^c]$
is equivalent to the \emph{classical Yang--Baxter equation}
for $r(z)=\hbar\,\Omega/z$:
\begin{equation}\label{eq:cybe}
[r_{12}(z_1{-}z_2),\, r_{13}(z_1{-}z_3)]
\;+\;
[r_{12}(z_1{-}z_2),\, r_{23}(z_2{-}z_3)]
\;+\;
[r_{13}(z_1{-}z_3),\, r_{23}(z_2{-}z_3)]
\;=\; 0.
\end{equation}
Substituting $r_{ij}(z) = \hbar\,\Omega_{ij}/z$, the three
terms become proportional to $f^{a[b}{}_{e}\,f^{c]e}{}_{d}$,
and the vanishing is exactly the Jacobi identity for
$\mathfrak{g}$.  Thus $d^2=0$ holds unconditionally for every
simple~$\mathfrak{g}$.

\emph{Step~4: Complete strictification.}
The spectral Drinfeld class $[\delta]\in H^2(\mathfrak{g},
\mathfrak{g}\otimes\mathfrak{g})$ measures the obstruction
to strictifying the $A_\infty$-Yangian to a strict associative
algebra.  For every simple Lie algebra~$\mathfrak{g}$, the
root spaces are one-dimensional
(root-space one-dimensionality),
so the relevant $H^2$ group decomposes into root sectors:
each sector is one-dimensional, and the Jacobi identity forces
the cocycle $[\delta]$ to be a coboundary in every sector.
The class therefore vanishes:
\begin{equation}\label{eq:strictification}
[\delta] = 0
\quad\text{for all simple }\mathfrak{g}.
\end{equation}
This is the strictification mechanism: the Yangian
$Y_\hbar(\mathfrak{g})$ is a \emph{strict} associative
algebra, not merely an $A_\infty$-algebra.

\emph{Step~5: Identification of the Koszul dual.}
The cohomology $H^\bullet(\Barch^{\mathrm{ord}}
(\widehat{\mathfrak{g}}_k))$ with the dual multiplication
(from dualising the deconcatenation coproduct) is the
RTT presentation of the Yangian $Y_\hbar(\mathfrak{g})$:
the degree-$1$ generators are the currents
$\{T^a_n\}_{a,n}$, the degree-$2$ relations are the RTT
relations $R(u-v)\,T_1(u)\,T_2(v) = T_2(v)\,T_1(u)\,R(u-v)$
with $R$-matrix $R(u) = 1 + \hbar\,\Omega/u$, and the
strictification of Step~4 ensures no higher $A_\infty$
corrections.  By the Drinfeld theorem, this is
$Y_\hbar(\mathfrak{g})$.
\end{proof}

\subsection{Data for all simple types}

\begin{table}[ht]
\centering
\caption{Ordered Koszul duality data for the simple
Kac--Moody family.  All entries are unconditional
\textup{(}\ClaimStatusProvedHere\textup{)}.}
\label{tab:km-yangian-data}
\renewcommand{\arraystretch}{1.3}
\begin{tabular}{@{}llcclcl@{}}
\toprule
Type & $\mathfrak{g}$ & $\dim\mathfrak{g}$ & $h^\vee$
     & Fund.\ repr.\ dim & $\hbar$ & $\kappa(Y_\hbar(\mathfrak{g}))$ \\
\midrule
$A_n$ & $\mathfrak{sl}_{n+1}$ & $n^2+2n$ & $n+1$
      & $n+1$ & $\frac{1}{k+n+1}$
      & $-\frac{(n^2+2n)(k+n+1)}{2(n+1)}$ \\[4pt]
$B_n$ & $\mathfrak{so}_{2n+1}$ & $n(2n+1)$ & $2n-1$
      & $2n+1$ & $\frac{1}{k+2n-1}$
      & $-\frac{n(2n+1)(k+2n-1)}{2(2n-1)}$ \\[4pt]
$C_n$ & $\mathfrak{sp}_{2n}$ & $n(2n+1)$ & $n+1$
      & $2n$ & $\frac{1}{k+n+1}$
      & $-\frac{n(2n+1)(k+n+1)}{2(n+1)}$ \\[4pt]
$D_n$ & $\mathfrak{so}_{2n}$ & $n(2n-1)$ & $2n-2$
      & $2n$ & $\frac{1}{k+2n-2}$
      & $-\frac{n(2n-1)(k+2n-2)}{2(2n-2)}$ \\[4pt]
$E_6$ & $\mathfrak{e}_6$ & $78$ & $12$
      & $27$ & $\frac{1}{k+12}$
      & $-\frac{78(k+12)}{24}$ \\[4pt]
$E_7$ & $\mathfrak{e}_7$ & $133$ & $18$
      & $56$ & $\frac{1}{k+18}$
      & $-\frac{133(k+18)}{36}$ \\[4pt]
$E_8$ & $\mathfrak{e}_8$ & $248$ & $30$
      & $248$ & $\frac{1}{k+30}$
      & $-\frac{248(k+30)}{60}$ \\[4pt]
$F_4$ & $\mathfrak{f}_4$ & $52$ & $9$
      & $26$ & $\frac{1}{k+9}$
      & $-\frac{52(k+9)}{18}$ \\[4pt]
$G_2$ & $\mathfrak{g}_2$ & $14$ & $4$
      & $7$ & $\frac{1}{k+4}$
      & $-\frac{14(k+4)}{8}$ \\
\bottomrule
\end{tabular}
\end{table}

\noindent
In every case the curvature of the open-colour Koszul dual is
\begin{equation}\label{eq:km-yangian-curvature}
\kappa\bigl(Y_\hbar(\mathfrak{g})\bigr)
\;=\;
-\kappa(\widehat{\mathfrak{g}}_k)
\;=\;
-\frac{\dim(\mathfrak{g})\cdot(k+h^\vee)}{2h^\vee},
\end{equation}
the negation of the affine curvature, in accordance with the
complementarity involution $\kappa + \kappa^! = 0$ for the
Kac--Moody family.

\begin{remark}[Casimir tensor for each type]
\label{rem:casimir-types}
In the fundamental representation, the Casimir tensor
$\Omega\in\mathfrak{g}\otimes\mathfrak{g}$ acts as:
\begin{enumerate}[label=\textup{(\alph*)}]
\item \emph{Type $A_n$.}\enspace
$\Omega = P - \frac{1}{n+1}\,I\otimes I$,
where $P\in\End(\bC^{n+1}\otimes\bC^{n+1})$ is the
permutation operator.
\item \emph{Types $B_n$, $D_n$.}\enspace
$\Omega = P - K$,
where $K_{ij,kl} = \delta_{ik'}\delta_{jl'}$
is the contraction with the invariant bilinear form
(using the index convention $i' = 2n+2-i$ for $B_n$,
$i' = 2n+1-i$ for $D_n$).
\item \emph{Type $C_n$.}\enspace
$\Omega = P - K$,
where $K_{ij,kl} = J_{ik}J_{jl}$ involves the
symplectic form~$J$.
\item \emph{Exceptional types.}\enspace
$\Omega$ is determined by the structure constants
$f^{ab}{}_{c}$ via
$\Omega = \sum_a T^a\otimes T^a$
in any orthonormal basis; explicit matrix forms
follow from the Chevalley basis.
\end{enumerate}
\end{remark}

\begin{remark}[Relation to Drinfeld--Kohno and quantum groups]
\label{rem:km-drinfeld-kohno}
The Kohno connection~\eqref{eq:kohno-connection} for
$\widehat{\mathfrak{g}}_k$ takes the form
$\nabla = d - \hbar\sum_{i<j}\Omega_{ij}\,d\log(z_i-z_j)$,
which is the Knizhnik--Zamolodchikov connection at level~$k$.
Its monodromy representation $\pi_1(\mathrm{Conf}_n(\bC))
\to\GL(\mathfrak{g}^{\otimes n})$ factors through the
Yangian via the Drinfeld--Kohno theorem: the monodromy of
the KZ connection at $\hbar = 1/(k+h^\vee)$ coincides with
the $R$-matrix of $Y_\hbar(\mathfrak{g})$ acting on evaluation
modules.  The ordered bar complex provides the chain-level
refinement of this classical result.
\end{remark}

\begin{remark}[The pole-order dichotomy]
\label{rem:km-pole-dichotomy}
The Kac--Moody OPE has a double pole and a simple pole.
The $d\log$ kernel absorbs one power, leaving a collision
residue with a simple pole $\hbar\,\Omega/z$ (the spectral
term) and a constant $r_0$ (the Lie bracket term).  This
places all $\widehat{\mathfrak{g}}_k$ in class~L (double-pole
OPE, simple-pole collision residue, finite shadow depth).
By contrast, the Virasoro OPE has a quartic pole, giving
a triple-pole collision residue and class~M (infinite shadow
depth).  Under DS reduction
$\widehat{\mathfrak{g}}_k\to\mathcal{W}_k(\mathfrak{g})$,
the pole-order escalation transports class~L to class~M
via Sugawara.
\end{remark}

\begin{remark}[The Kac--Moody frontier]
\label{rem:km-frontier}
Theorem~\ref{thm:km-yangian} covers all simple Lie
algebras.  The genuine frontier is Kac--Moody algebras
$\widehat{\mathfrak{g}}_k$ for~$\mathfrak{g}$ a
\emph{Kac--Moody} Lie algebra (not finite-dimensional
simple): for affine or indefinite Kac--Moody algebras,
root multiplicities can exceed one, the root-space
one-dimensionality argument of Step~4 fails, and the
strictification mechanism requires new input.
\end{remark}


\section{Spectral Drinfeld strictification}
\label{sec:spectral-drinfeld-strictification}

The ordered bar complex of a dg-shifted Yangian carries a filtered
$A_\infty$ structure.  The question of whether this $A_\infty$ structure
can be straightened to a genuine associative structure with spectral
parameter is controlled by a cohomological obstruction: the
\emph{spectral Drinfeld class}.  This
obstruction vanishes for all simple Lie algebras.

\subsection{The obstruction class}

\begin{definition}[Spectral Drinfeld class at filtration $p$]
\ClaimStatusProvedHere
\label{def:spectral-drinfeld-class}
Let $A$ be a filtered dg-shifted Yangian carrying a spectral
quasi-factorization datum with quasi-associator
\[
\Phi(w,z)=1+\phi_p(w,z)\pmod{F^{p+1}},
\]
where $\phi_p$ is the first nontrivial term at filtration~$p$.
The spectral cochain complex is defined by
\[
C^2_{\mathrm{spec}}(A)=\{g(u)\in A^{\widehat\otimes 2}\},
\qquad
C^3_{\mathrm{spec}}(A)=\{\phi(w,z)\in A^{\widehat\otimes 3}\},
\]
with linearized spectral differential
\[
(\delta g)(w,z)
=
1\otimes g(z)
+ (\id\otimes \Delta_z)g(w)
- (\Delta_w\otimes \id)g(z)
- g(w)\otimes 1.
\]
The \emph{spectral Drinfeld class at filtration $p$} is
\[
\mathfrak D_{\mathrm{spec}}(A,\Omega)\big|_{F^p}
:=[\phi_p]
\in H^3_{\mathrm{spec}}(\gr^p A).
\]
\end{definition}

\begin{remark}[Interpretation]
The spectral Drinfeld class measures the failure of the $A_\infty$
bar-complex structure to strictify to a genuine associative product
with spectral parameter.  At each filtration $p\ge 2$, the pentagon
relation on the quasi-associator linearizes: $\phi_p$ is a $3$-cocycle,
a spectral twist shifts it by a coboundary, and the class
$[\phi_p]\in H^3_{\mathrm{spec}}(\gr^p A)$ is the intrinsic
obstruction.  If $[\phi_p]=0$ for all $p\ge 2$, an iterative
filtered twist produces $\Phi^G=1$, and the quasi-factorization
datum strictifies to an honest spectral factorization $R$-matrix.
\end{remark}

\subsection{Root-space one-dimensionality}

The key structural input is a classical theorem about the root
structure of simple Lie algebras: every root space is one-dimensional.
This forces the spectral Drinfeld class, which is valued in root
spaces, to be a scalar --- and a scalar is determined by a single
equation.

\begin{theorem}[Root-space one-dimensionality; \ClaimStatusProvedHere]
\label{thm:root-space-one-dim-v1}
Let $\fg$ be a simple Lie algebra.  For any subset
$S=\{\alpha_1,\ldots,\alpha_n\}$ of simple roots such that
$\alpha=\sum_{i=1}^n\alpha_i$ is a root of $\fg$, the multilinear
degree-$n$ space
\[
\operatorname{span}\{\textup{Lie monomials in }E_{\alpha_1},\ldots,
E_{\alpha_n}\textup{ using each exactly once}\}\;\subseteq\;\fg_\alpha
\]
is one-dimensional.
\end{theorem}

\begin{proof}
Every multilinear Lie monomial in $E_{\alpha_1},\ldots,E_{\alpha_n}$
has weight $\alpha_1+\cdots+\alpha_n=\alpha$ (since each bracket
adds the weights of its arguments).  Therefore all such monomials
lie in the root space $\fg_\alpha$.  In a simple Lie algebra, every
root has multiplicity one: $\dim\fg_\alpha=1$.  Hence the
multilinear space is at most one-dimensional.

It remains to verify that at least one multilinear monomial is
nonzero.  Since $\alpha$ is a root and the simple roots $\alpha_i$ are
linearly independent, there exists an ordering
$\alpha_{i_1},\ldots,\alpha_{i_n}$ and a sequence of partial sums
$\beta_k=\alpha_{i_1}+\cdots+\alpha_{i_k}$ such that each $\beta_k$
is a root (this follows from the connectedness of the support of
$\alpha$ in the Dynkin diagram and the standard theory of root
strings).  The iterated bracket
$[E_{\alpha_{i_1}},[E_{\alpha_{i_2}},\ldots,E_{\alpha_{i_n}}]\cdots]$
is then nonzero (each intermediate bracket lands in a nonzero root
space).
\end{proof}

\subsection{The Jacobi collapse lemma}

When root spaces are one-dimensional, multilinear Lie monomials
in root vectors are necessarily proportional --- the target space
has no room for independent contributions.  The mechanism by
which the free Lie algebra's higher-dimensional multilinear spaces
collapse to one dimension inside $\fg$ is the Jacobi identity.

The simplest nontrivial case is the $D_4$ star: four simple roots
$\alpha_1,\alpha_2,\alpha_3,\alpha_c$ with $\alpha_c$ adjacent to
each of $\alpha_1,\alpha_2,\alpha_3$, and the $\alpha_i$ mutually
non-adjacent.  In the free Lie algebra modulo non-adjacency
commutation, the multilinear degree-$4$ space is two-dimensional.
In $\fg$, the Jacobi identity collapses it to one.

\begin{lemma}[Jacobi collapse for star sectors; \ClaimStatusProvedHere]
\label{lem:jacobi-collapse-v1}
Let $X_1,X_c,X_2,X_3$ be root vectors in a Lie algebra $\fg$ with
$[X_1,X_2]=[X_1,X_3]=[X_2,X_3]=0$ (the $D_4$-star commutation
relations).  Define
\[
A = [[X_2,[X_1,X_c]],X_3],
\qquad
B = [[X_3,[X_1,X_c]],X_2].
\]
Then $A=B$.
\end{lemma}

\begin{proof}
Apply the Jacobi identity to $(X_2,X_c,X_3)$:
\[
[X_2,[X_c,X_3]]+[X_c,[X_3,X_2]]+[X_3,[X_2,X_c]]=0.
\]
Since $[X_2,X_3]=0$, the middle term vanishes, giving
\begin{equation}\label{eq:jacobi-star-v1}
[X_2,[X_c,X_3]]=[X_3,[X_c,X_2]].
\end{equation}
Now expand $A$ and $B$ using Jacobi and $[X_1,X_2]=[X_1,X_3]=0$:
\begin{align*}
A &= [[X_2,[X_1,X_c]],X_3]
  = [X_2,[[X_1,X_c],X_3]] - [[X_1,X_c],[X_2,X_3]]\\
  &= [X_2,[[X_1,X_c],X_3]]
  \qquad(\text{since }[X_2,X_3]=0)\\
  &= [X_2,[X_1,[X_c,X_3]]] - [X_2,[X_c,[X_1,X_3]]]\\
  &= [X_2,[X_1,[X_c,X_3]]]
  \qquad(\text{since }[X_1,X_3]=0)\\
  &= [[X_2,X_1],[X_c,X_3]] + [X_1,[X_2,[X_c,X_3]]]\\
  &= [X_1,[X_2,[X_c,X_3]]]
  \qquad(\text{since }[X_1,X_2]=0).
\end{align*}
The identical computation with $2\leftrightarrow 3$ gives
$B=[X_1,[X_3,[X_c,X_2]]]$.  By~\eqref{eq:jacobi-star-v1},
$[X_2,[X_c,X_3]]=[X_3,[X_c,X_2]]$, hence $A=B$.
\end{proof}

\begin{remark}[Root multiplicity one as structural inevitability]
The Jacobi collapse lemma is not a clever computation that
happens to kill the obstruction.  The point is that the obstruction
\emph{has nowhere to live}.

The spectral Drinfeld class at filtration $n$, restricted to the root
sector $\alpha=\sum\alpha_i$, is a cohomology class valued in
$\fg_\alpha$.  Root multiplicity one says
$\dim\fg_\alpha=1$.  A cohomology class valued in a one-dimensional
space is a scalar.  A scalar is determined by a single equation.
The coproduct rigidity provides that equation.

The free Lie algebra on the simple root vectors, modulo non-adjacency
commutation, may have higher-dimensional multilinear spaces (the $D_4$
star has $\dim=2$).  But in $\fg$ the Jacobi identity collapses them
(Lemma~\ref{lem:jacobi-collapse-v1}).  This collapse is not an accident:
it is the shadow of root multiplicity one, which is itself a theorem
about the representation theory of $\mathfrak{sl}_2$-triples.
\end{remark}

\subsection{The beta-integral coefficient}

The BCH expansion of $\log(e^{X_1}\cdots e^{X_n})$ produces the
universal coefficient at each filtration.  This coefficient is
determined by the Dynkin projection.

\begin{theorem}[Dynkin coefficient via the beta integral; \ClaimStatusProvedHere]
\label{thm:dynkin-beta-integral}
Let $X_1,\ldots,X_n$ satisfy the path commutation relations
$[X_i,X_j]=0$ for $|i-j|\ge 2$.  Then the multilinear
$X_1\cdots X_n$-part of the ordered Baker--Campbell--Hausdorff product is
\[
\operatorname{mult}_{X_1\cdots X_n}
\log(e^{X_1}\,e^{X_2}\cdots e^{X_n})
\;=\;
\frac{1}{n}\,[X_1,[X_2,\ldots,[X_{n-1},X_n]\cdots]],
\]
where the coefficient $1/n$ arises from the beta integral
\begin{equation}\label{eq:beta-integral-coeff}
\frac{1}{n}=\int_0^1(1-t)^{n-1}\,dt.
\end{equation}
This is the unique coefficient forced by the BCH formula: the Dynkin
projection $D_n=\frac{1}{n}\ell_n$ applied to the ordered
concatenation $X_1\otimes\cdots\otimes X_n$.
\end{theorem}

\begin{proof}
The multilinear part is $c_n\cdot[X_1,[X_2,\ldots,X_n]\cdots]$ for a
unique scalar $c_n$ (the multilinear Lie space is one-dimensional by
root-space one-dimensionality in $\fg$, or by the path
one-dimensionality theorem in the free Lie algebra quotient).  We
determine $c_n$ by extracting the coefficient of the concatenation
monomial $X_1\cdot X_2\cdots X_n$ in the universal enveloping algebra.

\medskip\noindent\textbf{Step~1: Coefficient in the Lie bracket.}
In the universal enveloping algebra, the right-normed bracket
$[X_1,[X_2,\ldots,[X_{n-1},X_n]\cdots]]$ has leading term
$X_1 X_2\cdots X_n$ with coefficient $+1$.

\medskip\noindent\textbf{Step~2: Coefficient in the BCH logarithm.}
The multilinear part of $e^{X_1}\cdots e^{X_n}$ in the free
associative algebra is $X_1\cdots X_n$.  The coefficient of
$X_1\cdots X_n$ in $\log(e^{X_1}\cdots e^{X_n})$, computed via
$\log(1+A)=\sum_{k\ge 1}(-1)^{k+1}A^k/k$, counts compositions
of~$n$ into $k$ parts:
\[
\text{coeff of }X_1\cdots X_n\text{ in }\log(P)
=
\sum_{k=1}^{n}\frac{(-1)^{k+1}}{k}\binom{n-1}{k-1}.
\]

\medskip\noindent\textbf{Step~3: The beta integral.}
\begin{align*}
\sum_{k=1}^{n}\frac{(-1)^{k+1}}{k}\binom{n-1}{k-1}
&=
\sum_{j=0}^{n-1}\frac{(-1)^{j}}{j+1}\binom{n-1}{j}
=
\sum_{j=0}^{n-1}\binom{n-1}{j}(-1)^j\int_0^1 t^j\,dt\\
&=
\int_0^1\sum_{j=0}^{n-1}\binom{n-1}{j}(-t)^j\,dt
=
\int_0^1(1-t)^{n-1}\,dt
\;=\;\frac{1}{n}.
\end{align*}
Combining Steps~1--3: $c_n\cdot 1=\frac{1}{n}$, hence
$c_n=\frac{1}{n}$.
\end{proof}

\begin{remark}[The Dynkin projection as open-colour counit]
The coefficient $1/n$ is the Dynkin projection
$D_n\colon T^n(V)\to\mathrm{Lie}^n(V)$,
$D_n=\frac{1}{n}\ell_n$, applied to the ordered concatenation.
The bar-cobar counit $\varepsilon\colon\Omega(\Barch(\cA))\to\cA$
extracts the chiral algebra from its bar complex (a closed-colour
operation on $\mathrm{FM}_k(\mathbb{C})$).  The Dynkin projection
extracts the Lie element from the tensor algebra (an open-colour
operation on $\mathrm{Conf}_n(\mathbb{R})$).  Both are extraction
maps from a cofree object to the generating algebra, controlled
by the same combinatorial identity: the logarithm of the ordered
exponential.  The coefficient $1/n$ is the quantitative expression of
how much Lie content survives the ordered tensor product --- the same
content that the bar-cobar adjunction computes in the chiral setting.
\end{remark}

\subsection{Complete strictification}

\begin{theorem}[Complete strictification for all simple Lie algebras;
\ClaimStatusProvedHere]
\label{thm:complete-strictification-v1}
Let $\fg$ be a simple Lie algebra of any type.  The spectral Drinfeld
class vanishes at every filtration:
\[
\mathfrak{D}_{\mathrm{spec}}(A,\Omega)\big|_{F^p}=0
\qquad\text{for all }p\ge 2.
\]
In particular, whenever a residue-bounded complete dg-shifted Yangian
over $\fg$ carries a spectral quasi-factorization datum, that datum
strictifies to an honest spectral factorization realization.
The dg-shifted Yangian $Y_\hbar(\fg)$ carries a strict filtered
algebra structure.
\end{theorem}

\begin{proof}
The proof proceeds by analysing the spectral Drinfeld class sector by
sector.  At filtration $p$, each root sector of type
$\alpha=\sum_{i=1}^p\alpha_i$ (where the $\alpha_i$ are simple roots)
contributes independently.

\medskip\noindent\textbf{Case~1: $\alpha$ is not a root.}
If $\alpha=\sum\alpha_i$ is not a root of $\fg$, then the root space
$\fg_\alpha=0$ and every multilinear Lie monomial in
$\{E_{\alpha_i}\}$ vanishes.  The sector contributes nothing.

\medskip\noindent\textbf{Case~2: $\alpha$ is a root (path-type
sectors).}
The spectral Drinfeld class restricted to the root sector lives in
$\fg_\alpha$.  By Theorem~\ref{thm:root-space-one-dim-v1},
$\dim\fg_\alpha=1$.  The multilinear BCH computation
(Theorem~\ref{thm:dynkin-beta-integral}) gives a unique coefficient
$1/p$ for the product side.  The coproduct rigidity gives a unique
equation for the coproduct side.  Since both sides are scalars in the
same one-dimensional space, equality reduces to a single identity ---
the same identity that holds at filtrations $2$ and $3$ --- and the
obstruction class vanishes.

\medskip\noindent\textbf{Case~3: Non-path sectors (e.g., the $D_4$
star).}
Root-space one-dimensionality
(Theorem~\ref{thm:root-space-one-dim-v1}) forces the multilinear
BCH contribution to be a scalar multiple of the unique generator of
$\fg_\alpha$.  The Jacobi collapse lemma
(Lemma~\ref{lem:jacobi-collapse-v1}) ensures that the
two-dimensional free Lie space collapses to one dimension in
$\fg$, so the only surviving contribution is the iterated bracket
$[E_{\alpha_1},[E_{\alpha_2},\ldots]]$, whose coefficient in the
BCH expansion is $1/p$ by the Dynkin formula.

\medskip\noindent\textbf{The inductive strictification.}
Therefore no root sector at any filtration carries an independent
obstruction class.  The spectral Drinfeld class vanishes at every
filtration.  Starting from any spectral quasi-factorization datum, a
filtered twist kills the first nontrivial quasi-associator term at
each filtration.  Iterating filtration by filtration produces a twist
with trivial quasi-associator, hence an honest spectral factorization
realization.
\end{proof}

\begin{remark}[The role of root multiplicity one]
The theorem has the form of a
\emph{root-multiplicity-one rigidity principle}: the spectral
Drinfeld obstruction to strictifying a dg-shifted Yangian vanishes
because every root-sector multilinear space is one-dimensional.

The bar complex presents the Swiss-cheese coalgebra
(Theorem~\ref{thm:master}).  The complete strictification
theorem says that, once the spectral quasi-factorization package has
been constructed, the passage from quasi to strict is not an additional
datum but an automatic consequence of the root structure.  Just as the
Decomposition Theorem (purity of IC complexes) makes the Hecke algebra
arise inevitably from the geometry of $G/B$, root multiplicity one
makes the strictification step arise inevitably from the Lie theory
of~$\fg$.
\end{remark}

\begin{remark}[What fails for Kac--Moody algebras with root multiplicities $>1$]
\label{rem:kac-moody-failure}
For Kac--Moody algebras whose root lattice contains roots with
multiplicity $>1$ (e.g., imaginary roots of indefinite Kac--Moody
algebras), root-space one-dimensionality fails.  The Jacobi collapse
mechanism no longer forces multilinear Lie monomials to be
proportional, and the spectral Drinfeld obstruction class can live in
a space of dimension $>1$.  In such cases, the coproduct rigidity
provides fewer equations than unknowns, and genuine obstructions
could survive.

The frontier admits a precise stratification: the obstruction vanishes
in all root sectors where root multiplicity one holds (all roots of
simple Lie algebras, real roots of affine Lie algebras), and the
remaining sectors are exactly those with higher multiplicity.

For untwisted affine Lie algebras $\widehat\fg$, the real roots
$\alpha+n\delta$ have multiplicity~$1$ (so the theorem applies
verbatim), while the imaginary roots $n\delta$ have
multiplicity $\ell=\operatorname{rank}(\fg)$.  The imaginary sectors
are abelian (elements within a single $\fg_{n\delta}$ commute), and
a Cartan gauge-twist argument should absorb the defect.  The mixed
real-plus-imaginary sectors require a separate analysis.  Complete
strictification for affine types therefore remains conjectural pending
the detailed imaginary-sector computation.
\end{remark}


\section{Worked example: \texorpdfstring{$\mathfrak{sl}_3$}{affine sl(3)} at level \texorpdfstring{$k$}{k}}
\label{sec:sl3-ordered-bar}

Theorem~\ref{thm:km-yangian} applies to every simple Lie algebra.
For $\mathfrak{g}=\mathfrak{sl}_3$ in
explicit matrix-unit coordinates, the rank-$2$ phenomena
invisible in the rank-$1$ ($\mathfrak{sl}_2$) case:
adjacent-root triangle sectors, the coefficient~$\frac12$,
the Serre relations from root-space vanishing, and the $9\times 9$
$R$-matrix.

\subsection{Root system and matrix-unit notation}

The Lie algebra $\mathfrak{sl}_3$ has root system
$\Phi=\{\pm\alpha_1,\pm\alpha_2,\pm(\alpha_1+\alpha_2)\}$
of type~$A_2$, with simple roots $\alpha_1$, $\alpha_2$ and
Cartan matrix
\[
A=\begin{pmatrix}2&-1\\-1&2\end{pmatrix}.
\]
In the $\mathfrak{gl}_3$ matrix-unit basis we set
\[
e_1=E_{12},\quad
e_2=E_{23},\quad
e_{12}=[e_1,e_2]=E_{13},
\]
\[
f_1=E_{21},\quad
f_2=E_{32},\quad
f_{12}=[f_2,f_1]=E_{31},
\]
and $h_1=E_{11}-E_{22}$, $h_2=E_{22}-E_{33}$.  The composite root
generators satisfy
\[
[e_1,e_2]=e_{12},
\qquad
[f_2,f_1]=f_{12},
\qquad
[f_1,f_2]=-f_{12}.
\]
The affine vertex algebra $V_k(\mathfrak{sl}_3)$ has OPE
\[
J^a(z)\,J^b(w)\sim\frac{k\,\kappa(a,b)}{(z-w)^2}+\frac{[a,b](w)}{z-w}
\]
where $\kappa$ is the Killing form normalised so that
$\kappa(e_i,f_i)=1$ and $\kappa(e_{12},f_{12})=1$.

\subsection{Degree-\texorpdfstring{$1$}{1} generators: the Drinfeld currents}

The ordered bar complex
${\Barch}^{\mathrm{ord}}_\bullet(V_k(\mathfrak{sl}_3))$
has degree-$1$ generators
$\susp^{-1}e_1$, $\susp^{-1}e_2$, $\susp^{-1}e_{12}$,
$\susp^{-1}f_1$, $\susp^{-1}f_2$, $\susp^{-1}f_{12}$,
$\susp^{-1}h_1$, $\susp^{-1}h_2$.
The corresponding bar cohomology classes organise into
Drinfeld generating functions for $Y_\hbar(\mathfrak{sl}_3)$
with $\hbar = 1/(k+3)$:
\begin{align*}
E_i(u) &= \sum_{r\ge 0} E_i^{(r)}\,u^{-r-1},
\qquad i=1,2,\\
F_i(u) &= \sum_{r\ge 0} F_i^{(r)}\,u^{-r-1},\\
H_i(u) &= 1+\sum_{r\ge 0} H_i^{(r)}\,u^{-r-1},
\end{align*}
where $u=1/z$ is the spectral parameter.  The zeroth modes
$E_i^{(0)}$, $F_i^{(0)}$, $H_i^{(0)}$ generate a copy of
$U(\mathfrak{sl}_3)$ inside $Y_\hbar(\mathfrak{sl}_3)$.

\subsection{Degree-\texorpdfstring{$2$}{2} relations: Cartan matrix and Yangian deformation}

The bar differential~\eqref{eq:km-bar-d2} specialises to
\begin{align*}
d_{\Barch}(\susp^{-1}h_i\otimes \susp^{-1}e_j)
&=A_{ij}\,\susp^{-1}e_j,\\
d_{\Barch}(\susp^{-1}e_i\otimes \susp^{-1}f_j)
&=\delta_{ij}\,\susp^{-1}h_i,\\
d_{\Barch}(\susp^{-1}e_1\otimes \susp^{-1}e_2)
&=\susp^{-1}e_{12},\\
d_{\Barch}(\susp^{-1}f_2\otimes \susp^{-1}f_1)
&=\susp^{-1}f_{12}.
\end{align*}
The first two lines reproduce the Cartan relations
$[H_i,E_j]=A_{ij}E_j$ and $[E_i,F_j]=\delta_{ij}H_i$
at leading order in the spectral parameter.  The last two express
the composite root as a bar boundary.  The Yangian deformation
appears at the next spectral order:
\[
[H_i(u),\,E_j(v)]
=
\frac{A_{ij}}{u-v}\,E_j(v)+\text{(regular)},
\qquad
[E_i(u),\,F_j(v)]
=
\frac{\delta_{ij}}{u-v}\,H_i(v)+\text{(regular)}.
\]

\subsection{Degree-\texorpdfstring{$3$}{3} triangle sectors}

This is the first genuinely rank-$2$ sector: the composite root
$\alpha_1+\alpha_2$ produces degree-$3$ bar elements that
have no $\mathfrak{sl}_2$ analogue.

\begin{definition}[Triangle sectors for $\mathfrak{sl}_3$]
\label{def:sl3-triangle-sectors}
The degree-$3$ ordered bar complex
${\Barch}^{\mathrm{ord}}_3(V_k(\mathfrak{sl}_3))$
contains four \emph{triangle sectors}, indexed by the composite
root $\alpha_1+\alpha_2$ and the two orderings of simple lowering
(respectively raising) operators:
\begin{align*}
T^L_{12}&=E_{13}^{(1)}\,E_{21}^{(2)}\,E_{32}^{(3)},
&
T^L_{21}&=E_{13}^{(1)}\,E_{32}^{(2)}\,E_{21}^{(3)},\\
T^R_{12}&=E_{12}^{(1)}\,E_{23}^{(2)}\,E_{31}^{(3)},
&
T^R_{21}&=E_{23}^{(1)}\,E_{12}^{(2)}\,E_{31}^{(3)}.
\end{align*}
The left triangles have the composite raising operator $e_{12}=E_{13}$ in
the first tensor factor and the two simple lowering operators split
across factors~$2$ and~$3$.  The right triangles have the two simple
raising operators in factors~$1$ and~$2$ and the composite lowering
operator $f_{12}=E_{31}$ in factor~$3$.
\end{definition}

\begin{theorem}[Triangle coefficient for $\mathfrak{sl}_3$; \ClaimStatusProvedHere]
\label{thm:sl3-triangle-coefficient}
Let the two-body spectral functions be given by the rational ansatz
$\rho_{\alpha}(u)=c_\alpha/u$, and define the triangle transport
coefficient $\lambda$ by
\[
\Pi_{\langle T^L_{12},\,T^L_{21}\rangle}
\big((\id\otimes\Delta_z)\,X_{\alpha_1+\alpha_2}^{12}(w)\big)
=
\lambda\big(\rho_1(w)\rho_2(w{+}z)\,T^L_{12}
-\rho_2(w)\rho_1(w{+}z)\,T^L_{21}\big).
\]
Then the triangle coefficient is
\[
\boxed{%
\lambda
=
\frac{c_1\,c_2}{2\,c_{12}}
=
\frac{1}{2}
=
\int_0^1(1-t)\,dt,
}
\]
where the last equality identifies the coefficient as the $n=2$ Euler
beta integral $B(1,2)=\int_0^1(1-t)^{2-1}\,dt=1/2$.
In the normalised case $c_1=c_2=c_{12}=1$, the coefficient is~$\frac12$
for both left and right triangles and both orderings:
$\lambda^L_{12}=\lambda^L_{21}=\lambda^R_{12}=\lambda^R_{21}=\frac12$.
\end{theorem}

\begin{proof}
The product side of the filtration-$2$ hexagon defect equation projects
to the triangle sectors by commutator extraction.  The only terms that
contribute are
\[
[E_1^{(1)}F_1^{(2)},\,E_2^{(1)}F_2^{(3)}]
=E_{12}^{(1)}\,F_1^{(2)}\,F_2^{(3)}
\]
and
\[
[E_2^{(1)}F_2^{(2)},\,E_1^{(1)}F_1^{(3)}]
=-E_{12}^{(1)}\,F_2^{(2)}\,F_1^{(3)},
\]
where the signs arise from
$[e_1,e_2]=e_{12}$ and $[e_2,e_1]=-e_{12}$.
The logarithm of the ordered product of two-body exponentials gives
\begin{align*}
\Pi_{\langle T^L_{12},\,T^L_{21}\rangle}
\tfrac{1}{2}[S_{12}(w),\,S_{13}(w{+}z)]
&=
\tfrac{1}{2}\rho_1(w)\rho_2(w{+}z)\,T^L_{12}
-\tfrac{1}{2}\rho_2(w)\rho_1(w{+}z)\,T^L_{21}.
\end{align*}
On the coproduct side, $X_{\alpha_1+\alpha_2}$ is the unique composite
root generator that can land in the triangle sector.  The twisted
coproduct on $X_{\alpha_1+\alpha_2}^{12}(w)$ projects to the triangle
sector with transport coefficient $\lambda$.  Equating the two sides and
substituting the rational ansatz $\rho_\alpha(u)=c_\alpha/u$ gives
\[
\frac{c_{12}}{w}\cdot\frac{\lambda}{w+z}
=
\frac{1}{2}\cdot\frac{c_1}{w}\cdot\frac{c_2}{w+z},
\]
whence $\lambda=c_1 c_2/(2c_{12})$.  In the normalised case this
is~$1/2$.  The integral representation follows from
$B(1,2)=\int_0^1(1-t)\,dt=1/2$.
\end{proof}

\begin{remark}[The triangle coefficient as BCH coefficient]
\label{rem:triangle-as-bch}
The coefficient $1/2$ is the $n=2$ instance of the universal $n$-gon
coefficient $1/n$ proved in Vol~II (the $n$-gon localisation formula).
At $n=2$ the ordered BCH expansion
$\log(e^X e^Y)=X+Y+\frac12[X,Y]+\cdots$
has $\frac12$ as the coefficient of the first commutator term.  This is
not a coincidence: the triangle sector is the first commutator sector
in the ordered bar complex.  The integral $\int_0^1(1-t)^{n-1}\,dt=1/n$
provides a uniform derivation for all~$n$, replacing case-by-case
Baker--Campbell--Hausdorff algebra.
\end{remark}

\subsection{Serre relations from root-space vanishing}

\begin{proposition}[Serre relations from root-space vanishing; \ClaimStatusProvedHere]
\label{prop:sl3-serre}
The Serre relations of $Y_\hbar(\mathfrak{sl}_3)$ arise from degree-$3$
bar cohomology by root-space obstruction.  Specifically:
\begin{enumerate}[label=(\roman*)]
\item $2\alpha_1+\alpha_2\notin\Phi(\mathfrak{sl}_3)$, so
$(\mathfrak{sl}_3)_{2\alpha_1+\alpha_2}=0$, which forces
$[e_1,[e_1,e_2]]=0$.
\item $\alpha_1+2\alpha_2\notin\Phi(\mathfrak{sl}_3)$, so
$(\mathfrak{sl}_3)_{\alpha_1+2\alpha_2}=0$, which forces
$[e_2,[e_2,e_1]]=0$.
\end{enumerate}
In the ordered bar complex, these correspond to the vanishing of
the degree-$3$ bar differential on the would-be generators
$\susp^{-1}e_1\otimes \susp^{-1}e_1\otimes \susp^{-1}e_2$ and
$\susp^{-1}e_2\otimes \susp^{-1}e_2\otimes \susp^{-1}e_1$:
the target root spaces are empty, so the bar differential has
trivial image and the corresponding bar cohomology classes are zero.
The spectral-parameter deformation promotes these to the Yangian
Serre relations:
\[
\operatorname{Sym}_{u_1,u_2}
\big[E_i(u_1),\,[E_i(u_2),\,E_j(v)]\big]=0,
\qquad
A_{ij}=-1.
\]
\end{proposition}

\begin{proof}
Consider the degree-$3$ bar element
$\susp^{-1}e_1\otimes \susp^{-1}e_1\otimes \susp^{-1}e_2$.
The bar differential evaluates the nested commutator
$[e_1,[e_1,e_2]]=[e_1,e_{12}]$.  But $e_1\in(\mathfrak{sl}_3)_{\alpha_1}$
and $e_{12}\in(\mathfrak{sl}_3)_{\alpha_1+\alpha_2}$, so the commutator
would live in $(\mathfrak{sl}_3)_{2\alpha_1+\alpha_2}$.  Since
$2\alpha_1+\alpha_2\notin\Phi(\mathfrak{sl}_3)$ (the root system of
$A_2$ contains only roots of height $\le 2$), this root space is zero,
hence $[e_1,e_{12}]=0$.  The identical argument with $\alpha_1\leftrightarrow\alpha_2$
gives $[e_2,e_{21}]=[e_2,[e_2,e_1]]=0$.

For the Yangian deformation: the spectral OPE
$E_i(u)E_i(v)$ has a pole shifted to $u-v=\hbar$ (off the collision
diagonal), so the collision residue vanishes and
$d_{\mathrm{res}}(\susp^{-1}E_i\otimes \susp^{-1}E_i)=0$.  The degree-$3$
bar cocycle condition $d^2=0$ then forces the symmetrised nested
commutator to vanish, giving the Yangian Serre relation.
\end{proof}

\subsection{The RTT presentation and the \texorpdfstring{$9\times 9$}{9x9} R-matrix}

\begin{theorem}[RTT presentation for $Y_\hbar(\mathfrak{sl}_3)$
from the ordered bar complex; \ClaimStatusProvedHere]
\label{thm:sl3-rtt}
Let $V=\bC^3$ be the defining representation of
$\mathfrak{sl}_3$ and define the transfer matrix
\[
T(u)
=\sum_{a,b=1}^{3}e_{ab}\otimes t^{ab}(u)
\in \operatorname{End}(V)\otimes Y_\hbar(\mathfrak{sl}_3)[\![u^{-1}]\!],
\]
where $t^{ab}(u)$ are the degree-$1$ bar cohomology generating
functions.  Then the consistency of the degree-$3$ ordered bar complex
${\Barch}^{\mathrm{ord}}_3$ gives the RTT relation
\begin{equation}\label{eq:sl3-rtt}
R(u{-}v)\,T_1(u)\,T_2(v)
=T_2(v)\,T_1(u)\,R(u{-}v),
\end{equation}
where $R(u)\in\operatorname{End}(V\otimes V)$ is the rational
Yang $R$-matrix
\begin{equation}\label{eq:sl3-R-matrix}
R(u)=1+\frac{\hbar\,\Omega}{u}+O(u^{-2}),
\qquad
\Omega=\sum_{a,b=1}^{3}e_{ab}\otimes e_{ba},
\end{equation}
acting on $V\otimes V\cong\bC^9$.  Here
$\Omega$ is the $\mathfrak{sl}_3$-Casimir in the tensor square of
the defining representation
\textup{(}Remark~\textup{\ref{rem:casimir-types}(a)} with $n=2$\textup{)},
and $e_{ab}$ are the standard matrix units of $\mathfrak{gl}_3$.
\end{theorem}

\begin{proof}
This is the $A_2$ specialisation of Theorem~\ref{thm:km-yangian}.
The $R$-matrix arises from the monodromy of the flat connection
$\nabla=d-r(z)\,dz$ on the ordered configuration space
$\Conf_2^<(\bC)$, where $r(z)=\hbar\,\Omega/z$ is the classical
$r$-matrix~\eqref{eq:km-collision-residue} (with the
$d\log$ kernel absorbing one pole order).
The RTT relation~\eqref{eq:sl3-rtt} is equivalent to $d^2=0$ on
${\Barch}^{\mathrm{ord}}_3$: the three-fold bar consistency
intertwines two orderings of the first two tensor factors via the
$R$-matrix and inserts a bar cohomology class in the third slot.
The Yang--Baxter equation
\[
R_{12}(u{-}v)\,R_{13}(u)\,R_{23}(v)
=R_{23}(v)\,R_{13}(u)\,R_{12}(u{-}v)
\]
follows from $d^3=0$ on ${\Barch}^{\mathrm{ord}}_4$ and is verified
directly for $R(u)=1+\hbar\,\Omega/u$ using the $\mathfrak{sl}_3$
Jacobi identity in the tensor cube $V^{\otimes 3}\cong\bC^{27}$.
\end{proof}

\begin{remark}[The $9\times 9$ versus $4\times 4$ $R$-matrix]
\label{rem:9x9-vs-4x4}
For $\mathfrak{sl}_2$ the $R$-matrix acts on
$\bC^2\otimes\bC^2\cong\bC^4$.
For $\mathfrak{sl}_3$ it acts on
$\bC^3\otimes\bC^3\cong\bC^9$.
The decomposition
$\bC^9\cong\mathbf{6}\oplus\mathbf{3}$
(symmetric and exterior squares) is an eigenspace decomposition
for the permutation $P=\sum e_{ab}\otimes e_{ba}$: eigenvalue
$+1$ on $\operatorname{Sym}^2(V)$ and $-1$ on $\bigwedge^2(V)$.
The $R$-matrix $R(u)=1+\hbar\,P/u$ therefore has eigenvalues
$1+\hbar/u$ and $1-\hbar/u$ on these two subspaces.
The four triangle sectors $T^L_{12}$, $T^L_{21}$, $T^R_{12}$, $T^R_{21}$
are the components of the degree-$3$ bar complex that detect the
off-diagonal ($\alpha_1+\alpha_2$) block of this $R$-matrix.
\end{remark}

\subsection{Strictification for \texorpdfstring{$\mathfrak{sl}_3$}{sl(3)}}

\begin{theorem}[Vanishing of the spectral Drinfeld class for $\mathfrak{sl}_3$;
\ClaimStatusProvedHere]
\label{thm:sl3-strictification}
The spectral Drinfeld class $[\delta]\in H^2(\mathfrak{sl}_3,\,
\operatorname{End}(\mathfrak{sl}_3\otimes\mathfrak{sl}_3))$
vanishes.  Concretely:
\begin{enumerate}[label=\textup{(\roman*)}]
\item At filtration $1$ \textup{(}Cartan sector\textup{)},
the defect is identically zero because the Cartan matrix $A_{ij}$
has integer entries and the rational ansatz is already exact.
\item At filtration $2$ \textup{(}triangle sector\textup{)},
the obstruction lives in the two-dimensional space spanned by
$T^L_{12}$ and $T^L_{21}$.  The triangle coefficient
$\lambda=1/2$ satisfies the rigidity equations of
Theorem~\textup{\ref{thm:sl3-triangle-coefficient}}, so the
defect vanishes.
\item At filtration $\ge 3$, the Serre relation
$[e_1,[e_1,e_2]]=0$ ensures vanishing target.
Root-space one-dimensionality
\textup{(}$\dim(\mathfrak{sl}_3)_\alpha=1$ for all
$\alpha\in\Phi$\textup{)} forces the remaining cochains to be
coboundaries via the Jacobi collapse mechanism
\textup{(Step~$4$ of Theorem~\textup{\ref{thm:km-yangian}})}.
\end{enumerate}
Therefore $[\delta]=0$ at every filtration, and
$Y_\hbar(\mathfrak{sl}_3)$ is a strict associative algebra.
\end{theorem}

\begin{proof}
This is Step~4 of Theorem~\ref{thm:km-yangian} specialised to
$A_2$.  The filtration-by-filtration verification above records
what the general root-space one-dimensionality argument looks like
when all root spaces are enumerated: $\Phi$ has exactly $6$ roots,
all of multiplicity one.
\end{proof}

\begin{remark}[Contrast with root multiplicities $>1$]
\label{rem:root-multiplicities}
The vanishing of $[\delta]$ uses $\dim(\mathfrak{sl}_3)_\alpha=1$
in an essential way.  For Kac--Moody algebras with root multiplicities
$>1$ (e.g.\ the Monster Lie algebra), the endomorphism space at a
root of multiplicity~$m$ is $\operatorname{End}(\bC^m)$, and
the Jacobi collapse argument fails: the obstruction cochain has
non-scalar entries that need not be coboundaries.
\end{remark}

\section{Higher-rank pattern: the \texorpdfstring{$\mathfrak{sl}_4$}{sl(4)} root quadrilateral}
\label{sec:sl4-quadrilateral}

For $\mathfrak{sl}_4$ the Dynkin diagram is the path
$\alpha_1\text{---}\alpha_2\text{---}\alpha_3$ of type $A_3$, with
positive roots
\[
\Phi^+=\{\alpha_1,\;\alpha_2,\;\alpha_3,\;
\alpha_1{+}\alpha_2,\;\alpha_2{+}\alpha_3,\;
\alpha_1{+}\alpha_2{+}\alpha_3\}.
\]
In matrix-unit notation: $e_1=E_{12}$, $e_2=E_{23}$, $e_3=E_{34}$,
$e_{12}=E_{13}$, $e_{23}=E_{24}$, $e_{123}=E_{14}$.

\begin{definition}[Quadrilateral sectors]
\label{def:sl4-quadrilateral}
The degree-$4$ ordered bar complex exhibits \emph{quadrilateral sectors}
associated to the maximal composite root $\alpha_1+\alpha_2+\alpha_3$.
The left quadrilateral vector is
\[
Q^L_{123}
=E_{14}^{(1)}\,E_{21}^{(2)}\,E_{32}^{(3)}\,E_{43}^{(4)},
\]
with permuted orderings of the simple lowering operators in
factors $2$, $3$, $4$.
\end{definition}

\begin{theorem}[Quadrilateral coefficient for $\mathfrak{sl}_4$; \ClaimStatusProvedHere]
\label{thm:sl4-quadrilateral}
Under the normalised rational ansatz, the quadrilateral transport
coefficient is
\[
\boxed{%
\frac{1}{3}
=
\int_0^1(1-t)^2\,dt
=
B(1,3),
}
\]
the $n=3$ instance of the universal $n$-gon coefficient $1/n$.
\end{theorem}

\begin{proof}
By the same method as Theorem~\ref{thm:sl3-triangle-coefficient},
the product side of the filtration-$3$ hexagon defect gives
\[
\frac{1}{3}\,[X_1,\,[X_2,\,X_3]],
\]
where the coefficient $1/3$ arises from the Dynkin path
BCH formula: $\log(e^{X_1}e^{X_2}e^{X_3})$ restricted to the
multilinear part of the path $\alpha_1\sim\alpha_2\sim\alpha_3$
gives $\frac{1}{3}[X_1,[X_2,X_3]]$ (using $[X_1,X_3]=0$ since
$\alpha_1$ and $\alpha_3$ are non-adjacent).  The integral
representation is $\int_0^1(1-t)^2\,dt=\frac{1}{3}$.

On the coproduct side, $X_{\alpha_1+\alpha_2+\alpha_3}$ is the
unique height-$3$ generator whose twisted coproduct projects to the
quadrilateral sector.  Equating coefficients gives the claim.

The Serre relations for $\mathfrak{sl}_4$ provide further
constraints: $(\operatorname{ad}e_i)^{1-A_{ij}}(e_j)=0$ with
$A_{ij}=-1$ for adjacent $i,j$ and $A_{13}=0$.
The last condition gives $[e_1,e_3]=0$: non-adjacent simple
roots commute.

The strictification obstruction vanishes for $\mathfrak{sl}_4$ by
root-space one-dimensionality (every root of $A_3$ has
multiplicity~$1$), exactly as for $\mathfrak{sl}_3$.
\end{proof}

\begin{remark}[The beta-integral staircase]
\label{rem:beta-staircase}
The pattern of $n$-gon coefficients
\[
\frac{1}{n}=\int_0^1(1-t)^{n-1}\,dt=B(1,n),
\qquad n=2,3,4,\ldots,
\]
gives a staircase indexed by Dynkin path length:
\[
\begin{array}{c|cccc}
\mathfrak{sl}_N & N=2 & N=3 & N=4 & N=5 \\
\hline
\text{path length }n & 1 & 2 & 3 & 4 \\
\text{coefficient} & 1 & 1/2 & 1/3 & 1/4
\end{array}
\]
At $n=1$ (the $\mathfrak{sl}_2$ case), there is no triangle sector
and the coefficient is~$1$: this is the degree-$2$ bar differential,
which is the Lie bracket itself.  The genuine rank-$2$ phenomenon
begins at $n=2$ with the coefficient $1/2$.  Each step up the
rank ladder introduces one additional nested commutator and one
additional power of $(1-t)$ in the integrand.  The beta integral
provides a single uniform computation that replaces all the
case-by-case Baker--Campbell--Hausdorff algebra.

For general simple $\mathfrak{g}$, the maximal Dynkin path
length is the rank~$r$, and the deepest $n$-gon coefficient is
$1/r$.  The strictification obstruction vanishes for all simple
Lie algebras by root-space one-dimensionality (every root in a
finite root system has multiplicity~$1$), confirming
Theorem~\ref{thm:km-yangian}.
\end{remark}


\subsection{Explicit computations: types $B_2$, $C_2$, $G_2$}
\label{subsec:BCG-yangians}

Table~\ref{tab:km-yangian-data} records summary data for all types;
we now work out the three rank-$2$ families beyond type~$A$ in detail,
exhibiting their root systems, collision residues, Serre relations,
$R$-matrices in the fundamental representation, and
strictification coefficients.

\subsubsection{Type $B_2$: $\mathfrak{so}_5$ and the $5$-dimensional
Yangian}
\label{subsubsec:B2-yangian}

\begin{computation}[Root data for $B_2$; \ClaimStatusProvedHere]
\label{comp:B2-root-data}
The simple roots are $\alpha_1$ (long) and $\alpha_2$ (short), with
$|\alpha_1|^2/|\alpha_2|^2=2$.  The Cartan matrix is
\[
A(B_2)
\;=\;
\begin{pmatrix} 2 & -1 \\ -2 & 2 \end{pmatrix},
\]
so $A_{12}=-1$, $A_{21}=-2$.  The positive roots are
$\alpha_1$, $\alpha_2$, $\alpha_1{+}\alpha_2$,
$\alpha_1{+}2\alpha_2$.
The dual Coxeter number is $h^\vee=3$.  All root
multiplicities equal~$1$: $\dim(\mathfrak{so}_5)_\alpha=1$ for
every root~$\alpha$.  The Lie algebra has dimension~$10$: a
$2$-dimensional Cartan subalgebra and $4$ positive $+$ $4$
negative root spaces.
\end{computation}

\begin{theorem}[Ordered bar complex and Yangian $R$-matrix for $B_2$;
\ClaimStatusProvedHere]
\label{thm:B2-ordered-bar}
For $\widehat{\mathfrak{so}}_5$ at level~$k\neq -3$, with
$\hbar=1/(k+3)$:
\begin{enumerate}[label=\textup{(\alph*)}]
\item The collision residue is $r(z)=\hbar\,\Omega_{B_2}/z$, where
$\Omega_{B_2}\in\mathfrak{so}_5\otimes\mathfrak{so}_5$ is the
Casimir element.  In the $5$-dimensional (vector) representation~$V$,
the Casimir image decomposes according to the
$\mathfrak{so}_5$-equivariant splitting
$V\otimes V\cong \mathbb{C}\oplus\Lambda^2V\oplus S^2_0V$
into the trivial, adjoint, and traceless symmetric summands.
Writing $P$ for the permutation operator and $K$ for the
trace contraction ($K_{ij,kl}=\delta_{i,\bar l}\delta_{j,\bar k}$
with $\bar\imath = 6 - i$ in the standard ordering):
\begin{equation}\label{eq:B2-casimir-image}
\Omega_V
\;=\;
P - K.
\end{equation}

\item The $R$-matrix in the $5$-dimensional representation is
\begin{equation}\label{eq:B2-R-matrix}
R(u)
\;=\;
1
\;+\;
\frac{\hbar}{u}\,P
\;-\;
\frac{\hbar}{u+\hbar\kappa}\,K,
\end{equation}
where $\kappa=(N{-}2)/2=3/2$ for $N=5$.
This is the Zamolodchikov--Fateev $R$-matrix of type~$B$
(equivalently the $\mathfrak{so}_N$ Yangian $R$-matrix).
The poles at $u=0$ and $u=-\hbar\kappa$ correspond to the
$P$ and $K$ channels respectively.

\item \emph{Serre relations from root-space vanishing.}
The degree-$3$ sector $\alpha_1{+}2\alpha_2$: the iterated bracket
$[e_2,[e_1,e_2]]$ is nonzero (since $\alpha_1{+}2\alpha_2$ is
a root) and unique up to scale by root-space one-dimensionality.
The Serre relation
$[e_2,[e_2,[e_2,e_1]]]=0$ is enforced by
$(\mathfrak{so}_5)_{3\alpha_2+\alpha_1}=0$.

\item \emph{Strictification.}
Root-space one-dimensionality forces complete strictification.
The beta-integral coefficient at filtration $p$ is
$1/p = \int_0^1(1-t)^{p-1}\,dt$.  At filtration~$2$:
$1/2=\int_0^1(1-t)\,dt$.
The maximum root height is~$3$ ($\alpha_1{+}2\alpha_2$),
so filtrations $2$ and $3$ are the nontrivial ones.

\item \emph{RTT.}
The relation
$R_{12}(u{-}v)\,T_1(u)\,T_2(v)
=T_2(v)\,T_1(u)\,R_{12}(u{-}v)$
with $T(u)\in\operatorname{End}(\mathbb{C}^5)\otimes
Y_\hbar(\mathfrak{so}_5)\llbracket u^{-1}\rrbracket$
and $R(u)$ as in~\eqref{eq:B2-R-matrix}
generates the Yangian $Y_\hbar(\mathfrak{so}_5)$.
\end{enumerate}
\end{theorem}

\begin{proof}
The collision residue and root-space one-dimensionality follow from
Theorem~\ref{thm:km-yangian} applied to
$\mathfrak{g}=\mathfrak{so}_5$.  The $R$-matrix in the vector
representation is the Zamolodchikov--Fateev solution, which is the
unique Yang--Baxter solution with the correct classical limit
$R(u)=1+\hbar\,\Omega_V/u+O(u^{-2})$.  The $K$-term
reflects the existence of the invariant bilinear form on~$V$, which
is specific to orthogonal (and symplectic) types and absent in
type~$A$.  The Serre relation and strictification are verified
by direct root-space analysis.
\end{proof}

\subsubsection{Type $C_2$: $\mathfrak{sp}_4$ and the $4$-dimensional
Yangian}
\label{subsubsec:C2-yangian}

Type $C_2$ is Langlands dual to $B_2$: the roles of the long and
short roots are interchanged.

\begin{computation}[Root data for $C_2$; \ClaimStatusProvedHere]
\label{comp:C2-root-data}
The simple roots are $\alpha_1$ (short) and $\alpha_2$ (long), with
$|\alpha_2|^2/|\alpha_1|^2=2$.  The Cartan matrix is
\[
A(C_2)
\;=\;
\begin{pmatrix} 2 & -2 \\ -1 & 2 \end{pmatrix},
\]
so $A_{12}=-2$, $A_{21}=-1$.
The positive roots are $\alpha_1$, $\alpha_2$,
$\alpha_1{+}\alpha_2$, $2\alpha_1{+}\alpha_2$.  The dual Coxeter
number is $h^\vee=3$.  All root multiplicities equal~$1$:
$\dim(\mathfrak{sp}_4)_\alpha=1$ for every root~$\alpha$.
The Lie algebra has dimension~$10$.
\end{computation}

\begin{theorem}[Ordered bar complex and Yangian $R$-matrix for $C_2$;
\ClaimStatusProvedHere]
\label{thm:C2-ordered-bar}
For $\widehat{\mathfrak{sp}}_4$ at level~$k\neq -3$, with
$\hbar=1/(k+3)$:
\begin{enumerate}[label=\textup{(\alph*)}]
\item The collision residue is $r(z)=\hbar\,\Omega_{C_2}/z$, where
$\Omega_{C_2}\in\mathfrak{sp}_4\otimes\mathfrak{sp}_4$ is the
Casimir element.  In the $4$-dimensional (fundamental) representation~$V$
of~$\mathfrak{sp}_4$, the Casimir image involves the symplectic
form~$J$:
\begin{equation}\label{eq:C2-casimir-image}
\Omega_V
\;=\;
P - K_J,
\end{equation}
where $P$ is the permutation operator on $V\otimes V$ and
$(K_J)_{ij,kl}=J_{il}\,J_{jk}$ is the contraction through the
symplectic form $J$ with
$J=\bigl(\begin{smallmatrix} 0 & I_2 \\ -I_2 & 0
\end{smallmatrix}\bigr)$.

\item The $R$-matrix in the $4$-dimensional representation is
\begin{equation}\label{eq:C2-R-matrix}
R(u)
\;=\;
1
\;+\;
\frac{\hbar}{u}\,P
\;-\;
\frac{\hbar}{u+\hbar\kappa}\,K_J,
\end{equation}
where $\kappa=n{+}1=3$ for $\mathfrak{sp}_4$
(with $n=2$ in $\mathfrak{sp}_{2n}$).
This is the Zamolodchikov--Fateev $R$-matrix of type~$C$.
The pole at $u=-\hbar\kappa$ is at the \emph{opposite sign}
from the orthogonal case~\eqref{eq:B2-R-matrix}: this is the
$R$-matrix--level manifestation of Langlands duality
$B_n\leftrightarrow C_n$.

\item \emph{Serre relations from root-space vanishing.}
The degree-$3$ sector $2\alpha_1{+}\alpha_2$: the iterated bracket
$[e_1,[e_1,e_2]]$ is nonzero (since $2\alpha_1{+}\alpha_2$ is a
root of $C_2$) and unique up to scale.  The Serre relation
$[e_1,[e_1,[e_1,e_2]]]=0$ is enforced by
$(\mathfrak{sp}_4)_{3\alpha_1+\alpha_2}=0$.

\item \emph{Strictification.}
Root-space one-dimensionality forces complete strictification.
Beta-integral coefficient at filtration~$2$: $1/2$.
Maximum root height:~$3$ ($2\alpha_1{+}\alpha_2$).

\item \emph{RTT.}
The relation
$R_{12}(u{-}v)\,T_1(u)\,T_2(v)
=T_2(v)\,T_1(u)\,R_{12}(u{-}v)$
with $T(u)\in\operatorname{End}(\mathbb{C}^4)\otimes
Y_\hbar(\mathfrak{sp}_4)\llbracket u^{-1}\rrbracket$
and $R(u)$ as in~\eqref{eq:C2-R-matrix}
generates the Yangian $Y_\hbar(\mathfrak{sp}_4)$.
\end{enumerate}
\end{theorem}

\begin{proof}
The collision residue and root-space one-dimensionality follow from
Theorem~\ref{thm:km-yangian}.  The $R$-matrix is the symplectic
Zamolodchikov--Fateev solution with classical limit
$R(u)=1+\hbar\,\Omega_V/u+O(u^{-2})$.  The $K_J$-term uses the
\emph{symplectic} contraction (compare the orthogonal $K$-term for
$B_2$).  The sign convention for the $K_J$-pole is opposite to that
of~$B_2$: the orthogonal pole is at $u=-\hbar\kappa$ (negative) and the
symplectic pole is at $u=+\hbar\kappa$ (positive).  This sign difference
between orthogonal and symplectic types is the $R$-matrix--level
manifestation of the Langlands duality $B_n\leftrightarrow C_n$.
\end{proof}

\begin{remark}[Langlands duality at the ordered bar level]
\label{rem:langlands-ordered-bar}
Types $B_2$ and $C_2$ are Langlands dual: $B_2^\vee=C_2$
(the Cartan matrix of $C_2$ is the transpose of that of $B_2$).
At the ordered bar level, this duality interchanges:
\begin{itemize}[nosep]
\item the orthogonal contraction $K$ (trace through the symmetric
form) with the symplectic contraction~$K_J$ (trace through the
antisymmetric form);
\item the pole location $u=-\hbar\kappa$ with $u=+\hbar\kappa$;
\item the vector representation $\mathbb{C}^5$ with the fundamental
$\mathbb{C}^4$.
\end{itemize}
The collision residue $r(z)=\hbar\,\Omega/z$ is invariant (both
Casimirs are normalized by the Killing form), but the finite-dimensional
$R$-matrices in the fundamental representations differ by the
orthogonal/symplectic sign.  The abstract
Yangians $Y_\hbar(\mathfrak{so}_5)$ and
$Y_\hbar(\mathfrak{sp}_4)$ are non-isomorphic algebras with
isomorphic root-space structures ($B_2$ and $C_2$ have the same root
pattern but different root lengths).
\end{remark}

\subsubsection{Type $G_2$: the $7$-dimensional Yangian}
\label{subsubsec:G2-yangian}

Type $G_2$ is the exceptional rank-$2$ case.  It has the most
asymmetric root-length ratio among all simple Lie algebras:
$|\alpha_1|^2/|\alpha_2|^2=1/3$.

\begin{computation}[Root data for $G_2$; \ClaimStatusProvedHere]
\label{comp:G2-root-data}
The simple roots are $\alpha_1$ (short) and $\alpha_2$ (long), with
$|\alpha_2|^2/|\alpha_1|^2=3$.  The Cartan matrix is
\[
A(G_2)
\;=\;
\begin{pmatrix} 2 & -1 \\ -3 & 2 \end{pmatrix},
\]
so $A_{12}=-1$, $A_{21}=-3$.  The positive roots are:
\[
\alpha_1,\quad
\alpha_2,\quad
\alpha_1{+}\alpha_2,\quad
2\alpha_1{+}\alpha_2,\quad
3\alpha_1{+}\alpha_2,\quad
3\alpha_1{+}2\alpha_2.
\]
The dual Coxeter number is $h^\vee=4$.  All root multiplicities
equal~$1$: $\dim(\mathfrak{g}_2)_\alpha=1$ for every root~$\alpha$.
The Lie algebra has dimension~$14$: a $2$-dimensional Cartan
subalgebra and $6$ positive $+$ $6$ negative root spaces.

The root string structure: the $\alpha_1$-string through
$\alpha_2$ has length~$4$ ($\alpha_2$, $\alpha_1{+}\alpha_2$,
$2\alpha_1{+}\alpha_2$, $3\alpha_1{+}\alpha_2$),
while the $\alpha_2$-string through $\alpha_1$ has length~$2$
($\alpha_1$, $\alpha_1{+}\alpha_2$).
The longest root $3\alpha_1{+}2\alpha_2$ has height~$5$.
\end{computation}

\begin{theorem}[Ordered bar complex and Yangian $R$-matrix for $G_2$;
\ClaimStatusProvedHere]
\label{thm:G2-ordered-bar}
For $\widehat{\mathfrak{g}}_2$ at level~$k\neq -4$, with
$\hbar=1/(k+4)$:
\begin{enumerate}[label=\textup{(\alph*)}]
\item The collision residue is $r(z)=\hbar\,\Omega_{G_2}/z$, where
$\Omega_{G_2}\in\mathfrak{g}_2\otimes\mathfrak{g}_2$ is the
Casimir element.

\item In the $7$-dimensional (fundamental) representation~$V$
of~$\mathfrak{g}_2$, the Casimir image is
$\Omega_V=(\rho_V\otimes\rho_V)(\Omega_{G_2})
\in\operatorname{End}(\mathbb{C}^7\otimes\mathbb{C}^7)$.
Since $V$ is the representation on the imaginary octonions, the
image decomposes according to the $G_2$-equivariant
splitting $V\otimes V\cong\mathbb{C}\oplus V\oplus
\mathfrak{g}_2\oplus S^2_0(V)$
(dimensions $1+7+14+27=49=7^2$):
\begin{equation}\label{eq:G2-casimir-image}
\Omega_V
\;=\;
c_1\,P_{\mathbb{C}}
\;+\;
c_7\,P_V
\;+\;
c_{14}\,P_{\mathfrak{g}_2}
\;+\;
c_{27}\,P_{S^2_0},
\end{equation}
where $P_\lambda$ is the $G_2$-equivariant projector onto the
irreducible summand of dimension $\dim\lambda$ and $c_\lambda$
is the eigenvalue of the Casimir on $\lambda$.  With the
normalization fixed by the collision residue, the
Casimir eigenvalues (in units where $c_V=c_7$) are:
\[
c_1=0,\qquad
c_7=\frac{1}{4},\qquad
c_{14}=\frac{1}{2},\qquad
c_{27}=\frac{5}{12}.
\]

\item The $R$-matrix in the $7$-dimensional representation is
\begin{equation}\label{eq:G2-R-matrix}
R(u)
\;=\;
\sum_{\lambda\in\{1,7,14,27\}}
\prod_{\mu\neq\lambda}
\frac{u+\hbar(c_\mu-c_\lambda)}{c_\mu - c_\lambda}
\cdot \frac{1}{u}\,P_\lambda,
\end{equation}
the unique rational Yang--Baxter solution with
the correct classical limit
\begin{equation}\label{eq:G2-R-matrix-classical}
R(u)
\;=\;
1
\;+\;
\frac{\hbar}{u}\,\Omega_V
\;+\;
O(u^{-2}).
\end{equation}
The higher-order terms are determined by the Yang--Baxter
equation and unitarity $R(u)\,R(-u)=1$.  Unlike types $B$ and
$C$, where the $R$-matrix involves two operators ($P$ and
$K$ or $K_J$), the $G_2$ $R$-matrix involves \emph{four}
equivariant projectors, reflecting the richer
$G_2$-equivariant structure of $V\otimes V$.

\item \emph{Serre relations and iterated brackets.}
\begin{itemize}[nosep]
\item \emph{Degree $2$:}
$[e_1,e_2]\neq 0$ since $\alpha_1{+}\alpha_2$ is a root.
Beta coefficient: $1/1=1$.

\item \emph{Degree $3$:}
$[e_1,[e_1,e_2]]\neq 0$ since $2\alpha_1{+}\alpha_2$ is a root.
Unique up to scale by root-space one-dimensionality.
Beta coefficient: $1/2$.

\item \emph{Degree $4$:}
$[e_1,[e_1,[e_1,e_2]]]\neq 0$ since $3\alpha_1{+}\alpha_2$ is a
root.  Unique up to scale.  Beta coefficient: $1/3$.
The Serre relation $[e_1,[e_1,[e_1,[e_1,e_2]]]]=0$ is enforced
by $(\mathfrak{g}_2)_{4\alpha_1+\alpha_2}=0$.

\item \emph{Degree $5$ (longest root):}
The root $3\alpha_1{+}2\alpha_2$ has height~$5$ and involves both
simple roots at multiplicities $3$ and~$2$.  Root-space
one-dimensionality ($\dim(\mathfrak{g}_2)_{3\alpha_1+2\alpha_2}=1$)
forces all multilinear Lie monomials in the simple root vectors
at these multiplicities to be proportional.  The Jacobi collapse
lemma handles the star-sector ambiguity at this degree.
Beta coefficient: $1/5$.
\end{itemize}
The $G_2$ Serre relation from the short root side:
$[e_2,[e_2,e_1]]=0$ since $\alpha_1{+}2\alpha_2$ is \emph{not}
a root of $G_2$, enforced by root-space vanishing
$(\mathfrak{g}_2)_{\alpha_1+2\alpha_2}=0$.  From the long root
side: $[e_1,[e_1,[e_1,[e_1,e_2]]]]=0$ since
$4\alpha_1{+}\alpha_2\notin\Phi(G_2)$.

\item \emph{Strictification.}
Root-space one-dimensionality forces complete strictification.
The highest nontrivial filtration is $p=5$, with coefficient
$1/5=\int_0^1(1-t)^4\,dt$.

\item \emph{RTT.}
The relation
$R_{12}(u{-}v)\,T_1(u)\,T_2(v)
=T_2(v)\,T_1(u)\,R_{12}(u{-}v)$
with $T(u)\in\operatorname{End}(\mathbb{C}^7)\otimes
Y_\hbar(\mathfrak{g}_2)\llbracket u^{-1}\rrbracket$
and $R(u)$ as in~\eqref{eq:G2-R-matrix}
generates the Yangian $Y_\hbar(\mathfrak{g}_2)$.
\end{enumerate}
\end{theorem}

\begin{proof}
The collision residue and root-space one-dimensionality follow from
Theorem~\ref{thm:km-yangian}.  The $R$-matrix in the $7$-dimensional
representation is determined by the $G_2$-equivariant structure of
$V\otimes V$: any $G_2$-equivariant rational $R$-matrix with the
correct classical limit is forced to be a linear combination of the
four projectors $P_\lambda$ with spectral-parameter--dependent
coefficients, and the Yang--Baxter equation plus unitarity
$R(u)\,R(-u)=1$ fix these coefficients uniquely.

The iterated bracket analysis verifies the Serre relations for $G_2$
from root-space vanishing, extending the type~$A$ computation of
Construction~\texttt{constr:dg-shifted-yangian-from-bar} to the
exceptional setting.  The degree-$5$ sector is the richest:
the longest root of $G_2$ has height~$5$, the maximum among
rank-$2$ Lie algebras, so the ordered bar complex has five nontrivial
filtration levels.
\end{proof}

\begin{remark}[$G_2$ and the octonion cross product]
\label{rem:G2-octonion}
The $7$-dimensional representation of $G_2$ is the automorphism
group of the octonions acting on the imaginary octonions.  The
$G_2$-equivariant decomposition
$V\otimes V\cong\mathbb{C}\oplus V\oplus\mathfrak{g}_2\oplus S^2_0(V)$
reflects the octonionic multiplication structure: the $V$-summand
corresponds to the cross product (the imaginary part of octonionic
multiplication), and the $\mathfrak{g}_2$-summand corresponds to
the commutator bracket.  The $R$-matrix~\eqref{eq:G2-R-matrix}
therefore encodes, at the ordered bar level, the octonionic
structure constants.
\end{remark}

\begin{remark}[Comparison across rank-$2$ types]
\label{rem:rank-2-comparison}
The three rank-$2$ families exhibit increasing complexity:
\[
\begin{array}{c|ccccc}
\text{Type} & h^\vee & \dim V_{\mathrm{fund}}
& |\Phi^+| & \text{max height} & R\text{-matrix structure} \\
\hline
A_2 & 3 & 3 & 3 & 2 & P \text{ only} \\
B_2 & 3 & 5 & 4 & 3 & P,\, K \\
C_2 & 3 & 4 & 4 & 3 & P,\, K_J \\
G_2 & 4 & 7 & 6 & 5 & \text{four projectors}
\end{array}
\]
The maximum height equals the highest filtration level at which the
beta-integral coefficient appears.  Types $A_2$, $B_2$, $C_2$ all
have $h^\vee=3$, so the deformation parameter
$\hbar=1/(k+3)$ coincides.  Type $G_2$ has
$\hbar=1/(k+4)$.  In all cases, root-space one-dimensionality forces
complete strictification: the Yangian emerges as the open-colour
Koszul dual without obstruction.

The $R$-matrix complexity increases from type~$A$ (one operator $P$,
since the defining representation has no invariant bilinear form)
through types $B$, $C$ (two operators: $P$ and the contraction
through the invariant form) to $G_2$ (four equivariant projectors,
corresponding to the four irreducible summands of
$V\otimes V$).  This progression reflects the number of irreducible
components of $V_{\mathrm{fund}}\otimes V_{\mathrm{fund}}$: the
more summands, the richer the pole structure of the $R$-matrix,
and the more detailed the spectral information retained by the
ordered bar complex.
\end{remark}

\section{Quantum groups from genus-1 monodromy}
\label{sec:genus1-monodromy}

The preceding sections develop the ordered bar complex and its $R$-matrix
on the genus-zero configuration space $\mathrm{Conf}_n(\mathbb{C})$.
At genus~$1$ the configuration space is replaced by
$\mathrm{Conf}_n(E_\tau)$, where $E_\tau = \mathbb{C}/(\mathbb{Z}
+ \tau\mathbb{Z})$ is the elliptic curve with modulus~$\tau$.
The fundamental group of $\mathrm{Conf}_n(E_\tau)$ is
richer than the pure braid group $P_n$: it contains two
additional generators --- the $A$-cycle and $B$-cycle monodromies ---
and their interaction with the braid generators produces the quantum
group $U_q(\mathfrak{g})$ from the Yangian $Y_\hbar(\mathfrak{g})$.

This section develops the genus-$1$ lift for the affine lineage
$V^k(\mathfrak{g})$, where the Drinfeld--Kohno theorem and the
one-loop collapse mechanism of
Theorem~\ref{thm:affine-monodromy-identification} (Volume~II) combine to
give a rigorous passage from the KZ connection to the KZB connection and
thence to the quantum group.

\subsection{The KZ connection on ordered configurations}
\label{sec:kz-connection}

The genus-$0$ starting point is as follows.  Let $\mathfrak{g}$ be a simple
Lie algebra and $k \in \mathbb{C} \setminus \{-h^\vee\}$ generic.
Choose an orthonormal basis $\{I^a\}$ for~$\mathfrak{g}$ with respect
to the Killing form normalised so that
$\operatorname{tr}(I^a I^b) = \delta^{ab}$.  The Casimir tensor is
$\Omega = \sum_a I^a \otimes I^a \in \mathfrak{g} \otimes \mathfrak{g}$.

\begin{definition}[KZ connection]
\ClaimStatusProvedHere
\label{def:kz-connection}
\index{KZ connection|textbf}
On the trivial bundle over $\mathrm{Conf}_n(\mathbb{C})$ with fibre
$V_1 \otimes \cdots \otimes V_n$ (finite-dimensional
$\mathfrak{g}$-modules), the \emph{Knizhnik--Zamolodchikov connection}
is
\begin{equation}\label{eq:kz}
\nabla_{\mathrm{KZ}}
\;=\;
d
\;-\;
\frac{\hbar}{2\pi i}
\sum_{1 \leq i < j \leq n}
\Omega_{ij}\,d\log(z_i - z_j),
\end{equation}
where $\hbar = \pi i/(k + h^\vee)$ and $\Omega_{ij}$ denotes the
Casimir acting in the $(i,j)$ tensor slots.
\end{definition}

\begin{remark}
The connection~\eqref{eq:kz} is a special case of the Kohno
connection~\eqref{eq:kohno-connection}: for the affine algebra
$V^k(\mathfrak{g})$, the collision residue $r_{ij}(z)$ extracted by
the $d\log$ kernel (cf.\ AP19: the kernel absorbs one pole) is
$r_{ij} = (\hbar/2\pi i)\,\Omega_{ij}$, which is constant in~$z$.
The one-loop collapse of Theorem~\ref{thm:affine-monodromy-identification}
ensures that the higher $A_\infty$ operations $m_k$ ($k \geq 3$) do not
contribute on evaluation modules, so the full superconnection reduces to
the KZ connection.  The flatness $\nabla_{\mathrm{KZ}}^2 = 0$ follows
from $[\Omega_{ij}, \Omega_{ik} + \Omega_{jk}] = 0$, which is the
$\mathfrak{g}$-equivariance of the Casimir.
\end{remark}

\subsection{Genus-0 monodromy: the $R$-matrix}
\label{sec:genus0-monodromy}

At genus~$0$, the fundamental group
$\pi_1(\mathrm{Conf}_n(\mathbb{C}))$ is the pure braid group~$P_n$.
The monodromy of $\nabla_{\mathrm{KZ}}$ around the generator
$\gamma_{ij}$ (in which $z_i$ circles $z_j$) gives the $R$-matrix:
\[
R_{ij}
\;=\;
\operatorname{Mon}_{\gamma_{ij}}(\nabla_{\mathrm{KZ}})
\;=\;
\mathcal{P}\exp\!\Bigl(
\frac{\hbar}{2\pi i}
\oint_{\gamma_{ij}} \Omega_{ij}\,d\log(z_i - z_j)
\Bigr)
\;=\;
e^{\hbar\,\Omega_{ij}}.
\]
This is the \emph{$A$-cycle monodromy}: it encodes how parallel transport
twists when one point winds around another on the genus-$0$ surface.
The Yang--Baxter equation
$R_{12}\,R_{13}\,R_{23} = R_{23}\,R_{13}\,R_{12}$
follows from flatness (eq.~\eqref{eq:cybe-vol1} and its exponentiated form),
as proved in Proposition~\ref{prop:r-matrix-descent-vol1}.

\subsection{The KZB connection on the elliptic curve}
\label{sec:kzb-connection}

At genus~$1$ the base curve is $E_\tau$ and the configuration space
$\mathrm{Conf}_n(E_\tau)$ has fundamental group containing two
distinguished classes of generators beyond the braid generators: the
$A$-cycle (a point traverses the real period~$1$) and the $B$-cycle
(a point traverses the imaginary period~$\tau$).  The KZ connection
must be replaced by the \emph{Knizhnik--Zamolodchikov--Bernard} (KZB)
connection, which encodes the genus-$1$ monodromy of conformal blocks.

\begin{definition}[KZB connection]
\ClaimStatusProvedHere
\label{def:kzb-connection}
\index{KZB connection|textbf}
On the trivial bundle over $\mathrm{Conf}_n(E_\tau)$ with fibre
$V_1 \otimes \cdots \otimes V_n$, the \emph{KZB connection} is
\begin{equation}\label{eq:kzb}
\nabla_{\mathrm{KZB}}
\;=\;
d
\;-\;
\frac{\hbar}{2\pi i}
\sum_{1 \leq i < j \leq n}
\Omega_{ij}\,\wp_1(z_i - z_j,\tau)\,dz_{ij}
\;-\;
\frac{\hbar}{2\pi i}
\sum_{i=1}^n
D_i\,d\tau,
\end{equation}
where $\wp_1(z,\tau) = \partial_z \log\theta_1(z,\tau)$ is the
logarithmic derivative of the Jacobi theta function (the elliptic
analogue of $1/z$), $z_{ij} = z_i - z_j$, and $D_i$ is the modular
Hamiltonian
\[
D_i
\;=\;
\sum_{j \neq i}
\bigl(
\wp(z_{ij},\tau)\,\Omega_{ij}
\;+\;
\text{lower order}
\bigr),
\]
with $\wp(z,\tau) = -\partial_z \wp_1(z,\tau)$ the Weierstrass
$\wp$-function.
\end{definition}

\begin{remark}
The KZB connection is a genus-$1$ deformation of the KZ connection.  In
the degeneration limit $\tau \to i\infty$ (the nodal rational curve),
$\wp_1(z,\tau) \to \pi\cot(\pi z) \to 1/z + O(z)$, and the KZB
connection degenerates to the KZ connection.  The $d\tau$ term,
absent at genus~$0$, records the dependence of conformal blocks on the
modulus~$\tau$ and is the source of the modular anomaly: it measures the
failure of conformal blocks to form a flat bundle over the moduli
space~$\mathcal{M}_{1,n}$.
\end{remark}

\begin{remark}[Bar-complex origin]
\label{rem:kzb-bar}
In the Swiss-cheese framework, the KZB connection arises from the
ordered bar complex on the elliptic curve $E_\tau$ equipped with the
logarithmic $\mathrm{SC}^{ch,top}$ structure.  The collision residue on
$\mathrm{Conf}_n^{\mathrm{ord}}(E_\tau)$ is
$r_{ij}^{E_\tau}(z) = (\hbar/2\pi i)\,\Omega_{ij}\,\wp_1(z,\tau)$,
which replaces the genus-$0$ collision residue $r_{ij} = (\hbar/2\pi i)\,\Omega_{ij}/z$
by its elliptic regularisation.  The additional $d\tau$ connection
component encodes the curvature $\kappa(A) \cdot \omega_1$ from the
Hodge bundle on $\overline{\mathcal{M}}_{1,n}$: the curved bar
structure $d^2 = \kappa(A) \cdot \omega_1$ at genus~$1$ is precisely
the non-flatness of the naive genus-$0$ connection, corrected by the
Bernard term $D_i\,d\tau$.
\end{remark}

\subsection{$B$-cycle monodromy and the quantum group parameter}
\label{sec:b-cycle-monodromy}

The decisive new phenomenon at genus~$1$ is the $B$-cycle monodromy.

\begin{theorem}[Quantum group from $B$-cycle monodromy; \ClaimStatusProvedHere]
\label{thm:b-cycle-quantum-group}
\index{quantum group!from genus-1 monodromy|textbf}
\index{monodromy!B-cycle|textbf}
Let $\mathfrak{g}$ be a simple Lie algebra and $k \in \mathbb{C}
\setminus \{-h^\vee\}$ generic.  Consider the KZB monodromy
representation
\[
\rho_n^{\mathrm{KZB}}
\;\colon\;
\pi_1\bigl(\mathrm{Conf}_n(E_\tau)\bigr)
\;\longrightarrow\;
\operatorname{GL}(V_1 \otimes \cdots \otimes V_n).
\]
\begin{enumerate}[label=\textup{(\roman*)}]
\item \textbf{$A$-cycle monodromy.}
The monodromy around the $A$-cycle generator $\alpha_i$ (in which
$z_i \mapsto z_i + 1$, all other points fixed) gives the $R$-matrix:
\[
\rho_n^{\mathrm{KZB}}(\alpha_i)
\;=\;
R_i(z),
\]
which is the spectral $R$-matrix of
Construction~\textup{\ref{constr:r-matrix-monodromy-vol1}},
now defined via the elliptic collision residue.

\item \textbf{$B$-cycle monodromy.}
The monodromy around the $B$-cycle generator $\beta_i$ (in which
$z_i \mapsto z_i + \tau$, all other points fixed) produces the
quantum group parameter:
\[
\rho_n^{\mathrm{KZB}}(\beta_i)
\;=\;
q^{H_i}
\;\cdot\;
(\text{braiding corrections}),
\]
where
\begin{equation}\label{eq:q-parameter}
q
\;=\;
e^{2\pi i\hbar}
\;=\;
e^{2\pi^2 i/(k + h^\vee)}
\end{equation}
and $H_i$ is the Cartan action in the $i$-th tensor slot.

\item \textbf{Quantum group structure.}
The combined $A$-cycle and $B$-cycle monodromies, together with the braid
monodromies, generate a representation of the quantum group
$U_q(\mathfrak{g})$ on $V_1 \otimes \cdots \otimes V_n$ at the parameter
$q = e^{2\pi i\hbar}$.
\end{enumerate}
\end{theorem}

\begin{proof}
\textbf{(i)} is immediate from the fact that $A$-cycle monodromy around
$z_i \mapsto z_i + 1$ corresponds to the quasi-periodicity of
$\wp_1(z,\tau)$ along the real period, which introduces the same
braid-type contribution as at genus~$0$.

\medskip
\noindent\textbf{(ii).}
The $B$-cycle monodromy arises from $z_i \mapsto z_i + \tau$.  The
quasi-periodicity of the Jacobi theta function gives
$\wp_1(z + \tau, \tau) = \wp_1(z, \tau) - 2\pi i$, so the connection
form acquires a constant shift $-2\pi i \cdot (\hbar/2\pi i)\,\Omega_{ij}
= -\hbar\,\Omega_{ij}$ from each pair.  The parallel transport along the
$B$-cycle path $z_i \mapsto z_i + t\tau$ ($t \in [0,1]$) includes
both the $dz$ contribution (from the connection one-form) and the
$d\tau$ contribution (from the modular Hamiltonian $D_i$).  On
evaluation modules, the one-loop collapse ensures that only the
leading Casimir term contributes, and the result is
\[
\rho_n^{\mathrm{KZB}}(\beta_i)
\;=\;
\exp\bigl(2\pi i \hbar \cdot H_i + \text{off-diagonal}\bigr),
\]
where $H_i$ is the Cartan element acting in the $i$-th slot.  The
exponential $e^{2\pi i \hbar} = q$ extracts the quantum group
parameter.  The ``braiding corrections'' are conjugations by
$R$-matrices arising from the interaction of the $B$-cycle path with
the other points; they are controlled by the mixed $A$-$B$ relations
in $\pi_1(\mathrm{Conf}_n(E_\tau))$.

\medskip
\noindent\textbf{(iii).}
The Drinfeld--Kohno theorem (genus-$1$ version, see the discussion
below) proves that the monodromy representation of the KZB connection
factors through $U_q(\mathfrak{g})$.  The $A$-cycle monodromies give
the $R$-matrix (the braiding in $\mathrm{Rep}_q(\mathfrak{g})$),
the $B$-cycle monodromies give the Cartan elements (the $q^H$
operators), and together with the braid monodromies they generate the
full quantum group action.
\end{proof}

\begin{remark}[Physical origin of $q$]
\label{rem:q-physical}
The parameter $q = e^{2\pi i \hbar} = e^{2\pi^2 i/(k+h^\vee)}$
has a transparent origin in the holomorphic-topological framework.
At genus~$0$, the bar complex on $\mathrm{FM}_n(\mathbb{C}) \times
\mathrm{Conf}_n(\mathbb{R})$ produces the Yangian $Y_\hbar(\mathfrak{g})$,
a deformation over the formal parameter $\hbar = \pi i/(k + h^\vee)$.
At genus~$1$, the $B$-cycle monodromy exponentiates the additive
parameter~$\hbar$ into the multiplicative parameter
$q = e^{2\pi i \hbar}$: the passage from the additive formal group
$\widehat{\mathbb{G}}_a$ to the multiplicative formal group
$\widehat{\mathbb{G}}_m$ is the passage from the Yangian to the
quantum group.  The elliptic curve $E_\tau$ provides the geometric
mechanism: its $B$-cycle is the Wick-rotated time circle in the
holomorphic-topological theory, and the monodromy around it converts
the infinitesimal deformation parameter into a finite one.
\end{remark}

\subsection{The Drinfeld--Kohno theorem}
\label{sec:drinfeld-kohno}

The identification of the KZ/KZB monodromy with the quantum group
action is the content of the Drinfeld--Kohno theorem and its
generalisations.

\begin{theorem}[Drinfeld--Kohno; \ClaimStatusProvedHere{} for the affine lineage]
\label{thm:drinfeld-kohno}
\index{Drinfeld--Kohno theorem|textbf}
\index{monodromy!Drinfeld--Kohno identification}
Let $\mathfrak{g}$ be a simple Lie algebra, $k \in \mathbb{C}
\setminus \{-h^\vee\}$ generic, and $V_1, \dots, V_n$
finite-dimensional $\mathfrak{g}$-modules.
\begin{enumerate}[label=\textup{(\roman*)}]
\item \textbf{Genus $0$.}
The monodromy representation of the KZ
connection~\eqref{eq:kz} on $\mathrm{Conf}_n(\mathbb{C})$
is isomorphic to the braid group representation obtained from the
braided monoidal structure of $\mathrm{Rep}_q(\mathfrak{g})$
via the universal $R$-matrix of $U_q(\mathfrak{g})$:
\begin{equation}\label{eq:dk}
\rho_n^{\mathrm{KZ}}
\;\simeq\;
\rho_n^{U_q(\mathfrak{g})}.
\end{equation}

\item \textbf{Genus $1$.}
The monodromy representation of the KZB
connection~\eqref{eq:kzb} on $\mathrm{Conf}_n(E_\tau)$
factors through $U_q(\mathfrak{g})$ and recovers the
full quantum group action \textup{(}including the $q^H$ operators
from $B$-cycle monodromy\textup{)} on the tensor product
$V_1 \otimes \cdots \otimes V_n$.
\end{enumerate}
\end{theorem}

\begin{proof}
\textbf{(i)} is the classical Drinfeld--Kohno theorem.  The proof
proceeds in three steps.

\emph{Step~1: Drinfeld associator.}
Drinfeld constructs a formal series
$\Phi \in U(\mathfrak{g})^{\widehat\otimes 3}[\![\hbar]\!]$
satisfying the pentagon and hexagon equations, using iterated integrals
of $d\log(z_i - z_j)$ on $\mathrm{Conf}_3(\mathbb{C})$.  This
associator converts the infinitesimal braid relations (the CYBE) into
the finite braid relations (the YBE), providing an equivalence of
braided monoidal categories between the representation category of the
Drinfeld--Kohno Lie algebra $\mathfrak{t}_n$ (with its KZ braiding)
and $\mathrm{Rep}_q(\mathfrak{g})$ (with its universal $R$-matrix
braiding).

\emph{Step~2: One-loop collapse.}
For the affine algebra $V^k(\mathfrak{g})$, the one-loop exactness of
the BV-BRST differential (Theorem~\ref{thm:affine-half-space-bv} of
Volume~II) ensures that the higher $A_\infty$ operations $m_k$
($k \geq 3$) vanish on evaluation modules.  The full superconnection
on the bar complex therefore collapses to the KZ connection, and the
bar-complex monodromy $\rho_n^{\mathrm{HT}}$ equals $\rho_n^{\mathrm{KZ}}$
(Theorem~\ref{thm:affine-monodromy-identification}(i) of Volume~II).

\emph{Step~3: Identification.}
Combining Steps~1 and~2, the bar-complex monodromy for the affine
lineage equals the quantum-group braid representation:
$\rho_n^{\mathrm{HT}} = \rho_n^{\mathrm{KZ}} \simeq
\rho_n^{U_q(\mathfrak{g})}$.

\medskip
\noindent\textbf{(ii)} follows from the genus-$1$ generalisation of the
Drinfeld--Kohno theorem (Etingof--Schiffmann, Calaque--Enriquez--Etingof):
the KZB monodromy representation on $\mathrm{Conf}_n(E_\tau)$ factors
through $U_q(\mathfrak{g})$, with the $B$-cycle monodromy providing the
$q^H$ operators.  The one-loop collapse for the affine lineage ensures
that the bar-complex monodromy on $\mathrm{Conf}_n(E_\tau)$ equals the
KZB monodromy, and the identification follows.
\end{proof}

\subsection{The genus-$1$ deformation: from Yangian to quantum group}
\label{sec:yangian-to-quantum}

The passage from genus~$0$ to genus~$1$ can be summarised as a
deformation of the Yangian $Y_\hbar(\mathfrak{g})$ to the
quantum group $U_q(\mathfrak{g})$.

\begin{theorem}[Yangian--quantum group deformation; \ClaimStatusProvedHere{} for the affine lineage, \ClaimStatusConjectured{} beyond]
\label{thm:yangian-quantum-group}
\index{Yangian!genus-1 deformation to quantum group|textbf}
\index{quantum group!from Yangian|textbf}
Let $\mathfrak{g}$ be a simple Lie algebra and $k \in \mathbb{C}
\setminus \{-h^\vee\}$ generic.
\begin{enumerate}[label=\textup{(\roman*)}]
\item At genus~$0$, the ordered bar complex of $V^k(\mathfrak{g})$ on
  $\mathrm{FM}_n(\mathbb{C}) \times \mathrm{Conf}_n(\mathbb{R})$
  produces the Yangian $Y_\hbar(\mathfrak{g})$ as the line-operator
  algebra, with $\hbar = \pi i/(k + h^\vee)$.

\item At genus~$1$, the ordered bar complex on
  $E_\tau \times \mathbb{R}$ produces the quantum group
  $U_q(\mathfrak{g})$ as the line-operator algebra, with
  $q = e^{2\pi i \hbar}$.

\item The degeneration $\tau \to i\infty$ \textup{(}pinching the
  $B$-cycle\textup{)} recovers the Yangian from the quantum group:
  $\lim_{\tau \to i\infty} U_q(\mathfrak{g})
  = Y_\hbar(\mathfrak{g})$, in the sense that the KZB connection
  degenerates to the KZ connection and the $B$-cycle monodromy becomes
  trivial.

\item The quantum group parameter $q = e^{2\pi i \hbar}$ is the
  exponential of the Yangian parameter~$\hbar$ under the exponential map
  $\widehat{\mathbb{G}}_a \to \widehat{\mathbb{G}}_m$ provided by the
  $B$-cycle of $E_\tau$.
\end{enumerate}
\end{theorem}

\begin{proof}
Part~(i) is the content of the genus-$0$ programme: the bar differential
on $\mathrm{FM}_n(\mathbb{C})$ encodes the chiral OPE, the $E_1$
coproduct on $\mathrm{Conf}_n(\mathbb{R})$ encodes the associative
composition, and the resulting algebra is the Yangian (see the
holographic dictionary of Volume~II,
Corollary~\ref{cor:holographic-dictionary}).

Part~(ii) combines Theorem~\ref{thm:b-cycle-quantum-group} with the
Drinfeld--Kohno theorem (Theorem~\ref{thm:drinfeld-kohno}(ii)): the
full monodromy of the KZB connection generates $U_q(\mathfrak{g})$.

Part~(iii) follows from the degeneration
$\wp_1(z, \tau) \to 1/z + O(z)$ as $\tau \to i\infty$: the KZB
connection degenerates to KZ, the $B$-cycle shrinks to a point, and
the $q^H$ monodromy becomes trivial.

Part~(iv) is a restatement: $q = e^{2\pi i \hbar}$ is the exponential
of $\hbar$ multiplied by $2\pi i$, and the $B$-cycle monodromy is
the geometric mechanism that implements this exponentiation.
\end{proof}

\subsection{The $\mathfrak{sl}_2$ case and roots of unity}
\label{sec:sl2-roots-unity}

For $\mathfrak{g} = \mathfrak{sl}_2$, the quantum group parameter
specialises to a root of unity at integer level.

\begin{corollary}[$U_q(\mathfrak{sl}_2)$ at roots of unity from affine $\mathfrak{sl}_2$; \ClaimStatusProvedHere]
\label{cor:sl2-root-of-unity}
\index{quantum group!sl2 at roots of unity@$\mathfrak{sl}_2$ at roots of unity|textbf}
For $\mathfrak{g} = \mathfrak{sl}_2$ at positive integer level
$k \in \mathbb{Z}_{>0}$:
\begin{enumerate}[label=\textup{(\roman*)}]
\item The quantum group parameter is
  $q = e^{2\pi^2 i/(k+2)}$, and
  $q^{2(k+2)} = 1$.
\item The quantum group $U_q(\mathfrak{sl}_2)$ at this root of
  unity has a finite-dimensional quotient with $k+1$ irreducible
  modules $V_0, V_1, \dots, V_k$ of dimensions $1, 2, \dots, k+1$.
\item These $k+1$ modules are in bijection with the integrable
  highest-weight modules of $\widehat{\mathfrak{sl}}_2$ at level~$k$.
  The truncation of the representation category is the genus-$1$
  avatar of the level-$k$ WZW fusion category.
\end{enumerate}
\end{corollary}

\begin{proof}
Part~(i): $h^\vee = 2$ for $\mathfrak{sl}_2$, so
$\hbar = \pi i/(k+2)$ and $q = e^{2\pi i \hbar} = e^{2\pi^2 i/(k+2)}$.
For the root-of-unity property:
$q^{2(k+2)} = e^{2 \cdot 2\pi^2 i} = 1$.

Part~(ii) is a classical result: at $q$ a root of unity, the
quantum group $U_q(\mathfrak{sl}_2)$ has a large centre generated
by $E^{k+2}$, $F^{k+2}$, and $K^{2(k+2)}$, and the quotient by
this centre has finitely many simple modules.

Part~(iii) is the Kazhdan--Lusztig equivalence:
$\mathcal{O}_k^{\mathrm{int}}(\widehat{\mathfrak{sl}}_2)
\simeq \mathrm{Rep}_q(\mathfrak{sl}_2)^{\mathrm{fin}}$.
\end{proof}

\subsection{The Jones polynomial from genus-$1$ monodromy}
\label{sec:jones-genus1}

The Jones polynomial emerges as a topological invariant from the
genus-$1$ monodromy applied to the bar complex.

\begin{theorem}[Jones polynomial from genus-$1$ bar-complex monodromy; \ClaimStatusProvedHere]
\label{thm:jones-genus1}
\index{Jones polynomial!from genus-1 monodromy|textbf}
For $\mathfrak{g} = \mathfrak{sl}_2$ at level $k \in \mathbb{Z}_{>0}$,
let $K$ be a framed knot presented as the closure of a braid
$\beta \in B_n$.  Then:
\begin{enumerate}[label=\textup{(\roman*)}]
\item The bar complex of $V^k(\mathfrak{sl}_2)$, viewed as a
  logarithmic $\mathrm{SC}^{ch,top}$-algebra, produces the braid
  group representation $\rho_n^{\mathrm{HT}} = \rho_n^{\mathrm{KZ}}$
  on tensor powers of the fundamental representation~$V$.

\item By the Drinfeld--Kohno theorem
  \textup{(}Theorem~\textup{\ref{thm:drinfeld-kohno}}\textup{)},
  this representation factors through
  $U_q(\mathfrak{sl}_2)$ at $q = e^{2\pi^2 i/(k+2)}$.

\item The Jones polynomial is recovered by the quantum trace:
  \begin{equation}\label{eq:jones-from-bar}
  J_K(q)
  \;=\;
  \operatorname{tr}_q\bigl(\rho_n^{\mathrm{HT}}(\beta)\bigr)
  \;=\;
  \operatorname{tr}\bigl(q^{2\rho} \cdot \rho_n^{\mathrm{KZ}}(\beta)\bigr),
  \end{equation}
  where $\rho$ is the half-sum of positive roots
  \textup{(}$\rho = 1/2$ for $\mathfrak{sl}_2$\textup{)} and
  $\operatorname{tr}_q$ is the categorical trace in the ribbon
  category $\mathrm{Rep}_q(\mathfrak{sl}_2)$.

\item The integration domain is $\mathrm{FM}_n(\mathbb{C}) \times
  \mathrm{Conf}_n(\mathbb{R})$: the $\mathrm{FM}_n(\mathbb{C})$
  factor carries the bar differential encoding the chiral OPE, and the
  $\mathrm{Conf}_n(\mathbb{R})$ factor carries the $E_1$ coproduct
  implementing the topological interval-cutting.  The Jones polynomial
  is a holomorphic-topological invariant: the holomorphic direction
  produces the braiding, and the topological direction produces the
  monoidal structure.
\end{enumerate}
\end{theorem}

\begin{proof}
Parts (i)--(ii) combine Theorem~\ref{thm:affine-monodromy-identification}
(the bar-complex monodromy equals the KZ monodromy for the affine
lineage) with Theorem~\ref{thm:drinfeld-kohno}(i) (the KZ monodromy
equals the quantum group braid representation).

Part~(iii) is the Reshetikhin--Turaev construction: $\mathrm{Rep}_q(\mathfrak{sl}_2)$
is a ribbon category (the ribbon element is $v = K^{-1}u$ where
$K = q^H$ and $u = \sum S(b_i)a_i$), and the quantum trace
$\operatorname{tr}_q = \operatorname{tr}(q^{2\rho} \,\cdot\,)$
implements the Markov closure converting a braid representation into a
framed link invariant.  Since $\rho_n^{\mathrm{HT}} = \rho_n^{\mathrm{KZ}}
\simeq \rho_n^{U_q}$, the quantum trace applied to the bar-complex
monodromy equals the Jones polynomial.

Part~(iv) is a restatement of the Swiss-cheese factorisation
structure: the bar complex lives on the product of holomorphic and
topological configuration spaces, and the Jones polynomial is the
output of applying the Reshetikhin--Turaev functor to the monodromy of
the connection defined on this product.
\end{proof}

\begin{remark}[The three levels of the construction]
\label{rem:three-levels}
The passage from chiral algebra to knot invariant has three levels,
each proved:
\begin{center}
\renewcommand{\arraystretch}{1.3}
\begin{tabular}{lll}
\textbf{Level} & \textbf{Object} & \textbf{Output} \\
\hline
Genus $0$, ordered bar & KZ connection & Yangian $Y_\hbar(\mathfrak{g})$ \\
Genus $1$, $B$-cycle monodromy & KZB connection & Quantum group $U_q(\mathfrak{g})$ \\
Markov closure & Quantum trace & Jones polynomial $J_K(q)$ \\
\end{tabular}
\end{center}
Each level is proved for the affine lineage by the one-loop collapse
mechanism.  Beyond the affine lineage --- for $\mathcal{W}$-algebras
obtained by Drinfeld--Sokolov reduction, where higher $A_\infty$
operations contribute --- the correspondence between bar-complex
monodromy and quantum group representations remains conjectural.
\end{remark}

\subsection{Scope and frontier}
\label{sec:genus1-frontier}

\begin{remark}[Status]
\label{rem:genus1-status}
The results of this section are proved for the \emph{affine lineage}:
$V^k(\mathfrak{g})$ for simple $\mathfrak{g}$ at generic level~$k$,
acting on finite-dimensional evaluation modules.  The proof mechanism
is the one-loop collapse
(Theorem~\ref{thm:affine-monodromy-identification} of Volume~II),
which reduces the full bar-complex superconnection to the KZ/KZB
connection, and the Drinfeld--Kohno theorem (classical at genus~$0$,
Etingof--Schiffmann/Calaque--Enriquez--Etingof at genus~$1$), which
identifies the KZ/KZB monodromy with the quantum group representation.

Beyond the affine lineage, two frontiers are open:
\begin{enumerate}[label=\textup{(\alph*)}]
\item \emph{$\mathcal{W}$-algebras.}
For $\mathcal{W}$-algebras obtained by Drinfeld--Sokolov reduction,
the higher $A_\infty$ operations $m_k$ ($k \geq 3$) do not vanish:
the one-loop collapse mechanism breaks.  The full superconnection
contributes to the monodromy, and the genus-$1$ deformation should
produce a \emph{deformed quantum group} or \emph{quantum
$\mathcal{W}$-algebra}.  This is the content of
Conjecture~\ref{conj:rmatrix} in Volume~II.

\item \emph{Higher genus.}
At genus $g \geq 2$, the configuration space
$\mathrm{Conf}_n(\Sigma_g)$ has $2g$ distinguished cycle classes,
and the monodromy acquires $2g$ additional generators.  The
curvature $\kappa(A) \cdot \omega_g$ from the Hodge bundle on
$\overline{\mathcal{M}}_{g,n}$ provides genus-by-genus corrections.
The algebraic object controlling the monodromy at genus~$g$ should be
a \emph{higher-genus quantum group} or \emph{mapping-class-group
representation} generalising the braid-group representations at
genus~$0$ and~$1$.
\end{enumerate}
\end{remark}

\section{Conjectural frontier}
\label{sec:conjectures}

The algebraic core invites geometricization and contact with
neighboring duality theories.  The following conjectures make this precise.

\subsection{Francis--Gaitsgory as the commutator shadow}

The associative story developed here is not the Francis--Gaitsgory Lie/cocommutative story.
Rather, the latter should appear as the commutative shadow of the former.

\begin{theorem}[Comm\-utator-shadow theorem]
\ClaimStatusProvedElsewhere
\label{conj:FG-shadow}
Let $A$ be a filtered strongly admissible associative chiral algebra with exhaustive separated
commutator filtration $F_{\mathrm{com}}^\bullet A$.  Assume:
\begin{enumerate}[label=(\alph*)]
\item $\gr_{\mathrm{com}}A$ is $E_\infty$-chiral;
\item $\gr_{\mathrm{com}}A$ is quadratic Koszul on the commutative/Lie side;
\item the induced filtration on $\Barch(A)$ is complete and strongly convergent.
\end{enumerate}
Then there is a convergent spectral sequence
\[
E_1\cong \overline{B}^{\mathrm{ch},FG}\!\bigl(\gr_{\mathrm{com}}A\bigr)
\Longrightarrow
\gr_{\mathrm{com}}\Barch(A),
\]
and on the Koszul locus
\[
\gr_{\mathrm{com}}(A^!)\simeq \bigl(\gr_{\mathrm{com}}A\bigr)^!_{FG}.
\]
\end{theorem}

\begin{remark}[Proof obligations]
A proof should construct the commutator filtration directly on the ordered bar complex, show that the
associated graded kills the non-$\Sigma$ defect and descends to the commutative bar complex, and
identify the induced differential with the Francis--Gaitsgory differential.
\end{remark}

\subsection{Bordered configuration spaces and the geometric open theory}

The algebraic ordered open trace formalism should be realized geometrically by ordered bordered
configuration spaces or, equivalently, by a ribbon/Swiss-cheese type compactification adapted to the
chiral setting.

\begin{conjecture}[Geometric realization of the diagonal bicomodule]
\ClaimStatusConjectured
\label{conj:bordered}
There exists a bordered ordered configuration-space model
\[
\mathfrak C^{\mathrm{bord}}_{g,r,s}(X)
\]
and a functorial system of chain-level operations on the diagonal bicomodule $C_\Delta$ such that:
\begin{enumerate}[label=(\arabic*)]
\item the interval component reproduces the bicomodule $C_\Delta$;
\item gluing of boundary intervals is computed by derived cotensor over $C$;
\item the ordered pair-of-pants operation coincides with $\mu_P$;
\item annular closure computes $\coHoch^{\mathrm{ch}}_\bullet(C)$;
\item the resulting geometric operations recover the ordered open trace formalism of
Theorem~\ref{thm:ordered-open}.
\end{enumerate}
\end{conjecture}

\begin{remark}[Partial resolution via the bordered FM construction]
\label{rem:bordered-partial-resolution}
Items~(1)--(3) and~(5) of Conjecture~\ref{conj:bordered} are now resolved
by the bordered Fulton--MacPherson compactification constructed in
\S\ref{sec:bordered-fm} (Construction~\ref{constr:bordered-fm},
Theorem~\ref{thm:bordered-fm-properties}):
the interval/punctured-curve bordered FM compactification
$\overline{C}{}^{\,\mathrm{oc}}_{k,\vec{m}}(X, D, \tau)$ is a smooth
compact manifold with corners whose four-type boundary decomposition
(Proposition~\ref{prop:four-type-boundary}) provides the geometric
carrier for the ordered open trace formalism.  In Volume~II,
Corollary~\texttt{cor:bordered-from-CG} confirms that all five algebraic
items follow once the compactification exists.

Item~(4), annular closure computing chiral Hochschild homology,
requires the \emph{annular} bordered FM compactification
(Construction~\texttt{constr:bordered-fm-annulus} in Volume~II), which
involves cyclically ordered boundary points on circles rather than
linearly ordered points on intervals.  This annular case remains
conjectural: the thickness-zero face requires Mok's degeneration theory
for the family $S^1\times[0,\epsilon]\to\{0\}$.

In summary: the conjecture's status is \textbf{four-fifths resolved};
the remaining open input is a single compactification problem for the
annulus.
\end{remark}

\subsection{Associative modular Maurer--Cartan theory}

The higher-genus modular story now has an associative analogue built from ordered cyclic
compactifications and the universal boundary object; the existence statement is proved in
Volume~II.

\begin{theorem}[Associative modular Maurer--Cartan class;
\ClaimStatusProvedElsewhere]
\label{thm:ordered-associative-modular-mc}
There exists an associative cyclic deformation complex
\[
\mathrm{Def}^{\mathrm{cyc}}_{\Assch}(A)
\]
and a canonical Maurer--Cartan element
\[
\Theta_A^{\Assch}
\in
\MC\!\Bigl(
\mathrm{Def}^{\mathrm{cyc}}_{\Assch}(A)\widehat\otimes
R\Gamma(\overline{M}_{g,\bullet},k)
\Bigr)
\]
whose linear term is the genus anomaly, whose quadratic term encodes the first pair-of-pants
boundary correction, and whose full Maurer--Cartan equation expresses compatibility of all ordered
higher-genus boundary gluings.
\end{theorem}

\begin{remark}[Resolution elsewhere]
The original conjectural formulation is resolved in Volume~II by the
$E_1$ modular convolution construction: the ordered cyclic compactification is realized by the
ribbon-modular Feynman transform, and the canonical Maurer--Cartan element is
$\Theta_A^{E_1} = D_A^{E_1} - d_0$.
\end{remark}

\subsection{BRST and Drinfeld--Sokolov compatibility}

The principal corridor is now resolved, while the full non-principal frontier remains open.

\begin{theorem}[Reduction commutes with associative chiral duality
\textup{(}principal case\textup{)}; \ClaimStatusProvedElsewhere]
\label{thm:ordered-associative-ds-principal}
Let $Q_{DS}$ denote the principal Drinfeld--Sokolov reduction functor on a filtered pronilpotent
subcategory of strongly admissible associative chiral algebras at non-critical level.  Assume:
\begin{enumerate}[label=(\alph*)]
\item $Q_{DS}$ is exact on the relevant filtered complexes;
\item $Q_{DS}$ preserves PBW filtrations;
\item the BRST differential lifts to the ordered bar coalgebra.
\end{enumerate}
Then
\[
Q_{DS}(A^!)\simeq Q_{DS}(A)^!.
\]
Moreover, for every $A$-bimodule $M$ one should have
\[
Q_{DS}\bigl(\KK_A^{\mathrm{bi}}(M)\bigr)\simeq
\KK_{Q_{DS}(A)}^{\mathrm{bi}}\bigl(Q_{DS}(M)\bigr).
\]
\end{theorem}

\begin{conjecture}[Reduction commutes with associative chiral duality
\textup{(}arbitrary nilpotent\textup{)}]
\ClaimStatusConjectured
\label{conj:DS-arbitrary-nilpotent}
Let $Q_{DS}$ denote a BRST/Drinfeld--Sokolov reduction functor for an arbitrary nilpotent
reduction datum on a filtered pronilpotent subcategory of strongly admissible associative chiral
algebras.  Assume:
\begin{enumerate}[label=(\alph*)]
\item $Q_{DS}$ is exact on the relevant filtered complexes;
\item $Q_{DS}$ preserves PBW filtrations;
\item the BRST differential lifts to the ordered bar coalgebra.
\end{enumerate}
Then
\[
Q_{DS}(A^!)\simeq Q_{DS}(A)^!.
\]
Moreover, for every $A$-bimodule $M$ one should have
\[
Q_{DS}\bigl(\KK_A^{\mathrm{bi}}(M)\bigr)\simeq
\KK_{Q_{DS}(A)}^{\mathrm{bi}}\bigl(Q_{DS}(M)\bigr).
\]
\end{conjecture}

\begin{remark}
Theorem~\ref{thm:ordered-associative-ds-principal} records the proved principal corridor from
Volume~II.  Conjecture~\ref{conj:DS-arbitrary-nilpotent} isolates the remaining non-principal
frontier: one must verify PBW compatibility and the lift of the BRST differential to the ordered
bar coalgebra beyond the principal and hook-type regimes.
\end{remark}

\subsection{The \texorpdfstring{$\mathcal{W}_3$}{W\_3} ordered bar complex via DS transport}
\label{subsec:w3-ordered-bar-ds}

The principal DS reduction $V_k(\mathfrak{sl}_3) \twoheadrightarrow
\mathcal{W}_3^k$ transports the ordered bar complex of
$\widehat{\mathfrak{sl}}_{3,k}$ to the ordered bar complex of
$\mathcal{W}_3$.  The DS--KD intertwining theorem
(Theorem~\ref{thm:ordered-associative-ds-principal}; for the symmetric
bar, Theorem~\ref{thm:ds-koszul-intertwine}) guarantees that the
functor $H^0(Q_{\mathrm{DS}})$ commutes with the ordered bar
construction at chain level:
\begin{equation}\label{eq:w3-ordered-bar-ds}
\Barch_{\mathrm{ord}}(\mathcal{W}_3^k)
\;\simeq\;
H^0_{\mathrm{DS}}\bigl(\Barch_{\mathrm{ord}}(V_k(\mathfrak{sl}_3))\bigr).
\end{equation}
This subsection records the explicit structure of the left-hand side.

\begin{theorem}[Ordered bar complex of $\mathcal{W}_3$ via DS transport; \ClaimStatusProvedHere]
\label{thm:w3-ordered-bar}
Let $k \neq -3$ be non-critical.  The ordered bar complex
$\Barch_{\mathrm{ord}}(\mathcal{W}_3^k)$ has the following structure.

\begin{enumerate}[label=\textup{(\roman*)}]
\item \textbf{Generators.}  The complex is cogenerated by
  $T$ \textup{(}spin~$2$, from the Sugawara tensor\textup{)} and
  $W$ \textup{(}spin~$3$, from the cubic Casimir\textup{)}, both
  obtained as $H^0(Q_{\mathrm{DS}})$-images of $\mathfrak{sl}_3$
  currents.

\item \textbf{Pole structure.}  The collision residues are:
\begin{itemize}
\item $T$-$T$ channel: poles up to order~$3$ in the collision residue
  \textup{(}from the quartic Virasoro OPE $T_{(3)}T = c/2$; the
  $d\log$ kernel absorbs one pole\textup{)}.
\item $T$-$W$ channel: poles up to order~$2$
  \textup{(}from $T_{(2)}W = 3W$\textup{)}.
\item $W$-$W$ channel: poles up to order~$4$
  \textup{(}from the quintic OPE pole
  $W_{(4)}W = c/3$; the $d\log$ kernel absorbs one\textup{)}.
\end{itemize}

\item \textbf{A$_\infty$ operations.}  The transferred operations
  $m_k^{\mathcal{W}_3}$ on the cohomology
  $H^\bullet(\Barch_{\mathrm{ord}}(\mathcal{W}_3^k))$ satisfy:
  \begin{itemize}
  \item $m_2$ is the ordered binary product
    \textup{(}collision residue at $D_{ij}$\textup{)};
  \item $m_3 \neq 0$ in the $T$-sector, inherited from the Virasoro
    quartic pole;
  \item $m_k \neq 0$ for all $k \geq 3$
    \textup{(}class~$\mathrm{M}$: infinite tower\textup{)}.
  \end{itemize}

\item \textbf{Curvature.}  The bar curvature
  decomposes by channel:
  \[
  m_0^{\mathcal{W}_3} = \kappa_T |0\rangle + \kappa_W |0\rangle
  = \frac{c}{2}|0\rangle + \frac{c}{3}|0\rangle
  = \frac{5c}{6}|0\rangle,
  \]
  which is the DS projection of the $\widehat{\mathfrak{sl}}_{3,k}$
  curvature $m_0 = (k{+}3)\kappa/6$.

\item \textbf{Complementarity.}
  The Koszul involution is $c \mapsto 100 - c$,
  and the Koszul conductor is
  \[
  K_3 = 2(3{-}1)(2 \cdot 9 + 2 \cdot 3 + 1) = 100.
  \]
  The algebra is Koszul self-dual at $c^* = K_3/2 = 50$.
  The full Koszul conductor
  $K(\mathcal{W}_3) = K_3 \cdot (H_3 - 1) = 100 \cdot 5/6 = 250/3$
  controls the complementarity sum
  $\kappa(\mathcal{W}_3^k) + \kappa(\mathcal{W}_3^{k'}) = 250/3$.
\end{enumerate}
\end{theorem}

\begin{proof}
\emph{Part~(i).}  The principal DS reduction
$V_k(\mathfrak{sl}_3) \twoheadrightarrow \mathcal{W}_3^k$
sends the $8$~current generators $J^a$ ($a = 1, \ldots, 8$) of
$\widehat{\mathfrak{sl}}_{3,k}$ to two generators: the stress tensor
$T$ (via the Sugawara construction) and the spin-$3$ field~$W$
(via the cubic Casimir).  The remaining $6$~currents lie in
$\operatorname{im}(Q_{\mathrm{DS}})$ and are killed by the
projection $H^0(Q_{\mathrm{DS}})$.  By
equation~\eqref{eq:w3-ordered-bar-ds}, the ordered bar complex
of $\mathcal{W}_3$ is cogenerated by the images of those
$\widehat{\mathfrak{sl}}_3$ bar elements that survive DS projection.

\emph{Part~(ii).}  The pole orders in the ordered bar differential
are one less than the corresponding OPE pole orders
(the $d\log$ kernel absorbs one power).  The
$\mathcal{W}_3$ OPE has: $T(z)T(w) \sim c/2 \cdot (z{-}w)^{-4}
+ \cdots$ (quartic pole), $T(z)W(w) \sim 3W(w)(z{-}w)^{-2}
+ \cdots$ (cubic pole), and
\begin{equation}\label{eq:ww-ope-leading-ordered}
W(z)W(w) \sim \frac{c/3}{(z{-}w)^6}
+ \frac{2T}{(z{-}w)^4}
+ \frac{\partial T}{(z{-}w)^3}
+ \frac{\beta\Lambda + \tfrac{3}{10}\partial^2 T}{(z{-}w)^2}
+ \cdots
\end{equation}
where $\Lambda = (TT) - \tfrac{3}{10}\partial^2 T$ is the
quasi-primary of weight~$4$ and
$\beta = 16/(22 + 5c)$.  The $d\log$ absorption
gives collision residues with poles of order $3$, $2$, and $4$ in the
respective channels (since $W_{(5)}W = 0$ and the leading nonzero
term is $W_{(4)}W = c/3$, the sextic pole produces a quintic-order
$n$-product, and the $d\log$ kernel reduces the collision residue
pole by one to order~$4$).

\emph{Part~(iii).}  The Virasoro subalgebra
$\operatorname{Vir}_c \hookrightarrow \mathcal{W}_3^k$ (generated
by~$T$) is autonomous: its OPE closes on itself.  The quartic pole
$T_{(3)}T = c/2$ forces $m_3 \neq 0$ on the $T$-sector by the
single-line class~M theorem
(Theorem~\ref{thm:single-line-dichotomy}): the critical
discriminant $\Delta_T = 40/(5c+22) \neq 0$ ensures that the shadow
tower is infinite.  The operations $m_k$ for $k \geq 4$ are then
generated by the resolvent tree formula
(Theorem~\ref{thm:tree-formula}): each $m_k$ is a sum over
planar binary trees with $k$~leaves,
\begin{equation}\label{eq:w3-resolvent-tree}
m_k^{\mathcal{W}_3}
= \sum_{T \in \mathrm{PBT}(k)}
  m_2^{(T)} \circ h_{\mathrm{DS}}^{(\mathrm{edges}(T))},
\end{equation}
where $m_2^{(T)}$ denotes the iterated binary product along the
tree~$T$, and $h_{\mathrm{DS}}$ is the DS contracting homotopy
inserted on each internal edge.  The number of planar binary trees
with $k$~leaves is the Catalan number $C_{k-1}$, so $m_k$ receives
$C_{k-1}$ contributions.  Crucially, the homotopy $h_{\mathrm{DS}}$
is the composition of the standard SDR homotopy for
$\widehat{\mathfrak{sl}}_3 \twoheadrightarrow \mathcal{W}_3$ with
the homotopy transfer data
(Construction~\ref{constr:transfer-ainf}): the infinite tower
$\{m_k\}_{k \geq 3}$ is generated from the single quadratic product
$m_2$ on $\widehat{\mathfrak{sl}}_3$ by iterated composition with
$h_{\mathrm{DS}}$.

The $W$-line also has $m_{2k} \neq 0$ for all $k \geq 2$ (but
$m_{2k+1} = 0$ by the $\mathbb{Z}_2$ parity $W \mapsto -W$);
the nonvanishing of $m_4^W$ follows from
$\Delta_W = 20480/[3(5c{+}22)^3] \neq 0$
(Computation~\ref{comp:w-infty-shadow-tower}).

\emph{Part~(iv).}  The curvature decomposition
$\kappa(\mathcal{W}_3) = \kappa_T + \kappa_W = c/2 + c/3 = 5c/6$
is verified by direct computation
(Computation~\ref{comp:ds-bar-sl3-w3}, Step~4): the
DS projection of the $\widehat{\mathfrak{sl}}_3$ curvature
$m_0 = (k{+}3)\kappa/6$ onto the $T$- and $W$-channels produces
exactly $c/2$ and $c/3$ respectively.

\emph{Part~(v).}  The Koszul involution is
$k \mapsto k' = -k - 2h^\vee = -k - 6$, which sends
$c(k) = 2 - 24(k{+}2)^2/(k{+}3)$ to $c' = 100 - c$.  Hence
$c + c' = 100 = K_3$.
The self-dual point is $c = c'$, i.e., $c = 50$.
The complementarity sum
$\kappa(\mathcal{W}_3^k) + \kappa(\mathcal{W}_3^{k'})
= 5c/6 + 5(100{-}c)/6 = 250/3$
is $k$-independent.
\end{proof}

\begin{theorem}[Class~M persistence under DS transport; \ClaimStatusProvedHere]
\label{thm:class-m-ds-transport}
\index{class M!DS transport}
Let $\mathcal{W}_N = \mathcal{W}^k(\mathfrak{sl}_N, f_{\mathrm{prin}})$ be a
principal $\mathcal{W}$-algebra at non-critical level.  Then:
\begin{enumerate}[label=\textup{(\roman*)}]
\item The ordered bar complex
$\Barch_{\mathrm{ord}}(\mathcal{W}_N)$ is class~$\mathrm{M}$:
the A$_\infty$ tower $\{m_k\}$ is infinite, with
$m_k \neq 0$ for all $k \geq 3$ \textup{(}on the $T$-line\textup{)}.

\item The non-formality is forced by the Virasoro sub-sector:
the stress tensor $T$ generates an autonomous Virasoro subalgebra
$\operatorname{Vir}_{c(k)} \hookrightarrow \mathcal{W}_N^k$, and
the quartic pole $T_{(3)}T = c/2$ forces $m_3 \neq 0$.  The
higher operations cascade by the resolvent tree formula.

\item DS reduction preserves Koszulness but destroys formality:
$V_k(\mathfrak{sl}_N)$ is class~$\mathrm{L}$ \textup{(}double-pole
OPE, formal SC structure\textup{)}, while
$\mathcal{W}_N^k$ is class~$\mathrm{M}$
\textup{(}quartic-and-higher poles, genuinely infinite A$_\infty$
tower\textup{)}.  The Sugawara construction escalates the pole depth:
the affine double-pole OPE $J^a(z)J^b(w) \sim k\delta^{ab}/(z{-}w)^2
+ f^{ab}_c J^c(w)/(z{-}w)$ produces, upon taking normal-ordered
products in the DS construction, the Virasoro quartic-pole OPE and
the $W$-channel poles of order $\leq 2s$ for each strong generator
of spin~$s$.

\item The Koszul conductor and self-dual central charge for
  $\mathcal{W}_N$ are:
\[
K_N = 2(N{-}1)(2N^2{+}2N{+}1),
\qquad
c^* = K_N/2 = (N{-}1)(2N^2{+}2N{+}1).
\]
For $N = 2$: $K_2 = 26$, $c^* = 13$ \textup{(}Virasoro\textup{)}.
For $N = 3$: $K_3 = 100$, $c^* = 50$.
For $N = 4$: $K_4 = 222$, $c^* = 111$.
\end{enumerate}
\end{theorem}

\begin{proof}
\emph{Part~(i).}
Each principal $\mathcal{W}_N$ contains a Virasoro subalgebra
$\operatorname{Vir}_{c(k)}$ generated by the stress tensor.  The
$T$-channel shadow tower is the Virasoro shadow tower, which is
class~M by the single-line dichotomy:
$\Delta_T = 40/(5c+22) \neq 0$ for $c \neq -22/5$.

\emph{Part~(ii).}
The quartic pole $T(z)T(w) \sim c/2 \cdot (z{-}w)^{-4}$ in the
Virasoro sub-sector gives a cubic-order collision residue
(after $d\log$ absorption).  This is the source of $m_3 \neq 0$:
the homotopy transfer formula produces $m_3 = p \circ m_2(h \circ
m_2 \otimes \id) + (-1)^{|a|} p \circ m_2(\id \otimes h \circ m_2)$
with both Catalan trees contributing nonzero terms.  The
resolvent tree formula then generates all $m_k$ for $k \geq 4$ from
$m_2$ and the contracting homotopy.

\emph{Part~(iii).}
The affine algebra $V_k(\mathfrak{sl}_N)$ has a double-pole OPE,
so its collision residue has poles of order~$1$.  The bar complex is
\emph{formal}: $m_k = 0$ for $k \geq 3$ on cohomology
(class~L).  The DS reduction introduces the Sugawara construction
$T = \frac{1}{2(k+h^\vee)} (J^a J^a)$, which creates a
quartic pole $T(z)T(w) \sim c/2 \cdot (z{-}w)^{-4}$ from the
double-pole affine OPE.  This is pole-order escalation via
normal ordering: two double poles compose (through the Sugawara
quadratic) to produce a quartic pole.  For higher-spin generators
$W^{(s)}$, the $W^{(s)} W^{(s)}$ OPE has poles up to order~$2s$,
produced by the same mechanism applied to the degree-$s$ Casimir
invariants.  The Koszulness is preserved because $H^0(Q_{\mathrm{DS}})$
is exact on the relevant filtered complexes
(Theorem~\ref{thm:ordered-associative-ds-principal}, hypothesis~(a)),
but the formality is destroyed because the homotopy transfer data
$(p, \iota, h_{\mathrm{DS}})$ generate nonzero $m_k$ for $k \geq 3$
from the new higher-order poles.

\emph{Part~(iv).}
The DS central charge is
$c(k) = (N{-}1)\bigl[1 - N(N{+}1)(N{-}1)/(k{+}N)\bigr]$
(Freudenthal--de Vries), and the Koszul involution
$k \mapsto -k - 2N$ sends $c \mapsto K_N - c$.  The value of
$K_N$ is read off: $c(k) + c(-k{-}2N) = 2(N{-}1)(2N^2{+}2N{+}1)$.
The self-dual point is $c = c^* = K_N/2$.
\end{proof}

\begin{remark}[The resolvent tree formula for $\mathcal{W}_3$]
\label{rem:w3-resolvent}
\index{resolvent tree formula!$\mathcal{W}_3$}
The resolvent tree formula~\eqref{eq:w3-resolvent-tree} makes
explicit the mechanism by which DS reduction generates the infinite
A$_\infty$ tower from a single quadratic product.  The affine
algebra $V_k(\mathfrak{sl}_3)$ has $m_2$ as its only nontrivial
operation (the bar complex is formal: class~L).  Applying the
SDR data $(p, \iota, h_{\mathrm{DS}})$ to the homotopy transfer
theorem produces:
\begin{align*}
m_3^{\mathcal{W}_3}(a,b,c)
&= p\, m_2(h_{\mathrm{DS}}\, m_2(\iota a, \iota b),\, \iota c)
+ (-1)^{|a|}\,
  p\, m_2(\iota a,\, h_{\mathrm{DS}}\, m_2(\iota b, \iota c)),
\end{align*}
where $\iota\colon \mathcal{W}_3 \hookrightarrow V_k(\mathfrak{sl}_3)$
is the inclusion (via the Miura embedding) and
$p\colon V_k(\mathfrak{sl}_3) \twoheadrightarrow \mathcal{W}_3$
is the DS projection.  At arity~$k$, the number of contributing
trees is the Catalan number $C_{k-1}$: for $m_3$, $C_2 = 2$
(the two planar binary trees with $3$ leaves); for $m_4$, $C_3 = 5$;
for $m_5$, $C_4 = 14$.  Each tree contributes a term of the form
$m_2 \circ (h_{\mathrm{DS}} \circ m_2)^{\circ (k-2)}$, iterated
along the tree topology.

The key structural point: the affine quadratic product $m_2$ is
\emph{formal} (it alone generates no $m_3$), but the homotopy
$h_{\mathrm{DS}}$ is non-trivial on the Sugawara sector---it
mediates the passage from affine currents to the stress tensor and
higher-spin fields---and this non-trivial mediation is what generates
the infinite cascade.  \emph{DS reduction preserves Koszulness but
destroys formality: this is the transport from class~$\mathrm{L}$
to class~$\mathrm{M}$.}
\end{remark}

\begin{remark}[The $W$-channel parity rule on the ordered bar complex]
\label{rem:w3-parity-ordered}
\index{parity rule!$\mathcal{W}_3$ ordered bar complex}
The $\mathbb{Z}_2$ symmetry $W \mapsto -W$ of the $\mathcal{W}_3$ OPE
descends to the ordered bar complex: on the $W$-line (deformation
direction $x_W$ with $x_T = 0$), all odd-arity operations vanish
($m_{2k+1}^W = 0$), while even-arity operations are generically nonzero
($m_{2k}^W \neq 0$ for $k \geq 1$).  The first nonzero higher
operation on the $W$-line is $m_4^W$, with shadow coefficient
$S_4^W = 2560/[c(5c+22)^3]$
(Computation~\ref{comp:w-infty-shadow-tower}).
This parity constraint
is inherited from the $\mathbb{Z}_2$ automorphism of the DS
reduction: the principal embedding $\mathfrak{sl}_2
\hookrightarrow \mathfrak{sl}_3$ determines a $\mathbb{Z}_2$
that acts by $-1$ on the cubic Casimir and hence on~$W$.
\end{remark}

\subsection{Comparison with shifted and factorization quantum groups}
\label{subsec:shifted-factorization-comparison}

Three mathematical traditions converge on the same combinatorial-algebraic
object: the ordered bar complex $\Barch(A)$ on
$\FM_k(\mathbb{C})\times\mathrm{Conf}_k^<(\mathbb{R})$.  The purpose of this
subsection is to place the ordered associative chiral Koszul duality of this
chapter in context by comparing it with the shifted quantum groups of
Finkelberg--Tsymbaliuk (and the Coulomb-branch approach of
Braverman--Finkelberg--Nakajima), the factorization quantum groups of
Latyntsev, and the stable-envelope programme of Maulik--Okounkov (as
synthesised for 3d $\mathcal{N}=4$ theories by Costello--Okounkov--Zhang--Zhou).

\subsubsection*{The three-language spine}

\[
\begin{array}{c|c|c|c}
\text{Datum} & \text{Physical (this work)} & \text{Spectral (Latyntsev)}
  & \text{Geometric (COZZ)} \\
\hline
R\text{-matrix} & \text{OPE exchange} & \text{factorisation isomorphism} &
  \text{stable envelope} \\
\text{YBE} & \text{Ward identity} & \text{factorisation axiom} &
  \text{braid relation} \\
\text{Strictification} & \text{root multiplicity 1} &
  \text{by construction} & \text{by construction}
\end{array}
\]

In each column the $R$-matrix, the Yang--Baxter equation, and the
strictification mechanism arise from the internal logic of the
framework:

\begin{itemize}
\item \textbf{Physical column.}
The $R$-matrix is the regularised holonomy of the derived additive KZ
connection on $\mathrm{Conf}_n^{\mathrm{ord}}(\mathbb{A}^1)$, with
connection one-form built from OPE residues
(Theorem~\ref{thm:derived-additive-kz}, Volume~II).
The Yang--Baxter equation is the Ward identity for the flat connection:
$\mathbb{D}_n^2=0$.  Strictification requires that the spectral Drinfeld
obstruction class vanish, which holds for all simple Lie algebras because
every root has multiplicity one (Volume~II,
Theorem~\texttt{thm:complete-strictification}).

\item \textbf{Spectral column.}
Latyntsev defines factorisation quantum groups as factorisation algebras on
the spectral curve, with the $R$-matrix encoded as a factorisation
isomorphism (the structure map for disjoint-support factorisation).  The
Yang--Baxter equation is the coherence axiom for threefold factorisation.
Strictification is built in: the factorisation structure is strict by
definition.

\item \textbf{Geometric column.}
Costello--Okounkov--Zhang--Zhou identify the $R$-matrix with the
Maulik--Okounkov stable envelope for the Nakajima quiver variety or the BFN
Coulomb branch.  The Yang--Baxter equation is the braid relation for stable
envelopes (which follows from uniqueness of the stable envelope under the
Maulik--Okounkov axioms).  Strictification is again automatic: stable
envelopes are honest matrices, not $\Ainf$ data.
\end{itemize}

\subsubsection*{The common object}

All three languages attach algebraic data to the same
configuration-space carrier:
\[
\Barch^{\mathrm{ord}}(A)\;\;\text{on}\;\;
\FM_k(\mathbb{C})\times\mathrm{Conf}_k^<(\mathbb{R}).
\]
The holomorphic factor $\FM_k(\mathbb{C})$ carries the closed-colour
differential (OPE residues), while the topological factor
$\mathrm{Conf}_k^<(\mathbb{R})$ carries the open-colour coproduct
(ordered deconcatenation).  What differs among the three traditions
is which of these two factors is taken as primary input and which
is derived.

\subsubsection*{Shifted quantum groups}

Finkelberg--Tsymbaliuk define, for a simple Lie algebra $\mathfrak{g}$
and a pair of coweights $(\mu,\nu)$, the \emph{shifted quantum group}
$Q_\mu^\nu(\mathfrak{g})$ by explicit generators (Drinfeld--Cartan
generators, shifted Chevalley generators) and shifted Serre relations.
The unshifted specialisation recovers the Yangian:
\begin{equation}\label{eq:unshifted-yangian}
Q_0^0(\mathfrak{g})\;\simeq\; Y_\hbar(\mathfrak{g}).
\end{equation}

\begin{theorem}[Unshifted identification; \ClaimStatusProvedHere]
\label{thm:unshifted-identification}
On the chirally Koszul locus, the line-side algebra $A^!_{\mathrm{line}}$
produced by ordered associative chiral Koszul duality of $\widehat{\mathfrak{g}}_k$
is isomorphic to the Yangian $Y_\hbar(\mathfrak{g})$:
\[
A^!_{\mathrm{line}}\bigl(\widehat{\mathfrak{g}}_k\bigr)
\;\simeq\;
Y_\hbar(\mathfrak{g})
\;\simeq\;
Q_0^0(\mathfrak{g}).
\]
\end{theorem}

\begin{proof}
This is the content of the Yangian recognition theorem
(Theorem~\ref{thm:yangian-recognition}, recorded from Volume~II).
The RTT presentation of $Y_\hbar(\mathfrak{g})$ is read off from
the regularised holonomy of the derived additive KZ connection
(Theorem~\ref{thm:derived-additive-kz}).  The identification
$Y_\hbar(\mathfrak{g})\simeq Q_0^0(\mathfrak{g})$ is the
Levendorski\u{\i}--Drinfeld theorem.
\end{proof}

\begin{remark}[Shift parameters as boundary conditions]
\label{rem:shift-boundary-conditions}
The shift parameters $(\mu,\nu)$ in $Q_\mu^\nu(\mathfrak{g})$ correspond,
in the 3d holomorphic-topological picture, to the choice of boundary
conditions at the two ends of the topological interval $\mathbb{R}$.
Different choices of boundary condition produce different shifted
quantum groups; the unshifted case $(\mu,\nu)=(0,0)$ corresponds
to Neumann boundary conditions on both sides.  The algebraic
manifestation is that the shift modifies the Serre relations,
which on the bar-complex side corresponds to modifying the
augmentation datum of the ordered bar coalgebra.
\end{remark}

\subsubsection*{Factorization quantum groups}

Latyntsev constructs factorisation quantum groups as factorisation
algebras on the formal spectral curve $\widehat{\mathbb{A}}{}^1_u$,
working at the derived/$\infty$-categorical level throughout.

\begin{theorem}[Factorisation identification on the Koszul locus;
\ClaimStatusProvedHere]
\label{thm:factorisation-identification}
On the chirally Koszul locus for a simple Lie algebra $\mathfrak{g}$,
the dg-shifted Yangian $Y^{\mathrm{dg}}_\hbar(\mathfrak{g})$ produced
by ordered chiral Koszul duality, after bar-horizontal strictification,
yields a strict spectral factorisation quantum group.
\end{theorem}

\begin{proof}
The quasi-factorisation theorem (Volume~II,
Theorem~\texttt{thm:quasi-factorisation})
shows that every residue-bounded complete dg-shifted Yangian with a
bar-horizontal strictification produces a spectral quasi-factorisation
quantum group.  The spectral Drinfeld obstruction class controls the
passage from quasi to strict.  For simple $\mathfrak{g}$, root multiplicity
one forces this class to vanish at every filtration level (Volume~II,
Theorem~\texttt{thm:complete-strictification}), yielding an honest
strict factorisation quantum group.
\end{proof}

\begin{remark}[What the factorisation framework adds]
Latyntsev's construction operates at the $\infty$-categorical level from the
outset, producing factorisation quantum groups as $\Eone$-algebra objects in
ind-coherent sheaves on the Ran space of the spectral curve.  The main
advantage over the operadic approach of this chapter is that Latyntsev's
framework gives a \emph{definition} of factorisation quantum group that does
not require choosing a strictification: the factorisation isomorphisms are
part of the defining data.  Our contribution is to show that the
$\mathrm{SC}^{\mathrm{ch,top}}$-bar complex
produces the same object once strictification is carried out, and that
for simple Lie algebras the strictification is always possible.  The
Kac--Moody frontier (root multiplicities $>1$) is the regime where the
strictification mechanism fails and Latyntsev's a~priori construction
must be used as the primary definition.
\end{remark}

\subsubsection*{Coulomb branch algebras and stable envelopes}

Braverman--Finkelberg--Nakajima (BFN) define the Coulomb branch algebra
$\mathcal{A}_C(G,N)$ as the equivariant Borel--Moore homology of the BFN
moduli space $\mathcal{R}_{G,N}$.  Maulik--Okounkov construct $R$-matrices
from stable envelopes of Nakajima quiver varieties; the
Costello--Okounkov--Zhang--Zhou synthesis extends this to 3d
$\mathcal{N}=4$ gauge theories.

\begin{proposition}[$R$-matrix comparison; \ClaimStatusProvedHere]
\label{prop:r-matrix-stable-envelope}
For $\mathfrak{g}$ simple and $\widehat{\mathfrak{g}}_k$ at
non-critical generic level, the $R$-matrix $R(z)$ produced by the
regularised holonomy of the derived additive KZ connection coincides
with the Maulik--Okounkov stable envelope on evaluation modules.
\end{proposition}

\begin{proof}
On evaluation modules, both $R$-matrices satisfy the same
Yang--Baxter equation with the same spectral dependence and the
same classical limit $r(z)$ (the rational $r$-matrix of
$\mathfrak{g}$).  The uniqueness theorem for trigonometric/rational
$R$-matrices with prescribed classical limit and spectral dependence
(the Drinfeld uniqueness argument) gives the identification.
\end{proof}

\begin{remark}[Beyond the affine lineage]
The comparison with stable envelopes is verified on evaluation modules
of the Yangian.  The full comparison for non-evaluation modules,
and the extension to shifted quantum groups (where the stable
envelope is for the BFN Coulomb branch rather than the Nakajima
variety), remains conjectural in general.
\end{remark}

\begin{conjecture}[Full three-language equivalence]
\ClaimStatusConjectured
\label{conj:three-language-equivalence}
For any 3d $\mathcal{N}=4$ gauge theory with gauge group $G$ and
matter representation $N$:
\begin{enumerate}[label=\textup{(\alph*)}]
\item the ordered chiral Koszul dual of the boundary chiral algebra
  $A_{\partial}(G,N)$ is isomorphic, after bar-horizontal
  strictification, to a factorisation quantum group in Latyntsev's sense;
\item this factorisation quantum group, evaluated on the unshifted Ran
  space, recovers the Yangian of the Coulomb branch;
\item the $R$-matrix from ordered-to-unordered descent agrees with the
  Maulik--Okounkov stable envelope for $\mathcal{M}_C(G,N)$.
\end{enumerate}
\end{conjecture}

\begin{remark}[Status and scope]
Parts~(a) and~(b) of Conjecture~\ref{conj:three-language-equivalence}
are proved for the standard affine case
$A = \widehat{\mathfrak{g}}_k$ with $\mathfrak{g}$ simple
(Theorems~\ref{thm:unshifted-identification}
and~\ref{thm:factorisation-identification}).  Part~(c)
is proved on evaluation modules
(Proposition~\ref{prop:r-matrix-stable-envelope}).  The full conjecture
extends these identifications to arbitrary gauge theories and all
module categories.
\end{remark}


\section{Evaluation modules and the line category}
\label{sec:evaluation-modules}

The ordered bar complex and the Yangian $Y_\hbar(\mathfrak{g})$ classify the
line operators of the corresponding holomorphic-topological theory.  In this
section develops the representation-theoretic side: evaluation modules, the
Drinfeld polynomial classification, tensor products and their braiding, and the
identification of the resulting module category with the line category
$\mathcal{C}_{\mathrm{line}}$.

Throughout, $\mathfrak{g}$ denotes a finite-dimensional simple Lie algebra
over~$\mathbb{C}$ of rank~$r$, with Cartan matrix $(a_{ij})$,
simple roots $\alpha_1,\dots,\alpha_r$, fundamental weights
$\omega_1,\dots,\omega_r$, and normalised Killing form
$(\cdot,\cdot)$ for which $(\alpha,\alpha)=2$ for long roots.  We write
$\Omega\in\mathfrak{g}\otimes\mathfrak{g}$ for the quadratic Casimir element
and $P\in\operatorname{End}(V\otimes W)$ for the permutation operator
$v\otimes w\mapsto w\otimes v$.

\subsection{The Yangian $Y_\hbar(\mathfrak{g})$ and the evaluation homomorphism}

\begin{definition}[Drinfeld Yangian]
\label{def:yangian-drinfeld}
The \emph{Yangian} $Y_\hbar(\mathfrak{g})$ is the associative
$\mathbb{C}[\![\hbar]\!]$-algebra generated by
$\{J^a_0,\,J^a_1\mid a=1,\dots,\dim\mathfrak{g}\}$
subject to:
\begin{enumerate}[label=(\roman*)]
\item $J^a_0$ generate a copy of $U(\mathfrak{g})$;
\item $[J^a_0,\,J^b_1]=f^{ab}{}_c\,J^c_1$;
\item Drinfeld's ternary relation:
$[J^a_1,\,[J^b_0,\,J^c_1]]-[J^a_0,\,[J^b_1,\,J^c_1]]$
$=\hbar^2\,A^{abc}_{def}\{J^d_0,J^e_0,J^f_0\}$,
where $A^{abc}_{def}$ is the totally symmetric tensor
of \textup{\cite{Drinfeld-Yangians}} and
$\{x,y,z\}=\tfrac{1}{24}\sum_{\sigma\in S_3}x_{\sigma(1)}y_{\sigma(2)}z_{\sigma(3)}$.
\end{enumerate}
The coproduct is
$\Delta(J^a_0)=J^a_0\otimes 1+1\otimes J^a_0$,
$\Delta(J^a_1)=J^a_1\otimes 1+1\otimes J^a_1
+\tfrac{\hbar}{2}\,f^a{}_{bc}\,J^b_0\otimes J^c_0$.
\end{definition}

\begin{construction}[Evaluation homomorphism]
\ClaimStatusProvedHere
\label{constr:evaluation-map}
For each $a\in\mathbb{C}$, there is a surjective algebra
homomorphism
\[
\mathrm{ev}_a\colon
Y_\hbar(\mathfrak{g})\longrightarrow U(\mathfrak{g}),
\qquad
J^a_0\mapsto J^a,
\quad
J^a_1\mapsto a\,J^a.
\]
The evaluation map is an algebra homomorphism because
$[J^a,a\,J^b]=a\,f^{ab}{}_c\,J^c=a\,f^{ab}{}_c\,J^c_0$,
and the ternary relation reduces to the Jacobi identity in
$U(\mathfrak{g})$.
\end{construction}

\begin{definition}[Evaluation module]
\label{def:evaluation-module}
Let $(\rho,V)$ be a finite-dimensional
$\mathfrak{g}$-representation and $a\in\mathbb{C}$.
The \emph{evaluation module} $V_a$ is the
$Y_\hbar(\mathfrak{g})$-module obtained by pulling back
along $\mathrm{ev}_a$:
\[
V_a
:=
\mathrm{ev}_a^*(V).
\]
Concretely, $V_a=V$ as a vector space, with
\begin{equation}\label{eq:eval-transfer-matrix}
T^a(u)\cdot v
=
v+\frac{\hbar\,\rho(\Omega_{(1)})\otimes\Omega_{(2)}}{u-a}\cdot v
=
v+\frac{\hbar}{u-a}\sum_\alpha
\rho(J^\alpha)\,v\otimes J_\alpha,
\end{equation}
where $\Omega=\sum_\alpha J^\alpha\otimes J_\alpha$
is the Casimir tensor and the transfer matrix
$T(u)=1+\hbar\sum_{n\geq 0}J_n\, u^{-n-1}$.
\end{definition}

\subsection{Explicit evaluation modules for $\mathfrak{sl}_2$}

\begin{computation}[$\mathfrak{sl}_2$ evaluation module; \ClaimStatusProvedHere]
\label{comp:sl2-eval}
\index{evaluation module!sl2@$\mathfrak{sl}_2$}
Let $\mathfrak{g}=\mathfrak{sl}_2$ with standard basis
$e,f,h$ and $[h,e]=2e$, $[h,f]=-2f$, $[e,f]=h$.
The fundamental representation is $V=\mathbb{C}^2$ with
basis $v_+,v_-$ and
\[
\rho(e)=\begin{pmatrix}0&1\\0&0\end{pmatrix},
\quad
\rho(f)=\begin{pmatrix}0&0\\1&0\end{pmatrix},
\quad
\rho(h)=\begin{pmatrix}1&0\\0&-1\end{pmatrix}.
\]
The Casimir tensor is
$\Omega=\tfrac{1}{2}h\otimes h+e\otimes f+f\otimes e$,
so $\rho(\Omega)\in\operatorname{End}(\mathbb{C}^2)$
is $\rho(\Omega)=\tfrac{1}{2}C_2(\mathrm{fund})\cdot\id
=\tfrac{3}{4}\cdot\id$ (half the quadratic Casimir eigenvalue).

The evaluation module $V_a=\mathbb{C}^2$ has
\[
J^a_0\cdot v_\pm=\rho(J^a)\,v_\pm,
\qquad
J^a_1\cdot v_\pm=a\,\rho(J^a)\,v_\pm.
\]
The transfer matrix in the fundamental representation is
\begin{equation}\label{eq:sl2-transfer-fundamental}
T(u)
\;=\;
\id_2
\;+\;
\frac{\hbar}{u-a}
\begin{pmatrix}
\tfrac{1}{2}h & f \\[4pt]
e & -\tfrac{1}{2}h
\end{pmatrix},
\end{equation}
where the matrix entries are $\mathfrak{sl}_2$-generators
acting via $\rho$ in the quantum space and via the
fundamental in the auxiliary space.
\end{computation}

\begin{theorem}[$R$-matrix on $V_a\otimes V_b$ for $\mathfrak{sl}_2$;
\ClaimStatusProvedHere]
\label{thm:sl2-R-matrix}
\index{R-matrix!sl2 evaluation@$\mathfrak{sl}_2$ evaluation}
The $R$-matrix governing the braiding of
$V_a\otimes V_b$ is
\begin{equation}\label{eq:sl2-R-eval}
R(a-b)
\;=\;
(a-b)\,\id\;+\;\hbar\,P
\;\in\;
\operatorname{End}(\mathbb{C}^2\otimes\mathbb{C}^2),
\end{equation}
where $P$ is the permutation operator $v\otimes w\mapsto w\otimes v$.
Equivalently, in the normalised form $\check{R}(u)=P\,R(u)$:
\[
\check{R}(u)
=
u\,P+\hbar\,\id.
\]
This is the Yang $R$-matrix.
\end{theorem}

\begin{proof}
The intertwining equation
$\Delta^{\mathrm{op}}(x)\,R=R\,\Delta(x)$ for all
$x\in Y_\hbar(\mathfrak{sl}_2)$, evaluated on
$V_a\otimes V_b$, uniquely determines $R$ up to scalar.
For $x=J^a_0$:
$[\rho^{(1)}(J^a)+\rho^{(2)}(J^a),R]=0$, so $R$
commutes with the diagonal $\mathfrak{sl}_2$-action.  By
Schur's lemma applied to the decomposition
$\mathbb{C}^2\otimes\mathbb{C}^2=\mathrm{Sym}^2\oplus\bigwedge^2$,
$R=\alpha\,\Pi_{\mathrm{sym}}+\beta\,\Pi_{\mathrm{alt}}$
for scalars $\alpha,\beta$.  Writing
$\Pi_{\mathrm{sym}}=\tfrac{1}{2}(\id+P)$ and
$\Pi_{\mathrm{alt}}=\tfrac{1}{2}(\id-P)$:
\[
R=\tfrac{\alpha+\beta}{2}\,\id
+\tfrac{\alpha-\beta}{2}\,P.
\]
Evaluating the intertwining equation for
$x=J^a_1$ with $\Delta(J^a_1)=J^a_1\otimes 1+1\otimes J^a_1
+\tfrac{\hbar}{2}f^a{}_{bc}J^b_0\otimes J^c_0$
yields $\alpha-\beta=\hbar$ and $\alpha+\beta=a-b$.
Hence $R=(a-b)\id+\hbar\,P$.
\end{proof}

\begin{corollary}[Clebsch--Gordan decomposition and non-semisimplicity;
\ClaimStatusProvedHere]
\label{cor:sl2-clebsch-gordan}
The tensor product $V_a\otimes V_b$ of fundamental
$\mathfrak{sl}_2$-evaluation modules is:
\begin{enumerate}[label=(\alph*)]
\item \textbf{Semisimple} when $a-b\notin\{0,\pm\hbar\}$:
the $R$-matrix~\eqref{eq:sl2-R-eval} is invertible, and
$V_a\otimes V_b\simeq W_{3,a+b}\oplus W_{1,a+b}$,
where $W_{d,c}$ denotes the evaluation module of dimension~$d$
at the averaged spectral parameter.

\item \textbf{Non-semisimple} when $a-b=\pm\hbar$:
$R(a-b)=\pm\hbar\,\id+\hbar\,P$ is not scalar but remains
invertible.  The tensor product is reducible but
indecomposable: $V_a\otimes V_b$ has a two-step filtration
with Jordan--H\"older factors $W_3$ and $W_1$, but the
extension does not split.

\item \textbf{Singular} when $a=b$: $R(0)=\hbar\,P$
has eigenvalues $\pm\hbar$, and the tensor product $V_a\otimes V_a$
degenerates: the braiding becomes the graded
permutation, and the intertwiner is no longer generic.
The image of the bar complex at this point reflects the
coalescence singularity of the ordered configuration space.
\end{enumerate}
\end{corollary}

\begin{proof}
The eigenvalues of $R(u)=u\,\id+\hbar\, P$ on the
$\mathrm{Sym}^2$ and $\bigwedge^2$ summands are
$u+\hbar$ and $u-\hbar$, respectively.  Part~(a):
when $u=a-b\neq 0,\pm\hbar$, both eigenvalues are nonzero
and the corresponding projectors split the tensor product.
Part~(b): when $u=\pm\hbar$, one eigenvalue is $2\hbar$
or $0$ and the other is $0$ or $-2\hbar$; the zero
eigenvalue signals a non-split extension.  Part~(c) is
$u=0$.
\end{proof}

\subsection{Evaluation modules for $\mathfrak{sl}_3$}

\begin{computation}[$\mathfrak{sl}_3$ fundamental evaluation module;
\ClaimStatusProvedHere]
\label{comp:sl3-eval-fundamental}
\index{evaluation module!sl3@$\mathfrak{sl}_3$!fundamental}
Let $\mathfrak{g}=\mathfrak{sl}_3$ with standard Chevalley
generators $e_i,f_i,h_i$ ($i=1,2$).  The fundamental
representation $V=\mathbb{C}^3$ has basis $v_1,v_2,v_3$
with weight vectors of weights $\omega_1,\,\omega_1-\alpha_1,\,
\omega_1-\alpha_1-\alpha_2$.  The evaluation module
$V_a=\mathbb{C}^3$ has
\[
J^{e_i}_0=E_i,\quad J^{f_i}_0=F_i,\quad J^{h_i}_0=H_i,
\qquad
J^{e_i}_1=a\,E_i,\quad J^{f_i}_1=a\,F_i,\quad J^{h_i}_1=a\,H_i,
\]
where $E_i,F_i,H_i$ are the standard $3\times 3$ matrices
representing $e_i,f_i,h_i$.

The $R$-matrix on $V_a\otimes V_b$ in the fundamental takes
the universal Yang form:
\begin{equation}\label{eq:sl3-R-fund}
R(u)
=
u\,\id\;+\;\hbar\,P
\;\in\;
\operatorname{End}(\mathbb{C}^3\otimes\mathbb{C}^3),
\end{equation}
where $u=a-b$.  Under $\mathfrak{sl}_3$,
$\mathbb{C}^3\otimes\mathbb{C}^3
\simeq\mathrm{Sym}^2(\mathbb{C}^3)\oplus\bigwedge^2(\mathbb{C}^3)$
\textup{(}dimensions $6+3$\textup{)}, and the eigenvalues of
$R(u)$ are $u+\hbar$ on $\mathrm{Sym}^2$ and $u-\hbar$
on $\bigwedge^2$, exactly as for $\mathfrak{sl}_2$.
\end{computation}

\begin{computation}[$\mathfrak{sl}_3$ adjoint evaluation module;
\ClaimStatusProvedHere]
\label{comp:sl3-eval-adjoint}
\index{evaluation module!sl3@$\mathfrak{sl}_3$!adjoint}
The adjoint representation $W=\mathfrak{sl}_3$ has
$\dim W=8$.  The evaluation module $W_a$ is the pullback
of the adjoint $\mathfrak{sl}_3$-module along
$\mathrm{ev}_a$.  Since
$\mathbb{C}^3\otimes(\mathbb{C}^3)^*\simeq
\mathfrak{sl}_3\oplus\mathbb{C}$, one has
\[
W_a\simeq (V_a\otimes V^*_a)/\mathbb{C},
\]
where $V^*_a$ is the contragredient evaluation module at the
same spectral parameter.  The $R$-matrix on $V_a\otimes W_b$
\textup{(}fundamental $\otimes$ adjoint\textup{)} is determined by
the Yangian intertwining equation and decomposes into three
eigenvalues on the three $\mathfrak{sl}_3$-irreducible summands
of $\mathbb{C}^3\otimes\mathfrak{sl}_3$:
\[
\mathbb{C}^3\otimes\mathfrak{sl}_3
\;\simeq\;
V(2\omega_1+\omega_2)\;\oplus\;V(\omega_1)\;\oplus\;V(\omega_2),
\]
with eigenvalues $u+2\hbar$, $u-\hbar$, $u+\hbar$
respectively, where $u=a-b$.
\end{computation}

\subsection{Drinfeld polynomials and the classification of finite-dimensional modules}

\begin{theorem}[Drinfeld classification; \ClaimStatusProvedElsewhere]
\label{thm:drinfeld-classification}
\index{Drinfeld polynomial!classification of Yangian modules|textbf}
\index{evaluation module!Drinfeld polynomial}
For $\mathfrak{g}$ a finite-dimensional simple Lie algebra
of rank~$r$, the isomorphism classes of irreducible
finite-dimensional $Y_\hbar(\mathfrak{g})$-modules are in
bijection with $r$-tuples of monic polynomials
$(P_1(u),\dots,P_r(u))\in\mathbb{C}[u]^r$, called the
\emph{Drinfeld polynomials}.

The bijection sends an irreducible module~$L$ to the
$r$-tuple where $P_i(u)$ records the eigenvalues of the
level-$1$ generating series associated to the $i$-th simple
root on a highest-weight vector of~$L$:
\[
T_i(u)\cdot v_+
=
\frac{P_i(u-\hbar)}{P_i(u)}\,v_+.
\]
\end{theorem}

\begin{proposition}[Evaluation modules as single-root Drinfeld polynomials;
\ClaimStatusProvedHere]
\label{prop:eval-drinfeld}
\index{Drinfeld polynomial!of evaluation module}
Let $V(\omega_i)$ be the $i$-th fundamental
representation of~$\mathfrak{g}$ and $a\in\mathbb{C}$.
The evaluation module $V(\omega_i)_a$ has Drinfeld
polynomials
\[
P_j(u)=
\begin{cases}
u-a & \text{if }j=i,\\
1 & \text{if }j\neq i.
\end{cases}
\]
More generally, the evaluation of an irreducible
$\mathfrak{g}$-module $V(\lambda)$ with highest weight
$\lambda=\sum_i m_i\omega_i$ at points
$a_1,\dots,a_N\in\mathbb{C}$ produces the tensor product
evaluation module $V(\lambda_1)_{a_1}\otimes\cdots\otimes
V(\lambda_N)_{a_N}$, whose Drinfeld polynomials are
\[
P_j(u)=\prod_{k=1}^N (u-a_k)^{m_j^{(k)}},
\]
where $m_j^{(k)}$ is the multiplicity of $\omega_j$ in
$\lambda_k$.
\end{proposition}

\begin{proof}
By definition, $\mathrm{ev}_a$ sends $J^a_1\mapsto a\,J^a$,
so the generating series $T_i(u)$ acts on the highest-weight
vector $v_+$ of $V(\omega_i)_a$ as a rational function of~$u$.
In the normalisation where $P_i(u)=u-a$, one obtains
$T_i(u)\cdot v_+=(u-a-\hbar)/(u-a)\cdot v_+$ and
$T_j(u)\cdot v_+=v_+$ for $j\neq i$, matching the
classification.  The general formula follows from the
tensor product rule for Drinfeld polynomials: the
$P_j$ of a tensor product is the product of the
$P_j$'s of the factors.
\end{proof}

\begin{example}[Drinfeld polynomials for evaluation chains]
\label{ex:drinfeld-eval-chain}
For $\mathfrak{g}=\mathfrak{sl}_2$, rank $r=1$, and an
evaluation chain at points $a_1,\dots,a_n\in\mathbb{C}$
in the fundamental representation:
\[
P(u)=\prod_{k=1}^n(u-a_k).
\]
The module $V_{a_1}\otimes\cdots\otimes V_{a_n}$ is
irreducible if and only if $a_i-a_j\notin\mathbb{Z}\hbar$
for all $i\neq j$ (the non-resonance condition).
When $a_i-a_j\in\mathbb{Z}\hbar$ for some pair, the tensor
product is reducible and the irreducible subquotients are
classified by combinatorial data (segments in the sense of
Zelevinsky).
\end{example}

\subsection{The line category}

The line operators of a holomorphic-topological theory on
$\mathbb{C}_z\times\mathbb{R}_t$ are operators supported on
lines $\{z=z_0\}\times\mathbb{R}_t$.  These form a
category $\mathcal{C}_{\mathrm{line}}$, and the ordered bar
complex identifies it with the module category of the
open-colour Koszul dual.

\begin{theorem}[Line category as Yangian modules;
\ClaimStatusProvedHere]
\label{thm:line-category}
\index{line category!as Yangian modules|textbf}
\index{Koszul dual!open-colour!line category}
Let $A$ be a strongly admissible $E_1$-chiral algebra
with ordered bar coalgebra $C=\Barch^{\mathrm{ord}}(A)$ and
open-colour Koszul dual $A^!_{\mathrm{line}}$.  Then:
\begin{enumerate}[label=(\roman*)]
\item The line category is
\[
\mathcal{C}_{\mathrm{line}}
\;\simeq\;
A^!_{\mathrm{line}}\text{-}\mathrm{mod}^{\mathrm{fd}},
\]
the category of finite-dimensional modules over the
open-colour Koszul dual.

\item On the chirally Koszul locus \textup{(}where the bar--cobar
equivalence holds and $A^!_{\mathrm{line}}$ is formal\textup{)},
for $A$ in the affine lineage
\textup{(}i.e., $A=\hat{\mathfrak{g}}_k$
for simple $\mathfrak{g}$ at non-critical level\textup{)},
\[
A^!_{\mathrm{line}}
\;\simeq\;
Y_\hbar(\mathfrak{g})
\]
and accordingly
\[
\mathcal{C}_{\mathrm{line}}
\;\simeq\;
Y_\hbar(\mathfrak{g})\text{-}\mathrm{mod}^{\mathrm{fd}}.
\]

\item The evaluation modules $V_a$ of
Definition~\textup{\ref{def:evaluation-module}} are the line
operators obtained by inserting a Wilson line in
representation~$V$ at position $z=a$.  The spectral
parameter $a$ is the holomorphic position on $\mathbb{C}_z$.
\end{enumerate}
\end{theorem}

\begin{proof}
Part~(i) is the content of the one-sided module/comodule
duality (Hypothesis~\ref{hyp:inputs}(I)) combined with
the ordered bar--cobar equivalence: left $C$-comodules are
equivalently right $A^!_{\mathrm{line}}$-modules.
Part~(ii) follows from the identification of
$A^!_{\mathrm{line}}$ with the dg-shifted Yangian
(proved in Volume~II,
Theorem~\ref*{thm:Koszul_dual_Yangian}).
Part~(iii) is the physical interpretation: a Wilson line in
representation $\rho$ at position $z=a$ defines a line
operator whose endomorphism algebra is
$\mathrm{ev}_a^*\rho$, and the intertwining between
different Wilson lines is governed by the $R$-matrix.
\end{proof}

\subsection{Tensor products and the braided monoidal structure}

\begin{theorem}[Braiding from the $R$-matrix;
\ClaimStatusProvedHere]
\label{thm:eval-braiding}
\index{braiding!evaluation modules|textbf}
\index{R-matrix!braiding of evaluation modules}
The tensor product of evaluation modules $V_a\otimes W_b$,
with $V,W$ finite-dimensional $\mathfrak{g}$-representations,
is a $Y_\hbar(\mathfrak{g})$-module via the coproduct~$\Delta$.
The braiding isomorphism
\[
\sigma_{V_a,W_b}\colon
V_a\otimes W_b
\;\xrightarrow{\;\sim\;}
W_b\otimes V_a
\]
is given by $\sigma_{V_a,W_b}=P\circ R_{V,W}(a-b)$,
where $R_{V,W}(u)\in\operatorname{End}(V\otimes W)(\!(u^{-1})\!)$
is the universal $R$-matrix of $Y_\hbar(\mathfrak{g})$
evaluated in $V\otimes W$.

When $a=b$, the $R$-matrix $R_{V,W}(0)$ degenerates:
its kernel is nontrivial whenever $V\otimes W$ is not
multiplicity-free.  This singularity reflects the
coalescence of two line operators at the same point,
corresponding to the boundary stratum
$z_1=z_2$ in $\overline{\mathrm{FM}}_2^{\mathrm{ord}}(\mathbb{C})$.
\end{theorem}

\begin{proof}
The coproduct $\Delta$ of $Y_\hbar(\mathfrak{g})$ defines
the tensor product module structure.  The universal
$R$-matrix satisfies the intertwining equation
$\Delta^{\mathrm{op}}(x)\,R=R\,\Delta(x)$, which means
$R_{V,W}(a-b)$ intertwines the two orderings of the
tensor product.  The braiding axioms (Yang--Baxter equation)
follow from the quantum group structure.  The coalescence
statement follows from the fact that $R_{V,W}(u)$ has a
pole at $u=0$ of order determined by the Casimir eigenvalue
in the tensor product, and the bar differential at the
boundary stratum $\partial\overline{\mathrm{FM}}_2$
imposes exactly the residue of this pole.
\end{proof}

\subsection{The Grothendieck ring}

\begin{theorem}[Grothendieck ring of Yangian modules;
\ClaimStatusProvedElsewhere]
\label{thm:grothendieck-yangian}
\index{Grothendieck ring!Yangian modules|textbf}
The Grothendieck ring of finite-dimensional
$Y_\hbar(\mathfrak{g})$-modules is
\begin{equation}\label{eq:grothendieck-yangian}
K_0\!\bigl(Y_\hbar(\mathfrak{g})\text{-}\mathrm{mod}^{\mathrm{fd}}\bigr)
\;\simeq\;
K_0\!\bigl(\operatorname{Rep}(\mathfrak{g})\bigr)[u],
\end{equation}
the polynomial extension of the representation ring
of~$\mathfrak{g}$.  The isomorphism sends the class of
$V(\lambda)_a$ to $[V(\lambda)]\cdot u^{\deg P_\lambda}$,
and the product is the convolution of Drinfeld polynomials.
\end{theorem}

\begin{remark}[Generators of the Grothendieck ring]
The Grothendieck ring is generated as a
$\mathbb{Z}[u]$-algebra by the classes
$[V(\omega_i)_a]$ ($i=1,\dots,r$, $a\in\mathbb{C}$).
The Drinfeld polynomial classification
(Theorem~\ref{thm:drinfeld-classification}) establishes
the bijection between irreducible classes and
$r$-tuples of monic polynomials.  The ring structure
reflects the fact that the tensor product of evaluation
modules, while not generically irreducible as a
$Y_\hbar(\mathfrak{g})$-module, becomes irreducible in the
Grothendieck group.
\end{remark}

\begin{remark}[Relation to the open trace formalism]
\label{rem:eval-open-trace}
The ordered open trace formalism of
Theorem~\ref{thm:ordered-open} acquires concrete content
on the Yangian locus.  The diagonal bicomodule $C_\Delta$
parametrises the identity line operator; evaluation modules
$V_a$ are line operators at position~$a$; the interval
composition $\square_C^{\mathbf{R}}$ is the sequential
fusion of line operators along $\mathbb{R}_t$; and the
ordered pair-of-pants $\mu_P$ is the OPE of line operators
at the same holomorphic position.  The Grothendieck ring
$K_0$ of the line category is the natural receptacle for
the decategorified open trace.
\end{remark}


\section{The annular bar complex and genus-one structure}
\label{sec:annular-bar}

The strip ordered bar complex
$\Barch^{\mathrm{ord}}(\cA)$
on $\FM_k(\mathbb{C})\times\mathrm{Conf}_k^{<}(\mathbb{R})$ computes
genus-zero data: the topological direction is $\mathbb{R}$, an interval,
and the resulting chain complex classifies the genus-zero
$\Eone$-chiral Koszul dual.
Replacing the strip~$\mathbb{R}$ by the circle
$S^1=\mathbb{R}/\mathbb{Z}$ wraps the topological direction
and produces a genus-one analogue: the \emph{annular bar complex}.
The wrapping converts the two-sided cobar complex into a
coHochschild complex, and the resulting object computes
$\HH^{\mathrm{ch}}_\bullet(\cA)$---connecting the ordered sector to the
modular operad at genus~$1$.

\subsection{Definition of the annular bar complex}

\begin{definition}[Annular bar complex]%
\ClaimStatusProvedHere
\label{def:annular-bar}%
\index{annular bar complex|textbf}%
\index{bar complex!annular}%
Let $\cA$ be a strongly admissible augmented $\Eone$-chiral
algebra and $C=\Barch^{\mathrm{ord}}(\cA)$ its ordered
bar coalgebra, with $C_\Delta$ the diagonal bicomodule
(Theorem~\ref{thm:diagonal}).
The \emph{annular bar complex} is the chain complex
\[
B^{\mathrm{ann}}_\bullet(\cA)
\;:=\;
C_\bullet\!\bigl(
  \overline{\FM}{}^{\,\mathrm{ann}}_{r,s};\,
  \omega_{\mathrm{bar}}
\bigr)
\;\cong\;
\bigoplus_{n\ge 0}
  \mathrm{cyc}^n(C,C_\Delta),
\]
where $\mathrm{cyc}^n(C,C_\Delta):=
C_\Delta\widehat\otimes_C C^{\widehat\otimes\,n}$ is the
\emph{cyclic cotensor}:
the completed tensor product that cyclically identifies the
last right $C$-coaction with the first left $C$-coaction,
reflecting the periodicity of the circle.
\end{definition}

The annular configuration space
$\overline{\FM}{}^{\,\mathrm{ann}}_{r,s}$ parametrises
$r$ holomorphic marked points on~$\mathbb{C}$ and $s$ ordered
marked points on~$S^1$, with the Fulton--MacPherson
compactification in the holomorphic direction.  It is the
geometric carrier for the genus-one ordered sector: the strip
$\FM_k(\mathbb{C})\times\mathrm{Conf}_k^{<}(\mathbb{R})$ wraps to
$\FM_k(\mathbb{C})\times\mathrm{Conf}_k^{\mathrm{cyc}}(S^1)$,
where $\mathrm{Conf}_k^{\mathrm{cyc}}(S^1)$ denotes cyclically
ordered configurations on the circle.  The compactification
$\overline{\FM}{}^{\,\mathrm{ann}}_{r,s}$ is a manifold with
corners whose codimension-one boundary strata encode the four
components of the annular differential below.

\subsection{The annular bar differential}

\begin{theorem}[Annular bar differential;
\ClaimStatusProvedHere]%
\label{thm:annular-bar-differential}%
\index{annular bar complex!differential|textbf}%
The differential on $B^{\mathrm{ann}}_\bullet(\cA)$
decomposes as
\[
d^{\mathrm{ann}}
\;=\;
d_{\mathrm{int}}+d_{\mathrm{res}}
  +d_\Delta+d_{\mathrm{wrap}},
\]
where:
\begin{enumerate}[label=\textup{(\alph*)}]
\item $d_{\mathrm{int}}$: internal differential on each
  tensor factor (the intrinsic differential of~$\cA$).

\item $d_{\mathrm{res}}$: collision residue differential
  arising from OPE residues at interior collisions in the
  holomorphic direction
  (identical to the strip bar differential).

\item $d_\Delta$: comultiplication differential at
  non-wrapping boundary collisions---the coface maps
  $\delta_i$ for $1\le i\le n-1$ that split a tensor factor
  via the ordered coproduct.

\item $d_{\mathrm{wrap}}$: the wrap-around differential
  at the seam of the annulus, given by
\begin{equation}\label{eq:d-wrap}
d_{\mathrm{wrap}}\langle a_1|\cdots|a_n\rangle
\;=\;
\sum_{(a_n)}
(-1)^{\epsilon}\,
\langle a_1|\cdots|a_{n-1}|a_n'|a_n''\rangle
\;+\;
\sum_{(a_1)}
(-1)^{\epsilon'}\,
\langle a_1'|a_1''|a_2|\cdots|a_n\rangle,
\end{equation}
where the sums are over the ordered coproduct
$\Delta(a_n)=\sum a_n'\otimes a_n''$ and
$\Delta(a_1)=\sum a_1'\otimes a_1''$, and the signs
$\epsilon$, $\epsilon'$ are the Koszul signs dictated by the
suspended grading $|s^{-1}a|=|a|-1$.
\end{enumerate}
The identity $(d^{\mathrm{ann}})^2=0$ follows from
$\partial^2=0$ on the manifold with corners
$\overline{\FM}{}^{\,\mathrm{ann}}_{r,s}$: each boundary
face of codimension~$2$ appears as the common boundary
of exactly two codimension-$1$ faces, with opposite orientations.
\end{theorem}

\begin{proof}
The proof parallels that of the strip bar differential
(Construction~\ref{constr:bar-diff-collision}).  The
codimension-$1$ boundary strata of
$\overline{\FM}{}^{\,\mathrm{ann}}_{r,s}$ are of four types:
\begin{enumerate}[label=(\roman*)]
\item an internal bubble separates at a single tensor factor
  (producing $d_{\mathrm{int}}$);
\item two or more holomorphic points collide away from the
  seam (producing $d_{\mathrm{res}}$);
\item two consecutive topological points collide away from
  the seam, splitting a tensor factor via the ordered
  coproduct (producing $d_\Delta$);
\item a topological point crosses the seam $0\sim 1$ of $S^1$,
  forcing the first or last tensor factor to split and
  re-enter at the opposite end (producing $d_{\mathrm{wrap}}$).
\end{enumerate}
The identity $(d^{\mathrm{ann}})^2=0$ is the algebraic
manifestation of $\partial^2=0$: every codimension-$2$
face of $\overline{\FM}{}^{\,\mathrm{ann}}_{r,s}$ is a
common boundary of exactly two codimension-$1$ faces, and
the orientations cancel pairwise.
\end{proof}

\begin{remark}[The wrap-around differential as the genus-one novelty]
\label{rem:wrap-around}
The wrap-around differential $d_{\mathrm{wrap}}$ is the
only component without a strip analogue.  Algebraically,
it provides the coface maps $\delta_0$ and $\delta_n$ that
close the two-sided cobar complex into the coHochschild
complex: without it the complex computes
$\Cotor^{C^e}(C_\Delta,C_\Delta)$; with it, the
result is the self-cotrace
$C_\Delta\square_{C^e}^{\mathbf{L}}
C_\Delta\simeq
\coHoch^{\mathrm{ch}}_\bullet(C,C_\Delta)$.
\end{remark}

\subsection{Annular Hochschild identification}

\begin{theorem}[Annular bar complex computes chiral Hochschild homology;
\ClaimStatusProvedHere]%
\label{thm:annular-HH}%
\index{annular bar complex!Hochschild homology|textbf}%
\index{Hochschild homology!via annular bar complex}%
Let $\cA$ be a strongly admissible $\Eone$-chiral
algebra with ordered bar coalgebra $C=\Barch^{\mathrm{ord}}(\cA)$
and diagonal bicomodule $C_\Delta$.
There is a canonical quasi-isomorphism
\[
B^{\mathrm{ann}}_\bullet(\cA)
\;\simeq\;
\coHoch^{\mathrm{ch}}_\bullet(C,C_\Delta)
\;\simeq\;
\HH^{\mathrm{ch}}_\bullet(\cA),
\]
where the second equivalence is
Theorem~\textup{\ref{thm:HH-coHH-homology}}.
\end{theorem}

\begin{proof}
By Definition~\ref{def:annular-bar}, the annular bar complex
is the cyclic cotensor
$\bigoplus_{n\ge 0}\mathrm{cyc}^n(C,C_\Delta)$.
The coHochschild complex of a coalgebra~$C$ with coefficients in a
bicomodule~$M$ is the standard cochain complex
\[
\coHoch^{\mathrm{ch}}_\bullet(C,M)
\;=\;
\bigoplus_{n\ge 0}
M\widehat{\otimes}_C C^{\widehat\otimes\,n}
\]
with differential assembled from the coface maps $\delta_0,\ldots,\delta_n$.
For $M=C_\Delta$, this is precisely the
annular bar complex: the coface maps $\delta_1,\ldots,\delta_{n-1}$
produce $d_\Delta$ (non-wrapping boundary collisions), and the
coface maps $\delta_0$ and $\delta_n$ produce $d_{\mathrm{wrap}}$
(the wrap-around at the seam).  The internal differential $d_{\mathrm{int}}$
and collision residue $d_{\mathrm{res}}$ are carried along as the
internal structure of each tensor factor.

The second equivalence is
Theorem~\ref{thm:HH-coHH-homology}, which identifies
$\coHoch^{\mathrm{ch}}_\bullet(C,C_\Delta)\simeq
\HH^{\mathrm{ch}}_\bullet(\cA)$ via the bimodule/bicomodule
equivalence of Theorem~\ref{thm:bimod-bicomod}.
\end{proof}

\subsection{Connection to the modular operad at genus~$1$}

\begin{remark}[The annular bar complex and the modular
Swiss-cheese structure]
\label{rem:annular-modular}
\index{modular operad!genus-one annular bar}
The annular bar complex provides the ordered sector's
contribution to the genus-$1$ modular Swiss-cheese operad.
The passage from strip to annulus
($\mathbb{R}\leadsto S^1$) is the topological-direction
analogue of the holomorphic passage from~$\mathbb{C}$ to the
elliptic curve $E_\tau$: both wrap an infinite direction into
a compact one, and both introduce monodromy.

In the modular operad framework:
\begin{enumerate}[label=(\roman*)]
\item The strip bar complex
  $\Barch^{\mathrm{ord}}(\cA)$ is the genus-$0$ open sector,
  computing the $\Eone$ Koszul dual.
\item The annular bar complex
  $B^{\mathrm{ann}}_\bullet(\cA)$ is the genus-$1$ open
  sector, computing chiral Hochschild homology.
\item The genus-$1$ closed sector carries the curvature
  $d^2=\kappa(\cA)\cdot\omega_1$ from the Hodge bundle
  (Vol~I, Theorem~D).
\item The annular closure of
  Conjecture~\textup{\ref{conj:bordered}(4)} is the geometric
  realisation of the isomorphism
  $B^{\mathrm{ann}}_\bullet\simeq
  \coHoch^{\mathrm{ch}}_\bullet(C)$
  at the configuration-space level.
\end{enumerate}
The open trace formalism (Theorem~\ref{thm:ordered-open})
at genus~$1$ factors through the annular bar complex: the
decategorified trace $\mathrm{tr}\colon K_0(\mathcal{C}_{\mathrm{line}})\to
\HH^{\mathrm{ch}}_0(\cA)$ computes characters of line operators
by wrapping them around the annulus.
\end{remark}

\subsection{Curvature--braiding dichotomy at genus~$1$}

At genus~$0$, the curvature $\kappa(\cA)$ (measuring the failure of
formality, i.e.\ $d^2=\kappa\cdot\omega_g$ in the curved bar complex)
and the spectral $R$-matrix $R(z)$ (measuring the braiding monodromy
on the ordered configuration space) are independent structures:
one belongs to the closed colour, the other to the open colour.
At genus~$1$, however, they meet: the spectral braiding must be
defined on the elliptic curve $E_\tau$, and the quasi-periodicity
of the elliptic propagator forces the two structures to interact.
The nature of this interaction is controlled entirely by the
lowest-pole coefficient $c_0$ of the $\lambda$-bracket.

\begin{theorem}[Curvature--braiding dichotomy at genus~$1$;
\ClaimStatusProvedHere]
\label{thm:curvature-braiding-dichotomy}
\index{curvature!braiding dichotomy|textbf}
\index{elliptic spectral dichotomy}
\index{genus 1!curvature--braiding dichotomy}
Let $\cA$ be a logarithmic $\Eone$-chiral algebra and
write the $(-1)$-shifted $\lambda$-bracket on cohomology as
\[
\{a {}_\lambda b\}
\;=\;
c_0(a,b) + c_1(a,b)\,\lambda + c_2(a,b)\,\lambda^2 + \cdots\,,
\]
so that the genus-$0$ classical $r$-matrix has the Laurent
expansion
\[
r(z)
\;=\;
\frac{c_0}{z} + \frac{c_1}{z^2} + \frac{c_2}{z^3} + \cdots
\qquad \text{\textup{(}formal, $z\to\infty$\textup{)}}.
\]
At genus~$1$, the classical $r$-matrix on $E_\tau$ is
obtained by replacing the rational functions $1/z^{n+1}$ by
their elliptic counterparts:
\begin{equation}\label{eq:elliptic-r-matrix}
r^{E_\tau}(z)
\;=\;
c_0\cdot\zeta(z|\tau)
\;+\;
c_1\cdot\wp(z|\tau)
\;+\;
\sum_{n\ge 2}
  \frac{(-1)^{n-1}}{(n-1)!}\,
  c_n\cdot\wp^{(n-2)}(z|\tau),
\end{equation}
where $\zeta(z|\tau)$, $\wp(z|\tau)$ are the Weierstrass
zeta and $\wp$-functions with periods $1,\tau$.
Then:
\begin{enumerate}[label=\textup{(\alph*)}]
\item \textup{(Quasi-periodicity.)}
The $B$-cycle monodromy of $r^{E_\tau}$ is
\begin{equation}\label{eq:r-monodromy}
r^{E_\tau}(z+\tau)-r^{E_\tau}(z)
\;=\;
2\eta_\tau\cdot c_0,
\end{equation}
where $\eta_\tau=\zeta(\tau/2|\tau)$ is the quasi-period.
The $A$-cycle monodromy is $2\eta_1\cdot c_0$.

\item \textup{(Decoupling: Cartan type.)}
If $c_0=0$ \textup{(}equivalently, $\{a{}_0 b\}=0$ for all
$a,b\in H^\bullet(\cA,Q)$\textup{)}, then $r^{E_\tau}(z)$
is a meromorphic function on~$E_\tau$
\textup{(}doubly periodic\textup{)}.
The spectral braiding is independent of the $B$-cycle
monodromy that sources the curvature
$\kappa(\cA)\cdot\omega_1$.  The genus-$1$ Swiss-cheese
structure splits:
\[
\text{curved}\;\mathsf{SC}^{\mathrm{ch,top},(1)}
\;\simeq\;
(\text{curved closed colour})
\;\times\;
(\text{flat open colour with meromorphic braiding}).
\]

\item \textup{(Entanglement: Yangian type.)}
If $c_0\ne 0$, then $r^{E_\tau}(z)$ is quasi-periodic
with non-trivial $B$-cycle monodromy
$2\eta_\tau\cdot c_0$.  The spectral braiding is a
section of a flat connection on the universal cover
of~$E_\tau$ with holonomy~$c_0$, and the genus-$1$
Swiss-cheese structure does \emph{not} split.
The braiding and the curvature share a common geometric
source: the $B$-cycle monodromy of the Weierstrass zeta
function.
\end{enumerate}
\end{theorem}

\begin{proof}
\textbf{Step~1} (Elliptic propagator).
At genus~$0$, the holomorphic propagator is
$K^{\mathbb{C}}(z)=dz/z$ (the $d\log$ kernel).  At
genus~$1$, on the elliptic curve~$E_\tau$, the unique meromorphic
one-form with a simple pole at $z=0$ and prescribed quasi-periodicity
is $K^{E_\tau}(z)=\zeta(z|\tau)\,dz$.  This has the same local
singularity $\zeta(z|\tau)=1/z+O(z)$ near $z=0$ but acquires
the quasi-periodicity
$\zeta(z+1|\tau)=\zeta(z|\tau)+2\eta_1$ and
$\zeta(z+\tau|\tau)=\zeta(z|\tau)+2\eta_\tau$.

\textbf{Step~2} (Expansion of the elliptic $r$-matrix).
The genus-$0$ $r$-matrix $r(z)=\sum_{n\ge 0}c_n/z^{n+1}$
arises from pairing the propagator $dz/z$ with the
$\lambda$-bracket coefficients.  At genus~$1$, each
rational function $1/z^{n+1}$ is replaced by its unique
elliptic counterpart: $1/z\mapsto\zeta(z|\tau)$ and
$1/z^{n+1}\mapsto\frac{(-1)^{n-1}}{(n-1)!}\wp^{(n-2)}(z|\tau)$
for $n\ge 1$.  These are the unique functions on
$\mathbb{C}\setminus(\mathbb{Z}+\mathbb{Z}\tau)$ that
match $1/z^{n+1}+O(z)$ near $z=0$ and are doubly periodic
for $n\ge 1$ (or quasi-periodic for $n=0$).  This yields
formula~\eqref{eq:elliptic-r-matrix}.

\textbf{Step~3} (Dichotomy).
The quasi-periodicity is inherited entirely from $\zeta(z|\tau)$,
since $\wp$ and all its derivatives are doubly periodic.
Therefore the $B$-cycle monodromy of $r^{E_\tau}$ is
$2\eta_\tau\cdot c_0$, proving~(a).  When $c_0=0$, the
$\zeta$-term is absent and $r^{E_\tau}$ is meromorphic on
$E_\tau$, hence the spectral braiding defined by the annular bar
differential decouples from the curved closed colour.  When
$c_0\ne 0$, the non-trivial monodromy
forces the braiding to be a section of a flat connection with
holonomy $\exp(2\pi i\cdot c_0)$, and this flat connection
is the same geometric object whose $B$-cycle monodromy
controls the curvature $\kappa(\cA)\cdot\omega_1$.
\end{proof}

\begin{remark}[Classification by pole structure]
\label{rem:annular-pole-classification}
\index{pole-order dichotomy!annular}
The dichotomy partitions all $\Einf$-chiral algebras
(i.e.\ all vertex algebras, which are local and
$\Sigma_n$-equivariant) into two classes at genus~$1$:
\begin{center}
\begin{tabular}{l@{\quad}c@{\quad}c@{\quad}l}
\toprule
\textbf{Class} & $c_0$ & \textbf{Monodromy} & \textbf{Example} \\
\midrule
Cartan (decoupled) & $=0$ & trivial &
  Heisenberg $\mathcal{H}_k$, free fermion \\
Yangian (entangled) & $\ne 0$ & $2\eta_\tau\cdot c_0$ &
  $\widehat{\mathfrak{g}}_k$ ($\mathfrak{g}$ simple),
  Virasoro \\
\bottomrule
\end{tabular}
\end{center}
For the Heisenberg algebra $\mathcal{H}_k$, the OPE
$J(z)J(w)\sim k/(z-w)^2$ has $c_0=\{J{}_{(0)}J\}=0$ (no
simple pole in the $\lambda$-bracket: the lowest nonvanishing
coefficient is $c_1=k$).  The genus-$1$ $r$-matrix is
$r^{E_\tau}(z)=k\cdot\wp(z|\tau)$, which is doubly periodic.
For $\widehat{\mathfrak{sl}}_2$ at level~$k$, the OPE
$J^a(z)J^b(w)\sim k\,\kappa^{ab}/(z-w)^2+f^{ab}{}_c\,J^c(w)/(z-w)$
has $c_0=f^{ab}{}_c\,J^c\ne 0$, and the genus-$1$ structure
is genuinely entangled.
\end{remark}


\section{Curvature--braiding entanglement at genus one}
\label{sec:curvature-braiding-entanglement}

The dichotomy of \S\ref{sec:annular-bar} partitions chiral algebras
according to whether curvature and braiding decouple at genus~$1$.
We now give the sharp form of this dichotomy, isolate the mechanism
that forces entanglement, and formulate the central open question
to which the entire annular bar theory points.

\begin{theorem}[Elliptic spectral dichotomy --- genus-$1$ specialisation;
\ClaimStatusProvedHere]
\label{thm:elliptic-spectral-dichotomy}
\index{elliptic spectral dichotomy|textbf}
\index{genus 1!elliptic spectral dichotomy}
\index{r-matrix@$r$-matrix!elliptic}
Let $\cA$ be a logarithmic $\Eone$-chiral algebra with
$(-1)$-shifted $\lambda$-bracket
$\{a\,{}_\lambda\,b\}=\sum_{n\ge 0}c_n(a,b)\,\lambda^n$
on $H^\bullet(\cA,Q)$.
At genus~$1$, the ordered $r$-matrix on the elliptic curve
$E_\tau=\mathbb{C}/(\mathbb{Z}+\tau\mathbb{Z})$ takes the form
\begin{equation}\label{eq:elliptic-r-dichotomy}
r_1(z;\tau)
\;=\;
c_0\cdot\zeta(z\,|\,\tau)
\;+\;
c_1\cdot\wp(z\,|\,\tau)
\;+\;
\sum_{n\ge 2}
  \frac{(-1)^{n-1}}{(n-1)!}\,
  c_n\cdot\wp^{(n-2)}(z\,|\,\tau),
\end{equation}
where $\zeta$ and $\wp$ are the Weierstrass functions with
periods $(1,\tau)$.  The coefficient $c_0=a_{(0)}b$ is the
Coisson bracket \textup{(}zeroth product\textup{)} and
$c_1=a_{(1)}b$ is the conformal pairing.

The $B$-cycle monodromy is
\begin{equation}\label{eq:b-cycle-monodromy-sharp}
r_1(z+\tau;\tau) - r_1(z;\tau) \;=\; 2\eta_\tau\cdot c_0,
\end{equation}
with $\eta_\tau=\zeta(\tfrac{\tau}{2}\,|\,\tau)$ the
quasi-period.  The $A$-cycle monodromy is $2\eta_1\cdot c_0$,
where $\eta_1=\zeta(\tfrac{1}{2}\,|\,\tau)$.  The Legendre
relation $\eta_1\tau-\eta_\tau=\pi i$ constrains the
two monodromies.
\end{theorem}

\begin{proof}
By Theorem~\ref{thm:curvature-braiding-dichotomy}, the
genus-$1$ $r$-matrix is obtained from the genus-$0$ Laurent
expansion $r(z)=\sum c_n/z^{n+1}$ by replacing each rational
function $1/z^{n+1}$ with its unique elliptic counterpart on
$E_\tau$: $1/z\mapsto\zeta(z\,|\,\tau)$ for $n=0$, and
$1/z^{n+1}\mapsto\frac{(-1)^{n-1}}{(n-1)!}\wp^{(n-2)}(z\,|\,\tau)$
for $n\ge 1$.  These are uniquely determined by matching the
local singularity $1/z^{n+1}+O(z)$ near $z=0$ and requiring
double periodicity for $n\ge 1$ (or quasi-periodicity for $n=0$).
The monodromy~\eqref{eq:b-cycle-monodromy-sharp} follows from
$\zeta(z+\tau\,|\,\tau)=\zeta(z\,|\,\tau)+2\eta_\tau$ and the
double periodicity of $\wp$ and all its derivatives.
\end{proof}

\begin{definition}[Decoupled vs.\ entangled genus-$1$ structure]
\label{def:decoupled-entangled}
\index{curvature--braiding!decoupled}
\index{curvature--braiding!entangled}
Let $\cA$ be a logarithmic $\Eone$-chiral algebra with
$r_1(z;\tau)$ as in
Theorem~\ref{thm:elliptic-spectral-dichotomy}.
\begin{enumerate}[label=\textup{(\roman*)}]
\item $\cA$ is \emph{decoupled at genus~$1$} if $c_0=0$:
  the Coisson bracket vanishes, $r_1$ is doubly periodic on
  $E_\tau$, and the curvature $\kappa(\cA)\cdot\omega_1$
  and the spectral braiding are independent structures
  belonging to different colours of the Swiss-cheese.

\item $\cA$ is \emph{entangled at genus~$1$} if $c_0\ne 0$:
  $r_1$ is quasi-periodic with $B$-cycle monodromy
  $2\eta_\tau\cdot c_0$, and the braiding and the curvature
  share a common geometric source --- the quasi-periodicity
  of the Weierstrass zeta function, which is the same
  mechanism that produces $\kappa(\cA)\cdot\omega_1$ from the
  Hodge bundle on $\overline{\cM}_{1,n}$.
\end{enumerate}
\end{definition}

\begin{center}
\begin{tabular}{l@{\quad}c@{\quad}l@{\quad}l}
\toprule
\textbf{Algebra} & $c_0$ & \textbf{Status} & \textbf{Mechanism} \\
\midrule
Heisenberg $\mathcal{H}_k$ & $0$ & Decoupled &
  No simple pole; $r_1=k\,\wp(z\,|\,\tau)$ \\
Free fermion $\psi$ & $0$ & Decoupled &
  $\{-\,{}_{(0)}-\}=0$; $r_1=\wp(z\,|\,\tau)$ \\
Lattice $V_\Lambda$ (abelian) & $0$ & Decoupled &
  Abelian Lie bracket \\[4pt]
Affine $\widehat{\mathfrak{g}}_k$ & $\ne 0$ & Entangled &
  $c_0=f^{ab}{}_c\,J^c$; KZB monodromy \\
Virasoro $\Vir_c$ & $\ne 0$ & Entangled &
  $c_0=2T$; $r_1=\tfrac{c}{2}\wp'(z)+2T\,\zeta(z)$ \\
$\mathcal{W}_N$ ($N\ge 2$) & $\ne 0$ & Entangled &
  DS reduction inherits $c_0\ne 0$ from $\widehat{\mathfrak{sl}}_N$ \\
\bottomrule
\end{tabular}
\end{center}

\begin{remark}[The Eisenstein fan]
\label{rem:eisenstein-fan}
\index{Eisenstein series!fan structure}
\index{quasi-modular forms!and bar curvature}
The Laurent expansion of $\wp$ and its derivatives
near $z=0$ on $E_\tau$ takes the form
\[
\wp^{(n)}(z\,|\,\tau)
\;=\;
\frac{(-1)^n\,(n+1)!}{z^{n+2}}
\;+\;
\sum_{k\ge 1}
  \frac{(2k+n+1)!}{(2k)!\,(n+1)!}\,
  G_{2k+n+2}(\tau)\,z^{2k},
\]
where $G_{2k}(\tau)=\sum'_{(m,n)\in\mathbb{Z}^2}(m\tau+n)^{-2k}$
are the Eisenstein series.  Each OPE pole of order $2k$
in the $\lambda$-bracket thus generates a $G_{2k}(\tau)$
correction to the elliptic $r$-matrix, producing a fan of
modular corrections indexed by the pole spectrum.

The $G_2$ correction occupies a distinguished position:
it is \emph{quasi-modular}, transforming under
$\mathrm{SL}(2,\mathbb{Z})$ with the anomalous shift
\[
G_2(-1/\tau) \;=\; \tau^2\,G_2(\tau)-2\pi i\,\tau.
\]
This anomaly is the shadow of the Legendre relation
$\eta_1\tau-\eta_\tau=\pi i$, and it manifests in the
bar complex as the curvature $d^2=\kappa(\cA)\cdot\omega_1$
of the curved Swiss-cheese structure at genus~$1$.  More
precisely: the non-holomorphic completion
$\widehat{G}_2(\tau)=G_2(\tau)-\pi/\Imag\,\tau$
is modular, and the passage from $G_2$ to $\widehat{G}_2$
is the modular gauge transformation that renders the closed
colour well-defined on $\overline{\cM}_{1,n}$.  Whether
the analogous completion can be performed simultaneously
on the \emph{open} colour --- disentangling the braiding
from the curvature when $c_0\ne 0$ --- is the subject
of the following question.
\end{remark}

\begin{question}[The central genus-$1$ frontier]
\label{q:genus1-entanglement}
\index{genus 1!central open question}
\index{curvature--braiding!disentanglement problem}
\ClaimStatusOpen
Does the genus-$1$ coupling between the bar curvature
$d^2=\kappa(\cA)\cdot\omega_1$ and the spectral braiding
monodromy $2\eta_\tau\cdot c_0$ represent genuine entanglement
of the two Swiss-cheese colours, or can a modular gauge twist
$r_1(z;\tau)\mapsto r_1(z;\tau)+\phi(z;\tau)$ disentangle them?

More precisely: for $\cA$ with $c_0\ne 0$, does there exist
a doubly-periodic correction $\phi(z;\tau)$, meromorphic in~$z$
and modular in~$\tau$, such that
\[
\widetilde{r}_1(z;\tau)
\;:=\;
r_1(z;\tau)+\phi(z;\tau)
\]
is doubly periodic in~$z$ and the annular bar differential
$d^{\mathrm{ann}}$ computed from $\widetilde{r}_1$ remains
a differential \textup{(}$d^{\mathrm{ann}}\circ d^{\mathrm{ann}}
=\kappa\cdot\omega_1$, not something worse\textup{)}?

The obstruction lies in the Legendre relation.  The
quasi-period $\eta_\tau$ satisfies
$\eta_1\tau-\eta_\tau=\pi i$, which constrains the
$A$- and $B$-cycle monodromies simultaneously.  Any
correction $\phi$ that kills the $B$-cycle monodromy
$2\eta_\tau\cdot c_0$ must introduce an $A$-cycle defect of
magnitude $2\pi i\cdot c_0/\tau$, and this defect couples
to the modular parameter through the $\mathrm{SL}(2,\mathbb{Z})$
action.  The question reduces to whether the resulting
cocycle in
$H^1(\mathrm{SL}(2,\mathbb{Z}),\,
\mathcal{O}^{\mathrm{mer}}(E_\tau)\otimes\End(\cA))$
is trivial.

For $c_0=0$ the answer is trivially yes: no correction
is needed.  For $c_0\ne 0$, the Eisenstein fan suggests
the answer is no: the $G_2$ quasi-modularity anomaly
is intrinsic, and the entanglement between the two colours
at genus~$1$ is a genuine feature of the theory, not a
coordinate artifact.  But a proof --- or a counterexample
at some specific algebra --- remains open.
\end{question}


\section{Explicit $R$-matrices for the standard families}
\label{sec:explicit-r-matrices}

The general theory of \S\ref{sec:r-matrix-descent} constructs
the $R$-matrix as the monodromy of the Kohno connection
$\nabla=d-\sum_{i<j}r_{ij}(z_{ij})\,d\log(z_{ij})$.
For the Yangian $Y(\mathfrak{g})$ of a simple Lie algebra~$\mathfrak{g}$,
the rational $R$-matrix on the defining representation
$V=\mathbb{C}^N$ takes a closed form:
the \emph{Yang $R$-matrix} for type~$A$,
and the \emph{Zamolodchikov--Zamolodchikov $R$-matrix}
for types~$B$, $C$, $D$, and the exceptionals.
In each case, the classical limit
$R(u)=u\cdot\id_{V\otimes V}+\hbar\,r(u)+O(\hbar^2)$
recovers $r(u)=\Omega/u$, where $\Omega=\sum_a I^a\otimes I_a$
is the Casimir tensor (the image of the quadratic Casimir
under $\mathfrak{g}\hookrightarrow\operatorname{End}(V)^{\otimes 2}$).

All affine algebras $\hat\mathfrak{g}_k$ are $\Einf$-chiral
(every vertex algebra is local, hence $\Einf$).
The $R$-matrices below are derived from the local OPE
via the monodromy construction of
Construction~\ref{constr:r-matrix-monodromy-vol1};
they are \emph{not} independent input.

\subsection*{The universal formula: eigenvalue decomposition}

\begin{proposition}[Eigenvalue decomposition of the rational $R$-matrix]
\ClaimStatusProvedHere
\label{prop:r-matrix-eigenvalue}
Let $\mathfrak{g}$ be a simple Lie algebra, $V$ its defining
representation, and $R(u)\in\operatorname{End}(V\otimes V)(u)$ the
rational $R$-matrix of $Y(\mathfrak{g})$.  The tensor product
decomposes into irreducibles:
\[
V\otimes V \;=\; \bigoplus_\lambda V_\lambda.
\]
On each irreducible summand $V_\lambda$,
\begin{equation}\label{eq:r-eigenvalue}
R(u)\big|_{V_\lambda}
\;=\;
f_\lambda(u)\cdot\id_{V_\lambda},
\end{equation}
where $f_\lambda(u)$ is a rational function of~$u$
with $f_\lambda(u)\sim u$ as $u\to\infty$.
For type~$A$, $f_\lambda$ is linear: $f_\lambda(u)=u+c_\lambda$
with $c_\lambda=\tfrac{1}{2}(C_2(\lambda)-2\,C_2(V))$ and
$C_2(\mu)$ the eigenvalue of the quadratic Casimir on the
representation with highest weight~$\mu$.
For types $B$, $C$, $D$, and the exceptionals, $f_\lambda$
acquires rational corrections from the additional primitive
invariants in $\operatorname{End}_{\mathfrak{g}}(V\otimes V)$.
\end{proposition}

\begin{proof}
The $R$-matrix $R(u)$ commutes with the
diagonal $\mathfrak{g}$-action (it intertwines the $Y(\mathfrak{g})$
coproduct).  By Schur's lemma, $R(u)|_{V_\lambda}$ is a scalar
multiple of the identity on each irreducible $V_\lambda\subset V\otimes V$.
For type~$A$, the $R$-matrix $R(u)=u\cdot\id+P$ is polynomial of
degree one, so the scalar is linear in~$u$:
$f_\lambda(u)=u+c_\lambda$.
For types $B$/$C$/$D$, the trace/symplectic contraction~$Q$ or~$K$
introduces a single pole; for the exceptionals, multiple primitive
invariant tensors introduce multiple poles.  In every case, $f_\lambda$
is a rational function with $f_\lambda(u)\sim u$ as $u\to\infty$.
\end{proof}

\subsection*{Type $A_1$ ($\mathfrak{sl}_2$): the Yang $R$-matrix}

\begin{theorem}[Yang $R$-matrix for $\mathfrak{sl}_2$]
\ClaimStatusProvedHere
\label{thm:yang-r-sl2}
Let $V=\mathbb{C}^2$ be the defining representation of\/ $\mathfrak{sl}_2$.
The rational $R$-matrix of the Yangian $Y(\mathfrak{sl}_2)$ on
$V\otimes V$ is
\begin{equation}\label{eq:yang-sl2}
R(u)
\;=\;
u\cdot\id_{V\otimes V}
\;+\;
P,
\end{equation}
where $P\in\operatorname{End}(\mathbb{C}^2\otimes\mathbb{C}^2)$ is the
permutation operator $P(v\otimes w)=w\otimes v$.

The tensor product decomposes as
$\mathbb{C}^2\otimes\mathbb{C}^2=V_{\mathrm{sym}}\oplus V_{\mathrm{alt}}$,
where $V_{\mathrm{sym}}\cong\mathrm{Sym}^2(\mathbb{C}^2)$ is the adjoint
representation \textup{(}spin~$1$, dimension~$3$\textup{)} and
$V_{\mathrm{alt}}\cong\bigwedge^2(\mathbb{C}^2)$ is the trivial
representation \textup{(}spin~$0$, dimension~$1$\textup{)}.  The eigenvalues are:
\[
R(u)\big|_{V_{\mathrm{sym}}}=(u+1)\cdot\id,
\qquad
R(u)\big|_{V_{\mathrm{alt}}}=(u-1)\cdot\id.
\]

In the standard basis $\{e_1\otimes e_1,\,e_1\otimes e_2,\,
e_2\otimes e_1,\,e_2\otimes e_2\}$ of $\mathbb{C}^2\otimes\mathbb{C}^2$,
\begin{equation}\label{eq:yang-sl2-matrix}
R(u)
\;=\;
\begin{pmatrix}
u+1 & 0 & 0 & 0 \\
0 & u & 1 & 0 \\
0 & 1 & u & 0 \\
0 & 0 & 0 & u+1
\end{pmatrix}.
\end{equation}
\end{theorem}

\begin{proof}
The permutation operator has eigenvalue $+1$ on symmetric tensors and
$-1$ on antisymmetric tensors, giving the eigenvalue decomposition
directly from $R(u)=u\cdot\id+P$.

\emph{Yang--Baxter verification.}
The YBE
$R_{12}(u{-}v)\,R_{13}(u)\,R_{23}(v)
=R_{23}(v)\,R_{13}(u)\,R_{12}(u{-}v)$
in $\operatorname{End}((\mathbb{C}^2)^{\otimes 3})$ is verified by
direct expansion.  Both sides equal
\[
(u{-}v)\,u\,v\cdot\id
+(u{-}v)\,v\cdot P_{23}
+u(u{-}v)\cdot P_{13}
+u\,v\cdot P_{12}
+u\cdot P_{12}P_{13}
+v\cdot P_{13}P_{23}
+P_{12}P_{13}P_{23}.
\]
This uses only $P_{ij}^2=\id$ and
$P_{12}P_{13}P_{23}=P_{23}P_{13}P_{12}$ in
$\mathbb{C}[\Sigma_3]\subset\operatorname{End}(V^{\otimes 3})$, so
the YBE holds for any vector space~$V$, not just $\mathbb{C}^2$.

\emph{Classical limit.}
At large $|u|$, $R(u)/u=\id+P/u$, so the classical $r$-matrix is
$r(u)=P/u$.  Since $P$ acts as $\Omega$ (the Casimir tensor) on
$V\otimes V$ for $\mathfrak{sl}_N$, we recover $r(u)=\Omega/u$
satisfying the CYBE~\eqref{eq:cybe-vol1}.
\end{proof}

\subsection*{Type $A_2$ ($\mathfrak{sl}_3$): the Yang $R$-matrix}

\begin{theorem}[Yang $R$-matrix for $\mathfrak{sl}_3$]
\ClaimStatusProvedHere
\label{thm:yang-r-sl3}
Let $V=\mathbb{C}^3$ be the defining representation of\/ $\mathfrak{sl}_3$.
The rational $R$-matrix of\/ $Y(\mathfrak{sl}_3)$ on $V\otimes V$ is
\begin{equation}\label{eq:yang-sl3}
R(u)
\;=\;
u\cdot\id_{V\otimes V}
\;+\;
P.
\end{equation}

The decomposition
$\mathbb{C}^3\otimes\mathbb{C}^3
=\mathrm{Sym}^2(\mathbb{C}^3)\oplus\bigwedge^2(\mathbb{C}^3)$
gives the irreducible summands
$V_{(2,0)}\cong\mathbf{6}$ \textup{(}symmetric, highest weight
$(2,0)$\textup{)} and
$V_{(0,1)}\cong\overline{\mathbf{3}}$ \textup{(}alternating, highest weight
$(0,1)$\textup{)}.  The eigenvalues are:
\[
R(u)\big|_{\mathbf{6}}=(u+1)\cdot\id,
\qquad
R(u)\big|_{\overline{\mathbf{3}}}=(u-1)\cdot\id.
\]

In the standard basis $\{e_i\otimes e_j\}_{1\leq i,j\leq 3}$
of $\mathbb{C}^3\otimes\mathbb{C}^3$, ordered lexicographically,
$R(u)$ is the $9\times 9$ matrix with entries
$R(u)_{(ij),(kl)}
=u\,\delta_{ik}\delta_{jl}+\delta_{il}\delta_{jk}$:
diagonal entries~$u+1$ when $i=j$, and in each $2\times 2$ block
$\{e_i\otimes e_j,\,e_j\otimes e_i\}$ with $i\neq j$, diagonal
entries~$u$ and off-diagonal entries~$1$.
\end{theorem}

\begin{proof}
The Yang $R$-matrix for $\mathfrak{sl}_N$ on the defining
representation is universally $R(u)=u\cdot\id+P$, and the
eigenvalue decomposition on symmetric and alternating tensors is
$u+1$ and $u-1$ respectively, independent of~$N$.

The Yang--Baxter verification is identical to that
for $\mathfrak{sl}_2$: the formula $R(u)=u\cdot\id+P$ satisfies the YBE
in $\operatorname{End}(V^{\otimes 3})$ for any vector space~$V$,
since the proof uses only the symmetric group algebra relations.

The classical limit $r(u)=P/u=\Omega/u$ satisfies the
CYBE~\eqref{eq:cybe-vol1}.
\end{proof}

\subsection*{Type $A_3$ ($\mathfrak{sl}_4$): the Yang $R$-matrix}

\begin{theorem}[Yang $R$-matrix for $\mathfrak{sl}_4$]
\ClaimStatusProvedHere
\label{thm:yang-r-sl4}
Let $V=\mathbb{C}^4$ be the defining representation of\/ $\mathfrak{sl}_4$.
The rational $R$-matrix of\/ $Y(\mathfrak{sl}_4)$ on $V\otimes V$ is
\begin{equation}\label{eq:yang-sl4}
R(u)
\;=\;
u\cdot\id_{V\otimes V}
\;+\;
P,
\end{equation}
with the decomposition
$\mathbb{C}^4\otimes\mathbb{C}^4
=V_{(2,0,0)}\oplus V_{(0,1,0)}$,
where $V_{(2,0,0)}\cong\mathrm{Sym}^2(\mathbb{C}^4)\cong\mathbf{10}$
and $V_{(0,1,0)}\cong\bigwedge^2(\mathbb{C}^4)\cong\mathbf{6}$.
The eigenvalues are:
\[
R(u)\big|_{\mathbf{10}}=(u+1)\cdot\id,
\qquad
R(u)\big|_{\mathbf{6}}=(u-1)\cdot\id.
\]
\end{theorem}

\begin{proof}
The Yang $R$-matrix $R(u)=u\cdot\id+P$
is universal for all $\mathfrak{sl}_N$.  The eigenvalue $+1$ of~$P$
on symmetric tensors and $-1$ on alternating tensors gives eigenvalues
$u+1$ and $u-1$ of~$R(u)$.  The YBE holds identically in
$\operatorname{End}((\mathbb{C}^4)^{\otimes 3})$ by the same argument.
\end{proof}

\begin{remark}[The type-$A$ universality]
\label{rem:type-a-pattern}
For all $\mathfrak{sl}_N$, the Yang $R$-matrix on the defining
representation is $R(u)=u\cdot\id+P$, with two eigenvalues:
$u+1$ on $\mathrm{Sym}^2(\mathbb{C}^N)$ of dimension $\binom{N+1}{2}$,
and $u-1$ on $\bigwedge^2(\mathbb{C}^N)$ of dimension $\binom{N}{2}$.
The classical $r$-matrix is universally $r(u)=\Omega/u$.
The simplicity of the type-$A$ $R$-matrix --- a linear function
of the permutation operator alone --- reflects the
fact that for $\mathfrak{sl}_N$ the defining tensor product has exactly
two irreducible summands, so the centraliser algebra
$\operatorname{End}_{\mathfrak{sl}_N}(V\otimes V)$ is two-dimensional
(spanned by $\id$ and~$P$).
Types $B$, $C$, $D$, and the exceptionals have
$V\otimes V$ with three or more summands, and the centraliser
is correspondingly larger, necessitating additional correction terms.
\end{remark}

\subsection*{Type $B_2$ ($\mathfrak{so}_5$): the orthogonal $R$-matrix}

\begin{theorem}[Zamolodchikov $R$-matrix for $\mathfrak{so}_5$]
\ClaimStatusProvedHere
\label{thm:r-so5}
Let $V=\mathbb{C}^5$ be the defining \textup{(}vector\textup{)}
representation of\/ $\mathfrak{so}_5$, with invariant symmetric bilinear
form $(\cdot,\cdot)$.  Define:
\begin{itemize}
\item $P\in\operatorname{End}(V\otimes V)$: the permutation,
  $P(v\otimes w)=w\otimes v$;
\item $Q\in\operatorname{End}(V\otimes V)$: the trace contraction,
  $Q(v\otimes w)=(v,w)\sum_{i=1}^{5} e_i\otimes e_i$,
  where $\{e_i\}$ is an orthonormal basis.
  The key property is $Q^2=N\cdot Q$ with $N=5$.
\end{itemize}
The rational $R$-matrix of\/ $Y(\mathfrak{so}_5)$ on $V\otimes V$ is
\begin{equation}\label{eq:r-so5}
R(u)
\;=\;
u\cdot\id_{V\otimes V}
\;+\;
P
\;-\;
\frac{Q}{u+\kappa},
\qquad
\kappa=\tfrac{N-2}{2}=\tfrac{3}{2},
\end{equation}
where $N=5$ is the dimension of~$V$.

The tensor product decomposes into three irreducibles:
\begin{equation}\label{eq:so5-decomp}
\mathbb{C}^5\otimes\mathbb{C}^5
\;=\;
V_{(2,0)}\;\oplus\; V_{(0,1)}\;\oplus\; V_{(0,0)},
\end{equation}
where:
\begin{itemize}
\item $V_{(2,0)}\cong\mathbf{14}$: the traceless symmetric tensors
  \textup{(}highest weight $(2,0)$ of $B_2$\textup{)},
  $\dim=\tfrac{1}{2}N(N+1)-1=14$;
\item $V_{(0,1)}\cong\mathbf{10}$: the alternating tensors
  \textup{(}the adjoint $\cong\mathfrak{so}_5$, highest weight
  $(0,1)$\textup{)}, $\dim=\tfrac{1}{2}N(N-1)=10$;
\item $V_{(0,0)}\cong\mathbf{1}$: the trace scalar
  \textup{(}trivial representation\textup{)}, $\dim=1$.
\end{itemize}
The eigenvalues of $R(u)$ are:
\begin{align}
R(u)\big|_{\mathbf{14}}
&=(u+1)\cdot\id,
\label{eq:so5-eigen-14}\\
R(u)\big|_{\mathbf{10}}
&=(u-1)\cdot\id,
\label{eq:so5-eigen-10}\\
R(u)\big|_{\mathbf{1}}
&=\frac{(u+1)(u+\tfrac{3}{2})-5}{u+\tfrac{3}{2}}\cdot\id
=\frac{u^2+\tfrac{5}{2}u-\tfrac{7}{2}}{u+\tfrac{3}{2}}\cdot\id.
\label{eq:so5-eigen-1}
\end{align}
\end{theorem}

\begin{proof}
The operators $\id$, $P$, $Q$ span the centraliser algebra
$\operatorname{End}_{\mathfrak{so}_N}(V\otimes V)$ (the Brauer algebra
at parameter~$N$).  The most general $\mathfrak{so}_N$-invariant
$R$-matrix of Yangian form is $R(u)=u\cdot\id+\alpha P+\beta Q/(u+\gamma)$.
The normalisation fixes $\alpha=1$; the YBE determines
$\beta=-1$ and $\gamma=\kappa=(N-2)/2$.

\emph{Yang--Baxter verification.}
Expanding each $R_{ij}$ and using the Brauer algebra relations
$P_{ij}^2=\id$,
$P_{12}P_{13}=P_{23}P_{12}=P_{13}P_{23}$,
$Q_{ij}^2=N\,Q_{ij}$,
$Q_{12}Q_{13}=N\,Q_{13}$,
$P_{12}Q_{13}=Q_{13}P_{23}$,
$Q_{12}P_{13}=Q_{12}P_{23}$,
$P_{ij}Q_{ij}=Q_{ij}$,
all terms match provided $\kappa=(N-2)/2$.

\emph{Eigenvalue computation.}
$Q$ is rank-one with image $V_{(0,0)}$ and kernel
$V_{(2,0)}\oplus V_{(0,1)}$.
On $V_{(0,0)}$: $Q|_{V_{(0,0)}}=N\cdot\id$ and $P=+\id$.
So on $V_{(2,0)}$: $R(u)=(u+1)\cdot\id$;
on $V_{(0,1)}$: $R(u)=(u-1)\cdot\id$;
on $V_{(0,0)}$:
$R(u)=(u+1-N/(u+\kappa))\cdot\id
=((u+1)(u+\kappa)-N)/(u+\kappa)\cdot\id$.

\emph{Classical limit.}
$R(u)=u\cdot\id+P-Q/u+O(u^{-2})$ recovers $r(u)=\Omega/u$
with $\Omega=P-Q/N$.
\end{proof}

\begin{remark}[The rank-$2$ isogeny $\mathfrak{so}_5\cong\mathfrak{sp}_4$]
\label{rem:b2-c2-isogeny}
The algebras $\mathfrak{so}_5$ and $\mathfrak{sp}_4$ are isomorphic,
but the $R$-matrices in Theorems~\ref{thm:r-so5}
and~\ref{thm:r-sp4} act on different representations
($5$-dimensional vector vs.\ $4$-dimensional defining), yielding
different $R$-matrices.
\end{remark}

\subsection*{Type $C_2$ ($\mathfrak{sp}_4$): the symplectic $R$-matrix}

\begin{theorem}[Zamolodchikov $R$-matrix for $\mathfrak{sp}_4$]
\ClaimStatusProvedHere
\label{thm:r-sp4}
Let $V=\mathbb{C}^4$ be the defining representation of\/
$\mathfrak{sp}_4$, with invariant skew-symmetric bilinear form
$\omega\in\bigwedge^2(V^*)$.  Define:
\begin{itemize}
\item $P\in\operatorname{End}(V\otimes V)$: the permutation;
\item $K\in\operatorname{End}(V\otimes V)$: the symplectic trace
  contraction, $K(v\otimes w)=\omega(v,w)\,\omega^{-1}$, where
  $\omega^{-1}\in\bigwedge^2(V)$ is the dual bivector.
  The key property is $K^2=-2n\cdot K=-4K$ \textup{(}$2n=N=4$\textup{)},
  the sign from skew-symmetry of~$\omega$.
\end{itemize}
The rational $R$-matrix of\/ $Y(\mathfrak{sp}_4)$ on $V\otimes V$ is
\begin{equation}\label{eq:r-sp4}
R(u)
\;=\;
u\cdot\id_{V\otimes V}
\;+\;
P
\;+\;
\frac{K}{u-\kappa},
\qquad
\kappa=\tfrac{N+2}{2}=3,
\end{equation}
where $N=2n=4$.

The tensor product decomposes into three irreducibles:
\begin{equation}\label{eq:sp4-decomp}
\mathbb{C}^4\otimes\mathbb{C}^4
\;=\;
V_{(2,0)}\;\oplus\; V_{(0,1)}\;\oplus\; V_{(0,0)},
\end{equation}
where:
\begin{itemize}
\item $V_{(2,0)}\cong\mathbf{10}$: the full symmetric square
  $\mathrm{Sym}^2(\mathbb{C}^4)$, irreducible for
  $\mathfrak{sp}_{2n}$ \textup{(}the $\omega$-trace lives in
  $\bigwedge^2$ because $\omega$ is skew\textup{)}, $\dim=10$;
\item $V_{(0,1)}\cong\mathbf{5}$: the $\omega$-traceless alternating
  tensors \textup{(}adjoint of $C_2$\textup{)},
  $\dim=\tfrac{1}{2}N(N-1)-1=5$;
\item $V_{(0,0)}\cong\mathbf{1}$: the symplectic scalar
  $\mathrm{span}\{\omega^{-1}\}\subset\bigwedge^2(V)$, $\dim=1$.
\end{itemize}
The eigenvalues of $R(u)$ are:
\begin{align}
R(u)\big|_{\mathbf{10}}&=(u+1)\cdot\id,
\label{eq:sp4-eigen-10}\\
R(u)\big|_{\mathbf{5}}&=(u-1)\cdot\id,
\label{eq:sp4-eigen-5}\\
R(u)\big|_{\mathbf{1}}&=\frac{(u-1)(u-3)-4}{u-3}\cdot\id
=\frac{u^2-4u-1}{u-3}\cdot\id.
\label{eq:sp4-eigen-1}
\end{align}
\end{theorem}

\begin{proof}
The centraliser $\operatorname{End}_{\mathfrak{sp}_{2n}}(V\otimes V)$
is spanned by $\id$, $P$, $K$ (Brauer algebra at parameter $-2n$).
The orthogonal-symplectic duality $N\leftrightarrow -2n$ sends
$-Q/(u+(N-2)/2)$ to $+K/(u-(n+1))=+K/(u-\kappa)$
with $\kappa=(N+2)/2$.

\emph{Yang--Baxter verification.}
Same Brauer-algebra calculation as type~$B$, with $Q\to K$,
$N\to -2n$, and $K^2=-2n\cdot K$ replacing $Q^2=N\cdot Q$.

\emph{Eigenvalue computation.}
$K$ is rank-one with image $V_{(0,0)}$ and kernel
$V_{(2,0)}\oplus V_{(0,1)}$.
On $V_{(2,0)}=\mathrm{Sym}^2(V)$: $K=0$, $P=+\id$, so
$R(u)|_{\mathbf{10}}=(u+1)\cdot\id$.
On $V_{(0,1)}$: $K=0$, $P=-\id$, so
$R(u)|_{\mathbf{5}}=(u-1)\cdot\id$.
On $V_{(0,0)}$: $\omega^{-1}\in\bigwedge^2(V)$, so $P=-\id$;
from $K^2=-2nK$ and the rank-one structure,
$K|_{V_{(0,0)}}=-2n\cdot\id=-4\cdot\id$.
So $R(u)|_{\mathbf{1}}=(u-1-4/(u-3))\cdot\id
=(u^2-4u-1)/(u-3)\cdot\id$.

\emph{Classical limit.}
$R(u)=u\cdot\id+P+K/u+O(u^{-2})$ recovers $r(u)=\Omega/u$
with $\Omega=P+K/(2n)$.
\end{proof}

\subsection*{Type $G_2$: the exceptional $R$-matrix}

\begin{theorem}[Rational $R$-matrix for $G_2$]
\ClaimStatusProvedHere
\label{thm:r-g2}
Let $V=\mathbb{C}^7$ be the $7$-dimensional defining representation
of the exceptional Lie algebra~$G_2$.  The tensor product decomposes as
\begin{equation}\label{eq:g2-decomp}
\mathbb{C}^7\otimes\mathbb{C}^7
\;=\;
V_{\mathbf{27}}\;\oplus\; V_{\mathbf{14}}\;\oplus\;
V_{\mathbf{7}}\;\oplus\; V_{\mathbf{1}},
\end{equation}
where:
\begin{itemize}
\item $V_{\mathbf{27}}$: highest weight $(2,0)$, traceless
  symmetric square, $\dim=27$;
\item $V_{\mathbf{14}}$: highest weight $(0,1)$, adjoint
  $\cong\mathfrak{g}_2$, $\dim=14$;
\item $V_{\mathbf{7}}$: highest weight $(1,0)$, a copy of the
  defining representation inside $\bigwedge^2(V)$
  \textup{(}image of contraction with the
  $G_2$-invariant $3$-form $\varphi\in\bigwedge^3(V^*)$\textup{)},
  $\dim=7$;
\item $V_{\mathbf{1}}$: trivial, $\dim=1$.
\end{itemize}

The Yangian $R$-matrix is the unique $\mathfrak{g}_2$-invariant
rational YBE solution on $V\otimes V$ with
$R(u)\sim u\cdot\id$ as $u\to\infty$.  Let
$\Pi_\lambda$ denote the projector onto $V_\lambda$.  Then
\begin{equation}\label{eq:r-g2}
R(u)
\;=\;
\sum_{\lambda\in\{\mathbf{27},\mathbf{14},\mathbf{7},\mathbf{1}\}}
f_\lambda(u)\cdot\Pi_\lambda,
\end{equation}
with eigenvalues
\begin{equation}\label{eq:g2-eigenvalues}
\begin{aligned}
f_{\mathbf{27}}(u)&=u+1,\\
f_{\mathbf{14}}(u)&=u-1,\\
f_{\mathbf{7}}(u)&=\frac{(u+1)(u-2)}{u+2},\\
f_{\mathbf{1}}(u)&=\frac{(u+1)(u-4)}{u+4}.
\end{aligned}
\end{equation}
The poles at $u=-2$ and $u=-4$ correspond to the two primitive
$G_2$-invariant tensors in $\operatorname{End}(V\otimes V)$ beyond
the identity and permutation.
\end{theorem}

\begin{proof}
$\operatorname{End}_{G_2}(V\otimes V)$ is four-dimensional,
spanned by $\Pi_{\mathbf{27}}$, $\Pi_{\mathbf{14}}$,
$\Pi_{\mathbf{7}}$, $\Pi_{\mathbf{1}}$.  The most general
$\mathfrak{g}_2$-invariant $R$-matrix is
$R(u)=\sum_\lambda f_\lambda(u)\cdot\Pi_\lambda$.

\emph{Determination of eigenvalues.}
The $f_\lambda$ are determined by the YBE, unitarity
$R(u)R_{21}(-u)\propto\id$, and crossing symmetry.
The explicit computation was carried out by
Kuniba--Nakanishi--Suzuki (fusion procedure) and
Ogievetsky (projection method), yielding~\eqref{eq:g2-eigenvalues}.

The zeros $u=2$ in $f_{\mathbf{7}}$ and $u=4$ in
$f_{\mathbf{1}}$ are crossing points where the $R$-matrix
degenerates to a projector.

\emph{Yang--Baxter verification.}
On multiplicity-free sectors of $V^{\otimes 3}$, the YBE reduces
to $f_\lambda(u{-}v)f_\mu(u)f_\nu(v)
=f_\nu(v)f_\mu(u)f_\lambda(u{-}v)$, trivially satisfied.
$V_{\mathbf{7}}$ appears with multiplicity~$2$ in
$V^{\otimes 3}$; the resulting $2\times 2$ matrix YBE was verified
by Kuniba--Nakanishi--Suzuki via the $G_2$ Clebsch--Gordan
decomposition.

\emph{Classical limit.}
$f_\lambda(u)=u+c_\lambda+O(u^{-1})$ for each~$\lambda$,
so $R(u)=u\cdot\id+\Omega+O(u^{-1})$ with
$\Omega=\sum_\lambda c_\lambda\Pi_\lambda$ the Casimir tensor,
and $r(u)=\Omega/u$ satisfies the CYBE~\eqref{eq:cybe-vol1}.
\end{proof}

\begin{remark}[$G_2$ structural features]
\label{rem:g2-structure}
\begin{enumerate}[label=(\roman*)]
\item \emph{Four eigenvalues.}  Types $A$ have $2$;
  types $B$, $C$, $D$ have $3$; $G_2$ has $4$.
\item \emph{Two poles} at $u=-2$ and $u=-4$
  (one for types $B$/$C$/$D$, none for type~$A$).
\item \emph{Multiplicity in $V^{\otimes 3}$}: the $\mathbf{7}$
  has multiplicity~$2$, requiring a $2\times 2$ matrix YBE.
\item \emph{$\mathfrak{so}_7$ refinement.}
  $G_2\hookrightarrow\mathrm{SO}_7$ splits
  $\mathbf{27}_{\mathfrak{so}_7}\to
  \mathbf{27}_{G_2}\oplus\mathbf{7}_{G_2}$
  via the $3$-form $\varphi$; the $\mathfrak{so}_7$ eigenvalue
  $u+1$ on $\mathbf{27}_{\mathfrak{so}_7}$ splits into $u+1$
  on $\mathbf{27}_{G_2}$ and $(u+1)(u-2)/(u+2)$ on $\mathbf{7}_{G_2}$.
\end{enumerate}
\end{remark}

\begin{remark}[The universal classical $r$-matrix]
\label{rem:universal-classical}
For every simple $\mathfrak{g}$ and every finite-dimensional $V$,
the classical limit of the Yangian $R$-matrix is
\[
r(u)=\frac{\Omega}{u},
\qquad
\Omega=\sum_a I^a\otimes I_a
\in\mathfrak{g}\otimes\mathfrak{g}
\hookrightarrow
\operatorname{End}(V)\otimes\operatorname{End}(V),
\]
with $\{I^a\}$, $\{I_a\}$ dual bases for the Killing form.
The CYBE follows from the Jacobi identity:
$[\Omega_{12},\Omega_{13}]+[\Omega_{12},\Omega_{23}]
+[\Omega_{13},\Omega_{23}]=0$ is the Jacobi identity
in $\mathfrak{g}^{\otimes 3}$.

Quantum corrections distinguish the families:
type~$A$ has none; types $B$/$C$/$D$ have one rational correction
from the trace/symplectic contraction; the exceptionals have
multiple corrections from higher-order invariant tensors.
\end{remark}

\appendix

\section{Appendix: signs and conventions}

All formulas above are invariant under the usual choice of suspension conventions, provided one keeps
the same convention simultaneously in the ordered bar complex, the shuffle map, and the twisting
differential.  For clarity we summarize the minimal sign content.

\begin{enumerate}[label=(\arabic*)]
\item Reversal in Theorem~\ref{thm:opposite} uses the standard suspended Koszul sign
\[
(-1)^{\sum_{i<j}(|a_i|+1)(|a_j|+1)}.
\]

\item The left twisting term in Definition~\ref{def:Kbi} carries the sign needed to move the
coderivation on $C$ past the first coalgebra leg before applying the left $A$-action.

\item The right twisting term carries the additional sign needed to move the twisting cochain across
the middle bimodule factor.

\item In the strongly admissible regime the filtration is complete and bounded below, so all spectral
sequence and homotopy-transfer arguments used in the proofs converge.
\end{enumerate}

\section{Appendix: new results proved in this chapter}

For quick reference, the new results proved in this chapter are as follows.

\begin{enumerate}[label=(\arabic*)]
\item Ordered chiral shuffle theorem (Theorem~\ref{thm:shuffle}).
\item Opposite-duality and anti-involution transport (Theorem~\ref{thm:opposite},
Corollary~\ref{cor:anti}).
\item Enveloping duality (Corollary~\ref{cor:enveloping}).
\item Two-sided twisted bicomodule and its dg structure (Definition~\ref{def:Kbi},
Lemma~\ref{lem:Kbi-dg}).
\item Identification of the tangent ordered bar complex with the twisted bicomodule
(Theorem~\ref{thm:tangent=K}).
\item PBW-complete bimodule/bicomodule equivalence (Theorem~\ref{thm:bimod-bicomod}).
\item Diagonal correspondence (Theorem~\ref{thm:diagonal}).
\item Composition unit theorem (Corollary~\ref{cor:unit}).
\item Tensor--cotensor gluing (Corollary~\ref{cor:tensor-cotensor}).
\item Hochschild/coHochschild homology and cohomology equivalence
(Theorems~\ref{thm:HH-coHH-homology} and \ref{thm:HH-coHH-cohomology}).
\item Infinitesimal annular variation (Proposition~\ref{prop:infann}).
\item Ordered genus-zero open trace formalism (Theorem~\ref{thm:ordered-open}).
\item Calabi--Yau transport to an ordered Frobenius structure (Theorem~\ref{thm:CY}).
\item Open-colour double bar of $\hat{\mathfrak{sl}}_{2,k}$
(Proposition~\ref{prop:double-bar-sl2}).
\item Central extension invisible to the open-colour double bar
(Theorem~\ref{thm:central-extension-invisible}).
\item Two-colour double Koszul duality is involutive
(Theorem~\ref{thm:two-colour-double-kd}).
\item Non-redundancy of the two colours
(Corollary~\ref{cor:two-colours-non-redundant}).
\item Heisenberg ordered bar complex: trivial differential, deconcatenation
coproduct, collision residue $r^{\mathrm{coll}}(z) = k/z$,
$R$-matrix $R(z) = e^{k\hbar/z}$
(Theorems~\ref{thm:heisenberg-ordered-bar}
and \ref{thm:heisenberg-rmatrix}).
\item Abelian Yangian $Y(\mathfrak{u}(1))$ as ordered Koszul dual
of $\cH_k$ (Theorem~\ref{thm:heisenberg-yangian}).
\item Formality of $\Barch(\cH_k)$: class~G, shadow depth~$2$,
$m_k = 0$ for $k\ge 3$ (Theorem~\ref{thm:heisenberg-formality}).
\item Universal Kac--Moody Yangian theorem: for every simple
$\mathfrak{g}$ and non-critical $k\neq -h^\vee$,
$A^{!,\mathrm{ord}}_{\mathrm{line}}
(\widehat{\mathfrak{g}}_k)\simeq Y_\hbar(\mathfrak{g})$
at $\hbar = 1/(k+h^\vee)$; collision residue
$r(z) = \hbar\,\Omega/z$; CYBE from Jacobi; complete
strictification by root-space one-dimensionality; curvature
$\kappa(Y_\hbar(\mathfrak{g})) = -\dim(\mathfrak{g})\cdot
(k+h^\vee)/(2h^\vee)$
(Theorem~\ref{thm:km-yangian},
Table~\ref{tab:km-yangian-data}).
\item $\mathcal{W}_3$ ordered bar complex via DS transport: generators,
pole structure, A$_\infty$ operations (class~M), curvature $5c/6$,
complementarity $K_3 = 100$, self-dual at $c^* = 50$
(Theorem~\ref{thm:w3-ordered-bar}).
\item Class~M persistence under DS transport for all $\mathcal{W}_N$:
Koszulness preserved, formality destroyed, $K_N = 2(N{-}1)(2N^2{+}2N{+}1)$
(Theorem~\ref{thm:class-m-ds-transport}).
\item $\beta\gamma$ ordered bar complex: nilpotent collision residue
$\Theta = \beta\otimes\gamma - \gamma\otimes\beta$ with $\Theta^2 = 0$,
R-matrix $R = \id + 2\pi i\,\Theta$ (nilpotent, terminates at linear
order), Koszul dual a $4$-dimensional exterior algebra (class~C,
shadow depth~$4$)
(Theorems~\ref{thm:bg-ordered-bar}, \ref{thm:bg-rmatrix},
and \ref{thm:bg-koszul-dual}).
\item $bc$ ghost ordered bar complex: nilpotent collision residue
$\Theta_{bc} = b\otimes c - c\otimes b$ with $\Theta_{bc}^2 = 0$,
R-matrix $R = \id + 2\pi i\,\Theta_{bc}$ (nilpotent), Koszul dual
a symmetric algebra (class~C, shadow depth~$4$, complementarity
$K = 0$) (Theorem~\ref{thm:bc-ordered-bar}).
\item Wakimoto bar complex descent: BRST reduction of free-field
ordered bar complex recovers affine ordered bar complex; dg Yangian
$Y^{\mathrm{dg}}_\hbar(\fg)$ as BRST reduction of free-field Yangian
(Theorem~\ref{thm:wakimoto-ordered-bar}).
\item Lattice vertex algebras: collision residue diagonal in
$\alpha$-sector, R-matrix from OPE monodromy for symmetric cocycles
($\Einf$), independent braiding for non-symmetric cocycles (genuinely
$\Eone$), Koszul dual with inverse cocycle for self-dual lattices
(Theorems~\ref{thm:lattice-symmetric-ordered-bar},
\item Root-space one-dimensionality for simple Lie algebras
(Theorem~\ref{thm:root-space-one-dim-v1}).
\item Jacobi collapse for star sectors
(Lemma~\ref{lem:jacobi-collapse-v1}).
\item Dynkin coefficient via the beta integral: the universal BCH
coefficient $1/n = \int_0^1(1-t)^{n-1}\,dt$
(Theorem~\ref{thm:dynkin-beta-integral}).
\item Complete spectral Drinfeld strictification for all simple Lie
algebras: the spectral Drinfeld class vanishes at every filtration,
so the dg-shifted Yangian $Y_\hbar(\fg)$ carries a strict filtered
algebra structure (Theorem~\ref{thm:complete-strictification-v1}).
\ref{thm:lattice-nonsymmetric-ordered-bar},
and \ref{thm:lattice-ordered-koszul-dual}).
\item Evaluation homomorphism $\mathrm{ev}_a\colon Y_\hbar(\fg)\to U(\fg)$
for all simple types (Construction~\ref{constr:evaluation-map}).
\item $\mathfrak{sl}_2$ $R$-matrix on evaluation modules: $R(u)=u\,\id+\hbar P$
(Yang $R$-matrix), Clebsch--Gordan decomposition and non-semisimplicity
at resonance (Theorem~\ref{thm:sl2-R-matrix}, Corollary~\ref{cor:sl2-clebsch-gordan}).
\item $\mathfrak{sl}_3$ evaluation modules: fundamental and adjoint
(Computations~\ref{comp:sl3-eval-fundamental} and \ref{comp:sl3-eval-adjoint}).
\item Drinfeld polynomial classification of evaluation modules
(Proposition~\ref{prop:eval-drinfeld}).
\item Line category $\mathcal{C}_{\mathrm{line}}\simeq
Y_\hbar(\fg)\text{-}\mathrm{mod}^{\mathrm{fd}}$ on the affine Koszul locus
(Theorem~\ref{thm:line-category}).
\item Braided monoidal structure on evaluation modules from the universal
$R$-matrix, including coalescence singularity at $a=b$
(Theorem~\ref{thm:eval-braiding}).
\item Quantum group $U_q(\mathfrak{g})$ from genus-$1$ $B$-cycle monodromy
of the KZB connection, with $q = e^{2\pi i\hbar}$
(Theorem~\ref{thm:b-cycle-quantum-group}).
\item Drinfeld--Kohno theorem for the affine lineage: KZ/KZB monodromy
equals the quantum group braid representation
(Theorem~\ref{thm:drinfeld-kohno}).
\item Yangian--quantum group genus-$1$ deformation: $Y_\hbar(\mathfrak{g})$
at genus~$0$ deforms to $U_q(\mathfrak{g})$ at genus~$1$ via the
$B$-cycle exponential map $\widehat{\mathbb{G}}_a \to
\widehat{\mathbb{G}}_m$
(Theorem~\ref{thm:yangian-quantum-group}).
\item $U_q(\mathfrak{sl}_2)$ at roots of unity from affine
$\widehat{\mathfrak{sl}}_2$ at integer level; Kazhdan--Lusztig equivalence
(Corollary~\ref{cor:sl2-root-of-unity}).
\item Jones polynomial from genus-$1$ bar-complex monodromy via the
Reshetikhin--Turaev quantum trace
(Theorem~\ref{thm:jones-genus1}).
\item Annular bar complex: wrapping $\mathbb{R}\to S^1$ produces the
cyclic cotensor complex $B^{\mathrm{ann}}_\bullet(\cA)$ with four-component
differential $d^{\mathrm{ann}}=d_{\mathrm{int}}+d_{\mathrm{res}}+d_\Delta
+d_{\mathrm{wrap}}$ (Definition~\ref{def:annular-bar},
Theorem~\ref{thm:annular-bar-differential}).
\item Annular bar complex computes chiral Hochschild homology:
$B^{\mathrm{ann}}_\bullet(\cA)\simeq\HH^{\mathrm{ch}}_\bullet(\cA)$
(Theorem~\ref{thm:annular-HH}).
\item Curvature--braiding dichotomy at genus~$1$: the lowest
$\lambda$-bracket coefficient $c_0$ controls whether the genus-$1$
Swiss-cheese structure decouples ($c_0=0$, Cartan type) or entangles
($c_0\ne 0$, Yangian type) via elliptic quasi-periodicity
(Theorem~\ref{thm:curvature-braiding-dichotomy}).
\end{enumerate}

\bigskip

\noindent\textbf{Final principle.}
The diagonal is the boundary.  The bar coalgebra is the ambient ordered phase space.  The annulus is
the self-cotensor of the diagonal.  Once this is recognized, the associative chiral duality ceases
to be merely an equivalence between algebra and coalgebra and becomes a machine for open traces,
centers, and eventually bordered genus.

\begin{remark}[Non-commutative Euler product]%
\label{rem:e1-noncommutative-euler}%
The commutative sewing lift
$S_\cH(u) = \prod_p(1{-}p^{-u})^{-1}(1{-}p^{-u-1})^{-1}$
(Theorem~\ref{thm:dirichlet-weight-formula})
factors the Heisenberg sewing determinant over primes
via the unordered divisor function $\sigma_{-1}$.
The $\mathsf{E}_1$-ordered bar complex replaces
unordered by ordered factorizations; the analogous
ordered sewing lift should involve composition-weighted
divisor sums and factor in the ring of non-commutative
formal Dirichlet series.
\end{remark}


% ================================================================
\section{Rosetta stone: the $\mathsf{E}_1$ ordered bar complex across the landscape}
\label{sec:e1-rosetta-stone}
\index{Rosetta stone!ordered bar complex|textbf}
\index{ordered bar complex!landscape table|textbf}
\index{shadow depth!landscape classification|textbf}
% ================================================================

The ordered bar complex $\Barch(A)$ carries two layers of data:
the collision residue $r^{\mathrm{coll}}(z)$, which records the
holomorphic colour of the Swiss-cheese algebra, and
the $R$-matrix $R(z)$, which governs the passage from
ordered to unordered configurations.  The following table
collects this data for all standard families.

Every entry is an $\mathsf{E}_\infty$-chiral algebra
(a vertex algebra in the sense of Beilinson--Drinfeld:
the chiral product is defined on \emph{unordered}
configuration spaces with $\Sigma_n$-equivariant
factorization structure).  The ordered bar complex
$\Barch(A)$ retains the ordering as auxiliary bookkeeping;
the descent to the symmetric bar $B^\Sigma(A)$ is the
$R$-matrix-twisted quotient of Proposition~\ref*{prop:r-matrix-descent}.
None of the families below is $\mathsf{E}_1$-chiral in the sense
of Part~VII\@: genuinely $\mathsf{E}_1$-chiral algebras
(nonlocal, ordered configurations as primitive data) are the
Etingof--Kazhdan quantum vertex algebras, which lie outside
the Beilinson--Drinfeld framework.

\begin{theorem}[$\mathsf{E}_1$ ordered bar landscape;
\ClaimStatusProvedHere]%
\label{thm:e1-ordered-bar-landscape}%
\index{ordered bar complex!landscape theorem|textbf}%
\index{collision residue!landscape|textbf}%
\index{R-matrix!landscape|textbf}%
Let\/ $A$ be a strongly admissible augmented associative chiral
algebra from the standard landscape.  The ordered bar complex\/
$\Barch(A)$ has the following invariants, where\/ $r^{\mathrm{coll}}(z)$
denotes the collision residue
\textup{(}the OPE kernel after the $d\log$ bar propagator
absorbs one pole order\textup{)},
$R(z)$ is the spectral $R$-matrix on
$\mathrm{Conf}_2^{\mathrm{ord}}(\mathbb{C})$,
$r_{\max}$ is the shadow depth,
$A^!_{\mathrm{line}}$ is the ordered Koszul dual governing the
line category, and
``formal'' means the transferred $\Ainf$ operations
$m_k$ vanish on cohomology for all $k$ above the shadow depth:
\begin{center}
\small
\renewcommand{\arraystretch}{1.5}
\begin{tabular}{llcllcll}
\toprule
\textbf{Family}
  & \textbf{Class}
  & \makecell{\textbf{OPE} \\ \textbf{poles}}
  & \textbf{Collision residue}~$r^{\mathrm{coll}}(z)$
  & \textbf{$R$-matrix}~$R(z)$
  & $r_{\max}$
  & $A^!_{\mathrm{line}}$
  & \textbf{Formal?} \\
\midrule
$\mathcal{H}_k$
  & $\mathbf{G}$ & $2$
  & $k/z$
  & $e^{k\hbar/z}$
  & $2$
  & $\mathcal{H}_{-k}$\rlap{\textsuperscript{$\dagger$}}
  & yes \\[3pt]
$V_\Lambda$ \textup{(even)}
  & $\mathbf{G}$ & $2$
  & $\mathrm{diag}(\lambda_i)/z$
  & $\mathrm{diag}\bigl(e^{\lambda_i\hbar/z}\bigr)$
  & $2$
  & $V_{\Lambda^*}$\rlap{\textsuperscript{$\dagger$}}
  & yes \\[3pt]
$\widehat{\mathfrak{g}}_k$\; ($\mathfrak{sl}_2$)
  & $\mathbf{L}$ & $2$
  & $\hbar\,\Omega/z$
  & $u\,I + P/u$
  & $3$
  & $Y_\hbar(\mathfrak{sl}_2)$
  & yes \\[3pt]
$\widehat{\mathfrak{g}}_k$\; ($\mathfrak{sl}_3$)
  & $\mathbf{L}$ & $2$
  & $\hbar\,\Omega/z$
  & Yang ($9{\times}9$)
  & $3$
  & $Y_\hbar(\mathfrak{sl}_3)$
  & yes \\[3pt]
$\widehat{\mathfrak{g}}_k$\; (general)
  & $\mathbf{L}$ & $2$
  & $\hbar\,\Omega/z$
  & Yang ($\dim\mathfrak{g}^2$)
  & $3$
  & $Y_\hbar(\mathfrak{g})$
  & yes \\[3pt]
$\beta\gamma$
  & $\mathbf{C}$ & $1$
  & $\Theta$ \textup{(nilpotent)}
  & $I + 2\pi i\,\Theta$
  & $4$
  & $bc^!$
  & yes \\[3pt]
$bc$
  & $\mathbf{C}$ & $1$
  & $\Theta$ \textup{(nilpotent)}
  & $I - 2\pi i\,\Theta$
  & $4$
  & $\beta\gamma^!$
  & yes \\[3pt]
$\mathrm{Vir}_c$
  & $\mathbf{M}$ & $4$
  & $\dfrac{c/2}{z^3} + \dfrac{2T}{z}$
  & \textup{[class $\mathbf{M}$]}
  & $\infty$
  & $Y^{\mathrm{dg}}(\mathrm{Vir}_c)$
  & no \\[6pt]
$\mathcal{W}_3$
  & $\mathbf{M}$ & $6$
  & degree $5$ in $1/z$
  & \textup{[class $\mathbf{M}$]}
  & $\infty$
  & $Y^{\mathrm{dg}}(\mathcal{W}_3)$
  & no \\[3pt]
$\mathcal{W}_N$
  & $\mathbf{M}$ & $2N$
  & degree $2N{-}1$ in $1/z$
  & \textup{[class $\mathbf{M}$]}
  & $\infty$
  & $Y^{\mathrm{dg}}(\mathcal{W}_N)$
  & no \\
\bottomrule
\end{tabular}
\end{center}

\smallskip
\noindent
\textsuperscript{$\dagger$}The Koszul dual algebra $\mathcal{H}_k^!
= \mathrm{Sym}^{\mathrm{ch}}(V^*)$ is categorically distinct from
$\mathcal{H}_{-k}$ \textup{(}anti-pattern~AP33\textup{)}, but
both have $\kappa = -k$ and the line category
$\mathcal{C}_{\mathrm{line}} \simeq
\mathcal{H}_{-k}\text{-}\mathrm{mod}$
is equivalent.  Similarly, $V_\Lambda^!$ and $V_{\Lambda^*}$ share
the same modular characteristic but differ as chiral algebras.
\end{theorem}

\begin{proof}
The table assembles data from three sources.

\emph{OPE pole orders.}  The maximal OPE pole order on generators
is standard:
$J(z)J(w) \sim k/(z{-}w)^2$ for Heisenberg;
$J^a(z)J^b(w) \sim k\delta^{ab}/(z{-}w)^2 + f^{ab}{}_c J^c/(z{-}w)$
for affine Kac--Moody;
$\beta(z)\gamma(w) \sim 1/(z{-}w)$ and
$b(z)c(w) \sim 1/(z{-}w)$ for the free-field systems;
$T(z)T(w) \sim (c/2)/(z{-}w)^4 + 2T/(z{-}w)^2 + \partial T/(z{-}w)$
for Virasoro; and
$W_s(z)W_s(w)$ has poles up to order $2s$ for the
spin-$s$ generator of $\mathcal{W}_N$, giving
maximal pole $2N$ at the top spin.

\emph{Collision residues.}
The $d\log$ bar kernel absorbs one pole order
(Volume~I, Proposition~\ref*{V1-prop:dlog-absorption}):
the collision residue $r^{\mathrm{coll}}(z)$ from an OPE pole of
order~$d$ has pole of order~$d{-}1$.  For Heisenberg:
double-pole $k/(z{-}w)^2$ gives $r^{\mathrm{coll}} = k/z$.
For affine $\widehat{\mathfrak{g}}_k$: the Casimir
$\Omega = \sum_a t^a \otimes t_a$ produces
$r^{\mathrm{coll}} = \hbar\,\Omega/z$.
For $\beta\gamma$/$bc$: the simple-pole OPE $1/(z{-}w)$
gives a constant collision residue
$r^{\mathrm{coll}} = \Theta$, where
$\Theta \in \mathrm{End}(V)$ is the nilpotent
off-diagonal matrix encoding the pairing
($\Theta^2 = 0$).
For Virasoro: quartic pole gives
$(c/2)/z^3 + 2T/z$ (the $\partial T/(z{-}w)$ term of the OPE is
absorbed completely, producing no pole in the collision residue).

\emph{$R$-matrices.}
The $R$-matrix is the monodromy of the connection
$\nabla = d - r^{\mathrm{coll}} \cdot d\log(z)$ on
$\mathrm{Conf}_2^{\mathrm{ord}}(\mathbb{C})$.
For abelian algebras ($\mathcal{H}_k$, $V_\Lambda$):
the connection is scalar and
$R(z) = \exp(k\hbar/z)$
(Volume~II, Corollary~\ref*{cor:rosetta-heisenberg-projections}).
For affine $\widehat{\mathfrak{g}}_k$ at $\mathfrak{g}
= \mathfrak{sl}_2$: the $4 \times 4$ Yang $R$-matrix
$R(u) = uI + P$ (up to normalisation), where $P$ is
the permutation and $u = \hbar/z$ is the spectral parameter.
For $\beta\gamma$/$bc$: the nilpotent collision residue
$\Theta$ gives $R = \exp(2\pi i\,\Theta) = I + 2\pi i\,\Theta$
(the exponential terminates at the linear term since
$\Theta^2 = 0$).
For class~$\mathbf{M}$: the triple-pole collision residue
produces an $R$-matrix that is not expressible in closed form
as a finite product of elementary functions.

\emph{Shadow depth and formality.}
The shadow depth $r_{\max}$ is the highest arity at which the
transferred $\Ainf$ operation $m_k$ is non-vanishing on cohomology.
Class~$\mathbf{G}$: $m_{k \ge 3} = 0$ (abelian OPE, clean
Lagrangian intersection), so $r_{\max} = 2$ and the Swiss-cheese
structure is formal.
Class~$\mathbf{L}$: $m_3 \ne 0$ from the Lie bracket but
$m_{k \ge 4} = 0$ on cohomology (Jacobi closes the obstruction),
so $r_{\max} = 3$ and the structure is formal.
Class~$\mathbf{C}$: $m_4 \ne 0$ from the contact term but the
tower terminates at $r_{\max} = 4$ by rank-one abelian
rigidity, so the structure is formal.
Class~$\mathbf{M}$: $m_k \ne 0$ for all $k \ge 3$ on
cohomology (the quartic pole manufactures a propagating
obstruction), so $r_{\max} = \infty$ and the Swiss-cheese
structure is \emph{non-formal}.

\emph{Ordered Koszul duals.}
For affine $\widehat{\mathfrak{g}}_k$: the ordered Koszul dual is
the Yangian $Y_\hbar(\mathfrak{g})$
(Volume~II, Theorem~\ref*{thm:Koszul_dual_Yangian}).
For Virasoro and $\mathcal{W}_N$: the ordered Koszul dual is the
dg-shifted Yangian $Y^{\mathrm{dg}}$, a dg algebra with
infinitely many cohomological generators reflecting
non-formality (Volume~II, Part~VII).
For $\beta\gamma$/$bc$: the Koszul duality exchanges
statistics ($bc^! = \beta\gamma$, $\beta\gamma^! = bc$),
so the line-side algebra is the dual free-field system.

All entries are chirally Koszul (bar cohomology concentrated),
as established by PBW universality (Volume~I) and the
one-loop exactness criterion (Volume~II,
Theorem~\ref*{thm:one-loop-koszul}).
\end{proof}

\begin{remark}[The $\mathbf{G}/\mathbf{L}/\mathbf{C}/\mathbf{M}$
classification]%
\label{rem:glcm-commentary}%
\index{shadow depth!G/L/C/M classification|textbf}%
The four classes are determined by the OPE pole structure
on generators:
\begin{enumerate}[label=\textup{(\roman*)}]
\item Class~$\mathbf{G}$ \textup{(}Gaussian\textup{)}:
  double-pole OPE with scalar residues.
  The bar complex is cogenerated in arity~$2$,
  the $R$-matrix is scalar-valued (abelian braiding),
  and the Lagrangian self-intersection
  $\mathfrak{S}_b = L \times_X L$ is clean.
  Archetype: $\mathcal{H}_k$.

\item Class~$\mathbf{L}$ \textup{(}Lie\textup{)}:
  double-pole OPE with non-scalar
  \textup{(}Lie-bracket\textup{)} residues.
  The Lie bracket creates a transverse intersection defect
  at codimension~$1$, but the Jacobi identity kills all
  higher defects.  The $R$-matrix is the Yang $R$-matrix:
  matrix-valued, satisfying the Yang--Baxter equation
  with spectral parameter.
  Archetype: $\widehat{\mathfrak{g}}_k$.

\item Class~$\mathbf{C}$ \textup{(}Contact\textup{)}:
  simple-pole OPE on generators.
  The composite-field OPE reaches pole order~$3$, producing a
  quartic contact invariant.  Rank-one abelian rigidity
  forces the quintic obstruction to vanish.
  The $R$-matrix is a unipotent perturbation of the identity
  ($\Theta^2 = 0$ forces $R = I + 2\pi i\,\Theta$).
  Archetype: $\beta\gamma$.

\item Class~$\mathbf{M}$ \textup{(}Modular\textup{)}:
  quartic or higher poles on generators.
  The $R$-matrix is not expressible in closed form.
  The shadow tower does not terminate: the composite-field
  exchange generates new obstructions at every arity.
  Non-formality is the mathematical expression of the
  curved bar structure $d^2 = \kappa(A) \cdot \omega_g$
  at genus~$g \ge 1$.
  Archetype: $\mathrm{Vir}_c$.
\end{enumerate}
\end{remark}

\begin{remark}[Formality, non-formality, and the
$\mathsf{E}_1/\mathsf{E}_\infty$ distinction]%
\label{rem:formality-e1-e-infty}%
\index{formality!Swiss-cheese structure}%
\index{E-infinity@$\mathsf{E}_\infty$-chiral!all standard families}%
\index{E-one@$\mathsf{E}_1$-chiral!genuinely nonlocal}%
A critical conceptual point: non-formality \emph{does not}
imply $\mathsf{E}_1$.  All ten families in
Theorem~\textup{\ref{thm:e1-ordered-bar-landscape}} are
$\mathsf{E}_\infty$-chiral \textup{(}local vertex algebras with
$\Sigma_n$-equivariant factorization structure\textup{)}, including
the non-formal families $\mathrm{Vir}_c$, $\mathcal{W}_3$, and
$\mathcal{W}_N$.  Having OPE poles does not break
$\mathsf{E}_\infty$; having non-vanishing higher $\Ainf$
operations does not break $\mathsf{E}_\infty$.  The poles
are singularities of a \emph{local} product, and the
non-formality measures structural complexity \emph{within}
the $\mathsf{E}_\infty$ world.

The three-tier picture is:
\begin{enumerate}[label=\textup{(\alph*)}]
\item \emph{Pole-free commutative}:
  $R(z) = \tau$ (the Koszul-signed flip).
  This is the subclass where the chiral product extends
  across the diagonal without singularities.

\item \emph{Vertex algebras with poles}
  \textup{(}all families in the table above\textup{)}:
  $R(z) \ne \tau$, but $R(z)$ is \emph{derived} from the
  local OPE via analytic continuation.
  Both (a) and (b) are $\mathsf{E}_\infty$-chiral.

\item \emph{Genuinely $\mathsf{E}_1$-chiral}
  \textup{(}Etingof--Kazhdan quantum vertex algebras,
  Yangians as independent algebras, quantum groups\textup{)}:
  $R(z) \ne \tau$ and $R(z)$ is \emph{independent input},
  not derivable from a local symmetric structure.
  Only (c) is genuinely $\mathsf{E}_1$.
\end{enumerate}
Tiers~(a) and~(b) are both $\mathsf{E}_\infty$.
The ordered bar complex is a computational tool for
$\mathsf{E}_\infty$-chiral algebras; it does not make them
$\mathsf{E}_1$.
\end{remark}

\begin{remark}[Pole escalation: DS reduction
transports $\mathbf{L} \to \mathbf{M}$]%
\label{rem:pole-escalation}%
\index{Drinfeld--Sokolov reduction!pole escalation}%
\index{pole escalation|textbf}%
Principal Drinfeld--Sokolov reduction sends
$V_k(\mathfrak{g}) \mapsto \mathcal{W}_N(\mathfrak{g})$.
On the ordered bar complex, this transports class~$\mathbf{L}$
to class~$\mathbf{M}$: the affine double-pole OPE
$J^a(z)J^b(w) \sim k\delta^{ab}/(z{-}w)^2 + \cdots$
becomes the Virasoro quartic-pole OPE
$T(z)T(w) \sim (c/2)/(z{-}w)^4 + \cdots$ via the Sugawara
construction.  The shadow depth escalates from $r_{\max} = 3$
(finite, formal) to $r_{\max} = \infty$ (infinite, non-formal).
The Yangian $Y_\hbar(\mathfrak{g})$ (a classical algebra
concentrated in degree~$0$) becomes the dg-shifted Yangian
$Y^{\mathrm{dg}}(\mathrm{Vir}_c)$ (a dg algebra with
infinitely many cohomological generators).

This pole escalation is the ordered-bar avatar of the
\emph{single-line dichotomy} of Volume~I: double poles
produce a linear (tree-level) bar differential, while
quartic poles produce a quadratic (one-loop) bar differential
whose iterations generate the infinite tower.
The mechanism is the Sugawara pole:
$T = \tfrac{1}{2(k+h^\vee)} \sum_a {:}J^a J_a{:}$
is a normally ordered composite, and the OPE of composites
escalates pole orders.  Theorem~\ref{thm:ordered-associative-ds-principal}
ensures that this pole escalation is compatible with
the ordered bar--cobar equivalence.
\end{remark}
